\documentclass[11pt,a4paper]{article}
\usepackage[T1]{fontenc}
\usepackage{lmodern,amsmath,amssymb,amsthm,mathtools,mathrsfs,booktabs,longtable,enumitem}
\usepackage[margin=27mm]{geometry}
\usepackage{microtype}
\usepackage[hidelinks]{hyperref}
\hypersetup{pdftitle={Theta Prime Gamma Energy and its Spectral Structure},pdfauthor={Jongmin Choi},pdfsubject={Closed domains, global comparison bounds, scaled Schur reductions, weighted spectral synthesis and exact zero jet cores},pdfkeywords={theta kernel, prime powers, Gamma form, spectral synthesis, zero jets, Schur complement, variance bound}}
\setlength{\parindent}{1.2em}\setlength{\parskip}{0.2em}
\setlength{\emergencystretch}{2em}\allowdisplaybreaks[2]
\numberwithin{equation}{section}
\newtheorem{theorem}{Theorem}[section]
\newtheorem{proposition}[theorem]{Proposition}
\newtheorem{lemma}[theorem]{Lemma}
\newtheorem{corollary}[theorem]{Corollary}
\newtheorem{hypothesis}[theorem]{Hypothesis}
\theoremstyle{definition}\newtheorem{definition}[theorem]{Definition}
\theoremstyle{remark}\newtheorem{remark}[theorem]{Remark}
\newcommand{\R}{\mathbb R}\newcommand{\C}{\mathbb C}\newcommand{\N}{\mathbb N}
\newcommand{\E}{\mathcal E}\newcommand{\D}{\mathcal D}
\newcommand{\He}{L^2_{\mathrm{even}}(\mu)}\newcommand{\Le}{L^2_{\mathrm{even}}(dx)}
\newcommand{\Hh}{\mathcal H}\newcommand{\Hw}{\mathcal H}
\newcommand{\HH}{\mathcal H_W}\newcommand{\XX}{X_W}\newcommand{\EE}{\mathcal E_{N,\beta}}
\newcommand{\tg}{t_\Gamma}\newcommand{\cg}{c_\Gamma}
\newcommand{\qD}{\mathfrak q_D}\newcommand{\qth}{\mathfrak q_\theta}
\newcommand{\dom}{\operatorname{Dom}}\newcommand{\Var}{\operatorname{Var}}
\newcommand{\rank}{\operatorname{rank}}\newcommand{\ess}{\operatorname{ess}}\newcommand{\supp}{\operatorname{supp}}
\newcommand{\ip}[2]{\langle#1,#2\rangle}\newcommand{\norm}[1]{\lVert#1\rVert}\newcommand{\abs}[1]{\lvert#1\rvert}
\newcommand{\dd}{\mathop{}\!\mathrm d}\newcommand{\one}{\mathbf1}
\newcommand{\Zc}{\mathcal Z_\delta}\newcommand{\Rcal}{\mathcal R}
\newcommand{\RR}{\mathscr R}\newcommand{\SSS}{\mathcal S_{\!\exp}}\newcommand{\CE}{\mathcal C_{\exp}}
\newcommand{\Eloc}{\mathcal E_{\mathrm{loc}}}\newcommand{\Def}{\mathfrak D}
\newcommand{\PP}{P}\newcommand{\J}{J}\newcommand{\V}{V}\newcommand{\wo}{\mathcal W}
\newcommand{\Nn}{\mathcal N}\newcommand{\Ii}{\mathcal I}\newcommand{\Zz}{\mathcal Z}
\newcommand{\Kk}{\mathcal K}\newcommand{\Aa}{\mathcal A}\newcommand{\Ff}{\mathcal F}
\newcommand{\cl}{\overline{\operatorname{span}}}\newcommand{\Ran}{\operatorname{Ran}}\newcommand{\HS}{\mathfrak S_2}
\DeclareMathOperator{\Rea}{Re}\DeclareMathOperator{\Ima}{Im}
\title{Theta Prime Gamma Energy\\and its Spectral Structure}
\author{Jongmin Choi\\[3pt]\small Independent Researcher, Seoul, Korea\\\small ORCID iD: \href{https://orcid.org/0009-0008-7448-514X}{0009-0008-7448-514X}}
\date{November 2, 2026}
\begin{document}\maketitle
\begin{abstract}
We develop a single closed-form and spectral framework for the even theta-weighted energy containing the complete Gamma difference integral and every prime power. The maximal integral domain equals its compact smooth minimum closure. An explicit-formula identity gives its exact weighted physical realization and a compact-resolvent ordinary-$L^2$ pencil. Derivative radicals yield an infinite-dimensional half-eigenspace and essential spectral bottom $1/2$, without excluding lower discrete spectrum. A direct analytic variance bound is strengthened to $49307/100000$ using an accepted whole-local-domain sign input at radius $4/5$; an additional, explicitly conditional response contract gives $2861149/5760000$. All mean repairs, prime tails, touching Gamma kernels and response costs are retained. We construct delta-scaled finite Schur reductions and establish their error budgets, while identifying the unproved continuation sign. We also prove the quarter-curvature obstruction and the impossibility of a common nontrivial positive local-row factorization. Full-multiplicity weighted spectral synthesis follows from an actual-weight Abel endpoint estimate, and dual source recovery gives full-zero-jet form and graph cores. For the actual physical operator the complementary zero-jet sector is reducing; consequently the global half inequality is equivalent to positivity of every finite exact zero-jet matrix. Fixed-complexity high-frequency coercivity is uniform through complex collisions and arbitrary separations, but gives no complexity-independent sign. The global half inequality and the Riemann hypothesis remain unproved.
\end{abstract}
\noindent\textbf{Keywords.} Theta kernel; complete prime-power form; Gamma difference form; Schur complement; weighted spectral synthesis; zero-jet core.
\tableofcontents

\section{Introduction and logical structure}\label{sec:intro}
The object of this paper is one nonnegative jump energy, rather than a family of truncated prime or finite-matrix substitutes. Its centered variance gap can be studied in three complementary coordinates: profiles in $L^2(\mu)$, physical numerators in a reciprocal-theta weighted space, and an ordinary-$L^2$ confining pencil. The distinction between these spaces is essential. A form congruence is not an ordinary-$L^2$ unitary equivalence, a convergent response calculation does not determine a Schur sign, and approximation by full zero jets does not prove positivity on those jets.

The main analytic conclusions are the full-domain physical bridge, the exact spectral and radical identities, parameter-scaled finite complements, weighted spectral synthesis, stronger-topology zero-jet cores, and fixed-complexity coercivity. Two quantitative comparison branches have different computational premises. The first uses only an accepted sign theorem on an entire local minimum domain. The second assumes additional complete-operator residual and centered-response enclosures. These implications are stated separately, and neither input class is identified with a finite pivot computation.

\subsection{Common notation}
All physical spaces are complex and even unless a real profile is explicitly required. Inner products are linear in their first argument. The positive even kernel $\kappa$ is constructed in Section~\ref{sec:foundations}. Set
\begin{equation}\label{a:eq:weights}
 c(x)=\cosh(x/2),\qquad d\mu=2c\kappa\,dx,\qquad
 W=\frac{\kappa}{2c},\qquad \Hh=L^2_{\mathrm{even}}(\R,dx/W).
\end{equation}
The measure $\mu$ has mass one. Write $\norm{u}_{\Hh}^2=\int|u|^2/W$, and reserve $\norm{\cdot}_2$ for ordinary Lebesgue $L^2$. With $\Lambda$ the von Mangoldt function, define
\begin{equation}\label{eq:common-shifts}
 \rho(s)=\frac{e^{-s/2}}{1-e^{-2s}},\qquad
 a_n=\frac{\Lambda(n)}{\sqrt n},\qquad s_n=\log n.
\end{equation}
Thus $a_n$ retains every prime power, not merely primes. The exact numerator forms are
\begin{align}
 \tg[u]&=\int_0^\infty\rho(s)\norm{u(\cdot+s)-u}_2^2\,ds,\label{a:eq:physicalgamma}\\
 p[u]&=2\operatorname{Re}\sum_{n\ge2}a_n\int_\R u(x+s_n)\overline{u(x)}\,dx,\label{a:eq:physicalprime}\\
 q_0[u]&=\tg[u]-\cg\norm{u}_2^2-p[u],\qquad
 \cg=\log\pi-\psi(1/4),\quad\psi=\Gamma'/\Gamma.\label{a:eq:qzero}
\end{align}
The complete prime form is bounded on $\Hh$. The common numerator domain and centered space are
\begin{equation}\label{a:eq:X}
 \begin{gathered}
 X=\{u\in\Hh:\tg[u]<\infty\},\qquad
 \norm{u}_X^2=\norm{u}_{\Hh}^2+\tg[u],\\
 m(u)=\int_\R cu,\quad \Hh_c=\ker m,\quad X_c=X\cap\Hh_c.
 \end{gathered}
\end{equation}
The original energy $\E$ and its maximal domain $\D$ are defined in~\eqref{a:eq:Eg}--\eqref{a:eq:Edomain}. Their bridge is
\begin{equation}\label{eq:intro-bridge}
 \kappa\D=X,\qquad q_0[\kappa f]=\E[f]-\tfrac12\norm f_\mu^2,
 \qquad \mu(f)=2m(\kappa f).
\end{equation}
Multiplication $Uf=\kappa f$ is unitary from $\He$ to $\Hh$, but not to ordinary $\Le$. We use $U$ only for this unitary; the letters $J$ occurring in interface kernels or zero jets have their separately displayed meanings. The entire transform is $F_u(z)=\int u(x)\cos(zx)\,dx$; on even functions it equals $\mathcal Fu(z)=\int u(x)e^{izx}\,dx$. The ordinary unitary transform, used in multiplier estimates, is $\widehat u(t)=(2\pi)^{-1/2}\int u(x)e^{-itx}\,dx$.

\subsection{Results and their dependency classes}
\begin{center}\small
\begin{tabular}{@{}p{0.15\textwidth}p{0.52\textwidth}p{0.25\textwidth}@{}}
\toprule Class & Conclusion & Premise beyond the reproduced analysis\\\midrule
Analytic & Closed domains, actual operator, compact pencil, radicals, synthesis and zero-jet cores; fixed-complexity coercivity & Classical explicit formula and standard analytic facts\\[3pt]
Analytic & $\E\ge c_{\log2}\operatorname{Var}_\mu$, with $c_{\log2}>8939/425000$ & None computational\\[3pt]
Local sign & $\E\ge(49307/100000)\operatorname{Var}_\mu$ & Accepted whole-local-domain sign at $B_0=4/5$\\[3pt]
Response conditional & $\E\ge(2861149/5760000)\operatorname{Var}_\mu$ & Additional full residual and centered-response contracts\\[3pt]
Open & $\E\ge\tfrac12\operatorname{Var}_\mu$; all finite exact full-zero-jet matrices positive semidefinite & No proof is supplied\\\bottomrule
\end{tabular}
\end{center}
The analytic seed is useful because it is independent of the local certificate, not because it is the strongest comparison branch. The accepted local-sign coefficient is $0.49307$. The stronger coefficient $2861149/5760000$ is approximately $0.496727$ and has additional hypotheses. Their small remaining losses are positive; no limit argument removes them.

The proof order is as follows. Sections~\ref{sec:foundations}--\ref{sec:spectral} construct the common domains and actual spectral objects. Sections~\ref{sec:bounds} and~\ref{sec:schur-unified} give global comparisons and finite reductions. Section~\ref{sec:obstructions} establishes pointwise and local obstructions. Section~\ref{sec:synthesis-unified}, supported by the internal Abel and compact-model appendices, proves full weighted synthesis and stronger zero-jet closure. Section~\ref{sec:actual-jets} specializes those closures to the actual physical operator and proves the precise remaining sign equivalence. Section~\ref{sec:complexity-unified} establishes the fixed-complexity result. Long scalar, interface, response and cluster proofs are kept in internal appendices; no prior manuscript is used as a premise.

\section{Common analytic foundations}\label{sec:foundations}


\subsection{The theta kernel and the original closed energy}\label{a:sec:energy}

For $x\ge0$, define
\begin{equation}\label{a:eq:kappa}
 \kappa(x)=e^{x/2}\sum_{j\ge1}(4\tau_j(x)^2-6\tau_j(x))e^{-\tau_j(x)},
 \qquad \tau_j(x)=\pi j^2e^{2x},
\end{equation}
and extend evenly. We record the normalization and decay needed later.

\begin{lemma}[Smoothness and normalization]\label{a:lem:theta}
The kernel $\kappa$ is smooth, strictly positive, and even. There is a finite $C_\kappa>0$ such that
\begin{equation}\label{a:eq:envelopek}
 \kappa(x)\le C_\kappa\exp\!\left(\frac92|x|-\pi e^{2|x|}\right),
 \qquad
 C_\kappa=4\pi^2\sum_{j\ge1}j^4e^{-\pi(j^2-1)}.
\end{equation}
Furthermore $\int_\R2c\kappa\,dx=1$, and
\begin{equation}\label{a:eq:envelopeW}
 0<W(x)\le C_\kappa\exp(4|x|-\pi e^{2|x|}).
\end{equation}
\end{lemma}
\begin{proof}
The series in~\eqref{a:eq:kappa}, together with every derivative, converges uniformly on compact subsets of $[0,\infty)$. Each summand is positive because $\tau_j\ge\pi>3/2$. Dropping the negative polynomial term and using $e^{2x}\ge1$ gives~\eqref{a:eq:envelopek} for $x\ge0$.

To justify smooth evenness at zero, rather than assume it from reflection, set
\[
 H(x)=e^{x/2}\sum_{j\ge1}e^{-\pi j^2e^{2x}},\qquad x\in\R.
\]
The Jacobi transformation \cite[Eq.~20.7.32]{DLMF} implies
$H(-x)=H(x)+\sinh(x/2)$. Applying $D^2-1/4$ shows that $H''-H/4$ is even; direct differentiation identifies it with~\eqref{a:eq:kappa} on the positive half-line. Thus it gives the claimed smooth extension. Differentiating the transformation at zero yields $H'(0)=-1/4$. Since $c''=c/4$, integration by parts and the rapid decay of $H,H'$ at $+\infty$ give
\[
 \int_0^\infty c(x)\kappa(x)\,dx
  =[cH'-c'H]_0^\infty=\frac14.
\]
Evenness proves the normalization. Finally $2c(x)\ge e^{|x|/2}$ gives~\eqref{a:eq:envelopeW}.
\end{proof}

Put $\rho(s)=e^{-s/2}/(1-e^{-2s})$ for $s>0$, and $a_n=\Lambda(n)/\sqrt n$. The full positive energy is
\begin{align}
 \E_\Gamma[f]&=\int_0^\infty\rho(s)\int_\R
       \kappa(x)\kappa(x+s)|f(x+s)-f(x)|^2\,dx\,ds,\label{a:eq:Eg}\\
 \E_P[f]&=\sum_{n\ge2}a_n\int_\R
       \kappa(x)\kappa(x+s_n)|f(x+s_n)-f(x)|^2\,dx,\label{a:eq:Ep}\\
 \E[f]&=\E_\Gamma[f]+\E_P[f],\qquad
 \D=\{f\in\He:\E[f]<\infty\}.\label{a:eq:Edomain}
\end{align}
These are nonnegative integrals and sums, interpreted by Tonelli's theorem; no cancellation defines their domain. Write
\[
 \mu(f)=\int_\R f\,d\mu,\qquad
 \Var_\mu(f)=\norm{f}_\mu^2-|\mu(f)|^2.
\]

\begin{proposition}[Maximal domain and compact smooth form core]\label{a:prop:core}
The form $\E$ on $\He$ is densely defined, nonnegative, and closed. Its maximal integral domain $\D$ is the closure of even $C_c^\infty(\R)$ in the norm $(\norm{f}_\mu^2+\E[f])^{1/2}$. Constants belong to $\D$, and $\E[f]=0$ holds precisely for almost-everywhere constant $f$.
\end{proposition}
\begin{proof}
The theta envelope implies that for every $A_0>0$ there is a constant $C_{A_0}$ such that $\kappa(x)\le C_{A_0}e^{-A_0|x|}$. Consequently
\begin{equation}\label{a:eq:edgefinite}
 \int_\R\kappa(x)\kappa(x+s)\,dx
 \le C_{A_0}^2(s+A_0^{-1})e^{-A_0s},\qquad s\ge0.
\end{equation}
Since $a_n\le n$, the complete prime edge measure has finite mass, by taking $A_0>2$ in~\eqref{a:eq:edgefinite}. For every $L<\infty$,
\begin{equation}\label{a:eq:rhointegrable}
 \int_0^\infty\rho(s)\min\{1,L^2s^2\}\,ds<\infty.
\end{equation}
These observations admit compact smooth functions and constants to the domain.

For closedness, let $(f_j)$ be Cauchy in the form norm, with $L^2(\mu)$ limit $f$. A subsequence converges almost everywhere for $\mu$, hence for Lebesgue measure, because the density is strictly positive. A countable union of translated null sets is null, so all prime differences converge almost everywhere. Fubini's theorem gives convergence for the Gamma edge measure as well. For each fixed $k$, Fatou's lemma gives
\[
 \E[f-f_k]\le\liminf_{j\to\infty}\E[f_j-f_k].
\]
It follows that $f\in\D$ and $f_k\to f$ in form norm.

We prove the core assertion first for real profiles. Let $T_M(t)=\max\{-M,\min\{t,M\}\}$. The remainder $r_M(t)=t-T_M(t)$ is $1$-Lipschitz, tends pointwise to zero, and satisfies $|r_M(t)|\le|t|$. Thus dominated convergence in both edge measures and in $L^2(\mu)$ gives $T_M(f)\to f$ in form norm. It suffices to approximate bounded even profiles.

Choose smooth even cutoffs $0\le\chi_R\le1$ equal to one on $[-R,R]$, supported in $[-2R,2R]$, and with uniformly bounded Lipschitz constants. Equation~\eqref{a:eq:rhointegrable}, the finite prime-edge mass, and dominated convergence give $\E[\chi_R]\to0$. For bounded $f$, the identity for a product difference yields
\begin{align*}
 |[f(1-\chi_R)](y)-[f(1-\chi_R)](x)|^2
 &\le2(1-\chi_R(y))^2|f(y)-f(x)|^2\\
 &\quad+2\norm{f}_\infty^2|\chi_R(y)-\chi_R(x)|^2.
\end{align*}
Hence $f\chi_R\to f$ in form norm. It remains to smooth a compactly supported profile $g\in\D$.

Suppose $\supp g\subset[-K,K]$. Its ordinary $L^2$ norm is finite, since the density of $\mu$ is bounded below on this interval. For $0<s\le1$, a nonzero difference of $g$ has both endpoints in $[-K-1,K+1]$. The positive lower bound of $\kappa$ there controls the small-shift part of $\tg[g]$ by $\E_\Gamma[g]$. The part with $s>1$ is bounded by $4\norm{g}_2^2\int_1^\infty\rho(s)\,ds$. Thus $\tg[g]<\infty$.

Let $g_\epsilon=\zeta_\epsilon*g$, where $\zeta_\epsilon$ is an even compactly supported smooth approximate identity. Translations commute with convolution, and
\[
 \norm{\Delta_s(g_\epsilon-g)}_2\le2\norm{\Delta_sg}_2,
 \qquad\Delta_sg=g(\cdot+s)-g.
\]
Dominated convergence gives $\tg[g_\epsilon-g]\to0$. Since $\kappa$ is bounded,
$\E_\Gamma[g_\epsilon-g]\le\norm{\kappa}_\infty^2\tg[g_\epsilon-g]\to0$.
All the smoothed profiles have support in one fixed compact interval $J$. For any $h$ supported there,
\begin{equation}\label{a:eq:compactprimerate}
 \E_P[h]\le2\int_J\kappa(x)|h(x)|^2
     \sum_{n\ge2}a_n\bigl(\kappa(x-s_n)+\kappa(x+s_n)\bigr)\,dx
 \le C_J\norm{h}_2^2.
\end{equation}
Indeed, the complete rate in parentheses is uniformly bounded on $J$: use $\kappa(x\pm\log n)\le C_{A_0}e^{A_0\sup_J|x|}n^{-A_0}$ and $A_0>2$. Ordinary $L^2$ convergence therefore proves prime-energy convergence. The local upper bound of the density gives $L^2(\mu)$ convergence. Applying the argument to real and imaginary parts proves the complex core assertion.

Finally $\E[1]=0$. If $\E[f]=0$, the strictly positive Gamma density implies $f(y)=f(x)$ for almost every pair $x<y$. Fubini's theorem forces $f$ to be almost everywhere constant. Density of the core in $\He$ is immediate from the positive smooth density of $\mu$.
\end{proof}

The representation theorem for closed nonnegative forms \cite[Chapter~VI, Section~2]{Kato} defines $A$ by
\begin{equation}\label{a:eq:Adef}
 \begin{split}
 f\in\dom(A),\ Af=h\quad\Longleftrightarrow\quad
 f\in\D,\ h\in\He,\quad
 \E(f,g)=\ip{h}{g}_\mu\ \text{for all }g\in\D.
 \end{split}
\end{equation}
The operator is self-adjoint, $A1=0$, and $\ker A=\C1$. The centered space
$H_0=\{f\in\He:\mu(f)=0\}$ reduces $A$; its restriction is $A_E$. This construction uses only the nonnegative integral energy, before any identification with the physical form.


\subsection{Explicit bounds for the complete form}\label{a:sec:prime}

The bounds in this section account for all theta summands and all prime powers. A finite head is used only together with its complete remainder. We write $w_{\max}=\norm W_\infty$; the symbol $J$ will later denote a Neumann iteration index.

\begin{lemma}[Elementary theta bounds]\label{a:lem:explicitW}
For $r=|x|$,
\begin{equation}\label{a:eq:explicit_envelope}
 0<\kappa(x)\le42e^{(9/2)r-\pi e^{2r}},\qquad
 0<W(x)\le42e^{4r-\pi e^{2r}},\qquad w_{\max}<\frac12.
\end{equation}
For every integer $R\ge1$,
\begin{equation}\label{a:eq:rationalOmega}
 \Omega_R:=\sup_{|x|>R}W(x)
 <\overline\Omega(R):=42\,81^R(3/8)^{3\cdot7^R}.
\end{equation}
\end{lemma}
\begin{proof}
We use $3<\pi<22/7$, $\pi^2<10$, and the elementary exponential bounds
\begin{equation}\label{a:eq:elementarye}
 \frac{163}{60}=\sum_{j=0}^5\frac1{j!}<e
 <\frac{163}{60}+\frac1{720}\sum_{k\ge0}7^{-k}
 =\frac{11743}{4320}<\frac{87}{32}<\frac{11}{4}.
\end{equation}
In particular $e>27/10>8/3$, $e^2>7$, and $e^3>20$, since
$163^3=4330747>4320000=20\cdot60^3$.
The theta coefficient in Lemma~\ref{a:lem:theta} is less than $42$. Indeed, the tail of
$\sum_{k\ge1}k^4e^{-\pi(k^2-1)}$ starts below $16e^{-9}$, and its successive ratios, for $k\ge2$, are at most
$(3/2)^4e^{-15}<1/2$. Thus that sum is below
$1+32e^{-9}<26/25$, and
$4(22/7)^2(26/25)<42$.
This proves the two envelopes in~\eqref{a:eq:explicit_envelope}.

For the sharper supremum bound put $F(t)=(4t^2-6t)e^{-t}$. Since
\[
 F'(t)=-2(2t-1)(t-3)e^{-t},\qquad
 W(x)=\frac{\sum_{k\ge1}F(\pi k^2e^{2r})}{1+e^{-r}},
\]
$F$ decreases for $t\ge3$. The complete tail with $k\ge2$ satisfies
\[
 H(r):=\sum_{k\ge2}F(\pi k^2e^{2r})
 <\sum_{k\ge2}40k^4e^{-3k^2}<1280e^{-12}<\frac1{125}.
\]
The majorant starts at $640e^{-12}$ and has successive ratio at most
$(3/2)^4e^{-15}<1/2$; the last inequality uses $e^{12}>20^4=160000$.
Consider three fixed radial intervals. If $0\le r\le1/6$, then
$F(\pi e^{2r})\le F(3)=18e^{-3}<9/10$ and
$1+e^{-r}\ge2-r\ge11/6$, giving
\[
 W(x)<\frac{9/10+1/125}{11/6}=\frac{681}{1375}<\frac12.
\]
If $1/6\le r\le1/2$, then $\pi e^{2r}>4$,
$F(\pi e^{2r})<F(4)=40e^{-4}<20/27$, since $e^4>54$,
and $1+e^{-r}\ge3/2$. Therefore
\[
 W(x)<\frac{20/27+1/125}{3/2}=\frac{5054}{10125}<\frac12.
\]
Here $3/4-20/27=1/108>1/125$. Finally, if $r\ge1/2$, then
$\pi e^{2r}>3e>8$ and $e^8=e^6e^2>2800$, so
\[
 W(x)<208e^{-8}+\frac1{125}
 <\frac{13}{175}+\frac1{125}=\frac{72}{875}<\frac12.
\]
These uniform rational margins prove $w_{\max}<1/2$.

For~\eqref{a:eq:rationalOmega}, the function $4r-3e^{2r}$ decreases for $r\ge0$. Use $\pi>3$, $e^{4R}<81^R$, $e^{2R}>7^R$, and $e^{-1}<3/8$ in the envelope.
\end{proof}

\begin{lemma}[Complete weighted prime tail]\label{a:lem:prime}
For $N\ge2$, define
\begin{equation}\label{a:eq:rN}
 r_N=96(N+1)^3(27/512)^{N+1}.
\end{equation}
The full prime form is a bounded Hermitian form on $\Hw$, and
\begin{equation}\label{a:eq:headtail}
 |p[u]-p_N[u]|\le r_N\norm u_{\Hh}^2,
 \qquad |p_N[u]|\le2\sum_{2\le n\le N}a_n\norm u_2^2.
\end{equation}
In particular,
\begin{equation}\label{a:eq:head4}
 r_4=\frac{5380840125}{1099511627776}<\frac1{128},
 \qquad 0<\cg<6,
 \qquad \cg\norm u_2^2+|p_4[u]|\le9\norm u_2^2.
\end{equation}
Consequently the form
\begin{equation}\label{a:eq:Dform}
 d(u,z)=\cg\ip uz_2+p(u,z)
\end{equation}
is represented by a bounded self-adjoint operator $D$ on $\Hw$ with
\begin{equation}\label{a:eq:C0}
 \norm D\le C_0:=9+r_4<10.
\end{equation}
\end{lemma}
\begin{proof}
Let $\beta=\sqrt W$ and let $\tau_sF(x)=F(x+s)$ on ordinary $L^2$. With $z=|x|+|x+s|\ge s$, the envelope and the arithmetic--geometric mean inequality give
\begin{align*}
 \beta(x)\beta(x+s)
 &\le42\exp\left(2z-\frac\pi2(e^{2|x|}+e^{2|x+s|})\right)\\
 &\le42e^{2z-\pi e^z}\le42e^{2s-\pi e^s}.
\end{align*}
The last exponent decreases for $z\ge0$. Hence the self-adjoint series
\begin{equation}\label{a:eq:KP}
 K_P=\sum_{n\ge2}a_nM_\beta(\tau_{s_n}+\tau_{-s_n})M_\beta
\end{equation}
converges absolutely in operator norm: its $n$th term has norm at most
$84\Lambda(n)n^{3/2}e^{-\pi n}$.
Reflection interchanges the two shifts, so the sum preserves even functions. Substituting $u=\beta F$ gives
\begin{equation}\label{a:eq:prepresentation}
 p[u]=\ip{K_P(u/\beta)}{u/\beta}_2.
\end{equation}
This also proves absolute convergence of the original scalar series.

Using $\Lambda(n)\le\log n\le n$ and $n^{5/2}\le n^3$, its tail is bounded by
$84\sum_{n>N}n^3e^{-3n}$. For $n\ge3$ the successive ratio is below
$(4/3)^3(3/8)^3=1/8$. Summing the geometric majorant and using
$e^{-3}<27/512$ yields~\eqref{a:eq:rN}. Translation invariance and Cauchy--Schwarz give the finite-head estimate.

For the small head, the elementary bounds
\[
 \log2<7/10,\quad \log3<11/10,\quad
 \sqrt2>7/5,\quad \sqrt3>22/13
\]
give $a_2<1/2$, $a_3<13/20$, $a_4<7/20$, and thus
$2(a_2+a_3+a_4)<3$. The logarithm inequalities follow, for example, by keeping the first five terms of $e^{7/10}$ and the first six terms of $e^{11/10}$; these exceed $2$ and $3$, respectively. The square-root comparisons follow by squaring.
The digamma reflection and duplication formulas, together with $\psi(1)=-\gamma$, give
\[
 \cg=\log\pi+\gamma+\frac\pi2+3\log2.
\]
These classical formulas are recorded in \cite[Sections 5.4--5.5]{DLMF}. In particular $\cg>0$. For the upper bound, use $0<\gamma<1$, $\pi/2<11/7$, $\log2<7/10$, and $\log\pi<6/5$. The last inequality follows from $\pi<22/7<e^{6/5}$, using the first five terms of the exponential series. Their sum gives
$\cg<6/5+1+11/7+21/10<6$.
Substitution $N=4$ in~\eqref{a:eq:rN} yields the fraction in~\eqref{a:eq:head4}; comparison with $1/128$ is immediate by cross multiplication.
Finally~\eqref{a:eq:head4} and $w_{\max}<1$ imply
$|d[u]|\le(9+r_4)\norm u_{\Hh}^2$. Polarization and the bounded-form representation theorem give $D$ and~\eqref{a:eq:C0}.
\end{proof}

The operator $K_P$ is not claimed to be compact. A decaying multiplication operator composed with a translation need not be compact on a nonatomic $L^2$ space. The bounded operator $D$ acts in the weighted space; its existence does not assert that the full unweighted prime row belongs to ordinary $L^2$.


\subsection{The explicit formula and the compact physical bridge}\label{a:sec:bridge}

We now prove the identity that connects the positive energy to the physical form. The classical explicit formula is used as a theorem, without any assumption about the location of nontrivial zeta zeros inside the critical strip. We first fix its normalization, then justify its use with the noncompact theta kernel, and finally perform the ground-state calculation with all endpoints retained.

In this section use the nonunitary Fourier transform
\[
 \mathcal Fu(z)=\int_\R u(x)e^{izx}\,dx,
 \qquad m(u)=\int_\R c(x)u(x)\,dx.
\]
Let $\rho_\zeta$ range over the distinct nontrivial zeros of $\zeta$, with multiplicity $\nu_{\rho_\zeta}$, and put $w_{\rho_\zeta}=(\rho_\zeta-1/2)/i$. Thus $|\operatorname{Im}w_{\rho_\zeta}|<1/2$; no reality of $w_{\rho_\zeta}$ is asserted. We use the classical zero-counting bound
\begin{equation}\label{a:eq:zerocount}
 \sum_{|\operatorname{Im}\rho_\zeta|\le T}\nu_{\rho_\zeta}
       =O(T\log(T+2))\qquad(T\ge2).
\end{equation}
This weak form of the classical Riemann--von Mangoldt estimate \cite{vonMangoldt} is sufficient for every zero sum below.

\begin{theorem}[Classical Guinand explicit formula]\label{a:thm:guinand}
For $g\in C_c^\infty(\R)$,
\begin{align}\label{a:eq:guinand}
 \sum_{\rho_\zeta}\nu_{\rho_\zeta}\mathcal Fg(w_{\rho_\zeta})
 &=\mathcal Fg(i/2)+\mathcal Fg(-i/2)\notag\\
 &\quad+\frac1{2\pi}\int_\R
   \bigl(\operatorname{Re}\psi(1/4+it/2)-\log\pi\bigr)
       \mathcal Fg(t)\,dt\notag\\
 &\quad-\sum_{n\ge2}a_n\bigl(g(s_n)+g(-s_n)\bigr).
\end{align}
The zero sum is absolutely convergent. This is the classical formula in the Fourier convention displayed here \cite{Guinand}, with prime powers collected by the von Mangoldt function.
\end{theorem}

The statement is quoted, not proposed as a new result. Absolute convergence follows from the rapid decay of $\mathcal Fg$ in the closed strip $|\operatorname{Im}z|\le1/2$ and~\eqref{a:eq:zerocount}.

\begin{lemma}[Bilinear physical normalization]\label{a:lem:bilinear}
For real even $u,v\in C_c^\infty(\R)$,
\begin{equation}\label{a:eq:bilinearzeros}
 q_0(u,v)+2m(u)m(v)
   =\sum_{\rho_\zeta}\nu_{\rho_\zeta}
       \mathcal Fu(w_{\rho_\zeta})\mathcal Fv(w_{\rho_\zeta}).
\end{equation}
\end{lemma}
\begin{proof}
The digamma integral \cite[Eq.~5.9.16]{DLMF} gives, after subtracting its value at $1/4$ and setting its integration variable equal to $2s$,
\begin{equation}\label{a:eq:digammamultiplier}
 \operatorname{Re}\psi(1/4+it/2)-\psi(1/4)
    =2\int_0^\infty\rho(s)(1-\cos(ts))\,ds=:d_\Gamma(t).
\end{equation}
The difference integral is convergent at zero and infinity. Plancherel and polarization identify
\[
 \tg(u,v)=\frac1{2\pi}\int_\R
       d_\Gamma(t)\mathcal Fu(t)\mathcal Fv(t)\,dt
\]
for real even tests. Apply Theorem~\ref{a:thm:guinand} to $g=u*v$. Its Fourier transform is the product of the two transforms, and $g(0)=\int uv$. Evenness gives
\[
 \mathcal Fu(i/2)=\mathcal Fu(-i/2)=m(u),\qquad
 g(s_n)+g(-s_n)
  =\int_\R\bigl(u(x+s_n)v(x)+u(x)v(x+s_n)\bigr)\,dx.
\]
These are respectively the pole and the polarized prime term. The constant in the Gamma term is $\psi(1/4)-\log\pi=-\cg$, proving~\eqref{a:eq:bilinearzeros}.
\end{proof}

For complex even tests the zero-side product is
$\mathcal Fu(w)\overline{\mathcal Fv(\overline w)}$. It is not generally $|\mathcal Fu(w)|^2$ when $u=v$ and $w$ is nonreal. We need only the real bilinear identity and its conjugate-linear extension in the second argument; no positivity is inferred from this zero sum.

\begin{lemma}[Entire theta transform]\label{a:lem:thetaxi}
Let
\[
 \xi(s)=\tfrac12s(s-1)\pi^{-s/2}\Gamma(s/2)\zeta(s).
\]
Then the entire continuation of this function satisfies
\begin{equation}\label{a:eq:thetaxi}
 \mathcal F\kappa(z)=\xi(1/2+iz),\qquad z\in\C.
\end{equation}
In particular, $\mathcal F\kappa(w_{\rho_\zeta})=0$ for every nontrivial zero, whether or not it is on the critical line.
\end{lemma}
\begin{proof}
Write $\Theta(t)=\sum_{j\ge1}e^{-\pi j^2t}$, so that the function in Lemma~\ref{a:lem:theta} is $H(x)=e^{x/2}\Theta(e^{2x})$. For $\operatorname{Re}s>1$, termwise integration gives
\[
 \pi^{-s/2}\Gamma(s/2)\zeta(s)
   =\int_0^\infty\Theta(t)t^{s/2-1}\,dt.
\]
The Jacobi transformation is
$\Theta(t)=t^{-1/2}\Theta(1/t)+(t^{-1/2}-1)/2$.
Split the integral at one and substitute $t\mapsto1/t$ in its lower part. The elementary remainder integrates to $1/[s(s-1)]$, and hence
\begin{align}\label{a:eq:mellintheta}
 \pi^{-s/2}\Gamma(s/2)\zeta(s)
 &=\frac1{s(s-1)}+I(s),\notag\\
 I(s)&=\int_1^\infty\Theta(t)
       \bigl(t^{s/2-1}+t^{(1-s)/2-1}\bigr)\,dt.
\end{align}
The integral $I$ is entire by local uniform double-exponential domination after $t=e^{2x}$. Thus
$\xi(s)=1/2+s(s-1)I(s)/2$ is entire. For $s=1/2+iz$, the same substitution gives
\begin{equation}\label{a:eq:xiH}
 \xi(1/2+iz)=\frac12
       -2(z^2+1/4)\int_0^\infty H(x)\cos(zx)\,dx.
\end{equation}
Because $\kappa=H''-H/4$ on the positive half-line, evenness and two integrations by parts give
\[
 \mathcal F\kappa(z)
  =-2H'(0)-2(z^2+1/4)\int_0^\infty H(x)\cos(zx)\,dx.
\]
All boundary terms at infinity vanish for each complex $z$. The value $H'(0)=-1/4$ from Lemma~\ref{a:lem:theta} proves~\eqref{a:eq:thetaxi}. The defining prefactors of $\xi$ are finite and nonzero at a nontrivial zero of $\zeta$, so the last assertion follows.
\end{proof}

\begin{proposition}[Weak theta source identity]\label{a:prop:thetanull}
For every even $h\in C_c^\infty(\R)$,
\begin{equation}\label{a:eq:thetanull}
 q_0(\kappa,h)=-\overline{m(h)}.
\end{equation}
All terms on the left exist separately as convergent form pairings. Equivalently, $(D_\Gamma-\cg-P)\kappa=-c$ on even compact smooth tests. This is a weak identity and does not assert that $c$ belongs to ordinary $L^2$.
\end{proposition}
\begin{proof}
We first take $h$ real. Let $\chi_R(x)=\chi(x/R)$, $R\ge1$, where $\chi$ is even and smooth, equals one on $[-1,1]$, has support in $[-2,2]$, and satisfies $0\le\chi\le1$. Put $k_R=\kappa\chi_R$. Differentiating the theta series on each half-line shows that $\kappa$ and its derivatives decrease faster than every exponential. Consequently $k_R\to\kappa$ in $H^1(\R)$ and in $L^1(e^{|x|/2}dx)$. For $v\in H^1(\R)$,
\[
 \tg[v]\le\norm{v'}_2^2\int_0^1\rho(s)s^2\,ds
              +4\norm{v}_2^2\int_1^\infty\rho(s)\,ds.
\]
It follows that the Gamma pairing, the ordinary $L^2$ pairing, and the moment in~\eqref{a:eq:bilinearzeros} converge as $R\to\infty$.

The complete prime pairing also converges without a spatially uniform incoming-rate assumption. If $\supp h\subset[-K,K]$, changing variables in one of its two terms gives
\[
 p(k_R,h)=\sum_{n\ge2}a_n\int_{-K}^K h(x)
           \bigl(k_R(x+s_n)+k_R(x-s_n)\bigr)\,dx.
\]
For any fixed $A_0>2$,
\begin{equation}\label{a:eq:bridgeprimerate}
 \sup_{|x|\le K}\sum_{n\ge2}a_n
       \bigl(\kappa(x+s_n)+\kappa(x-s_n)\bigr)
 \le2C_{A_0}e^{A_0K}\sum_{n\ge2}a_nn^{-A_0}<\infty.
\end{equation}
Here $a_n\le n$ suffices. Since $|k_R|\le\kappa$, dominated convergence proves the assertion and absolute convergence of the limiting pairing.

For the zero side, uniformly in $R$ and in $|\operatorname{Im}z|\le1/2$,
\[
 |\mathcal Fk_R(z)|\le\int_\R\kappa(x)e^{|x|/2}\,dx<\infty.
\]
Repeated integration by parts in the transform of the fixed compact test gives
$|\mathcal Fh(z)|\le C_{h,N}(1+|\operatorname{Re}z|)^{-N}$ throughout this strip. Equation~\eqref{a:eq:zerocount} therefore supplies a summable majorant for the products in~\eqref{a:eq:bilinearzeros}, for example with $N=3$. Pointwise transform convergence and Lemma~\ref{a:lem:thetaxi} make their limiting sum zero. Thus
\[
 q_0(\kappa,h)+2m(\kappa)m(h)=0.
\]
Lemma~\ref{a:lem:theta} gives $m(\kappa)=1/2$, proving~\eqref{a:eq:thetanull} for real $h$. Splitting an arbitrary $h$ into real and imaginary parts and using conjugate linearity in the second argument proves the stated version.
\end{proof}

\begin{theorem}[Compact physical bridge]\label{a:thm:compactbridge}
For every even $f\in C_c^\infty(\R)$,
\begin{equation}\label{a:eq:provedbridge}
 q_0[\kappa f]=\E[f]-\frac12\norm{f}_\mu^2.
\end{equation}
\end{theorem}
\begin{proof}
Set $h=\kappa|f|^2$, which is a real even compact smooth test. For $a=\kappa(x)$, $b=\kappa(x+s)$, $F=f(x)$, and $G=f(x+s)$, the pointwise identity
\[
 |bG-aF|^2-(b-a)(b|G|^2-a|F|^2)=ab|G-F|^2
\]
gives, after integration,
\begin{equation}\label{a:eq:gammaground}
 \tg[\kappa f]-\tg(\kappa,h)=\E_\Gamma[f].
\end{equation}
This is an equality between finite quantities: $\kappa$, $h$, and $\kappa f$ are in $H^1$, their Gamma pairings are finite by the estimate in the preceding proof, and $\E_\Gamma[f]$ is finite by Proposition~\ref{a:prop:core}.

For the primes the two endpoint terms are both essential. Their full expansion gives
\begin{align}\label{a:eq:primeground}
 p[\kappa f]-p(\kappa,h)
 &=\sum_{n\ge2}a_n\int_\R\kappa(x)\kappa(x+s_n)
       \bigl(2\operatorname{Re}(f(x+s_n)\overline{f(x)})\notag\\
 &\hspace{44mm}-|f(x+s_n)|^2-|f(x)|^2\bigr)\,dx\notag\\
 &=-\E_P[f].
\end{align}
The compact-profile prime energy is finite, the compact correlation $p[\kappa f]$ has only finitely many nonzero shifts, and $p(\kappa,h)$ is absolutely convergent by~\eqref{a:eq:bridgeprimerate}. Thus no difference of divergent prime expressions is used. Finally,
$\norm{\kappa f}_2^2=\ip{\kappa}{h}_2$ cancels the constant term. Combining~\eqref{a:eq:gammaground}, \eqref{a:eq:primeground}, and~\eqref{a:eq:thetanull} yields
\[
 q_0[\kappa f]=\E[f]-m(h)
             =\E[f]-\int_\R c\kappa|f|^2\,dx,
\]
which is~\eqref{a:eq:provedbridge}.
\end{proof}

The theorem uses the classical explicit formula and the theta transform, not a zero-location hypothesis. It establishes an identity of forms, not a positivity statement for $q_0$ or for the centered energy shifted by $1/2$.

\subsection{The exact numerator domain}\label{a:sec:domain}

We first analyze $X$ directly from its two defining nonnegative terms. Because $W$ is positive and bounded,
\begin{equation}\label{a:eq:L2bound}
 \norm{u}_2^2\le\norm{W}_\infty\norm{u}_{\Hh}^2.
\end{equation}
The Gamma form is closed on ordinary $L^2$, for example by the same difference-and-Fatou argument used for $\E$.

\begin{lemma}\label{a:lem:Xcore}
The space $X$ of~\eqref{a:eq:X} is Hilbert, and even $C_c^\infty(\R)$ is dense in $X$.
\end{lemma}
\begin{proof}
A sequence Cauchy in $X$ has a weighted $L^2$ limit and an ordinary Gamma-form limit. Equation~\eqref{a:eq:L2bound} identifies their ordinary $L^2$ limits, proving completeness.

Use the cutoffs $\chi_R$ from Proposition~\ref{a:prop:core}. For $h_R=(1-\chi_R)u$ the product identity gives
\begin{align*}
 |h_R(x+s)-h_R(x)|^2
 &\le2(1-\chi_R(x+s))^2|u(x+s)-u(x)|^2\\
 &\quad+2|u(x)|^2|\chi_R(x+s)-\chi_R(x)|^2.
\end{align*}
The first term tends to zero in the Gamma integral by dominated convergence. The second is dominated by
$2|u(x)|^2\min\{1,L^2s^2\}$, which is integrable by~\eqref{a:eq:rhointegrable}. Weighted convergence is another application of dominated convergence. Thus compactly supported vectors are dense in $X$.

For a fixed compactly supported vector, even smooth mollification converges in ordinary $L^2$ and in $\tg$, by the translation argument in Proposition~\ref{a:prop:core}. The supports remain in one compact interval, where $1/W$ is bounded. Ordinary $L^2$ convergence consequently gives weighted convergence as well.
\end{proof}

For original profiles and physical numerators the coordinates are
\begin{equation}\label{a:eq:coordinates}
 u=\kappa f,\qquad v=2cf=\frac uW,\qquad
 F=\frac u{\sqrt W}=\sqrt{2c\kappa}\,f.
\end{equation}
Only $f\mapsto F$ is the displayed unitary map from $\He$ to ordinary $\Le$. The map $f\mapsto u$ is an isometry to the weighted space with norm $\norm{\cdot}_{\Hh}$, since
\begin{equation}\label{a:eq:weightedidentity}
 \norm{\kappa f}_{\Hh}=\norm{f}_\mu.
\end{equation}

\begin{proposition}[Extension of the compact bridge]\label{a:prop:domain}
Multiplication by $\kappa$ gives a bijection $\D\to X$, with equivalent graph norms, and
\begin{equation}\label{a:eq:extendedbridge}
 q_0(\kappa f,\kappa g)=\E(f,g)-\tfrac12\ip{f}{g}_\mu,
 \qquad f,g\in\D.
\end{equation}
\end{proposition}
\begin{proof}
Put $C_{\mathrm b}=|\cg|\norm{W}_\infty+\norm{K_P}$. On the compact smooth core, Theorem~\ref{a:thm:compactbridge}, Lemma~\ref{a:lem:prime}, and~\eqref{a:eq:weightedidentity} imply
\begin{align}
 \tg[\kappa f]&\le\E[f]+(C_{\mathrm b}+\tfrac12)\norm{f}_\mu^2,\label{a:eq:graph1}\\
 \E[f]&\le\tg[\kappa f]+(C_{\mathrm b}+\tfrac12)\norm{f}_\mu^2.\label{a:eq:graph2}
\end{align}
Taking compact-profile approximants from Proposition~\ref{a:prop:core} and using completeness of $X$ proves that $\kappa\D\subset X$.

Conversely, for $u\in X$ choose even $u_j\in C_c^\infty$ converging in $X$, by Lemma~\ref{a:lem:Xcore}. Since $\kappa$ is smooth and strictly positive, $f_j=u_j/\kappa$ are legitimate even compact smooth profiles. Equation~\eqref{a:eq:graph2} makes $(f_j)$ Cauchy in the original energy graph norm. Its weighted limit is $f=u/\kappa$, so closedness of $\E$ gives $f\in\D$. This proves surjectivity and graph-norm equivalence. The Gamma form and the prime form are continuous in $X$, so the quadratic identity passes to the limit; polarization gives~\eqref{a:eq:extendedbridge}.
\end{proof}

Recall $m(u)=\int_\R c(x)u(x)\,dx$. The integral exists absolutely for finite weighted norm, because
\begin{equation}\label{a:eq:meanbound}
 |m(u)|\le\left(\int_\R c^2W\,dx\right)^{1/2}\norm{u}_{\Hh}
       =\tfrac12\norm{u}_{\Hh}.
\end{equation}
Here $\int c^2W=\tfrac12\int c\kappa=1/4$. Thus $\mu(f)=2m(\kappa f)$, and the full pole form obeys
\begin{equation}\label{a:eq:polebridge}
 q[u]:=q_0[u]+2|m(u)|^2
       =\E[f]-\tfrac12\Var_\mu(f),\qquad u=\kappa f\in X.
\end{equation}
The coefficient $2$ in the pole is fixed by normalization. It is not an additional positive contribution available after centering.


\subsection{Complex theta continuation and bounded Fourier evaluations}\label{d:sec:theta}
The real-line normalization has already been proved. We record the additional complex estimates needed for spectral synthesis.


\begin{lemma}[Theta continuation and tails]\label{d:lem:theta-tails}
The function $\kappa$ has an even holomorphic continuation to $|\operatorname{Im}z|<\pi/4$.
There are constants $0<A_0\le A_1<\infty$ such that
\begin{equation}\label{d:eq:W-two-sided}
 A_0 e^{4x-\pi e^{2x}}\le W(x)\le A_1 e^{4x-\pi e^{2x}}
 \qquad(x\ge0).
\end{equation}
For each nonnegative integer $j$ and each $\beta_0<\pi/4$,
\begin{equation}\label{d:eq:strip-upper}
 |\kappa^{(j)}(x+i\beta)|\le C_{j,\beta_0}
 \exp\!\left((9/2+2j)|x|-\pi\cos(2\beta_0)e^{2|x|}\right)
 \quad(|\beta|\le\beta_0).
\end{equation}
For fixed $b$ in the strip, as $x\to+\infty$,
\begin{equation}\label{d:eq:complex-asymptotic}
 \kappa(x+b)=4\pi^2e^{9(x+b)/2}e^{-\pi e^{2x}e^{2b}}
 \bigl(1+O(e^{-2x})\bigr).
\end{equation}
The last error estimate is uniform when $b$ ranges over a compact subset of the strip.
\end{lemma}
\begin{proof}
Let
\[
 \Psi(t)=\sum_{n\ge1}e^{-\pi n^2t}\quad(\operatorname{Re}t>0),\qquad
 h(z)=e^{z/2}\Psi(e^{2z}).
\]
The Jacobi transformation~\cite[Eq.~20.7.32]{DLMF} gives, first on the real line,
$h(-x)=h(x)+\sinh(x/2)$. Both sides are holomorphic in the strip, so the identity holds there. Direct differentiation gives
$\kappa=h''-h/4$ on $x\ge0$. Since $(D^2-1/4)\sinh(x/2)=0$, this formula defines an even holomorphic extension and agrees with the stipulated reflection. All derivatives of the defining series converge locally uniformly in the strip.

For a lower bound in~\eqref{d:eq:W-two-sided}, keep the $n=1$ term and use
$4\tau^2-6\tau\ge(4-6/\pi)\tau^2$ for $\tau\ge\pi$, together with
$e^{x/2}\le2c(x)\le2e^{x/2}$. For the upper bound, discard the negative term and use
\[
 \sum_{n\ge1} n^4e^{-\pi n^2e^{2x}}
 \le e^{-\pi e^{2x}}\sum_{n\ge1}n^4e^{-\pi(n^2-1)}.
\]
The series on the right converges. Termwise differentiation introduces only fixed polynomial powers of $n^2e^{2z}$ and fixed exponential factors. The same estimate with $\pi\cos(2\beta_0)$ in the exponential proves~\eqref{d:eq:strip-upper}; evenness handles negative $x$.

For~\eqref{d:eq:complex-asymptotic}, the lower-degree part of the $n=1$ polynomial has relative size $O(e^{-2x})$. The terms $n\ge2$ have, on a compact set of $b$, an additional factor bounded by $C\exp(-c e^{2x})$ with $c>0$. This proves the stated relative estimate. In particular,
\begin{equation}\label{d:eq:W-asymptotic}
 W(x)=4\pi^2e^{4x-\pi e^{2x}}\bigl(1+O(e^{-x})\bigr)
 \quad(x\to+\infty).
\end{equation}
\end{proof}


\begin{lemma}[Entire transform and bounded evaluations]\label{d:lem:Fourier}
For $u\in\Hh$, the function $F_u$ is even entire, the map $u\mapsto F_u$ is injective, and
\begin{equation}\label{d:eq:eval-bound}
 |F_u^{(k)}(z)|\le\norm u_{\Hh}
 \left(\int_\R W(x)|x|^{2k}e^{2B|x|}\,dx\right)^{1/2}
 \quad (|\operatorname{Im}z|\le B).
\end{equation}
Moreover
\begin{equation}\label{d:eq:Xi-xi}
 \Xi(z):=F_\kappa(z)=\xi(1/2+iz),\qquad
 \Xi(0)>0,\quad \Xi(i/2)=\Xi(-i/2)=\tfrac12.
\end{equation}
Its zero set is discrete and invariant under $w\mapsto-w$ and $w\mapsto\overline w$, preserving multiplicity.
\end{lemma}
\begin{proof}
Cauchy--Schwarz gives~\eqref{d:eq:eval-bound}; the two-sided theta estimate makes every displayed integral finite and justifies locally uniform differentiation. It also puts $u$ in $L^1(e^{B|x|}dx)$ for every $B$. Ordinary Fourier uniqueness proves injectivity. Lemma~\ref{a:lem:thetaxi} already proves the exact completed transform, including its Mellin calculation. Positivity, evenness, real coefficients and normalization give the remaining assertions.
\end{proof}


\subsection{The exact Gamma form and its weighted representative}\label{d:sec:Gamma-form}

We now construct the weighted Gamma representative and its ambient graph core. These facts will serve both the physical realization and the later stronger-topology synthesis theorem.

Retain the same theta weight $W$ and put
\begin{align}
 \rho(s)&=\frac{e^{-s/2}}{1-e^{-2s}},\qquad s>0,\label{d:eq:Gamma-rho}\\
 \tg[u]&=\int_0^\infty\rho(s)
       \norm{u(\cdot+s)-u}_{L^2(\R)}^2\,ds,\label{d:eq:Gamma-form}\\
 {X}&=\{u\in\Hh:\tg[u]<\infty\},\qquad
 \norm{u}_{{X}}^2=\norm{u}_{\Hh}^2+\tg[u].
 \label{d:eq:Gamma-X}
\end{align}
The form is polarized with the linear-first convention. Although its vectors are even, its translation differences are taken on the whole real line. All positive shifts in~\eqref{d:eq:Gamma-form} are retained.

Write $M=\norm{W}_\infty<\infty$. The inclusion
$\norm{u}_2\le\sqrt M\norm{u}_{\Hh}$ is continuous. The kernel satisfies $\rho(s)=O(s^{-1})$ at zero and $O(e^{-s/2})$ at infinity, so
\begin{equation}\label{d:eq:Gamma-J}
 J_j=\int_0^\infty s^j\rho(s)\,ds<\infty\qquad(j=1,2).
\end{equation}
Use the unitary ordinary Fourier transform in the next two sections. Plancherel and Tonelli give, including infinite values,
\begin{equation}\label{d:eq:Gamma-multiplier}
 \tg[u]=\int_\R d_\Gamma(\xi)|\widehat u(\xi)|^2\,d\xi,
 \qquad d_\Gamma(\xi)=2\int_0^\infty\rho(s)(1-\cos(\xi s))\,ds.
\end{equation}
The inequalities $1-\cos y\le |y|$ and
$\norm{u(\cdot+s)-u}_2\le s\norm{u'}_2$ imply
\begin{equation}\label{d:eq:Gamma-first-order}
 0\le d_\Gamma(\xi)\le2J_1|\xi|,\qquad
 \tg[u]\le J_2\norm{u'}_2^2
 \quad(u\in H^1(\R)).
\end{equation}
Here $H^1(\R)$ always means the ordinary unweighted Sobolev space. These are upper estimates, not replacements of the exact Gamma form by a first-order form.

\begin{proposition}[Closed form and compact smooth form core]\label{d:prop:Gamma-form-core}
The form~\eqref{d:eq:Gamma-form} is densely defined, nonnegative, and closed in $\Hh$. The even compact smooth functions are dense in ${X}$. It has a nonnegative self-adjoint representative $T_\Gamma$ in $\Hh$.
\end{proposition}
\begin{proof}
The domain and compact smooth form-core assertion are Lemma~\ref{a:lem:Xcore}. The multiplier identity above proves nonnegativity and closedness equivalently. The representation theorem gives $T_\Gamma$; no graph-core conclusion is inferred from form density.
\end{proof}

Let $D_\Gamma$ be the ordinary nonnegative Fourier multiplier with symbol $d_\Gamma$ and its maximal ordinary $L^2$ multiplier domain. Fix any bounded self-adjoint operator $\mathsf D$ on $\Hh$, and define
\begin{equation}\label{d:eq:Gamma-graph}
 B_{\mathsf D}=T_\Gamma-\mathsf D,\qquad
 \mathcal G_{\mathsf D}=\operatorname{Dom}T_\Gamma,
 \qquad
 \norm{u}_{\mathcal G_{\mathsf D}}^2
 =\norm{u}_{\Hh}^2+\norm{B_{\mathsf D}u}_{\Hh}^2.
\end{equation}
The bounded self-adjoint perturbation is self-adjoint on $\operatorname{Dom}T_\Gamma$, and the two graph norms are equivalent, since
\[
 \norm{B_{\mathsf D}u}_{\Hh}\le\norm{T_\Gamma u}_{\Hh}
       +\norm{\mathsf D}\norm{u}_{\Hh},\qquad
 \norm{T_\Gamma u}_{\Hh}\le\norm{B_{\mathsf D}u}_{\Hh}
       +\norm{\mathsf D}\norm{u}_{\Hh}.
\]
No relation between $\mathsf D$ and $\Kk$ is assumed.

\begin{lemma}[Actual graph admission from one ordinary derivative]\label{d:lem:Gamma-H1}
If $u\in\Hh\cap H^1(\R)$ is even, then $u\in\operatorname{Dom}T_\Gamma$ and
\begin{equation}\label{d:eq:Gamma-H1}
 T_\Gamma u=W D_\Gamma u,\qquad
 \norm{B_{\mathsf D}u}_{\Hh}
 \le2\sqrt M J_1\norm{u'}_2+\norm{\mathsf D}\norm{u}_{\Hh}.
\end{equation}
Consequently, for either $E={X}$ or $E=\mathcal G_{\mathsf D}$,
\begin{equation}\label{d:eq:Gamma-E-upper}
 \norm{u}_{\Hh}\le\norm{u}_E
 \le C_E\bigl(\norm{u}_{\Hh}+\norm{u}_{H^1(\R)}\bigr)
 \qquad(u\in\Hh\cap H^1(\R)).
\end{equation}
\end{lemma}
\begin{proof}
Equation~\eqref{d:eq:Gamma-first-order} first gives $u\in{X}$ and $D_\Gamma u\in L^2$. The vector $g=WD_\Gamma u$ belongs to $\Hh$ because
\[
 \norm{g}_{\Hh}^2=\int_\R W|D_\Gamma u|^2
 \le M\norm{D_\Gamma u}_2^2.
\]
For every $z\in{X}$, the ordinary multiplier form identity yields
$\tg(u,z)=\ip{D_\Gamma u}{z}_2=\ip{g}{z}_{\Hh}$.
The representation criterion therefore admits $u$ to the actual weighted operator domain and identifies $T_\Gamma u=g$. Now~\eqref{d:eq:Gamma-first-order} and boundedness of $\mathsf D$ prove~\eqref{d:eq:Gamma-H1}. The form and graph estimates imply~\eqref{d:eq:Gamma-E-upper}. This is an admission argument, not an assumed maximal-domain identity for a formal product.
\end{proof}

\subsection{A direct compact smooth graph core}\label{d:sec:Gamma-graph-core}

We need compact smooth density in the graph norm itself. The form core from Proposition~\ref{d:prop:Gamma-form-core} alone is insufficient.

\begin{lemma}[Weighted cutoff commutator]\label{d:lem:Gamma-commutator}
If $\chi$ is real, bounded, even, and Lipschitz, multiplication by $\chi$ preserves ${X}$ and $\operatorname{Dom}T_\Gamma$. There is a bounded operator $C_\chi$ on $\Hh$ such that
\begin{align}
 \tg(\chi u,z)-\tg(u,\chi z)
 &=\ip{C_\chi u}{z}_{\Hh},\label{d:eq:Gamma-comm-form}\\
 T_\Gamma(\chi u)&=\chi T_\Gamma u+C_\chi u,
 \qquad
 \norm{C_\chi}\le2M\operatorname{Lip}(\chi)J_1.
 \label{d:eq:Gamma-comm-row}
\end{align}
The first identity holds for $u,z\in{X}$, and the second for $u\in\operatorname{Dom}T_\Gamma$.
\end{lemma}
\begin{proof}
The product rule gives
\[
 \tg[\chi u]
 \le2\norm{\chi}_\infty^2\tg[u]
       +2\operatorname{Lip}(\chi)^2J_2\norm{u}_2^2,
\]
proving form-domain preservation. Put $\beta=\sqrt W$ and use the isometry $U_Wu=u/\beta$ to ordinary even $L^2$. Expansion of the combined form difference in~\eqref{d:eq:Gamma-comm-form} gives the kernel
\begin{equation}\label{d:eq:Gamma-comm-kernel}
 (U_W C_\chi U_W^{-1}f)(x)
 =\int_\R\beta(x)\beta(y)\rho(|x-y|)
                 (\chi(x)-\chi(y))f(y)\,dy.
\end{equation}
Its absolute row and column integrals are at most
$2M\operatorname{Lip}(\chi)\int_0^\infty s\rho(s)\,ds$.
The elementary Schur estimate gives the asserted norm bound. For precision, first delete $|x-y|<\epsilon$; the omitted operator norm is at most
$2M\operatorname{Lip}(\chi)\int_0^\epsilon s\rho(s)\,ds\to0$.
The form difference converges to~\eqref{d:eq:Gamma-comm-form} by the difference-product identity and form Cauchy--Schwarz. The separately divergent diagonal rate is never split off. The sign in~\eqref{d:eq:Gamma-comm-kernel} is the sign of $[T_\Gamma,\chi]=T_\Gamma\chi-\chi T_\Gamma$.

For $u\in\operatorname{Dom}T_\Gamma$ and all $z\in{X}$,
\[
 \tg(\chi u,z)
 =\ip{T_\Gamma u}{\chi z}_{\Hh}+\ip{C_\chi u}{z}_{\Hh}
 =\ip{\chi T_\Gamma u+C_\chi u}{z}_{\Hh}.
\]
The representation criterion proves operator-domain preservation and~\eqref{d:eq:Gamma-comm-row}, rather than presupposing them.
\end{proof}

For the scaled cutoffs from the preceding section, every $u\in\operatorname{Dom}T_\Gamma$ therefore satisfies
\begin{align}
 \norm{u-\chi_Ru}_{\Hh}&\le\norm{\boldsymbol1_{|x|>R}u}_{\Hh},\notag\\
 \norm{T_\Gamma u-T_\Gamma(\chi_Ru)}_{\Hh}
 &\le\norm{\boldsymbol1_{|x|>R}T_\Gamma u}_{\Hh}
       +\frac{2M\norm{\chi'}_\infty J_1}{R}\norm{u}_{\Hh}.
 \label{d:eq:Gamma-cutoff-graph}
\end{align}
Both right sides tend to zero for the fixed graph vector. No uniform tail rate on a graph unit ball is claimed.

\begin{lemma}[Compact graph vectors and their noncompact rows]\label{d:lem:Gamma-compact-row}
If $u\in\operatorname{Dom}T_\Gamma$ is supported in $[-A,A]$, then $u\in\operatorname{Dom}D_\Gamma$ in ordinary $L^2$. For $h>0$ and $|x|>A+h$,
\begin{equation}\label{d:eq:Gamma-exterior}
 D_\Gamma u(x)=-\int_{-A}^A\rho(|x-y|)u(y)\,dy,
 \qquad
 |D_\Gamma u(x)|\le
 \frac{e^{A/2}\norm{u}_1e^{-|x|/2}}{1-e^{-2h}}.
\end{equation}
\end{lemma}
\begin{proof}
For $u\in{X}$, define $D_\Gamma u$ as a tempered distribution through the multiplier; the at-most-linear bound~\eqref{d:eq:Gamma-first-order} suffices. For $u\in\operatorname{Dom}T_\Gamma$, compact testing gives
\begin{equation}\label{d:eq:Gamma-local-row}
 D_\Gamma u=(T_\Gamma u)/W\quad\hbox{in }L^2_{\rm loc}.
\end{equation}
This follows by comparing the distributional form pairing with the representative pairing. Division by $W$ is legitimate locally only.

On a test support at positive distance from $[-A,A]$, the diagonal term vanishes. Expanding the form there and applying Fubini gives the exterior convolution in~\eqref{d:eq:Gamma-exterior}. The distance $h$ controls the denominator of $\rho$, while $|x-y|\ge|x|-A$ controls its numerator. If
$w_{A+h}=\min_{|x|\le A+h}W(x)>0$, the local row and exterior bound imply
\[
 \norm{D_\Gamma u}_2^2
 \le w_{A+h}^{-1}\norm{T_\Gamma u}_{\Hh}^2
       +\frac{2e^{-h}\norm{u}_1^2}{(1-e^{-2h})^2}<\infty.
\]
Thus the global distributional row is ordinary $L^2$, which is exactly the ordinary multiplier-domain criterion. Compact support of the input was essential; the output row is not asserted to be compactly supported.
\end{proof}

\begin{theorem}[Compact smooth graph core]\label{d:thm:Gamma-graph-core}
The even compact smooth functions form a graph core for $T_\Gamma$ and for every $B_{\mathsf D}$ in~\eqref{d:eq:Gamma-graph}.
\end{theorem}
\begin{proof}
First use~\eqref{d:eq:Gamma-cutoff-graph}. For a fixed compact graph vector $u$ supported in $[-A,A]$, let $u_\epsilon=\eta_\epsilon*u$, where $\eta$ is a nonnegative even compact smooth mollifier of mass one supported in $[-1,1]$. Lemma~\ref{d:lem:Gamma-compact-row} allows ordinary multiplier commutation:
\[
 D_\Gamma u_\epsilon=\eta_\epsilon*(D_\Gamma u).
\]
Both input and row converge in ordinary $L^2$. For $\epsilon\le1$ the input difference has support in $[-A-1,A+1]$, so
\[
 \norm{u_\epsilon-u}_{\Hh}
 \le w_{A+1}^{-1/2}\norm{u_\epsilon-u}_2,\qquad
 \norm{T_\Gamma(u_\epsilon-u)}_{\Hh}
 \le\sqrt M\norm{D_\Gamma(u_\epsilon-u)}_2.
\]
Each $u_\epsilon$ is already admitted by Lemma~\ref{d:lem:Gamma-H1}; the second estimate compares actual rows. A diagonal cutoff--mollifier sequence proves the graph core. Equivalent graph norms under bounded perturbation prove the assertion for $B_{\mathsf D}$. This theorem concerns the full ambient space $\Hh$, not a restriction to $\Kk$.
\end{proof}


\section{Actual realization and the spectral threshold}\label{sec:spectral}
The weighted realization and the ordinary-$L^2$ pencil are distinct operators with a common form numerator. The first is unitarily equivalent to the original energy operator; the second is a confining parameter family used to count signs.


\subsection{The weighted physical representative}\label{sec:weighted-physical}
Let $D_b$ be the bounded self-adjoint correction defined by
\begin{equation}\label{eq:Db-physical}
 \ip{D_bu}{v}_{\Hh}=\cg\ip uv_2+p(u,v),\qquad
 B=T_\Gamma-D_b,\quad \dom B=\dom T_\Gamma.
\end{equation}
Lemma~\ref{a:lem:prime} proves boundedness, including every prime power; in its notation $D_b=D$ and $\norm{D_b}<10$. The closed-form representation theorem and~\eqref{eq:intro-bridge} imply
\begin{equation}\label{eq:actual-unitary}
 UAU^{-1}=B+\tfrac12I,\qquad Uf=\kappa f.
\end{equation}
Indeed both sides represent the same closed form on $X$, so uniqueness of the self-adjoint representative proves the operator and domain identity. This does not identify the ordinary-$L^2$ representative of the same algebraic form.


\subsection{Derivative radicals and finite interpolation}\label{a:sec:radical}

The scaled response uses actual radicals of $q_0$, meaning vectors that annihilate the entire form domain under polarization. A vector with only $q_0[u]=0$ would not suffice. We also need ordinary-$L^2$ density, but not weighted density; the two statements differ for this kernel.

\begin{lemma}[Full-domain derivative radicals]\label{a:lem:radical}
For $j\ge1$, put
\begin{equation}\label{a:eq:radical}
 r_j=\kappa^{(2j)}-4^{-j}\kappa,\qquad
 \Rcal=\operatorname{span}\{r_j:j\ge1\}.
\end{equation}
Then $r_j\in X_c$ and
\begin{equation}\label{a:eq:radicalidentity}
 q_0(r_j,z)=0\qquad(z\in X).
\end{equation}
The family $(r_j)_{j\ge1}$ is linearly independent.
\end{lemma}
\begin{proof}
For each integer $m\ge0$, differentiating the theta series on the positive half-line gives
\[
 \kappa^{(m)}(x)=e^{x/2}\sum_{n\ge1}P_m(\pi n^2e^{2x})e^{-\pi n^2e^{2x}},
\]
where $P_0(t)=4t^2-6t$ and
$P_{m+1}(t)=(1/2-2t)P_m(t)+2tP_m'(t)$.
Thus $\deg P_m\le m+2$. Positivity and the first theta term also give
\begin{align}
 |\kappa^{(m)}(x)|&\le C_m
        e^{(9/2+2m)|x|-\pi e^{2|x|}},\label{a:eq:derivativeupper}\\
 \kappa(x)&\ge(4\pi^2-6\pi)
        e^{(9/2)|x|-\pi e^{2|x|}}.\label{a:eq:thetalower}
\end{align}
The upper bound follows by summing
$n^{2m+4}e^{-\pi(n^2-1)}$; the lower bound uses
$4\pi^2e^{4x}-6\pi e^{2x}\ge(4\pi^2-6\pi)e^{4x}$ for $x\ge0$.
Evenness extends both estimates to the full line. Therefore
\begin{equation}\label{a:eq:derivativeweighted}
 \frac{|\kappa^{(m)}(x)|^2}{W(x)}
 \le C_m'e^{(5+4m)|x|-\pi e^{2|x|}}.
\end{equation}
Every derivative has finite weighted norm and belongs to $H^1(dx)$, hence has finite Gamma form. In particular $d_j=\kappa^{(2j)}\in X$.

Integration by parts gives
\[
 \mathcal Fd_j(z)=(-z^2)^j\xi(1/2+iz),\qquad
 m(d_j)=4^{-j}m(\kappa)=4^{-j}/2.
\]
Use $d_j\chi_R$ in the bilinear explicit formula exactly as in the proof of Proposition~\ref{a:prop:thetanull}. For clarity, all three cutoff controls persist:~\eqref{a:eq:derivativeupper} gives convergence in $H^1$ and in $L^1(e^{|x|/2}dx)$; on each compact test support the complete shifted prime row is dominated by a summable $n^{-A_0}$ majorant with $A_0>2$; and the zero side is dominated by
\[
 \left(\int_\R|d_j(x)|e^{|x|/2}\,dx\right)
 C_{h,N}(1+|\operatorname{Re}w|)^{-N}.
\]
The zero-counting bound makes this summable. Every limiting product vanishes because $\mathcal Fd_j(w_{\rho_\zeta})=0$. Hence, on compact even tests,
\[
 q_0(d_j,h)=-4^{-j}\overline{m(h)}.
\]
Continuity of $q_0$ on $X\times X$ and the compact smooth form core extend the identity to all $h\in X$. Subtracting $4^{-j}$ times the source identity for $\kappa$ proves~\eqref{a:eq:radicalidentity}; the moment formula proves $r_j\in X_c$.

If $\sum_{j=1}^Jb_jr_j=0$, Fourier transformation gives
\[
 \left(\sum_{j=1}^Jb_j\bigl((-z^2)^j-4^{-j}\bigr)\right)
                 \xi(1/2+iz)=0\qquad(z\in\R).
\]
The second factor is nonzero on a neighborhood of zero, since
$\mathcal F\kappa(0)=\int\kappa>0$. The polynomial must vanish identically. Comparing its highest degrees successively gives all $b_j=0$.
\end{proof}

\begin{lemma}[Holomorphic translations]\label{a:lem:strip}
The real kernel has an even holomorphic continuation to
$\mathcal S=\{z:|\operatorname{Im}z|<\pi/4\}$. For every $m\ge0$ and every $0\le\vartheta<\pi/4$,
\begin{equation}\label{a:eq:stripbound}
 |\kappa^{(m)}(r+iy)|\le C_{m,\vartheta}
 e^{(9/2+2m)|r|}
 e^{-\pi\cos(2\vartheta)e^{2|r|}}
 \quad(|y|\le\vartheta).
\end{equation}
The map $b\mapsto\kappa(\cdot-b)$ is holomorphic from $\mathcal S$ into ordinary $L^2(dx)$.
\end{lemma}
\begin{proof}
The continuation and derivative estimate are Lemma~\ref{d:lem:theta-tails}. It remains to prove ordinary-$L^2$-valued holomorphy. Put $d_\vartheta=\pi\cos(2\vartheta)>0$.

On a compact set of translation parameters choose $B<\infty$ and $\vartheta<\pi/4$ with $|\operatorname{Re}b|\le B$ and $|\operatorname{Im}b|\le\vartheta$. Formula~\eqref{a:eq:stripbound} supplies the common square-integrable envelope
\[
 C_{m,\vartheta}e^{(9/2+2m)(|x|+B)}
       e^{-d_\vartheta e^{2\max\{|x|-B,0\}}}.
\]
Taylor's integral remainder, with the bound for the second derivative, proves complex differentiability of $b\mapsto\kappa(\cdot-b)$ in $L^2$. Repetition gives derivative $(-1)^m\kappa^{(m)}(\cdot-b)$ of every order.
\end{proof}

\begin{proposition}[Ordinary density and weighted incompleteness]\label{a:prop:density}
The closure of $\Rcal$ in ordinary $\Le$ is all of $\Le$. Its closure in $\Hh_c$ with the weighted norm is a proper subspace.
\end{proposition}
\begin{proof}
Suppose $h\in\Le$ is orthogonal to every $r_j$, and define the bilinear convolution pairing
\[
 F(b)=\int_\R h(x)\kappa(x-b)\,dx,\qquad b\in\mathcal S.
\]
Lemma~\ref{a:lem:strip} makes $F$ holomorphic. Since $h$ and $\kappa$ are even, $F$ is even. Orthogonality to the real functions $r_j$ gives
$F^{(2j)}(0)=4^{-j}F(0)$ for all $j\ge1$.
All Taylor coefficients therefore agree with those of $F(0)\cosh(b/2)$, and the identity theorem yields
\[
 F(b)=F(0)\cosh(b/2)\qquad(b\in\mathcal S).
\]
For real $b$, Cauchy--Schwarz bounds $|F(b)|$ by
$\norm h_2\norm\kappa_2$, independently of $b$. Thus $F(0)=0$, so $F=0$. Since $\kappa\in L^1$, the convolution belongs to $L^2$ and Plancherel gives
$\mathcal F\kappa(t)\widehat h(t)=0$ almost everywhere. The first factor is a nonzero entire function, whose real zeros are discrete and hence have measure zero. It follows that $h=0$, proving ordinary density.

For the different weighted assertion, choose any nontrivial zero of $\zeta$, and write it as $1/2+iw$. Such zeros exist by the classical counting theorem \cite{vonMangoldt}; their position on or off the critical line is irrelevant. Put
\[
 u_w(x)=W(x)\cos(wx).
\]
The theta envelope gives $u_w\in\Hw$. The entire transform identity and integration by parts imply
\[
 m(u_w)=\frac12\mathcal F\kappa(w)=0,\qquad
 \ip{u_w}{r_j}_{\Hh}
 =\int_\R\cos(wx)r_j(x)\,dx
 =\bigl((-w^2)^j-4^{-j}\bigr)\mathcal F\kappa(w)=0.
\]
Thus $u_w\in \Hh_c$ is orthogonal to $\Rcal$, and it is nonzero because $u_w(0)=W(0)>0$. This proves weighted incompleteness. It does not identify $u_w$ as an eigenstate or as a counterexample to any variance inequality.
\end{proof}

The ordinary density in Proposition~\ref{a:prop:density} will be applied only to finitely many ordinary-$L^2$ continuous moment functionals. It gives no bound on the weighted norm of an interpolant. In particular it cannot be substituted for density in the original centered profile space under $u=\kappa f$.


\subsection{A source map and its actual invariant domain}\label{b:sec:source}
In this subsection let $a(x)=(2c(x))^{-1}$, $L=-A$, and
\begin{equation}\label{b:eq:nu}
 d\nu(s)=\rho(s)\,ds+\sum_{n\ge2}a_n\delta_{s_n}(ds),\qquad s>0.
\end{equation}
The function $a(x)$ is distinct from the coefficients $a_n$ and from the complex scalar $a=i/2$ used later only in the synthesis sections.


\begin{equation}\label{b:eq:source-main}
 Q\kappa=-c,\qquad Q=D_\Gamma-\cg-P,
\end{equation}
where, on smooth rapid inputs,
\begin{align}
 D_\Gamma h(x)&=\int_0^\infty\rho(s)
       [2h(x)-h(x+s)-h(x-s)]\,ds,\label{b:eq:Drow}\\
 Ph(x)&=\sum_{n\ge2}a_n[h(x+s_n)+h(x-s_n)].\label{b:eq:Prow}
\end{align}
The compensated difference in~\eqref{b:eq:Drow} is never split into divergent diagonal rates.

Let $\RR$ be the space of smooth rapid-exponential functions with seminorms
\[
 p_{B,m}(h)=\sup_x e^{B|x|}|h^{(m)}(x)|<\infty
 \quad(B>0,\ m\ge0).
\]
Let $\SSS$ be the smooth small-exponential class with seminorms
\[
 q_{\epsilon,m}(v)=\sup_xe^{-\epsilon|x|}|v^{(m)}(x)|<\infty
 \quad(\epsilon>0,\ m\ge0).
\]
Finally, $\CE$ consists of all smooth functions satisfying, for each fixed $m$, a bound $|f^{(m)}(x)|\le C_m e^{B_m|x|}$ for some finite $C_m,B_m$. The exponent may depend on $m$. This is an algebra, $\SSS\subset\CE$, and the theta derivative estimates~\eqref{a:eq:derivativeupper} imply $\kappa\CE\subset\RR$.

\begin{lemma}[Continuous full source map]\label{b:lem:T}
The map $T=aQ$ sends $\RR$ continuously into $\SSS$. The maps $h\mapsto WQh$, $h\mapsto W(D_\Gamma-\cg)h$, and $h\mapsto WPh$ send $\RR$ continuously into $\RR$.
\end{lemma}
\begin{proof}
For $0<s\le1$ the $m$th differentiated Gamma difference is bounded by $s^2\norm{h^{(m+2)}}_\infty$, and for $s>1$ by $4\norm{h^{(m)}}_\infty$. These bounds are integrable against $\rho(s)\,ds$ and justify every differentiation. Thus every derivative of $D_\Gamma h$ is bounded.

For $B>1/2$, the elementary inequalities $\Lambda(n)\le\log n$ and $|x\pm\log n|\ge\log n-|x|$ give
\begin{equation}\label{b:eq:primepoint}
 |(Ph)^{(m)}(x)|\le
  2p_{B,m}(h)e^{B|x|}\sum_{n\ge2}(\log n)n^{-B-1/2}.
\end{equation}
The scalar sum is finite by the integral test, and includes every prime power. It proves locally uniform absolute convergence of the series and every differentiated series. Since $|a^{(j)}(x)|\le C_j e^{-|x|/2}$, choosing $B=1/2+\epsilon$ and applying Leibniz's rule proves the $\SSS$ bounds. The same argument gives the full tail estimate
\begin{equation}\label{b:eq:primetailpoint}
 q_{\epsilon,m}\bigl(a(P-P_N)h\bigr)
 \le C_{m,\epsilon}\max_{0\le j\le m}p_{1/2+\epsilon,j}(h)
           \sum_{n>N}(\log n)n^{-1-\epsilon}\longrightarrow0.
\end{equation}
Here $P_N$ retains the terms with $n\le N$ in~\eqref{b:eq:Prow}. This is a weighted-growth seminorm estimate, not a uniform unweighted incoming-rate estimate. Every derivative of $Qh$ grows at most as $e^{(1/2+\epsilon)|x|}$. Multiplication by the rapid theta factor $W$ gives all target $\RR$ seminorms. The same proof applies separately to its Gamma and prime parts.
\end{proof}

\begin{proposition}[Actual generator and invariant algebra]\label{b:prop:generator}
Every even $f\in\CE$ belongs to $\dom(A)$, and
\begin{equation}\label{b:eq:generator-source}
 Af=\tfrac12f+T(\kappa f),\qquad
 Lf=-\tfrac12f-T(\kappa f).
\end{equation}
The right sides lie in $\CE$, so every such profile belongs to $\dom(A^m)$ for each finite $m$. Pointwise,
\begin{equation}\label{b:eq:generator-jump}
 Lf(x)=a(x)\int_{s>0}\!
 \left\{\kappa(x+s)[f(x+s)-f(x)]
       +\kappa(x-s)[f(x-s)-f(x)]\right\}\,d\nu(s).
\end{equation}
\end{proposition}
\begin{proof}
Put $u=\kappa f\in\RR$. The theta bounds imply $u\in H^1(\R)$ and
$\int |u|^2/W=\int |f|^2\,d\mu<\infty$. Thus $u\in X$, and the full bridge admits $f$ to $\D$. Against an even compact smooth $h$, the compensated Gamma row can be paired with $\kappa h$ by the Taylor estimates above; the prime row is locally absolutely convergent by~\eqref{b:eq:primepoint}. Hence
\[
 q_0(u,\kappa h)=\int Q(u)\kappa\overline h\,dx
                =\ip{T(u)}h_\mu.
\]
Lemma~\ref{b:lem:T} gives $T(u)\in\SSS\subset\He$, because $\mu$ has all exponential moments. The bridge yields
$\E(f,h)=\ip{f/2+T(u)}h_\mu$ on the compact smooth form core. Both sides are continuous in the form norm of $h$, so the equality extends to all $h\in\D$. The representation criterion~\eqref{a:eq:Adef} proves~\eqref{b:eq:generator-source}. Its right side is in $\CE$, proving invariance and iteration.

The weak equation~\eqref{b:eq:source-main} upgrades to a pointwise equality because Lemma~\ref{b:lem:T} makes $Q\kappa$ smooth and the residual is even. In the jump row, near zero the profile difference is $O(s)$ and $\rho(s)=O(1/s)$; its Gamma tail is integrable. The prime row converges locally absolutely by~\eqref{b:eq:primepoint}, applied to $\kappa f$ and $\kappa$. Subtracting the two compensated source rows gives
\[
 Q(\kappa f)-fQ\kappa
  =-\int_{s>0}\!\left\{\kappa(x+s)[f(x+s)-f(x)]
             +\kappa(x-s)[f(x-s)-f(x)]\right\}\,d\nu(s).
\]
Now use $Q\kappa=-c$ and $ac=1/2$. This proves~\eqref{b:eq:generator-jump} with the coefficient $1/(2c)$.
\end{proof}


For every $\delta>0$, set
\begin{equation}\label{eq:ell-unified}
 \ell_\delta[u]=q_0[u]+\delta\norm u_{\Hh}^2,\qquad \dom\ell_\delta=X.
\end{equation}


\subsection{The closed pencil and compact resolvent}\label{a:sec:pencil}

\begin{proposition}[Closed ordinary-$L^2$ realization]\label{a:prop:closedpencil}
For every $\delta>0$, the form $\ell_\delta$ of~\eqref{eq:ell-unified} is closed and lower bounded on $\Le$. Its self-adjoint representative is characterized by
\begin{equation}\label{a:eq:Ldef}
 \begin{split}
 u\in\dom(L_\delta),\ L_\delta u=h\quad\Longleftrightarrow\quad
 u\in X,\ h\in\Le,\quad
 \ell_\delta(u,z)=\ip{h}{z}_2\ \text{for all }z\in X.
 \end{split}
\end{equation}
The construction here uses the direct numerator estimates.
\end{proposition}
\begin{proof}
Choose an integer $N$ for which $r_N\le\delta/2$ and set
\[
 C_N=|\cg|+2\sum_{2\le n\le N}a_n.
\]
The complete head--tail estimate gives
\begin{equation}\label{a:eq:coercive}
 \ell_\delta[u]\ge\tg[u]+\frac\delta2\norm{u}_{\Hh}^2-C_N\norm{u}_2^2.
\end{equation}
Its upper absolute bound is a finite multiple of $\norm{u}_X^2$, by Lemma~\ref{a:lem:prime} and~\eqref{a:eq:L2bound}. Therefore
$\ell_\delta[u]+(C_N+1)\norm{u}_2^2$ defines a positive norm squared equivalent to $\norm{u}_X^2$. Completeness and density follow from Lemma~\ref{a:lem:Xcore}. The closed-form representation theorem gives the self-adjoint operator and~\eqref{a:eq:Ldef}.
\end{proof}

The choice of $N$ in this proof is existential and analytic. The form remains the complete form for every such choice. In distributional notation the operator is
\begin{equation}\label{a:eq:formalL}
 L_\delta=D_\Gamma-\cg-P+\frac\delta W,
 \qquad
 Pu=\sum_{n\ge2}a_n\bigl(u(\cdot+s_n)+u(\cdot-s_n)\bigr),
\end{equation}
where $D_\Gamma$ is represented by $\tg$. For $u\in X$, the prime row converges in $L^2$ on every fixed compact interval: each shifted local norm is bounded by the local supremum of $\sqrt W$ times $\norm{u}_{\Hh}$, and the theta envelope makes the resulting tails summable. Also $u/W$ is locally $L^2$. The Gamma term is defined distributionally by its form. Operator-domain membership requires that the \emph{sum} in~\eqref{a:eq:formalL} be represented by an ordinary $L^2$ vector; it does not assert that each term is globally $L^2$. Compact smooth functions are a form core. No operator-core claim for this ordinary-$L^2$ realization is being made.

\begin{lemma}[Logarithmic Fourier coercivity]\label{a:lem:fourier}
For the unitary Fourier transform,
\begin{equation}\label{a:eq:multiplier}
 \tg[u]=\int_\R d_\Gamma(\xi)|\widehat u(\xi)|^2\,d\xi,
 \quad
 d_\Gamma(\xi)=2\int_0^\infty\rho(s)(1-\cos(\xi s))\,ds
 \ge\frac12\log(1+4\xi^2).
\end{equation}
\end{lemma}
\begin{proof}
Plancherel's theorem and Tonelli's theorem give the multiplier identity, allowing infinite values. The inequality $1-e^{-2s}\le2s$ gives
\[
 d_\Gamma(\xi)\ge\int_0^\infty e^{-s/2}\frac{1-\cos(\xi s)}s\,ds.
\]
For general $a>0$, differentiating this last integral with $e^{-as}$ in place of $e^{-s/2}$ yields $\int_0^\infty e^{-as}\sin(\xi s)\,ds=\xi/(a^2+\xi^2)$. Integrating back from $\xi=0$ gives $\tfrac12\log(1+\xi^2/a^2)$. Taking $a=1/2$ proves the bound.
\end{proof}

\begin{proposition}[Compactness]\label{a:prop:compact}
The embedding $X\hookrightarrow\Le$ is compact. Consequently $L_\delta$ has compact resolvent for every $\delta>0$.
\end{proposition}
\begin{proof}
The unit ball of $X$ is uniformly spatially tight, since
\begin{equation}\label{a:eq:spatialtight}
 \int_{|x|>R}|u(x)|^2\,dx
 \le\sup_{|x|>R}W(x)\,\norm{u}_{\Hh}^2\longrightarrow0
\end{equation}
uniformly on that ball. Lemma~\ref{a:lem:fourier} gives uniform Fourier tightness:
\begin{equation}\label{a:eq:fouriertight}
 \int_{|\xi|>M}|\widehat u(\xi)|^2\,d\xi
 \le\frac{\tg[u]}{\tfrac12\log(1+4M^2)}\longrightarrow0.
\end{equation}
Let $P_M$ be the Fourier projection onto $[-M,M]$. The operator
$\one_{[-R,R]}P_M$ is Hilbert--Schmidt: its space--frequency kernel has constant modulus $(2\pi)^{-1/2}$ on a finite rectangle and vanishes elsewhere. Moreover
\[
 \norm{u-\one_{[-R,R]}P_Mu}_2
 \le\norm{\one_{\{|x|>R\}}u}_2+\norm{(I-P_M)u}_2.
\]
Thus the inclusion is a uniform limit on the $X$ unit ball of compact operators. The cutoffs and Fourier projection preserve evenness.

By~\eqref{a:eq:coercive}, a resolvent $(L_\delta+aI)^{-1}$ with $a>C_N$ maps ordinary $L^2$ boundedly into $X$. Composing this map with the compact inclusion proves compactness of the resolvent. Hence the eigenvalues are real, have finite multiplicity, are bounded below, and tend to $+\infty$.
\end{proof}

This compactness concerns $L_\delta$ on ordinary $L^2$, not the original operator $A_E$ on $L^2(\mu)$. The two Hilbert norms are different.


\subsection{Actual eigenstates and the negative index}\label{f:sec:spectrum}

Proposition~\ref{a:prop:domain} gives, for $\lambda=1/2-\delta$,
\begin{equation}\label{f:eq:congruence}
 \ell_\delta(\kappa f,\kappa g)
       =\E(f,g)-\lambda\ip{f}{g}_\mu,
       \qquad f,g\in\D.
\end{equation}
This is a form congruence through a bijection of form domains, not a unitary equivalence of operators on their displayed Hilbert spaces.

\begin{proposition}[Fixed-distance spectral finiteness]\label{f:prop:finite}
For every $t<1/2$, the spectral projection $\one_{(-\infty,t]}(A)$ has finite rank. The centered restriction satisfies $A_E\ge\varepsilon_0 I$ for some $\varepsilon_0>0$ and
\begin{equation}\label{f:eq:essential}
 \inf\sigma_{\ess}(A_E)\ge\frac12,
\end{equation}
where the infimum of the empty essential spectrum is $+\infty$.
\end{proposition}
\begin{proof}
For $t<0$ the assertion follows from $A\ge0$. Let $0\le t<1/2$, and take $\delta=1/2-t$. Every $f$ in the range of $\one_{[0,t]}(A)$ lies in $\dom(A)$ and satisfies $\E[f]\le t\norm{f}_\mu^2$. Equation~\eqref{f:eq:congruence} sends this range injectively into a subspace of $X$ on which $\ell_\delta\le0$.

A lower-bounded compact-resolvent operator has only finitely many nonpositive eigenvalues. Its nonpositive spectral subspace is therefore finite dimensional. Any subspace of its form domain on which the form is nonpositive has dimension at most that finite number: otherwise a nonzero vector in the subspace would be orthogonal to the entire nonpositive spectral subspace, where the form is strictly positive. Applying this observation to $L_\delta$ proves finite rank.

The centered restriction has no zero eigenvector, by Proposition~\ref{a:prop:core}. Apply the finite-rank statement at any one $t_0\in(0,1/2)$, for example $t_0=1/4$. On its centered finite-dimensional spectral range, if nonzero, the smallest eigenvalue is strictly positive; on the complementary range the operator is bounded below by $t_0$. This proves a positive lower bound. Finally, finite rank below every $t<1/2$ excludes essential spectrum below $1/2$ and permits accumulation of subthreshold eigenvalues only at $1/2$.
\end{proof}

The proposition is a consequence of the compact pencil and the bridge. It does not require an additional estimate for primes in short intervals. The next result establishes equality in~\eqref{f:eq:essential}.


\begin{proposition}[Actual half eigenstates and essential bottom]\label{f:prop:halfeigen}
The profiles
\begin{equation}\label{f:eq:halfprofiles}
 g_j=\kappa^{(2j)}/\kappa-4^{-j}\quad(j\ge1)
\end{equation}
belong to $\dom A_E$, satisfy $A_Eg_j=g_j/2$, and are linearly independent. Hence
\begin{equation}\label{f:eq:exactessential}
 \dim\ker(A_E-\tfrac12I)=\infty,\qquad
 \inf\sigma_{\mathrm{ess}}(A_E)=\tfrac12.
\end{equation}
\end{proposition}
\begin{proof}
Lemma~\ref{a:lem:radical} proves $r_j\in X_c$ and $q_0(r_j,z)=0$ for every $z\in X$, with linear independence. The representation criterion for $B$ gives $r_j\in\dom B$, $Br_j=0$. Apply~\eqref{eq:actual-unitary} and $U^{-1}r_j=g_j$. The resulting infinite-dimensional eigenspace puts $1/2$ in the essential spectrum. Proposition~\ref{f:prop:finite} excludes essential spectrum below it.
\end{proof}
The equality states do not exclude lower discrete spectrum. No completeness of the derivative family in the centered weighted space has been asserted; its complement is the subject of Section~\ref{sec:synthesis-unified}.


The compact-resolvent eigenvalues are enumerated increasingly, with multiplicity, as $\eta_1(\delta)\le\eta_2(\delta)\le\cdots$. For $0<\delta<1/2$ and $\lambda=1/2-\delta$, the kernel correspondence to be proved is
\begin{equation}\label{eq:pencil-kernel}
 \ker L_\delta=\{\kappa f:f\in\dom A_E,\ A_Ef=\lambda f\}.
\end{equation}


\begin{proposition}[Kernel correspondence]\label{f:prop:kernel}
For $0<\delta<1/2$ and $\lambda=1/2-\delta$, equation~\eqref{eq:pencil-kernel} holds as an exact linear correspondence, including multiplicities.
\end{proposition}
\begin{proof}
If $u\in\ker L_\delta$, then $u\in X$ and $f=u/\kappa\in\D$. Testing~\eqref{a:eq:Ldef} against every $z=\kappa g\in X$ and using~\eqref{f:eq:congruence} gives
\[
 \E(f,g)=\lambda\ip{f}{g}_\mu\qquad(g\in\D).
\]
The representation criterion~\eqref{a:eq:Adef} proves $f\in\dom(A)$ and $Af=\lambda f$. Since $\lambda>0$ and $A1=0$, self-adjointness implies $\mu(f)=0$. Thus $f\in\dom(A_E)$ is a genuine centered eigenstate.

Conversely, for such an eigenstate $f$, the vector $u=\kappa f$ belongs to $X$ and~\eqref{f:eq:congruence} vanishes against every $z\in X$. Its ordinary-$L^2$ representing vector is zero. Hence $u\in\dom(L_\delta)$ and $L_\delta u=0$.
\end{proof}

No pointwise eigen-row or additional regularity of an eigenstate is used here. The correspondence identifies kernels at the indicated parameter; it does not identify nonzero eigenvalues of $L_\delta$ with eigenvalues of $A_E$.

\begin{proposition}[Exact form index]\label{f:prop:index}
For $0<\delta<1/2$ and $\lambda=1/2-\delta$,
\begin{align}
 n_-(L_\delta)&=1+\rank\one_{(0,\lambda)}(A_E),\label{f:eq:index}\\
 \dim\ker L_\delta&=\dim\ker(A_E-\lambda).\label{f:eq:nullity}
\end{align}
For $u\in X$, put $a=2m(u)$ and $u_0=u-a\kappa$. Then
\begin{align}
 m(u_0)&=0,\qquad
 \norm{u}_{\Hh}^2=|a|^2+\norm{u_0}_{\Hh}^2,\label{f:eq:decompnorm}\\
 \ell_\delta[u]&=-\lambda|a|^2+
              \E[u_0/\kappa]-\lambda\norm{u_0}_{\Hh}^2.\label{f:eq:decompform}
\end{align}
\end{proposition}
\begin{proof}
Write $f=u/\kappa=a+g$, where $a=\mu(f)$ and $\mu(g)=0$. Constants have zero energy, and polarization gives $\E(1,g)=0$. Also $\norm{1}_\mu=1$. The orthogonal decomposition in $L^2(\mu)$ and~\eqref{f:eq:congruence} prove~\eqref{f:eq:decompnorm}--\eqref{f:eq:decompform}.

A bijective linear change of form domain preserves the maximal dimension of a strictly negative form subspace. On the constant line the right side of~\eqref{f:eq:decompform} is strictly negative. On the centered subspace its negative index is the rank of the spectral projection of $A_E$ below $\lambda$, by the spectral theorem. Proposition~\ref{f:prop:finite} and the positive gap make this the finite rank in~\eqref{f:eq:index}. For a compact-resolvent operator its form negative index equals its number of negative eigenvalues. The nullity identity is Proposition~\ref{f:prop:kernel}.
\end{proof}

\begin{corollary}[Failure of the half bound]\label{f:cor:failure}
The global inequality $\E[f]\ge\tfrac12\Var_\mu(f)$ on $\D$ fails if and only if there exists an actual eigenpair
\[
 A_Ef=\lambda f,\qquad f\in\dom(A_E)\setminus\{0\},\qquad0<\lambda<1/2.
\]
\end{corollary}
\begin{proof}
Subtract the mean of a violating profile to obtain $g\ne0$ with $\E[g]/\norm{g}_\mu^2<1/2$. Choose $t$ strictly between this quotient and $1/2$. The spectral projection $\one_{(-\infty,t]}(A_E)$ is nonzero, since otherwise the spectral theorem would give $\E[g]\ge t\norm{g}_\mu^2$. It is finite rank by Proposition~\ref{f:prop:finite}, so its invariant range contains an eigenvector. Its eigenvalue is positive by the gap and at most $t<1/2$. The converse follows by evaluating the energy on the eigenvector.
\end{proof}

\subsection{The mean and the operator-domain obstruction}\label{f:sec:mean}

The decomposition in Proposition~\ref{f:prop:index} is orthogonal in the weighted norm, but not in ordinary physical $L^2$. Two concrete facts make the distinction unavoidable.

\begin{proposition}[The mean is unbounded in ordinary $L^2$]\label{f:prop:mean}
The functional $m$ is continuous on $X$ and on the weighted space, with bound~\eqref{a:eq:meanbound}, but is not bounded in the ordinary $L^2$ norm, even on even $C_c^\infty(\R)$. Its zero-mean subspace in $X$ is closed in $X$ and dense in ordinary $\Le$.
\end{proposition}
\begin{proof}
Only unboundedness and density remain to be shown. Fix a nonzero, nonnegative, even $\phi\in C_c^\infty(\R)$ and put $M=m(\phi)>0$. For sufficiently large $R$, the two translates in
\[
 \psi_R(x)=\phi(x-R)+\phi(x+R),\qquad
 b_R(x)=\frac{\psi_R(x)}{2\cosh(R/2)M}
\]
have disjoint supports. The addition formula for the hyperbolic cosine and evenness of $\phi$ give $m(\psi_R)=2\cosh(R/2)M$. Hence
\begin{equation}\label{f:eq:unboundedmean}
 m(b_R)=1,\qquad
 \norm{b_R}_2^2=\frac{\norm{\phi}_2^2}{2\cosh^2(R/2)M^2}\longrightarrow0.
\end{equation}
Every $b_R$ is in the compact smooth numerator core. Thus neither $m$ nor the quadratic pole $2|m|^2$ has an ordinary-$L^2$ bound. For any even compact smooth $u$, the corrected function $u-m(u)b_R$ has zero mean and converges to $u$ in ordinary $L^2$. Density of the compact smooth core in $\Le$ proves the last claim. No smallness of $\norm{b_R}_X$ or $\norm{b_R}_{\Hh}$ is asserted.
\end{proof}

\begin{proposition}[The theta kernel is generally not in the operator domain]\label{f:prop:kdomain}
For every $\delta>0$, $\kappa\in X$. If $\delta\ne1/2$, then $\kappa\notin\dom(L_\delta)$. At $\delta=1/2$,
\begin{equation}\label{f:eq:halfkernel}
 L_{1/2}\ge0,\qquad \ker L_{1/2}=\C\kappa.
\end{equation}
For $0<\delta<1/2$, $\kappa$ is nevertheless a strictly negative form direction.
\end{proposition}
\begin{proof}
Since $1\in\D$, Proposition~\ref{a:prop:domain} admits $\kappa$ to $X$ and gives $\norm{\kappa}_{\Hh}=1$, $m(\kappa)=1/2$. Put $\lambda=1/2-\delta$, now allowing either sign. Equation~\eqref{f:eq:congruence} gives the exact polarized identity
\begin{equation}\label{f:eq:sourcerow}
 \ell_\delta(\kappa,z)=-2\lambda\overline{m(z)},\qquad z\in X.
\end{equation}
If $\kappa\in\dom(L_\delta)$, then compact smooth even tests in~\eqref{a:eq:Ldef} force its ordinary-$L^2$ representative to equal $-2\lambda c$ as an even distribution. For $\lambda\ne0$ this is impossible, since $c\notin L^2(\R)$. When $\lambda=0$, the form is $\E[u/\kappa]\ge0$ and its zero vectors are precisely multiples of $\kappa$, by Proposition~\ref{a:prop:core}. Finally $\ell_\delta[\kappa]=-\lambda<0$ for $0<\delta<1/2$.
\end{proof}

Thus a negative \emph{form direction} is not necessarily an operator-domain vector, much less an ordinary-$L^2$ ground eigenvector. The mean-zero space is not an ordinary-$L^2$ closed codimension-one hyperplane. Moreover $\int u\overline\kappa\,dx=0$ is different from $m(u)=0$: the two functionals on compact tests would be proportional only if $\kappa$ were a scalar multiple of $c$, contrary to their decay. Removing $\kappa$ by an ordinary-$L^2$ orthogonal projection therefore does not implement centering of the original problem.

\subsection{The second eigenvalue and the remaining sign}\label{f:sec:sign}

The compact-resolvent construction gives a precise family of sign questions, but it supplies no positive sign for the second eigenvalue throughout the required parameter interval.

\begin{proposition}[Equivalent sign conditions]\label{f:prop:sign}
The following statements are equivalent:
\begin{enumerate}
\item $\E[f]\ge\tfrac12\Var_\mu(f)$ for every $f\in\D$.
\item For every $0<\delta<1/2$, $\eta_2(\delta)>0$.
\item For every $0<\delta<1/2$ and every nonzero $u\in X$ with $m(u)=0$, one has $\ell_\delta[u]>0$.
\item $\ker L_\delta=\{0\}$ for every $0<\delta<1/2$.
\end{enumerate}
\end{proposition}
\begin{proof}
If the first statement holds, then on centered vectors~\eqref{f:eq:congruence} gives $\ell_\delta[u]\ge\delta\norm{u}_{\Hh}^2>0$. It also gives no centered spectrum of $A_E$ at or below $\lambda=1/2-\delta$. Equations~\eqref{f:eq:index}--\eqref{f:eq:nullity} then show that $L_\delta$ has exactly one negative eigenvalue and no zero eigenvalue; compact resolvent gives $\eta_2(\delta)>0$.

Conversely, if $\eta_2(\delta)>0$, the index and nullity identities exclude centered spectrum at or below $\lambda$, proving strict positivity on the centered form subspace. If the third statement holds for all $\delta$, letting $\delta\downarrow0$ for each fixed centered $u$ gives $q_0[u]\ge0$. By~\eqref{a:eq:polebridge} and subtraction of the original mean, this is the first statement. Finally, the equivalence of the first and fourth statements follows from Corollary~\ref{f:cor:failure} and Proposition~\ref{f:prop:kernel}.
\end{proof}

Two valid estimates illustrate the limits of the construction. First, min--max on the common domain $X$ and~\eqref{a:eq:L2bound} give
\begin{equation}\label{f:eq:monotonicity}
 \eta_j(\delta_2)\ge\eta_j(\delta_1)
       +\frac{\delta_2-\delta_1}{\norm{W}_\infty},
       \qquad\delta_2>\delta_1>0.
\end{equation}
At $\delta=1/2$, Proposition~\ref{f:prop:kdomain} and compact resolvent give $\eta_2(1/2)>0$. However, decreasing $\delta$ decreases the form. Inequality~\eqref{f:eq:monotonicity} does not prohibit a second eigenvalue from crossing zero at a smaller positive parameter.

Second, suppose $u\in\ker L_\delta$ with $0<\delta<1/2$ and normalize $\norm{u}_{\Hh}=1$. Then
\begin{equation}\label{f:eq:stateidentity}
 \tg[u]+\delta=\cg\norm{u}_2^2+p[u].
\end{equation}
Choose $N$ with $r_N\le\delta/2$ as in Proposition~\ref{a:prop:closedpencil}. The same complete-tail estimate yields
\begin{equation}\label{f:eq:statebound}
 \tg[u]+\frac\delta2\le C_N\norm{u}_2^2,
 \qquad \norm{u}_2^2\ge\frac{\delta}{2C_N}.
\end{equation}
Thus a normalized fixed-gap candidate cannot lose all its ordinary physical $L^2$ mass. The estimate is not contradictory: reducing the tail by enlarging $N$ also enlarges the head constant $C_N$. No estimate reversing this balance is proved. Testing the kernel equation against $\kappa$ gives only $m(u)=0$, by~\eqref{f:eq:sourcerow}; this recovers centering and provides no second independent sign constraint.


\section{Global variance comparisons and exact mean repair}\label{sec:bounds}
There are three comparison levels. The first is wholly analytic. The second uses only a whole-local-domain sign. The third has additional response hypotheses. None of the proofs replaces the complete energy by its finite head.


\begin{theorem}[Analytic variance and inertia bounds]\label{a:thm:mainseed}
For every $f\in\D$,
\begin{equation}\label{a:eq:mainseed}
 \E[f]\ge c_{\log 2}\Var_\mu(f),\qquad
 c_{\log 2}=\frac4{85}\left(\frac{63}{128}-7238e^{-12}\right)
 >\frac{8939}{425000}>\frac1{48}.
\end{equation}
For each $\delta>0$, the form $\ell_\delta$ on $X$ has a lower-bounded self-adjoint representative $L_\delta$ in ordinary $\Le$, with compact resolvent. If
\[
 \frac12-c_{\log 2}<\delta<\frac12,
\]
then $L_\delta$ has exactly one negative eigenvalue, counted with multiplicity, and zero operator nullity. At $\delta_3=23/48$, its exact mean Schur scalar is $-1/48$, and on $X_c=\{u\in X:\int_\R cu=0\}$,
\begin{equation}\label{a:eq:mainmargin}
 \ell_{\delta_3}[u]\ge\frac{509}{2550000}\norm{u}_{\Hh}^2.
\end{equation}
The strongest open endpoint furnished by this particular bound is $1/2-c_{\log 2}$, not its rational rounding $1/2-8939/425000$.
\end{theorem}


\subsection{A radius-sensitive analytic variance bound}\label{a:sec:seed}

We now prove the first quantitative conclusion. The estimate uses a fixed small prime head, but the weighted remainder from Lemma~\ref{a:lem:prime} remains present throughout. Combining the full-domain bridge with~\eqref{a:eq:head4} gives, for every $f\in\D$,
\begin{equation}\label{a:eq:masscontrol}
 \E[f]\ge\tg[\kappa f]+\frac{63}{128}\norm f_\mu^2
                         -9\norm{\kappa f}_2^2.
\end{equation}
Indeed the coefficient of $\norm f_\mu^2$ before rounding is $1/2-r_4>63/128$. The head through $4$ includes the prime square $4$; every later prime correlation is charged to $r_4$.

\subsubsection{Even folding and the global mean}

For $r>0$, write $I_r=[-r,r]$ and set
\begin{equation}\label{a:eq:radiusdata}
 \varepsilon_r=\mu(I_r^c),\qquad
 \Omega_r=\sup_{I_r^c}W,\qquad
 \alpha_r=\frac{\rho(r)}{4c(r)}
 =\frac1{2(e^r+1)(1-e^{-2r})}.
\end{equation}

\begin{lemma}[Folded local variance with global centering]\label{a:lem:fold}
If $f\in\D$ and $\mu(f)=0$, put $T=\norm f_\mu^2$ and
$A_r=\int_{I_r}|f|^2\,d\mu$. Then
\begin{equation}\label{a:eq:foldbound}
 \E[f]\ge\alpha_r(A_r-\varepsilon_rT),\qquad
 A_r\le\frac{\E[f]}{\alpha_r}+\varepsilon_rT.
\end{equation}
\end{lemma}
\begin{proof}
On $[0,r]$, let $d\nu(a)=4c(a)\kappa(a)\,da$. Evenness gives
\[
 \nu([0,r])=\mu(I_r),\qquad
 \int_0^r|f|^2\,d\nu=A_r,\qquad
 \int_0^rf\,d\nu=\int_{I_r}f\,d\mu.
\]
Restrict the full ordered-pair Gamma integral to $I_r^2$. The two same-sign regions and the opposite-sign region fold to
\begin{equation}\label{a:eq:foldidentity}
 2\int_{0<a<b<r}[\rho(b-a)+\rho(b+a)]
       \kappa(a)\kappa(b)|f(b)-f(a)|^2\,da\,db.
\end{equation}
For $0<a<b$, the complete positive exponential series gives
\begin{align*}
 \rho(b-a)+\rho(b+a)
 &=2\sum_{j\ge0}e^{-(2j+1/2)b}\cosh((2j+1/2)a)\\
 &\ge2c(a)\rho(b).
\end{align*}
Relative to $d\nu(a)d\nu(b)$, the coefficient in~\eqref{a:eq:foldidentity} is therefore at least
\[
 \frac{\rho(b-a)+\rho(b+a)}{8c(a)c(b)}
 \ge\frac{\rho(b)}{4c(b)}\ge\alpha_r,
\]
since $\rho$ decreases and $c$ increases on the positive half-line. Tonelli's theorem justifies both the restriction and the positive series. The ordered-pair variance identity yields
\[
 \E[f]\ge\alpha_r\left((1-\varepsilon_r)A_r
                  -\left|\int_{I_r}f\,d\mu\right|^2\right).
\]
The original global mean is retained:
\[
 \left|\int_{I_r}f\,d\mu\right|^2
 =\left|\int_{I_r^c}f\,d\mu\right|^2
 \le\varepsilon_r(T-A_r).
\]
Substitution proves~\eqref{a:eq:foldbound}. This proof applies to complex profiles on the entire maximal domain; it makes no local-mean substitution.
\end{proof}

Because $w_{\max}<1/2$,
\begin{equation}\label{a:eq:masswindow}
 \norm{\kappa f}_2^2=\int_\R W|f|^2\,d\mu
 \le\frac12A_r+\Omega_rT.
\end{equation}
Dropping only the nonnegative $\tg$ term in~\eqref{a:eq:masscontrol}, and then using Lemma~\ref{a:lem:fold}, gives the general template
\begin{equation}\label{a:eq:seedtemplate}
 \E[f]\ge
 \frac{63/128-9\Omega_r-(9/2)\varepsilon_r}
      {1+9/(2\alpha_r)}\,T
\end{equation}
whenever the numerator is positive. The restriction of a nonnegative integral supplies a lower bound; it does not replace the underlying energy by a local model.

\subsubsection{The radius \texorpdfstring{$\log2$}{log 2}}

At $r_0=\log2$, formula~\eqref{a:eq:radiusdata} becomes
\begin{equation}\label{a:eq:alphalog2}
 \alpha_{r_0}=\frac29,\qquad
 A_{r_0}\le\frac92\E[f]+\varepsilon_{r_0}T.
\end{equation}
The complete theta envelope holds at every nonnegative radius. Since $4r-3e^{2r}$ decreases,
\begin{equation}\label{a:eq:Omegalog2}
 \Omega_{r_0}\le42e^{4r_0-3e^{2r_0}}=672e^{-12}.
\end{equation}
For the mass tail, use $2c(r)\le2e^{r/2}$ and the envelope for $\kappa$ itself. Evenness and $t=e^{2r}$ give
\begin{align}
 \varepsilon_{r_0}
 &\le168\int_{r_0}^\infty e^{5r-3e^{2r}}\,dr
 =84\int_4^\infty t^{3/2}e^{-3t}\,dt\notag\\
 &\le42\int_4^\infty t^2e^{-3t}\,dt
 =\frac{2380}{9}e^{-12}.\label{a:eq:epslog2}
\end{align}
Here $t^{3/2}\le t^2/2$ for $t\ge4$, and integration by parts gives
\[
 \int_4^\infty t^2e^{-3t}\,dt
 =e^{-12}\left(\frac{16}{3}+\frac89+\frac2{27}\right)
 =\frac{170}{27}e^{-12}.
\]
Equations~\eqref{a:eq:Omegalog2}--\eqref{a:eq:epslog2} are full exterior estimates, not values sampled at the boundary.

\begin{proof}[Proof of the variance bound in Theorem~\ref{a:thm:mainseed}]
For centered $f$, combine~\eqref{a:eq:masscontrol}, \eqref{a:eq:masswindow}, and~\eqref{a:eq:alphalog2} to get
\[
 \E[f]\ge\left(\frac{63}{128}-9\Omega_{r_0}
                       -\frac92\varepsilon_{r_0}\right)T
                         -\frac{81}{4}\E[f].
\]
Move the last term to the left. The coefficient of $e^{-12}$ in the tail cost is
$9\cdot672+(9/2)(2380/9)=6048+1190=7238$. Therefore
\[
 \E[f]\ge\frac4{85}\left(\frac{63}{128}-7238e^{-12}\right)T
 =c_{\log2}T.
\]
By~\eqref{a:eq:elementarye}, $e^{12}>160000$, whence
\begin{align*}
 c_{\log2}&>\frac4{85}\left(\frac{63}{128}-\frac{7238}{160000}\right)
 =\frac4{85}\frac{8939}{20000}=\frac{8939}{425000}=:c_3,\\
 c_3-\frac1{48}
 &=\frac{8939\cdot48-425000}{425000\cdot48}
 =\frac{509}{2550000}>0.
\end{align*}
For general $f\in\D$, replace $f$ by $f-\mu(f)$. Constants belong to $\D$, the energy is unchanged, and the squared norm becomes $\Var_\mu(f)$. This proves~\eqref{a:eq:mainseed} on the full original domain.
\end{proof}

No optimality is asserted for $c_{\log2}$. It is the explicit coefficient proved by this particular radius and these estimates. In particular the coarse template~\eqref{a:eq:seedtemplate}, with~\eqref{a:eq:masscontrol} fixed, is bounded above by $63/128<1/2$: its loss terms are nonnegative and its denominator is at least one. That observation limits this template, not other possible estimates for the original energy.

\subsubsection{The exact mean and the inertia interval}

Let $\Hh_c=\{u\in\Hw:m(u)=0\}$ and $X_c=X\cap \Hh_c$. Since
$\norm\kappa_{\Hh}=1$ and $\ip u\kappa_{\Hh}=2m(u)$, the map
\begin{equation}\label{a:eq:centering}
 \mathsf C u=u-2m(u)\kappa
\end{equation}
is the exact weighted orthogonal projection onto $\Hh_c$, and it preserves $X$. The full-domain bridge gives, for $u=a\kappa+w$ with $w\in X_c$ and $a=2m(u)$,
\begin{align}
 \norm u_{\Hh}^2&=|a|^2+\norm w_{\Hh}^2,\label{a:eq:meansplitnorm}\\
 \ell_\delta[u]&=(\delta-1/2)|a|^2+\ell_\delta[w].\label{a:eq:meansplit}
\end{align}
This is weighted centering, not ordinary-$L^2$ projection away from $\kappa$.

\begin{proof}[Proof of the inertia conclusion in Theorem~\ref{a:thm:mainseed}]
Write $g_\delta=c_{\log2}+\delta-1/2$. The variance bound implies
\begin{equation}\label{a:eq:seedcoercivity}
 \ell_\delta[w]\ge g_\delta\norm w_{\Hh}^2\qquad(w\in X_c).
\end{equation}
For $1/2-c_{\log2}<\delta<1/2$, this margin is positive. The head--tail bound and $w_{\max}<1/2$ also give
\[
 \ell_\delta[w]\ge\tg[w]-C\norm w_{\Hh}^2,
 \qquad C=9/2+r_4.
\]
Combine this estimate with~\eqref{a:eq:seedcoercivity}, assigning weight
$\lambda=g_\delta/[2(g_\delta+C)]$ to the first one. Then
\begin{equation}\label{a:eq:fullseedcoercivity}
 \ell_\delta[w]\ge\lambda\tg[w]
                         +\frac{g_\delta}{2}\norm w_{\Hh}^2.
\end{equation}
Thus the centered infinite complement is coercive in the full form norm. By~\eqref{a:eq:meansplit}, its exact elimination leaves the scalar $\delta-1/2<0$. The maximal strictly negative form dimension is one and the bilinear form kernel is zero. For a lower-bounded compact-resolvent operator, projection onto its negative spectral subspace is injective on every strictly negative form subspace, while that spectral subspace attains the bound. Thus the form negative index equals the number of negative eigenvalues with multiplicity. By~\eqref{a:eq:Ldef}, a form-kernel vector is precisely an operator-domain vector with representing row zero, so the nullities also agree.

At $\delta_3=23/48$, the scalar is $-1/48$ and
$g_{\delta_3}>c_3-1/48=509/2550000$, proving~\eqref{a:eq:mainmargin}.
The endpoint $1/2-c_{\log2}$ is not certified by the non-strict variance inequality alone. The rational endpoint $1/2-c_3$ lies strictly inside the certified interval, with margin $c_{\log2}-c_3>0$.
\end{proof}

Zero operator nullity does not mean that an indefinite quadratic form has no isotropic vectors. Likewise, the mean Schur scalar $-1/48$ is a congruence entry, not the lowest ordinary-$L^2$ eigenvalue. Indeed the source row gives
$\ell_\delta(\kappa,z)=(2\delta-1)\overline{m(z)}$ on $X$; for $\delta\ne1/2$ its distributional representative is $(2\delta-1)c\notin L^2(dx)$. Hence $\kappa$ is a negative form direction for $0<\delta<1/2$, but not an operator-domain eigenvector of $L_\delta$.


\subsection{The local interval and the two computational premises}\label{c:sec:premises}
Set $B_0=4/5$, $I=(-B_0,B_0)$ and $O=\R\setminus[-B_0,B_0]$. This scalar radius is distinct from the weighted operator $B$ in~\eqref{eq:Db-physical}. The local minimum domain $\D_I$ is the closure of even $C_c^\infty(I)$, zero extended, in $(\tg+\norm{\cdot}_2^2)^{1/2}$. The form
\[
 q[u]=q_0[u]+2|m(u)|^2
\]
has restriction $q_I=q|_{\D_I}$ and ordinary physical $L^2(I)$ representative $L_I$. The sharp-cut and exterior minimum-domain statements used below are proved in Appendix~\ref{c:app:cuts}; they are not additional computational premises.


\begin{hypothesis}[Accepted local input theorem]\label{c:hyp:local}
On the whole domain $\D_I$, $q_I[u]\ge0$.
\end{hypothesis}
The accepted computer-assisted input states the stronger inequality $q_I[u]\ge10^{-22}\norm u_2^2$. Only its sign is used in the first theorem. Its physical $L^2$ coefficient is not transferred unchanged to a weighted global gap. Appendix~\ref{c:app:provenance} identifies the accepted evidence; it does not claim a fresh replay.

\begin{theorem}[Local-sign branch]\label{c:thm:localglobal}
Under Hypothesis~\ref{c:hyp:local}, every even $f\in\D$ satisfies
\begin{equation}\label{c:eq:49307}
 \E[f]\ge\frac{49307}{100000}\Var_\mu(f).
\end{equation}
The same comparison permits the endpoint
\begin{equation}\label{c:eq:lambdaendpoint}
 \begin{gathered}
 \lambda_{\rm cmp}=\frac12-\frac12\left(a_1+\sqrt{a_1^2+4L_1^2}\right)
 >\frac{49307}{100000},\\
 a_1=\frac{39703}{56448000},\qquad L_1=\frac{2099}{320000},\qquad
 \E[f]\ge\lambda_{\rm cmp}\Var_\mu(f).
 \end{gathered}
\end{equation}
No local inverse-response enclosure is needed. The displayed rational coefficient improves $493/1000$ by exactly $7/100000$; it does not approach $1/2$ through a parameter family.
\end{theorem}

\begin{theorem}[Additional-response branch]\label{c:thm:conditionalglobal}
Under the complete-operator and centered-response contracts in Hypothesis~\ref{c:hyp:response}, every even $f\in\D$ satisfies
\begin{equation}\label{c:eq:newglobal}
 \E[f]\ge\frac{2861149}{5760000}\Var_\mu(f).
\end{equation}
These contracts imply local positivity but are not implied by Hypothesis~\ref{c:hyp:local}. Their numerical enclosure assertions are inputs, not computations established here.
\end{theorem}
The latter coefficient exceeds $71501077/144000000$, obtained from a larger source cost under the same contracts, by exactly $24/125000$. This comparison concerns conditional implications, not the endpoint $1/2$. Arguments are given for real-even functions; decomposition into real and imaginary parts gives the complex statements. Inner products are linear in the first argument. 


\subsection{Analytic estimates and the two source coordinates}\label{c:sec:estimates}
Write $\alpha=\sqrt W$, and set
\[
 A(x)=\frac{(4t^2-6t)e^{-t}}{1+e^{-x}}+128t^2e^{-4t},
 \quad t=\pi e^{2x},\qquad A=A({B_0}).
\]
Appendix~\ref{c:app:scalar} proves
\begin{gather}
 A<1/9216,\quad W({B_0}+h)\le Ae^{-25h},\quad
 \int_{B_0}^\infty W<1/230400,\label{c:eq:envelope}\\
 \sup W<1/2,\quad \int_0^\infty W\le1/8,\quad
 p_*:=\mu(O)<3/49000,\quad \cg<383/70<11/2,\label{c:eq:basiccaps}\\
 W(1/10)<21/50,\quad W(1/4)<1/4,\quad W(9/20)<3/40,\label{c:eq:pointW}\\
 \kappa(2/5)<1/4,\quad \kappa(1/2)<7/100,\quad
 \kappa(3/5)<1/50,\quad \kappa(7/10)<3/1000,\quad
 \kappa(3/4)<1/1000.\label{c:eq:pointkappa}
\end{gather}
Both $\kappa$ and $W$ decrease on $[1/10,\infty)$. Also,
\begin{equation}\label{c:eq:rhoelementary}
 \rho(s)\le\frac1{2s}+\frac14,\qquad
 \frac56<\rho({B_0})<\frac9{10},\qquad \pi^2<10.
\end{equation}
These are full-series analytic inequalities.

There are two source norms. The weighted profile interface $\mathcal J=\mathcal J_p+\mathcal J_\Gamma$ is the inside-to-outside operator in
\begin{equation}\label{c:eq:energycross}
 \E[g+h]=\E[g]+\E[h]-2\ip{\mathcal Jg}{h}_\mu,
 \qquad \supp g\subset\overline I,\quad\supp h\subset\overline O.
\end{equation}
Under the unitary map
\begin{equation}\label{c:eq:foldU}
 (U_+ f)(x)=\sqrt{4c(x)\kappa(x)}f(x),\qquad x>0,
\end{equation}
its Gamma kernel is $\sqrt{W(x)W(y)}[\rho(x-y)+\rho(x+y)]$, $x>{B_0}$, $0<y<{B_0}$.

The physical exterior source in the second branch is
\begin{equation}\label{c:eq:physicalsource}
 (Ju)(y)=\int_O\rho(|x-y|)u(x)\,dx+
 \sum_{n\ge2}a_n[u(y+s_n)+u(y-s_n)],\qquad y\in I,
\end{equation}
where exterior inputs are extended by zero. Its input norm is
$\wo(u)=\int_O|u|^2/W\,dx$, and its output norm is physical $L^2(I)$.
With $v(x)=\sqrt2u(x)/\alpha(x)$ on $x>{B_0}$, the input norm becomes $\norm v_2^2=\wo(u)$ and the output fold is multiplication by $\sqrt2$. The extra factor $\sqrt{W(y)}$ belongs only to the profile interface.

\begin{lemma}[Carleman bound]\label{c:lem:carleman}
The kernel $(s+t)^{-1}$ on the positive half-line has norm at most $\pi$, as does every rectangular restriction.
\end{lemma}
\begin{proof}
For the Schur weight $w(s)=s^{-1/2}$,
\[
 \int_0^\infty\frac{t^{-1/2}}{s+t}\,dt
 =2s^{-1/2}\int_0^\infty\frac{dr}{1+r^2}=\pi s^{-1/2}.
\]
The transposed integral is identical. Weighted Cauchy--Schwarz proves the assertion on compact tests away from zero, and density completes it.
\end{proof}

\subsection{Four weighted prime channels and all their overlaps}\label{c:sec:weightedprime}
For the folded adjoint of the profile interface the branches are
\begin{align*}
 (P_nv)(y)&=a_n\sqrt{W(y)W(s_n+y)}v(s_n+y)\one_{s_n+y>{B_0}},\\
 (M_nv)(y)&=a_n\sqrt{W(y)W(s_n-y)}v(s_n-y)\one_{s_n-y>{B_0}}.
\end{align*}
If $y-s_n>0$, it is still inside; a negative exterior input gives exactly the reflected-minus branch. Retain $H_p=P_2+M_3+M_4+M_5$. Its domains, up to null endpoints, are
\begin{center}\small
\begin{tabular}{@{}lll@{}}\toprule
Channel & Inside output $y$ & Exterior input $x$\\\midrule
$P_2$ & $({B_0}-\log2,{B_0})$ & $({B_0},{B_0}+\log2)$\\
$M_3$ & $(0,\log3-{B_0})$ & $({B_0},\log3)$\\
$M_4$ & $(0,\log4-{B_0})$ & $({B_0},\log4)$\\
$M_5$ & $(0,{B_0})$ & $(\log5-{B_0},\log5)$\\\bottomrule
\end{tabular}
\end{center}
Here $M_2=0$ and $a_6=0$. Appendix~\ref{c:app:scalar} gives
\begin{equation}\label{c:eq:coefficientcaps}
 a_2<99/200,\quad a_3<16/25,\quad a_4<7/20,\quad a_5<29/40,
 \qquad a_n\le1.
\end{equation}
In particular $a_4=(\log2)/2$, not $(\log4)/2$.

The four channel-square weights, in units of $A$, are bounded by
$21/50,1/4,3/40,1/100$. For $P_2$, $y>{B_0}-\log2>1/10$.
For $M_3$ split at $y=1/4$: below it $x-{B_0}>6/125$ gives
$W(x)W(y)<Ae^{-6/5}/2<A/4$.
For $M_4$ split at $9/20$: below it $x-{B_0}>17/125$ gives $Ae^{-17/5}/2<A/40$.
For $M_5$ split at $3/5$: below it $x-{B_0}>209/1000$ gives $A/280<A/100$.
Above each split use \eqref{c:eq:pointW}--\eqref{c:eq:pointkappa}.
The common input diagonal is at most
\begin{align}
 A\bigl[(99/200)^2(21/50)&+(16/25)^2/4
 +(7/20)^2(3/40)+(29/40)^2/100\bigr]\notag\\
 &=A\frac{879017}{4000000}<\frac{11A}{50}.\label{c:eq:weighteddiag}
\end{align}
No output orthogonality is assumed.

A cross operator has the form $Cv(x)=\one_D(x)w(x)v(\phi(x))$,
where $\phi$ is injective and $|\phi'|=1$, so $\norm C\le\sup_D|w|$.
The pair $C+C^*$ costs twice this bound, even when $D$ is not reflection invariant.
Every pair is listed below; the last column bounds $\sqrt{W(x)W(x')}/A$.
\begin{center}\small
\begin{tabular}{@{}llll@{}}\toprule
Pair & First input domain & Second input $x'$ & Cap\\\midrule
$P_2,M_3$ & $({B_0},\log6-{B_0})$ & $\log6-x$ & $1/10$\\
$P_2,M_4$ & $({B_0},\log8-{B_0})$ & $\log8-x$ & $1/300$\\
$P_2,M_5$ & $({B_0},{B_0}+\log2)$ & $\log10-x$ & $1/6000$\\
$M_3,M_4$ & $({B_0},\log3)$ & $x+\log(4/3)$ & $1/32$\\
$M_3,M_5$ & $({B_0},\log3)$ & $x+\log(5/3)$ & $1/500$\\
$M_4,M_5$ & $({B_0},\log4)$ & $x+\log(5/4)$ & $1/15$\\\bottomrule
\end{tabular}
\end{center}
These rational comparisons follow from the envelope and are proved in Appendix~\ref{c:app:scalar}. In the first three rows use $W(y)<21/50$, and in the next two use $W(y)<1/2$.
In the last row the sharper full product is below $A/200$:
for $y\ge9/20$ use $(3/40)(1/15)$; otherwise the sum of exterior distances is
$\log20-2{B_0}-2y>495/1000$, its exponent exceeds $6$, and the product is below $A/800$.
The six paired costs in units of $A$ are at most
\[
 27/1000,\quad1/2000,\quad3/50000,\quad7/1000,\quad1/1000,\quad13/5000.
\]
Their sum is $1908/50000<1/25$. Thus $\norm{H_p}^2<(13/50)A<(51/100)^2A$.

The exact complement is $\sum_{n\ge3}P_n+\sum_{n\ge7}M_n$.
From $e^{B_0}<9/4$, $e^{2{B_0}}<5$ and $25/2>12$,
\begin{gather}
 \norm{P_n}\le\sqrt{A/2}(9/(4n))^{12},\qquad
 \norm{M_n}\le\sqrt{A/2}(5/n)^{12},\notag\\
 \sum_{n\ge3}(9/(4n))^{12}\le\frac{14}{11}(3/4)^{12}<1/24,\qquad
 \sum_{n\ge7}(5/n)^{12}\le\frac{18}{11}(5/7)^{12}<3/100.\label{c:eq:bothtails}
\end{gather}
The last comparisons are $178564176<184549376$ and
$439453125000<456762477633$. Both complete tails converge in operator norm.
Their triangle inequalities pay all head--tail and tail--tail products.
Since $\sqrt2>7/5$,
\begin{equation}\label{c:eq:weightedprimenorm}
 \norm{\mathcal J_p}<\sqrt A\left(\frac{51}{100}+\frac{43}{840}\right)
 <\frac9{16}\sqrt A<\frac3{512}<\frac{59}{10000}.
\end{equation}
The scalar slack before multiplication by $\sqrt A$ is $11/8400$.

\subsection{The exact mean repair and six Gamma strips}\label{c:sec:meanrepair}
For $x>{B_0}$ and $0<y<{B_0}$ define
\begin{equation}\label{c:eq:remainderkernel}
 B_\Gamma(x,y)=\rho(x-y)+\rho(x+y)-2c(y)\rho(x).
\end{equation}
The full positive series for $\rho$ gives
\[
 B_\Gamma(x,y)=2\sum_{j\ge1}e^{-(2j+1/2)x}
 [\cosh((2j+1/2)y)-\cosh(y/2)]\ge0.
\]
It decreases in $x$ and increases in $y$. Only the $j=0$ term cancels.
Let $\mathcal T_\Gamma$ have kernel $\sqrt{W(x)W(y)}B_\Gamma(x,y)$ and set
\[
 e(y)=U_+ 1(y)=2\sqrt2c(y)\sqrt{W(y)},\qquad
 z(x)=\frac{\sqrt{W(x)}\rho(x)}{\sqrt2}.
\]
For an inside profile $g$ the exact identity is
\begin{equation}\label{c:eq:rankrepair}
 U_+ \mathcal J_\Gamma g=\mathcal T_\Gamma U_+ g+z\mu(g),\qquad
 \mu(g)=\ip{U_+ g}{e}.
\end{equation}
Thus the rank-one term has not been discarded.

For $y\le b<{B_0}$, monotonicity and \eqref{c:eq:rhoelementary} give
\[
 0\le B_\Gamma(x,y)\le B_\Gamma({B_0},b)
 <\frac{{B_0}}{{B_0}^2-b^2}-\frac76=:F(b).
\]
The subtraction is made at $x={B_0}$, not by imposing a false lower bound for $\rho(x)$ in the exterior.
Divide the input at $2/5,1/2,3/5,7/10,3/4$.
The first five strips give valid Hilbert--Schmidt estimates:
\begin{center}\small
\begin{tabular}{@{}llll@{}}\toprule
Strip & $\int W(y)\,dy$ cap & $F(b)$ & Squared norm cap\\\midrule
$(0,2/5)$ & $1/8$ & $1/2$ & $1/7372800$\\
$(2/5,1/2)$ & $1/80$ & $23/26$ & $529/12460032000$\\
$(1/2,3/5)$ & $7/2000$ & $71/42$ & $5041/116121600000$\\
$(3/5,7/10)$ & $1/1000$ & $25/6$ & $625/8294400000$\\
$(7/10,3/4)$ & $3/40000$ & $1703/186$ & $2900209/106278912000000$\\\bottomrule
\end{tabular}
\end{center}
Each row uses the exterior integral $1/230400$.
The squared norms are respectively below $(140,43,44,76,28)/10^9$.

On the touching strip $(3/4,4/5)$, $W(y)<1/2000$ and
\[
 B_\Gamma(x,y)\le\rho(x-y)+\rho(x+y)
 \le\frac1{2(x-y)}+\frac{51}{62}.
\]
The singular term is half a Carleman kernel after $s=x-{B_0}$, $t={B_0}-y$.
Its norm is below
$\frac\pi2(9216\cdot2000)^{-1/2}<37/100000$.
Its squared cap is $1/7372800$, and $1369\cdot7372800>10^{10}$.
The bounded term has norm at most
\[
 \frac{51}{62}(230400\cdot40000)^{-1/2}
 =\frac{17}{1984000}<\frac9{1000000}.
\]
Thus the last strip norm is below $379/10^6<19/50000$, whose square is below $145/10^9$.
No Hilbert--Schmidt claim is made for its touching singularity.
The input projections are orthogonal, giving
$\mathcal T_\Gamma \mathcal T_\Gamma^*=\sum_iT_iT_i^*$, without output orthogonality.
Consequently
\begin{equation}\label{c:eq:gammaremainnorm}
 \norm{\mathcal T_\Gamma}^2<\frac{140+43+44+76+28+145}{10^9}
 =\frac{476}{10^9}<\left(\frac7{10000}\right)^2.
\end{equation}
Both physical boundaries are already included in the even fold.

\subsection{The global comparison from local sign}\label{c:sec:firstproof}
For an exterior profile $h$, the physical bridge gives
\[
 \E[h]=\tfrac12\norm h_\mu^2+\tg[\kappa h]
 -\cg\norm{\kappa h}_2^2-p[\kappa h].
\]
The full-prime signed identity in Appendix~\ref{c:sec:signed} is analytic and needs neither response data nor an exterior moment constraint. With $u=\kappa h$ it gives
\begin{equation}\label{c:eq:exteriorenergy}
 \E[h]\ge\left(\tfrac12-\frac{31}{46080}\right)\norm h_\mu^2
 +\tg[\kappa h]+D_p(\kappa h/\alpha),
\end{equation}
where $D_p$ is the complete nonnegative prime slack defined in \eqref{c:eq:slack}.
Thus the cap $31/46080$, not the earlier coarser $7/10240$, is available in the local-sign branch.

For a real centered profile $v=g+h$ cut on $I$ and $O$, one has
$\mu(g)=-\mu(h)$, not $\mu(g)=0$. Put
\[
 H=U_+ h,\qquad e_O=U_+ \one_O=2\sqrt2c\alpha\big|_{({B_0},\infty)},\qquad
 p_*:=\norm{e_O}^2=\mu(O).
\]
Then $\mu(h)=\ip H{e_O}$.
The removed term in \eqref{c:eq:rankrepair} enters the energy with sign
\[
 -2\mu(g)\ip zH=2\mu(h)\ip zH.
\]
This product of distinct exterior functionals is not assigned a sign on its own. Local positivity simultaneously supplies the penalty
$-\tfrac12|\mu(g)|^2=-\tfrac12|\mu(h)|^2$.

\begin{lemma}[The combined mean form]\label{c:lem:coupledmean}
Let $e_O,z$ be real vectors in the folded exterior space with
$0\le z\le e_O/4$ pointwise and $p_*=\norm{e_O}^2$.
For every real exterior vector $H$,
\begin{equation}\label{c:eq:combinedmean}
 -\tfrac12|\ip H{e_O}|^2+2\ip H{e_O}\ip zH
 \ge-\tfrac12p_*\norm H^2.
\end{equation}
More precisely, the optimal nonnegative loss for this form is
$\norm{e_O}\norm w-\ip{e_O}w$, where $w=z-e_O/4$.
\end{lemma}
\begin{proof}
With $(a\otimes b)H=\ip Hb a$ on the real space, the combined operator is
\[
 K=e_O\otimes z+z\otimes e_O-\tfrac12e_O\otimes e_O
   =e_O\otimes w+w\otimes e_O.
\]
Its two possible nonzero eigenvalues are
$\ip{e_O}w\pm\norm{e_O}\norm w$; the formula includes dependent vectors by continuity and the remaining spectrum is zero.
Since $-e_O/4\le w\le0$, one has
$\norm w\le\norm{e_O}/4$ and $\ip{e_O}w\ge-p_*/4$.
Consequently $K\succeq-p_*I/2$, proving the inequality and the stated exact loss.
\end{proof}
Here the pointwise hypothesis is automatic:
\[
 \frac{z(x)}{e_O(x)}=\frac{\rho(x)}{4c(x)}<\frac14\qquad(x>{B_0}),
\]
since $\rho$ decreases, $\rho({B_0})<9/10$ and $c(x)\ge1$.
The estimate removes the need for the separate absolute rank charge; it does not discard the local mean penalty.

Local sign gives $\E[g]\ge\tfrac12\norm g_\mu^2-\tfrac12|\mu(g)|^2$.
Use \eqref{c:eq:exteriorenergy}, Lemma~\ref{c:lem:coupledmean},
the complete weighted prime bound $\norm{\mathcal J_p}<3/512$, and
$\norm{\mathcal T_\Gamma}<7/10000$ to obtain
\begin{equation}\label{c:eq:twobytwo}
 \E[v]\ge\tfrac12\norm g_\mu^2+(\tfrac12-a_1)\norm h_\mu^2
 -2L_1\norm g_\mu\norm h_\mu,
\end{equation}
where
\begin{equation}\label{c:eq:improvedconstants}
 a_1=\frac{31}{46080}+\frac3{98000}=\frac{39703}{56448000},
 \qquad L_1=\frac3{512}+\frac7{10000}=\frac{2099}{320000}.
\end{equation}
Only $p_*<3/49000$ was used for the combined mean loss. The nonnegative terms in \eqref{c:eq:exteriorenergy} were dropped only after the signed identity had been applied.

For $\delta_1=693/100000$, the residual comparison matrix is
\[
 M_{\delta_1}=\begin{pmatrix}\delta_1&-L_1\\-L_1&\delta_1-a_1\end{pmatrix}.
\]
The exact principal-minor checks are
\begin{equation}\label{c:eq:margin}
 \delta_1-a_1=\frac{8787041}{1411200000}>0,\qquad
 \det M_{\delta_1}=\frac{2244433}{17920000000000}>0.
\end{equation}
For the determinant, on denominator $1128960000000000$, the numerators of
$\delta_1^2$, $\delta_1a_1$ and $L_1^2$ are respectively
$54218191104$, $5502835800$ and $48573956025$.
Their signed difference is $141399279=63\cdot2244433$.
Thus \eqref{c:eq:twobytwo} proves the coefficient
$1/2-\delta_1=49307/100000$.
More generally $M_\delta\succeq0$ precisely when
\[
 \delta\ge\frac12\left(a_1+\sqrt{a_1^2+4L_1^2}\right),
\]
which gives \eqref{c:eq:lambdaendpoint}. This threshold is strictly positive.
It is an exact endpoint for this comparison, not an assertion of the optimal constant of the actual energy, and not a family of signs tending to zero loss.

First apply the argument to $v=f-\mu(f)$ with $f$ even compact smooth.
Constants and sharp cuts belong to the corresponding minimum domains by Appendix~\ref{c:app:cuts}.
The bounded interface operators and norm-convergent tails justify the cross identity on these cuts.
Finally use the global smooth form core, along which energy and variance converge.
This proves Theorem~\ref{c:thm:localglobal} on the full maximal domain.
Real and imaginary parts give the complex statement without assigning a sign to the rank-one cross term separately.

\subsection{A joint Gram estimate for the physical source}\label{c:sec:joint}
The output coordinate in this section is physical $L^2(I)$, not weighted profile space.
For $v(x)=\sqrt2u(x)/\alpha(x)$ on $x>{B_0}$ the folded Gamma source is
\[
 (\mathcal T_\Gamma^{\rm phys}v)(y)=\int_{B_0}^\infty
 [\rho(x-y)+\rho(x+y)]\alpha(x)v(x)\,dx,\qquad0<y<{B_0}.
\]
There is no further factor $\sqrt2$.
Because $x+y>{B_0}$, \eqref{c:eq:rhoelementary} implies the positive majorant
\[
 0\le\rho(x-y)+\rho(x+y)\le\frac1{2(x-y)}+b,\qquad b=\frac{23}{20}.
\]
Define
\[
 (Cv)(y)=\int_{B_0}^\infty\frac{\alpha(x)v(x)}{2(x-y)}\,dx,\qquad
 (Rv)(y)=b\int_{B_0}^\infty\alpha(x)v(x)\,dx.
\]
Pointwise $|\mathcal T_\Gamma^{\rm phys}v|\le(C+R)|v|$.
This is domination, not equality with the actual Gamma source.
The touching kernel is half a Carleman kernel, hence
\[
 \norm C\le\frac\pi2\sqrt A.
\]
Set
\[
 I_0=\int_{B_0}^\infty\alpha(x)^2\,dx,\qquad
 I_2=\int_{B_0}^\infty\alpha(x)^2\log^2\frac{x}{x-{B_0}}\,dx.
\]
Then $I_0\le A/25$ and
\begin{equation}\label{c:eq:regulargram}
 \norm R^2=b^2B_0I_0\le\frac{529}{12500}A,\qquad
 \norm{C^*R}=\frac b2\sqrt{I_2I_0}.
\end{equation}
The second identity follows from the exact rank-one formula
\[
 (C^*Rv)(x)=\frac b2\alpha(x)\log\frac{x}{x-{B_0}}
                  \int_{B_0}^\infty\alpha(t)v(t)\,dt.
\]

\subsubsection{The logarithmic cross integral}
With $t=25(x-{B_0})$,
\[
 I_2/A\le\frac1{25}\int_0^\infty e^{-t}\log^2(1+20/t)\,dt.
\]
Take $a=61/20>\log21$, since $e^3>20$ and $e^{1/20}>21/20$.
Writing $r(t)=\max\{-\log t,0\}$, one has
$\log(1+20/t)\le a+r(t)$.
The inequality $e^{-t}\le1-t+t^2/2$ on $[0,1]$ and
\[
 \int_0^1t^k(-\log t)\,dt=(k+1)^{-2},\qquad
 \int_0^1t^k\log^2t\,dt=2(k+1)^{-3}
\]
give
\[
 \int_0^1e^{-t}(-\log t)\,dt\le29/36,\qquad
 \int_0^1e^{-t}\log^2t\,dt\le193/108.
\]
Consequently
\begin{equation}\label{c:eq:I2}
 I_2/A\le\frac1{25}\left(\frac{3721}{400}+
       \frac{1769}{360}+\frac{193}{108}\right)
 =\frac{172837}{270000}<\left(\frac{81}{100}\right)^2.
\end{equation}
Thus $2\norm{C^*R}<(1863/10000)A$ and
\begin{equation}\label{c:eq:gammajoint}
 \norm{J_\Gamma}_{\wo\to L^2(I)}^2/A
 <\frac{121}{49}+\frac{529}{12500}+\frac{1863}{10000}
 <\frac{529}{196}.
\end{equation}
The last margin is $9524/9800000$.
This proves the standalone bound
$\norm{J_\Gamma}_{\wo\to L^2(I)}<(23/14)\sqrt A$.
The calculation concerns $C^*R$, not the square integral of a touching Gamma kernel.

\subsubsection{The physical prime head and complete tails}
In the same folded spaces the physical channels are
\[
 (P_nv)(y)=a_n\alpha(s_n+y)v(s_n+y)\one_{s_n+y>{B_0}},\quad
 (M_nv)(y)=a_n\alpha(s_n-y)v(s_n-y)\one_{s_n-y>{B_0}}.
\]
In this section $H_p=P_2+M_3+M_4+M_5$ refers to these physical channels.
They have the same domains and pair maps as in Section~\ref{c:sec:weightedprime}, but without $\sqrt{W(y)}$.
Their common input diagonal is at most $(422091/400000)A$.
Indeed below $x_5=\log5-{B_0}$ the square sum is at most $(6217/8000)A$;
above $x_5$, $\alpha(x)/\sqrt A<e^{-9/80}<9/10$, giving
\[
 D(x)/A\le\frac{81}{100}\bigl[(99/200)^2+(16/25)^2
 +(7/20)^2+(29/40)^2\bigr]=\frac{422091}{400000}.
\]
The six cross products use the last column of the table in Section~\ref{c:sec:weightedprime}.
Their paired sum is
\[
 \frac{198}{3125}+\frac{231}{200000}+\frac{957}{8000000}
 +\frac7{500}+\frac{116}{62500}+\frac{203}{6000}
 =\frac{109751}{960000}.
\]
Hence
\begin{equation}\label{c:eq:physicalhead}
 \norm{H_p}^2\le\frac{5613847}{4800000}A<\frac{169}{144}A.
\end{equation}
The exact complement $T_p=\sum_{n\ge3}P_n+\sum_{n\ge7}M_n$ obeys
\begin{equation}\label{c:eq:physicaltail}
 \norm{T_p}<\left(\frac1{24}+\frac3{100}\right)\sqrt A
\end{equation}
by the same complete integer tails \eqref{c:eq:bothtails}, now without $1/\sqrt2$.
Thus no prime power or head--tail cross term is omitted.

\subsubsection{Every singular Gamma and prime-head product}
Let $\sigma=25/2$. Direct substitution gives the \emph{input}-Gram kernels
\begin{align}
 (C^*P_n)(z,t)&=\frac{a_n\alpha(z)\alpha(t)}
                      {2(z-t+s_n)},&
 z>{B_0},\quad\max({B_0},s_n)<t<{B_0}+s_n,\label{c:eq:inputgramplus}\\
 (C^*M_n)(z,t)&=\frac{a_n\alpha(z)\alpha(t)}
                      {2(z+t-s_n)},&
 z>{B_0},\quad\max({B_0},s_n-{B_0})<t<s_n.\label{c:eq:inputgramminus}
\end{align}
They vanish outside the indicated intervals. They are not the output-Gram products $CP_n^*$ or $CM_n^*$.

For $P_2$, put $\ell=\log2$, $s=z-{B_0}$, $r=t-{B_0}\in(0,\ell)$.
The kernel is bounded by
$(a_2A/2)e^{-\sigma(s+r)}/(s+\ell-r)$.
On $r\le\ell/2$ the denominator is at least $\ell/2$, so an actual Hilbert--Schmidt majorant gives
\[
 \norm{(C^*P_2)_{\rm low}}/A
 \le\frac{a_2}{2\sigma\ell}<\frac{297}{10000}.
\]
On $r>\ell/2$, set $h=\ell-r$. The denominator is $s+h$, so use Carleman, retaining the factor $1/2$:
\[
 \norm{(C^*P_2)_{\rm high}}/A
 \le\frac{a_2\pi}{2}\,2^{-25/4}<\frac{1089}{89600}.
\]
This corner $(z,t)=({B_0},{B_0}+\log2)$ is not Hilbert--Schmidt.

For $M_3,M_4$, let $d_n=2{B_0}-s_n$. The bounds
$d_3>1/2$, $d_4>1/5$ genuinely separate their kernels.
The exponential square integrals give
\[
 \norm{C^*M_n}/A\le\frac{a_n}{4\sigma d_n},\qquad
 \norm{C^*M_3}/A<16/625,\quad \norm{C^*M_4}/A<7/200.
\]
For $M_5$, however, set $d=s_5-2{B_0}>0$,
$s=z-{B_0}$ and $h=t-(s_5-{B_0})$.
The kernel majorant is
\[
 \frac{a_5A}{2}e^{-\sigma d}\frac{e^{-\sigma(s+h)}}{s+h}.
\]
Although its two input locations are separated, the corner is touching.
Carleman, not a denominator gap or Hilbert--Schmidt estimate, gives
$\norm{C^*M_5}/A\le a_5\pi/2<319/280$.
All four products are therefore paid by
\begin{equation}\label{c:eq:CHp}
 \norm{C^*H_p}/A<
 \frac{297}{10000}+\frac{1089}{89600}+\frac{16}{625}
 +\frac7{200}+\frac{319}{280}
 =\frac{2781497}{2240000}<\frac54.
\end{equation}

\subsubsection{The regular-majorant product and the full source}
On the common exterior input coordinate,
\begin{align*}
 H_p^*1(t)&=\alpha(t)w(t),\\
 w(t)&=a_2\one_{({B_0},{B_0}+\log2)}(t)+a_3\one_{({B_0},\log3)}(t)
 +a_4\one_{({B_0},\log4)}(t)+a_5\one_{(\log5-{B_0},\log5)}(t).
\end{align*}
Since $w<221/100$ and $R=b\,1\otimes\alpha$, the exact rank-one product satisfies
\begin{equation}\label{c:eq:HpR}
 H_p^*R=b(H_p^*1)\otimes\alpha,\qquad
 2\norm{H_p^*R}\le2b(221/100)I_0\le\frac{5083}{25000}A.
\end{equation}
Expanding the bounded-operator Gram of $C+R+H_p$ now charges all three cross pairs:
\begin{align}
 \norm{C+R+H_p}^2/A
 &<\frac{121}{49}+\frac{169}{144}+\frac52+
   \frac{529}{12500}+\frac{1863}{10000}+\frac{5083}{25000}\notag\\
 &=\frac{121}{49}+\frac{169}{144}+\frac52+\frac{21597}{50000}
 <\frac{324}{49}.\label{c:eq:jointgram}
\end{align}
For the last comparison use $21597/50000<9/20$ and
$169/144+59/20=2969/720<203/49$, whose cross products are
$145481<146160$.
Positive-kernel domination gives
$|Jv|\le(C+R+H_p+T_p)|v|$.
Combining \eqref{c:eq:physicaltail} and \eqref{c:eq:jointgram} proves
\begin{theorem}[Complete analytic source bound]\label{c:thm:source}
The physical source extends from compact exterior tests to every even finite-$\wo$ input and satisfies
\begin{equation}\label{c:eq:sourcej}
 j:=\norm J_{\wo\to L^2(I)}
 <\left(\frac{18}{7}+\frac1{24}+\frac3{100}\right)\sqrt A
 =\frac{11101}{4200}\sqrt A<\frac{11101}{403200}<\frac{69}{2500}.
\end{equation}
This theorem needs neither local sign nor response enclosures.
\end{theorem}
The extension follows by density and boundedness; both full prime tails converge in operator norm. The domination statement is not a Loewner order between rectangular operators.

\subsection{The additional response contracts and inverse cost}\label{c:sec:response}
Let $e_0=(2{B_0})^{-1/2}$ and
$e_n={B_0}^{-1/2}\cos(n\pi x/{B_0})$, $n\ge1$, on $I$.
Let $E:\R^{64}\to L^2_{\rm even}(I)$ insert $e_0,\ldots,e_{63}$ and set
\[
 H_{64}=L_I+9EE^*,\qquad \mathsf U_{64}=H_{64}^{-1}E.
\]
Appendix~\ref{c:sec:lift} proves analytically that
\begin{equation}\label{c:eq:Hcoercivity}
 H_{64}\succeq61I/28,\qquad H_{64}^{-1}\preceq D_{\rm old},
\end{equation}
where $D_{\rm old}$ is the explicit cosine-diagonal majorant defined there.
These facts are not numerical enclosures.

Choose fixed real trial coefficients $C_{\rm trial}\in\R^{1024\times64}$,
defining columns $\widetilde{\mathsf U}_{64}$ in the first 1024 local cosine modes.
These columns belong to the local operator domain: their killed Gamma row has logarithmic endpoint singularities in $L^2$, and their remaining rows are bounded.
Set
\[
 R_{64}=E-H_{64}\widetilde{\mathsf U}_{64},\qquad
 F=I_{64}/9-E^*H_{64}^{-1}E.
\]
The regular-majorant operator $R$ in Section~\ref{c:sec:joint} is unrelated to $R_{64}$.

\begin{hypothesis}[Full residual and centered-response contracts]\label{c:hyp:response}
For the fixed trial matrix there are symmetric matrices $F_*,V_{\rm up}$,
and its positive-diagonal Cholesky factor $M$ with $F_*=MM^*$, such that:
\begin{enumerate}[label=\textup{(\alph*)}]
\item The actual complete-operator inequalities hold:
\begin{equation}\label{c:eq:contracta}
 0\prec F_*\preceq F,\qquad
 R_{64}^*H_{64}^{-1}R_{64}
 \preceq R_{64}^*D_{\rm old}R_{64}
 \preceq V_{\rm up}\preceq F_*/100.
\end{equation}
All residual rows, including those beyond 1024 and the infinite remainder, are included.
\item Set $V=\widetilde{\mathsf U}_{64} M^{-*}$ and $Z=C_{\rm trial}M^{-*}$.
For one fixed rational centering vector $a\in\R^{64}$, the exact finite-row, spatial-tail, Gamma-tail, and prime-tail quantities specified in Appendix~\ref{c:app:rows} have enclosures satisfying
\begin{equation}\label{c:eq:contractb}
 \sqrt{I_{\rm fin}}+e_X<\frac{81}{2500},\qquad
 e_\Gamma<\frac{39}{10000},\qquad
 b_{16}\norm{Z_{<64}}_{\mathrm{HS}}<10^{-21},\qquad
 \norm{Z_{64:1023}}_{\mathrm{HS}}<\frac13.
\end{equation}
Here $b_{16}=500/20^{25}$ bounds the entire omitted prime source after 16.
The rows and their centering refer to the same $V$, not to independently chosen response vectors.
\end{enumerate}
\end{hypothesis}
The theorem is valid for any fixed data satisfying these precise contracts.
An accepted finite matrix sign or a finite replay alone does not establish their full analytic enclosure interpretation.
Neither \eqref{c:eq:contracta} nor \eqref{c:eq:contractb} is inferred from Hypothesis~\ref{c:hyp:local}.
No new enclosure run is asserted.

To see local coercivity from (a), put $f_*=\lambda_{\min}(F_*)>0$ and
$K_0=3H_{64}^{-1/2}E$. Then
$K_0^*K_0=I-9F\preceq(1-9f_*)I$, so the same upper bound holds for
$K_0K_0^*$. Therefore
\begin{equation}\label{c:eq:Lcoercive}
 L_I=H_{64}^{1/2}(I-K_0K_0^*)H_{64}^{1/2}
 \succeq9f_*H_{64}\succeq(549f_*/28)I>0.
\end{equation}
The mean pole remains in $L_I$ throughout.

For $m_O(u)=\int_Ocu=0$, the fixed centered-row decomposition in Appendix~\ref{c:app:rows} gives
\begin{equation}\label{c:eq:responsebound}
 \norm{V^*Ju}\le r\sqrt{\wo(u)},\qquad
 r<\frac{363}{10000}+\frac j3+10^{-21}<\frac{91}{2000}.
\end{equation}
It is essential to insert the sharper fraction from \eqref{c:eq:sourcej} before rounding:
\[
 \frac{91}{2000}-\frac{363}{10000}-\frac{11101}{1209600}
 =\frac{683}{30240000}>10^{-21}.
\]
The coarser bound $j<69/2500$ alone leaves no displayed room for this last tail.

Woodbury gives the actual inverse
\[
 T=L_I^{-1}=H_{64}^{-1}+\mathsf U_{64}F^{-1}\mathsf U_{64}^*.
\]
For $s\in L^2_{\rm even}(I)$ set $h=\ip s{H_{64}^{-1}s}$.
Since $\mathsf U_{64}=\widetilde{\mathsf U}_{64}+H_{64}^{-1}R_{64}$, inverse order and (a) imply
\[
 \ip s{Ts}\le h+
 \left[\norm{V^*s}+\frac1{10}\sqrt h\right]^2.
\]
Indeed
$\norm{M^{-1}R_{64}^*H_{64}^{-1}s}\le\sqrt h/10$
follows by conjugating \eqref{c:eq:contracta}.
Young's inequality with parameter $1/20$ yields
\[
 \ip s{Ts}\le\frac{121}{100}h+\frac{21}{20}\norm{V^*s}^2.
\]
For $s=Ju$, $h\le(28/61)j^2\wo(u)$.
Using \eqref{c:eq:responsebound} and then $j<69/2500$,
\begin{align}
 \frac{\ip{Ju}{TJu}}{\wo(u)}
 &<\frac{847}{1525}\left(\frac{69}{2500}\right)^2
 +\frac{21}{20}\left(\frac{91}{2000}\right)^2\notag\\
 &=\frac{4032567}{9531250000}+\frac{173901}{80000000}
 <\frac{17}{40000}+\frac{87}{40000}=\frac{13}{5000}.\label{c:eq:inversecost}
\end{align}
For the first comparison use $847/1525<5/9$ and
$17/40000-4761/11250000=81/45000000>0$; the second uses $173901<174000$.
Including $u=0$, the universal statement is
$\ip{Ju}{L_I^{-1}Ju}\le(13/5000)\wo(u)$ on centered exterior inputs.

\subsection{The signed exterior form and the conditional global bound}\label{c:sec:conditionalproof}
Let $\D_O^0$ be the minimum closure of even $C_c^\infty(O)$ with
$m_O(u)=0$ in $(\tg+\wo)^{1/2}$.
For $u\in\D_O^0$ define
\[
 S[u]=\tg[u]-\cg\norm u_2^2-p[u]-\ip{Ju}{L_I^{-1}Ju}.
\]
This is the actual signed Schur form of the mean-inclusive local operator.
Appendix~\ref{c:sec:signed} proves, without a numerical premise,
\begin{equation}\label{c:eq:signedcapmain}
 \cg\norm u_2^2+p[u]\le\frac{31}{46080}\wo(u)-D_p(u/\alpha),
\end{equation}
where the full nonnegative slack is
\[
 D_p(g)=\sum_{n\ge2}a_n\int_{\{x,x+s_n\in O\}}
  \alpha(x)\alpha(x+s_n)|g(x+s_n)-g(x)|^2\,dx.
\]
Every prime power and both row orientations are retained.
Combining \eqref{c:eq:inversecost} and \eqref{c:eq:signedcapmain},
\begin{equation}\label{c:eq:Slower}
 S[u]\ge\tg[u]+D_p(u/\alpha)-\varepsilon_8\wo(u),\qquad
 \varepsilon_8=\frac{31}{46080}+\frac{13}{5000}
 =\frac{18851}{5760000}.
\end{equation}

For $f\in\D$ use the exterior-mean gauge
\[
 \beta_O=\mu(O)^{-1}\int_Of\,d\mu,\qquad
 u=\kappa\one_O(f-\beta_O),\quad
 v=\kappa\one_I(f-\beta_O).
\]
Appendix~\ref{c:app:cuts} puts $u$ in $\D_O^0$ and $v$ in $\D_I$.
Moreover
\[
 \wo(u)=\int_O|f-\beta_O|^2\,d\mu
 \le\int_O|f-\mu(f)|^2\,d\mu\le\Var_\mu(f).
\]
The full bridge and local square completion give
\begin{equation}\label{c:eq:squaremain}
 \E[f]-\tfrac12\Var_\mu(f)=q_I[v-L_I^{-1}Ju]+S[u].
\end{equation}
The local square is nonnegative by \eqref{c:eq:Lcoercive}.
Dropping only that square, $\tg$ and $D_p$ in \eqref{c:eq:Slower} proves
Theorem~\ref{c:thm:conditionalglobal}.
The exact comparison is
\[
 \frac{2861149}{5760000}-\frac{71501077}{144000000}
 =\frac{24}{125000}>0.
\]
The full-domain claim uses the proved equality of minimum and maximal energy domains, not a change of realization.

\subsection{The remaining sign and the absence of a uniform weighted gap}\label{c:sec:limits}
The exact remaining target is a spatially faithful inverse comparison, not a uniform estimate obtained by replacing its two costs by constants. Assume here that $L_I$ is strictly positive, as supplied by the accepted stronger local theorem $q_I\ge10^{-22}\norm{\cdot}_2^2$ or by Hypothesis~\ref{c:hyp:response}. Then the global half bound on $\D$ is equivalent to
\begin{equation}\label{c:eq:remainingS}
 S[u]\ge0\quad\hbox{for every }u\in\D_O^0.
\end{equation}
Sufficiency follows from \eqref{c:eq:squaremain}.
For necessity, take any $u\in\D_O^0$, set $v=L_I^{-1}Ju\in\D_I$ and
$f=(u+v)/\kappa$. The domain identification in Section~\ref{a:sec:domain} gives $f\in\D$; its exterior mean is zero and its local square vanishes, so its half-defect is $S[u]$.
If one assumes only $L_I\succeq0$, an ordinary inverse requires an additional coercivity or generalized-inverse and range argument.

By \eqref{c:eq:signedq0}, the target \eqref{c:eq:remainingS} is equivalently
\begin{equation}\label{c:eq:spatialtarget}
 \int_O(\cg W+V_p)|u/\alpha|^2\,dx+\ip{Ju}{L_I^{-1}Ju}
 \le\tg[u]+D_p(u/\alpha),\qquad u\in\D_O^0.
\end{equation}
Here $V_p$ is the exact complete row potential \eqref{c:eq:primepotential}.
The following analytic obstruction explains why the spatial potential and actual inverse response in this inequality cannot simply be replaced by the fixed upper cost $\varepsilon_8\wo$ and absorbed by a uniform weighted gap.

\begin{proposition}[No uniform weighted gap for the full positive slack]\label{c:prop:nogap}
On the original centered exterior minimum domain,
\begin{equation}\label{c:eq:nogap}
 \inf_{0\ne u\in\D_O^0}
 \frac{\tg[u]+D_p(u/\alpha)}{\wo(u)}=0.
\end{equation}
The same holds with $D_p$ replaced by its $n=2$ summand $D_2$.
Neither local positivity nor response contracts are required for this proposition.
\end{proposition}
\begin{proof}
Choose a nonnegative $\eta\in C_c^\infty(0,1)$ with ordinary $L^2$ norm one.
For $R>{B_0}$, translate it to $(R,R+1)$ and $(R+2,R+3)$, reflect each evenly to the negative half-line, and divide each reflected pair by $\sqrt2$.
Denote the resulting full-line orthonormal functions by $\phi_{1,R},\phi_{2,R}$.
Their derivative norms are independent of $R$.
Put
\[
 m_i=\int_O c\alpha\phi_{i,R}\,dx>0,\qquad
 g_R=\frac{m_2\phi_{1,R}-m_1\phi_{2,R}}{\sqrt{m_1^2+m_2^2}},
 \qquad u_R=\alpha g_R.
\]
Then $\norm{g_R}_2=1$, $m_O(u_R)=0$, and $\wo(u_R)=1$ exactly.
The coefficient vector has Euclidean norm one, so $\norm{g_R'}_2$ is bounded uniformly even when the moments are small.
Every $u_R$ is itself in the centered compact smooth exterior core.

We justify decay of $\alpha'$ from the positive theta series, not from an upper bound for $W$ alone.
On the positive exterior write
\[
 W=\sum_{j\ge1}w_j,\qquad
 w_j=\frac{(4t_j^2-6t_j)e^{-t_j}}{1+e^{-x}},\qquad
 t_j=\pi j^2e^{2x}>3.
\]
Termwise differentiation gives
\[
 (\log w_j)'=2+\frac{4t_j}{2t_j-3}-2t_j+
 \frac{e^{-x}}{1+e^{-x}},\qquad
 |(\log w_j)'|\le2t_j+7.
\]
Cauchy--Schwarz in this positive series yields
\begin{equation}\label{c:eq:alphaprimedecay}
 |\alpha'|^2=\frac{|W'|^2}{4W}
 \le\frac14\sum_{j\ge1}w_j(2t_j+7)^2.
\end{equation}
The polynomial-times-double-exponential series on the right has a summable majorant at any fixed exterior radius and tends uniformly to zero on escaping tails. The same holds for $\alpha$; evenness supplies the negative half-line.
Consequently
\[
 \norm{u_R}_2\le\sup_{|x|>R}\alpha\longrightarrow0,\qquad
 \norm{u_R'}_2\le\sup_{|x|>R}|\alpha'|
 +\sup_{|x|>R}\alpha\,\norm{g_R'}_2\longrightarrow0.
\]
The full nonlocal Gamma estimate
\[
 \tg[u]\le\norm{u'}_2^2\int_0^1\rho(s)s^2\,ds
 +4\norm u_2^2\int_1^\infty\rho(s)\,ds
\]
proves $\tg[u_R]\to0$.

For completeness, the full prime potential also vanishes uniformly on escaping tails.
For $x>{B_0}$, each existing same-side product in \eqref{c:eq:primepotential} is at most $An^{-25/2}$; the reflected product is at most $Ae^{25{B_0}}n^{-25/2}$.
Since $a_n\le1$, a uniform summable row majorant is
\[
 A(2+e^{25{B_0}})\sum_{n\ge2}n^{-25/2}.
\]
For each fixed shift, the row product tends uniformly to zero as $x\to+\infty$; its reflected branch eventually disappears.
A finite-head and summable-tail split, followed by evenness, therefore gives
$\sup_{|x|>R}V_p(x)\to0$.
The complete edge-square inequality gives
\[
 0\le D_p(g_R)\le2\int_OV_p|g_R|^2\,dx
 \le2\sup_{|x|>R}V_p(x)\longrightarrow0.
\]
All prime powers and both physical half-lines are retained.
Together with $\wo(u_R)=1$, this proves \eqref{c:eq:nogap}.
Finally $0\le D_2\le D_p$ proves its dyadic version.
\end{proof}
The proposition is compatible with the trivial kernel of $D_p$ proved in Appendix~\ref{c:sec:signed}: injectivity is not a positive uniform gap.
It rules out $\tg+D_p\ge\varepsilon\wo$ for every fixed $\varepsilon>0$, hence in particular absorption of $\varepsilon_8\wo$ by that route.
It gives no negative direction of $S$ and does not obstruct the relative estimate \eqref{c:eq:spatialtarget}.

\subsubsection{What a regularized sign would and would not imply}
Let $A_2[u]=\tg[u]+D_2(u/\alpha)$ on the same original centered exterior domain.
Under the additional response contracts, \eqref{c:eq:Slower} implies only
\[
 A_2[u]\le S[u]+\varepsilon_8\wo(u).
\]
It does not imply $S[u]+A_2[u]\ge0$.
If that sign were independently proved on the original domain, averaging the two inequalities would give
$S\ge-(\varepsilon_8/2)\wo$ and the coefficient $5741149/11520000$.
That number has an unproved premise and is not an achieved bound.
More generally a separate sign $S+\eta A_2\ge0$ with $\eta>0$ yields only
\[
 S\ge-\frac{\eta\varepsilon_8}{1+\eta}\wo.
\]
One case $\eta=1$ cannot be iterated to zero loss by reusing the same inequalities.
A family with $\eta\downarrow0$ on the same original domain, with all residual, source-tail and domain errors controlled, would prove \eqref{c:eq:remainingS}; no such family is established here.
A weaker $A_2^{1/2}$ completion would require its own faithful-domain and continuous-extension proofs. The full slack is not extended to a weaker topology by assertion.
A rank-seven lift, if used in a future sign analysis, would be $H_7=L_I+7E_7E_7^*$, not the rank-64 lift $H_{64}=L_I+9EE^*$ used here.
Neither the improved fixed comparison nor the conditional response cost supplies a uniform continuation to zero loss.
The global half bound and the Riemann hypothesis remain unproved.

\subsubsection{A positive loss interval for a congruent global pencil}
For $0<\delta<1/2$ define the profile form on $\D$ by
\[
 Q_\delta[f]=\E[f]-(\tfrac12-\delta)\norm f_\mu^2.
\]
If $\E\ge\lambda\Var_\mu$, then on the centered subspace
\[
 Q_\delta[f]\ge(\lambda-\tfrac12+\delta)\norm f_\mu^2,
 \qquad\mu(f)=0.
\]
Thus Theorem~\ref{c:thm:localglobal} excludes centered negativity and nullity whenever
\begin{equation}\label{c:eq:pencilinterval}
 \frac{693}{100000}<\delta<\frac12.
\end{equation}
Under the response contracts, Theorem~\ref{c:thm:conditionalglobal} also gives
$18851/5760000<\delta<1/2$.
The constant direction has value $Q_\delta[1]=-(1/2-\delta)<0$ and is form-orthogonal to the centered subspace. Hence the profile form has negative index one and zero nullity on either applicable strict interval.

The closed form congruence proved in Sections~\ref{a:sec:pencil} and~\ref{f:sec:spectrum} transfers this negative index and nullity to the actual ordinary-$L^2$ pencil $L_\delta$.
This $L_\delta$ is distinct from the fixed local operator $L_I$ inverted above.
The transfer requires the stated congruence and domain identification; it is not inferred from a finite matrix.
These are sufficient positive-loss intervals. Neither endpoint inequality, a single lifted calculation, nor the location of an essential spectral threshold supplies positivity for every $\delta>0$ or a cofinal certified sequence tending to zero.


\section{Delta scaled finite Schur reductions}\label{sec:schur-unified}
The preceding variance bounds give actual signs on specified parameter intervals. At every positive parameter one can also construct a positive infinite complement and exact scaled response, but the finite sign is then a separate question. The radical interpolation cost must remain in the reference form.


\begin{theorem}[Scaled complement and exact Schur response]\label{a:thm:mainschur}
For every $\delta>0$ and $0<\theta<1$, finitely many explicitly prescribed interior moments define a weighted orthogonal decomposition $X_c=V\oplus Z$ such that
\begin{equation}\label{a:eq:maincomparison}
 (1-\theta)P_\delta[z]\le\ell_\delta[z]
 \le(1+\theta)P_\delta[z],\qquad
 P_\delta[z]=\tg[z]+\delta\norm{z}_{\Hh}^2.
\end{equation}
There is a finite-dimensional radical lift $H:V\to X_c$ satisfying $P_VHv=v$ and $q_0(Hv,z)=0$ for all $z\in X$. Write
\[
 h(v)=Hv,\quad p(v)=h(v)-v\in Z,\quad G[v]=\norm{h(v)}_{\Hh}^2.
\]
If $A_{\delta,Z}$ is the positive weighted operator represented by $\ell_\delta|_Z$, then the exact Schur form is
\begin{equation}\label{a:eq:mainschur}
 S_\delta[v]=\delta G[v]
    -\delta^2\ip{p(v)}{A_{\delta,Z}^{-1}p(v)}_{\Hh}.
\end{equation}
Finite Neumann sums $z_J$ of exact positive reference solves satisfy
\begin{equation}\label{a:eq:mainerror}
 0\le\ell_\delta[h(v)-z_J(v)]-S_\delta[v]
 \le\frac{\delta\theta^{2J+2}}{1-\theta}G[v].
\end{equation}
They also have a weighted residual bound $C_0\theta^J\sqrt{G[v]}$ with $C_0<10$. The norm of $H$ is charged through $G$; no uniform conditioning or weighted completeness is assumed.
\end{theorem}

The proof of Theorem~\ref{a:thm:mainseed} is in Section~\ref{a:sec:seed}. Sections~\ref{a:sec:radical}--\ref{a:sec:response} prove Theorem~\ref{a:thm:mainschur}. On the interval in Theorem~\ref{a:thm:mainseed}, the exact finite form also satisfies
$S_\delta[v]\ge(c_{\log 2}+\delta-1/2)\norm v_{\Hh}^2$.
This finite sign follows from the full-domain variance estimate. It is not a uniform continuation theorem as $\delta$ tends to zero.


\subsection{Finite moment complements for positive parameters}\label{a:sec:moments}

The full correction $d=\cg\ip\cdot\cdot_2+p$ is bounded in the weighted norm, but need not be small globally. Interior moments make its ordinary-$L^2$ part small relative to the Gamma form, while the theta weight controls spatial leakage. The construction below quantifies both effects.

Write
\[
 d(T)=\frac12\log(1+4T^2),\qquad
 \epsilon_m=\frac{2RT}{\pi}\frac{(RT)^{4m}}{((2m)!)^2}.
\]
The scalar function $d(T)$ is distinct from the bilinear correction $d(u,z)$.

\begin{theorem}[Explicit relative complement]\label{a:thm:moment}
Fix $\delta>0$ and $0<\theta<1$. Choose $N\ge2$, $C>0$, $T>0$, an integer $R\ge1$, and an integer $m\ge1$ so that
\begin{align}
 r_N&\le\theta\delta/2,&
 C&\ge\cg+2\sum_{2\le n\le N}a_n,\label{a:eq:choices1}\\
 d(T)&\ge2C/\theta,&
 \Omega_R&\le\theta\delta/(8C),&
 \epsilon_m&\le1/4.\label{a:eq:choices2}
\end{align}
If $u\in X$ and
\begin{equation}\label{a:eq:moments}
 \int_{-R}^Rx^{2k}u(x)\,dx=0\qquad(0\le k<m),
\end{equation}
then
\begin{align}
 \norm u_2^2&\le\frac2{d(T)}\tg[u]+4\Omega_R\norm u_{\Hh}^2,
                   \label{a:eq:Fouriermoment}\\
 |d[u]|&\le\theta P_\delta[u],\qquad
 (1-\theta)P_\delta[u]\le\ell_\delta[u]
                \le(1+\theta)P_\delta[u],\label{a:eq:relative}
\end{align}
where $P_\delta[u]=\tg[u]+\delta\norm u_{\Hh}^2$.
\end{theorem}
\begin{proof}
Put $I=[-R,R]$, $u_i=\one_Iu$, and $u_o=\one_{I^c}u$. Evenness and~\eqref{a:eq:moments} remove the first $m$ terms of the cosine Taylor expansion. For real arguments the remainder is bounded by $|\xi x|^{2m}/(2m)!$. Hence the unitary transform satisfies, for $|\xi|\le T$,
\begin{align*}
 |\widehat{u_i}(\xi)|
 &\le(2\pi)^{-1/2}\int_I|u(x)|\frac{|\xi x|^{2m}}{(2m)!}\,dx\\
 &\le\sqrt{R/\pi}\,\norm{u_i}_2\frac{(RT)^{2m}}{(2m)!}.
\end{align*}
If $F_T$ denotes Fourier projection onto $[-T,T]$, then
$\norm{F_Tu_i}_2^2\le\epsilon_m\norm{u_i}_2^2$.
Only ordinary $L^2$ membership of the cut piece is used; no Gamma-domain assertion about $u_i$ is required. The high-frequency part of the original $u$ is controlled by Lemma~\ref{a:lem:fourier}. Thus
\begin{align*}
 \norm u_2^2
 &\le2\epsilon_m\norm{u_i}_2^2+2\norm{u_o}_2^2
                     +\tg[u]/d(T)\\
 &\le2\epsilon_m\norm u_2^2+2\Omega_R\norm u_{\Hh}^2
                     +\tg[u]/d(T).
\end{align*}
Absorbing $2\epsilon_m\le1/2$ proves~\eqref{a:eq:Fouriermoment}.
The complete head--tail estimate then gives
\begin{align*}
 |d[u]|&\le C\norm u_2^2+r_N\norm u_{\Hh}^2\\
 &\le\frac{2C}{d(T)}\tg[u]+(4C\Omega_R+r_N)\norm u_{\Hh}^2
 \le\theta P_\delta[u].
\end{align*}
Since $\ell_\delta=P_\delta-d$, the two-sided estimate follows.
\end{proof}

\begin{corollary}[A terminating explicit prescription]\label{a:cor:choices}
The choices in Theorem~\ref{a:thm:moment} can be made without a spectral sign assumption. One sufficient prescription is:
\begin{enumerate}
\item Choose an integer $N\ge2$ with $r_N\le\theta\delta/2$.
\item Put $C=2N+4$, $L=\lceil2C/\theta\rceil$, and $T=(11/4)^L/2$.
\item Choose an integer $R\ge1$ with
$42\,81^R(3/8)^{3\cdot7^R}\le\theta\delta/(8C)$.
\item Put $m=\lceil2(RT)^2\rceil$.
\end{enumerate}
For rational $\delta$ and $\theta$, all sufficient tests in this prescription use integer and rational operations.
\end{corollary}
\begin{proof}
The bound $a_n\le\log n/\sqrt n\le2/e<1$ follows by maximizing $\log x/\sqrt x$ for $x\ge1$. Together with $\cg<6$ this gives the stated $C$. Since $\log(11/4)>1$ by~\eqref{a:eq:elementarye},
$d(T)>\log(2T)=L\log(11/4)>L\ge2C/\theta$.
The remainder $r_N$ decreases to zero geometrically: its successive ratio for $N\ge2$ is at most $1/8$. The rational envelope in~\eqref{a:eq:rationalOmega} tends to zero, so the choices of $N$ and $R$ terminate.

Write $A_0=RT\ge1$. Since $(2m)!\ge m^m$, $m\ge2A_0^2$, and $4^m\ge4m$,
\[
 \epsilon_m<\frac{2A_0}{3}\left(\frac{A_0^2}{m}\right)^{2m}
 \le\frac{2A_0}{3}4^{-m}
 \le\frac{A_0}{6m}\le\frac1{12A_0}<\frac14.
\]
This verifies the final test.
\end{proof}

The prescription is deliberately coarse. It proves finite existence with explicit sufficient inequalities, not a practically small matrix dimension. No parameter selection or finite matrix is evaluated here.

\subsubsection{Exact weighted moment representatives}

With $I=[-R,R]$ and the choices above, define
\begin{equation}\label{a:eq:momentbasis}
 b_k=W\one_Ix^{2k},\qquad
 v_k=\mathsf Cb_k=b_k-2m(b_k)\kappa\quad(0\le k<m),
\end{equation}
and let
\begin{equation}\label{a:eq:VZ}
 V=\operatorname{span}\{v_0,\ldots,v_{m-1}\},\qquad
 Z=\{u\in X_c:\text{\eqref{a:eq:moments} holds}\}.
\end{equation}

\begin{lemma}[Admitted basis and exact decomposition]\label{a:lem:VZ}
Each $v_k$ belongs to $X_c$, the family is independent, and
\begin{equation}\label{a:eq:orthdecomp}
 X_c=V\oplus Z,\qquad V\perp_{\Hh} Z.
\end{equation}
The subspace $Z$ is form closed and has codimension $m$ in $X_c$. Its weighted closure is
$\Zc=\Hh_c\cap V^{\perp_{\Hh}}$, and $X\cap\Zc=Z$.
\end{lemma}
\begin{proof}
Each $b_k$ is bounded, compactly supported, and piecewise $C^1$, with at most two boundary jumps. For $0<s\le1$ the squared translation increment is $O(s)$: the boundary strips have length $O(s)$ and bounded values, while the smooth common interior contributes $O(s^2)$. Since $\rho(s)=O(1/s)$ at zero and is integrable away from zero, $\tg[b_k]<\infty$. Its weighted norm is finite on the compact interval. Since $\kappa\in X$, centering gives $v_k\in X_c$.

For $u\in X_c$,
\[
 \ip u{v_k}_{\Hh}=\ip u{b_k}_{\Hh}
                    =\int_Ix^{2k}u(x)\,dx.
\]
If $\sum\alpha_kv_k=0$, then $W\one_IP=a\kappa$ for an even polynomial $P$. Outside $I$, positivity of $\kappa$ forces $a=0$. Inside $I$, positivity of $W$ gives $P=0$, so all coefficients vanish. Thus the weighted Gram matrix is positive definite. Projection onto its finite span preserves $X_c$ and proves~\eqref{a:eq:orthdecomp}. The moment functionals are weighted-norm continuous, so their common kernel is form closed and has the asserted codimension.

Finally $X_c$ is dense in $\Hh_c$: apply $\mathsf C$ to a weighted-norm dense sequence in $X$. Applying $I-P_V$ to this approximation proves density of $Z$ in $\Zc$, and the characterization $X\cap\Zc=Z$ follows from the moment identities.
\end{proof}

In particular the infinite complement has genuine full-form coercivity,
\begin{equation}\label{a:eq:Zcoercive}
 \ell_\delta[z]\ge(1-\theta)P_\delta[z]
 \ge(1-\theta)\min\{1,\delta\}\norm z_X^2,
\end{equation}
not merely positivity on a finite sample. Its weighted gap is at least $(1-\theta)\delta$.

\subsection{Exact finite Schur reduction and radical lifts}\label{a:sec:schur}

For $v\in V$, define the usual exact response $z_v\in Z$ by
\begin{equation}\label{a:eq:usualresponse}
 \ell_\delta(z_v,z)=\ell_\delta(v,z)\quad(z\in Z),\qquad
 S_\delta[v]=\ell_\delta[v-z_v].
\end{equation}
Lax--Milgram in the full form norm gives its existence and uniqueness by~\eqref{a:eq:Zcoercive}. Completing the square gives
\begin{equation}\label{a:eq:usualsquare}
 \ell_\delta[v+z]=S_\delta[v]+\ell_\delta[z+z_v]
                         \qquad(v\in V,\ z\in Z).
\end{equation}
Consequently the negative index and bilinear nullity of $\ell_\delta|_{X_c}$ equal those of the finite form $S_\delta$. Together with~\eqref{a:eq:meansplit}, the full form is congruent to the mean scalar, the finite Schur form, and the positive infinite complement. For $0<\delta<1/2$, positivity of $S_\delta$ is exactly the remaining finite condition for negative index one and nullity zero.

The bare finite block has the useful identity
\begin{equation}\label{a:eq:bareblock}
 \ell_\delta(v_i,v_j)=\ell_\delta(b_i,b_j)
       +(2-4\delta)m(b_i)\overline{m(b_j)}.
\end{equation}
It follows by expanding the weighted centering and using
$q_0(\kappa,z)=-\overline{m(z)}$ and
$\ip\kappa z_{\Hh}=2\overline{m(z)}$.
For compact $b_i,b_j$, prime shifts with $\log n>2R$ have disjoint supports and contribute exactly zero to this bare block. This support fact does not truncate the full noncompact response; all its tails remain controlled by Lemma~\ref{a:lem:prime}.

\begin{lemma}[Finite radical interpolation]\label{a:lem:lift}
There are a finite-dimensional subspace $F\subset\Rcal$ and a linear map $H:V\to F$ such that $P_VHv=v$. Writing
\begin{equation}\label{a:eq:liftdata}
 h(v)=Hv=v+p(v),\qquad p(v)=P_{\Zc}h(v)\in Z,
\end{equation}
the finite form
\begin{equation}\label{a:eq:G}
 G[v]=\norm{h(v)}_{\Hh}^2
     =\norm v_{\Hh}^2+\norm{p(v)}_{\Hh}^2
\end{equation}
is positive definite.
\end{lemma}
\begin{proof}
If a nonzero linear combination of the $m$ interior moments annihilated $\Rcal$, its ordinary-$L^2$ representative
$\one_I\sum_{k<m}\alpha_kx^{2k}$ would be orthogonal to the ordinary-dense space $\Rcal$. Proposition~\ref{a:prop:density} would force that representative to vanish, and then all coefficients would be zero. Therefore the moment map $\Rcal\to\C^m$ is onto. Its image already has full rank on a finite span $F$ of derivative radicals. On $V$, the same moment map is the invertible weighted Gram matrix. A linear right inverse on $F$ followed by this matrix defines $H$. The moment identities give $P_VHv=v$ and $p(v)\in Z$. Orthogonality gives~\eqref{a:eq:G}; $P_VHv=v$ makes $H$ injective, so $G$ is positive definite.
\end{proof}

Only finite interpolation has been proved. There is no uniform bound on $\norm H$ as $\delta$, $\theta$, or the moment choices change. Nor has a particular full-rank derivative minor been located and enclosed. An implementation must find and certify that minor and charge its conditioning. The reference form $G$ explicitly retains this cost, rather than replacing it silently by $\norm v_{\Hh}^2$.

\subsection{Delta-scaled response and uniform preconditioning}\label{a:sec:response}

All operators in this section act in the weighted complement Hilbert space $\Zc$, not in ordinary $L^2$. Let $T_Z$ be the nonnegative self-adjoint representative of the Gamma form restricted to $Z$, whose form domain is $Z$. Define
\begin{equation}\label{a:eq:weightedops}
 P_{\delta,Z}=T_Z+\delta I,\qquad
 D_Z=P_{\Zc}D|_{\Zc},\qquad
 A_{\delta,Z}=P_{\delta,Z}-D_Z.
\end{equation}
The Gamma form on $\Zc$ is densely defined and closed: its norm after adding the weighted squared norm is precisely the restricted $X$ norm, and Lemma~\ref{a:lem:VZ} supplies density. Since $D_Z$ is bounded with norm at most $C_0$, the bounded-perturbation theorem gives
\begin{equation}\label{a:eq:opdomains}
 \dom A_{\delta,Z}=\dom P_{\delta,Z}=\dom T_Z.
\end{equation}
The relative complement estimate gives
$A_{\delta,Z}\ge(1-\theta)\delta I$ and hence
$\norm{A_{\delta,Z}^{-1}}\le[(1-\theta)\delta]^{-1}$.
These statements concern a weighted constrained operator. They construct no ordinary-$L^2$ endpoint operator at $\delta=0$.

\begin{theorem}[Exact scaled Schur formula]\label{a:thm:scaled}
For the lift in Lemma~\ref{a:lem:lift}, put
\begin{equation}\label{a:eq:liftresponse}
 z_*(v)=\delta A_{\delta,Z}^{-1}p(v),\qquad
 e(v)=h(v)-z_*(v).
\end{equation}
Then
\begin{equation}\label{a:eq:scaledschur}
 S_\delta[v]=\ell_\delta[e(v)]
 =\delta G[v]-\delta^2\ip{p(v)}{A_{\delta,Z}^{-1}p(v)}_{\Hh}.
\end{equation}
The usual response in~\eqref{a:eq:usualresponse} is $z_v=z_*(v)-p(v)$.
\end{theorem}
\begin{proof}
The full-domain radical identity gives, for every $z\in Z$,
\[
 \ell_\delta(h(v),z)=\delta\ip{h(v)}z_{\Hh}
                    =\delta\ip{p(v)}z_{\Hh}.
\]
Thus $z_*$ is exactly the response from the affine representative $h(v)$. The affine sets $h(v)+Z$ and $v+Z$ coincide. Moreover $q_0[h(v)]=0$, so $\ell_\delta[h(v)]=\delta G[v]$. Completing the square gives~\eqref{a:eq:scaledschur}. Substituting $h(v)=v+p(v)$ into the response equation gives the last assertion.
\end{proof}

The factor $\delta$ in the forcing is exact. It does not arise by subtracting two separately uncontrolled terms of order one. Nevertheless the second term in~\eqref{a:eq:scaledschur} can be comparable to or larger than the first, and its sign consequence needs a separate estimate.

Regard $Z$ as the energy Hilbert space with inner product $P_\delta$. Set
\begin{equation}\label{a:eq:Kx}
 K_\delta=P_{\delta,Z}^{-1}D_Z,\qquad
 x(v)=\delta P_{\delta,Z}^{-1}p(v).
\end{equation}
The relative bound~\eqref{a:eq:relative} implies that $K_\delta$ is self-adjoint in this energy inner product and
\begin{equation}\label{a:eq:contraction}
 \norm{K_\delta}_{P_\delta\to P_\delta}\le\theta,
 \qquad \norm{x(v)}_{P_\delta}^2\le\delta\norm{p(v)}_{\Hh}^2.
\end{equation}
Indeed $P_\delta(K_\delta u,z)=d(u,z)$, so the quadratic bound yields the operator bound by the spectral theorem for bounded self-adjoint operators. Testing the reference equation for $x$ against $x$, and using
$\norm x_{\Hh}\le\delta^{-1/2}\norm x_{P_\delta}$, proves the second estimate.

\begin{theorem}[Exact-reference response series and charged errors]\label{a:thm:neumann}
For each integer $J\ge0$, define
\begin{equation}\label{a:eq:iterates}
 z_J(v)=\sum_{k=0}^JK_\delta^kx(v),\qquad
 e_J(v)=h(v)-z_J(v),\qquad T_J[v]=\ell_\delta[e_J(v)].
\end{equation}
Then $z_J\to z_*$ in the full form norm and
\begin{equation}\label{a:eq:Schurerror}
 0\le T_J[v]-S_\delta[v]
 \le\frac{\delta\theta^{2J+2}}{1-\theta}\norm{p(v)}_{\Hh}^2
 \le\frac{\delta\theta^{2J+2}}{1-\theta}G[v].
\end{equation}
For the residual
\begin{equation}\label{a:eq:residual}
 r_J(v;z)=\delta\ip{p(v)}z_{\Hh}-\ell_\delta(z_J(v),z),
 \qquad z\in Z,
\end{equation}
one has
\begin{align}
 r_J(v;z)&=P_\delta(K_\delta^{J+1}x(v),z),\label{a:eq:energyrow}\\
 \norm{r_J(v;\cdot)}_{P_\delta^*}
 &\le\sqrt\delta\,\theta^{J+1}\norm{p(v)}_{\Hh},\label{a:eq:energyresidual}\\
 r_J(v;z)&=\ip{D_ZK_\delta^Jx(v)}z_{\Hh},\label{a:eq:weightedrow}\\
 |r_J(v;z)|&\le C_0\theta^J\sqrt{G[v]}\norm z_{\Hh}.\label{a:eq:weightedresidual}
\end{align}
For the usual response, the corresponding approximate vector is
$\widetilde z_v=z_J(v)-p(v)$; its residual is exactly~\eqref{a:eq:residual}.
\end{theorem}
\begin{proof}
The response equation is $(I-K_\delta)z_*=x$. Since $\norm K_\delta\le\theta<1$ in the energy norm, the Neumann series converges to $z_*$ in that norm, which is equivalent to the $X$ norm for fixed $\delta>0$. Telescoping gives~\eqref{a:eq:energyrow} and~\eqref{a:eq:energyresidual}.

Put $r=K_\delta^{J+1}x$. Then
$z_*-z_J=(I-K_\delta)^{-1}r$. Exact completion of the square gives
\begin{align*}
 T_J-S_\delta
 &=\ell_\delta[z_*-z_J]
 =P_\delta(r,(I-K_\delta)^{-1}r)\\
 &\le\frac1{1-\theta}\norm r_{P_\delta}^2.
\end{align*}
Use~\eqref{a:eq:contraction} and~\eqref{a:eq:G} to obtain~\eqref{a:eq:Schurerror}.

For the weighted residual, a domain argument is essential. Every application of $P_{\delta,Z}^{-1}$ maps $\Zc$ into $\dom T_Z$, and $D_Z$ is bounded on $\Zc$. Hence $x$ and all finite terms $K_\delta^kx$ belong to the common operator domain in~\eqref{a:eq:opdomains}. The geometric identity now holds as an actual weighted operator equation:
\[
 A_{\delta,Z}z_J=\delta p-D_ZK_\delta^Jx.
\]
This proves~\eqref{a:eq:weightedrow}. Finally,
\[
 \norm{D_ZK_\delta^Jx}_{\Hh}
 \le C_0\delta^{-1/2}\norm{K_\delta^Jx}_{P_\delta}
 \le C_0\theta^J\norm p_{\Hh}
 \le C_0\theta^J\sqrt{G[v]}.
\]
For $\widetilde z_v=z_J-p$, the radical identity gives
\[
 \ell_\delta(v,z)-\ell_\delta(\widetilde z_v,z)
 =\delta\ip pz_{\Hh}-\ell_\delta(z_J,z)=r_J(v;z),
\]
which verifies the stated relation to the usual response.
\end{proof}

This weighted residual estimate is not inferred from mere form membership of an arbitrary trial vector. It uses the bounded correction $D_Z$ and equality of the two operator domains. No form-core-to-operator-core upgrade is made.

\subsubsection{Residual budgets and the implementation boundary}

For $\theta=1/2$, the energy error is
\begin{equation}\label{a:eq:halferror}
 0\le T_J-S_\delta\le2\delta\,4^{-J-1}G.
\end{equation}
For any $\eta>0$, choosing $J$ so that
$\theta^{2J+2}/(1-\theta)\le\eta$ makes the error at most $\eta\delta G$, uniformly in $\delta$. The weighted-dual estimate has a different parameter budget. If
\begin{equation}\label{a:eq:weightedbudget}
 C_0^2\theta^{2J}\le(1-\theta)\eta\delta^2,
\end{equation}
then its residual constant $\tau=C_0\theta^J$ obeys
$\tau^2\le(1-\theta)\eta\delta^2$.
To see the relevance, a residual satisfying
$|r(v;z)|\le\tau\sqrt{G[v]}\norm z_{\Hh}$ and a complement gap
$\ell_\delta[z]\ge b\delta\norm z_{\Hh}^2$ gives Schur error at most
$\tau^2G/(b\delta)$. This follows by testing the exact response error equation against the error and applying Cauchy--Schwarz. With $b=1-\theta$, condition~\eqref{a:eq:weightedbudget} gives the same target $\eta\delta G$.

These formulas are analytic constructions involving finitely many \emph{exact infinite-dimensional} positive reference solves $P_{\delta,Z}^{-1}$. They are not computed finite-matrix certificates. Replacing a reference solve by a Galerkin vector requires additional certified residual or enclosure estimates. Form density does not itself provide those estimates. Neither an evaluated reference inverse nor a quantitative enclosure of the interpolation matrix is supplied here.

\subsection{Fixed-interval sign and signed continuation}\label{a:sec:continuation}

\subsubsection{The actual finite sign supplied by the variance bound}

\begin{theorem}[Finite Schur sign on the proved interval]\label{a:thm:finiteseed}
Fix $1/2-c_{\log2}<\delta<1/2$ and any one of the moment decompositions above. With $g_\delta=c_{\log2}+\delta-1/2$,
\begin{align}
 S_\delta[v]&\ge g_\delta\norm v_{\Hh}^2,\label{a:eq:Sseed}\\
 \delta^2\ip p{A_{\delta,Z}^{-1}p}_{\Hh}
 &\le\delta G[v]-g_\delta\norm v_{\Hh}^2.\label{a:eq:signedresponse}
\end{align}
In particular $S_{23/48}\ge(509/2550000)I_V$ in the inherited weighted metric. For the same fixed lift, put
\begin{equation}\label{a:eq:sigmaH}
 \sigma_H=\inf_{v\ne0}\frac{\norm v_{\Hh}^2}{G[v]}>0.
\end{equation}
Every exact-reference iterate satisfies
\begin{equation}\label{a:eq:Tseed}
 T_J[v]\ge g_\delta\norm v_{\Hh}^2\ge g_\delta\sigma_HG[v].
\end{equation}
\end{theorem}
\begin{proof}
Since $z_*\in Z$, weighted projection gives $P_Ve(v)=v$. Thus
$\norm{e(v)}_{\Hh}^2\ge\norm v_{\Hh}^2$. Apply the full-domain bound~\eqref{a:eq:seedcoercivity} to $e(v)$ and use~\eqref{a:eq:scaledschur}. The same projection argument applies to $e_J(v)=h(v)-z_J(v)$. Finally $\sigma_H>0$ follows from finite dimensionality and injectivity of $H$.
\end{proof}

This is a sign theorem for the actual finite Schur form with its full infinite response retained. Its source is the global energy bound, not residual smallness. It gives no lower bound for $\sigma_H$ uniform over changing moment spaces or parameters.
For this fixed $V,Z,H,\delta$, choose
$\eta=g_\delta\sigma_H/(2\delta)$. Then~\eqref{a:eq:Tseed} is
$T_J\ge2\eta\delta G$. A finite $J$ meets both
\begin{equation}\label{a:eq:seedbudgets}
 \frac{\theta^{2J+2}}{1-\theta}\le\frac{g_\delta\sigma_H}{2\delta},
 \qquad C_0^2\theta^{2J}\le(1-\theta)\eta\delta^2.
\end{equation}
This existence statement does not compute $\sigma_H$, locate the interpolation minor, or turn the exact positive solves into numerical matrices.

\begin{corollary}[The stronger comparison branches in the same Schur coordinates]\label{cor:stronger-schur}
Suppose a branch of Section~\ref{sec:bounds} proves $\E\ge c_*\Var_\mu$. For every $\delta>1/2-c_*$ and every moment decomposition at that parameter,
\[
 S_\delta[v]\ge(c_*+\delta-1/2)\norm v_{\Hh}^2.
\]
If also $\delta<1/2$, the ordinary-$L^2$ pencil has negative index one and zero nullity. In particular the local-sign input gives this conclusion for $693/100000<\delta<1/2$, and the additional response contracts give it for $18851/5760000<\delta<1/2$.
\end{corollary}
\begin{proof}
The bridge gives $\ell_\delta[u]\ge(c_*+\delta-1/2)\norm u_{\Hh}^2$ on $X_c$. The exact response vector has projection $v$, so its weighted norm is at least $\norm v_{\Hh}$. Evaluating this inequality at that vector proves the Schur estimate. The weighted mean decomposition then leaves one strictly negative scalar and a positive centered form, proving the index and nullity statements exactly as in Theorem~\ref{a:thm:mainseed}.
\end{proof}
This corollary transfers already established branch bounds; it adds no new signed estimate and does not extend any branch to $\delta=0$.

\subsubsection{The signed loss under a parameter decrease}

Fix one $V,Z,H$ on an interval where the complement remains positive, and write
\begin{equation}\label{a:eq:renormalized}
 B_Z=T_Z-D_Z,\qquad A_{\delta,Z}=B_Z+\delta I,
 \qquad R_\delta=S_\delta/\delta.
\end{equation}
The same weighted constrained operator $B_Z$ occurs at each parameter, with a common operator domain. It is not an ordinary-$L^2$ endpoint realization.

\begin{proposition}[Signed renormalized continuation]\label{a:prop:continuation}
Suppose $A_{a,Z}\ge\nu I$, $0<h<\nu$, and $b=a-h>0$. Then
\begin{align}
 R_b[v]-R_a[v]
 &=h\ip{p(v)}{B_Z(B_Z+b)^{-1}(B_Z+a)^{-1}p(v)}_{\Hh},
                  \label{a:eq:signedcontinuation}\\
 R_b[v]&\ge R_a[v]
 -\frac{h(a-\nu)_+}{\nu(\nu-h)}\norm{p(v)}_{\Hh}^2.
                  \label{a:eq:continuationloss}
\end{align}
If the complement was built at $a$ with relative parameter $\theta$, then for $b=qa$ and $\theta<q<1$,
\begin{equation}\label{a:eq:ratio_loss}
 R_{qa}[v]\ge R_a[v]
 -\frac{(1-q)\theta}{(1-\theta)(q-\theta)}\norm{p(v)}_{\Hh}^2.
\end{equation}
\end{proposition}
\begin{proof}
Divide~\eqref{a:eq:scaledschur} by $\delta$ and use the resolvent identity
\[
 a(B_Z+a)^{-1}-b(B_Z+b)^{-1}
 =(a-b)B_Z(B_Z+b)^{-1}(B_Z+a)^{-1}.
\]
It is a bounded operator identity by the spectral theorem and proves~\eqref{a:eq:signedcontinuation}. The spectrum of $B_Z$ lies in $[\nu-a,\infty)$. Nonnegative spectral values contribute nonnegatively. At a negative value $-d$, where $0\le d\le a-\nu$, the negative magnitude is
\[
 \frac{hd}{(a-d)(b-d)}.
\]
It increases with $d$ for $0\le d<b$. Its endpoint is
$h(a-\nu)/[\nu(\nu-h)]$, proving~\eqref{a:eq:continuationloss}. For the relative complement use $\nu=(1-\theta)a$ and substitute $h=(1-q)a$.
\end{proof}

Differentiating the same spectral identity gives
\begin{equation}\label{a:eq:derivativeR}
 \frac{d}{d\delta}R_\delta[v]
 =-\frac{q_0[z_*(v)]}{\delta^2}.
\end{equation}
The response belongs to $\dom B_Z$, so the expression is defined. Formula~\eqref{a:eq:derivativeR} is not a monotonicity theorem: no favorable sign of $q_0[z_*]$ is supplied beyond the proved interval. The continuation loss is charged only to the possible negative spectrum of $B_Z$, but that signed contribution has not been removed.

For example, if at one parameter $a$ the full centered form obeys
$\ell_a\ge g\norm{\cdot}_{\Hh}^2$, let
\begin{equation}\label{a:eq:beta}
 \beta_H=\sup_{v\ne0}\frac{\norm{p(v)}_{\Hh}^2}{G[v]}=1-\sigma_H<1.
\end{equation}
The same projection argument gives
$R_a\ge(g/a)(1-\beta_H)G$. Therefore a sufficient one-step sign criterion is
\begin{equation}\label{a:eq:onestep}
 \frac ga(1-\beta_H)
 >\frac{(1-q)\theta}{(1-\theta)(q-\theta)}\beta_H.
\end{equation}
Every quantity must refer to the same $V,Z,H$. The inequality has not been verified here at any parameter below $1/2-c_{\log2}$.

\subsection{The unresolved sign and the scope of the construction}\label{a:sec:scope}

At a general positive parameter, the positive complement and the scaled response formula give only
\begin{equation}\label{a:eq:anglebound}
 S_\delta[v]\ge\delta\left(G[v]
                 -\frac{\norm{p(v)}_{\Hh}^2}{1-\theta}\right).
\end{equation}
Its sufficient angle condition is $\beta_H<1-\theta$. Ordinary density does not establish that condition. Decreasing $\theta$ changes the required moments and generally changes $H$ and $\beta_H$; one cannot keep a favorable interpolation cost while independently improving the complement parameter.

A finite-dimensional model makes the logical limitation explicit. On $\C^2$ take
\[
 \begin{gathered}
 B=\operatorname{diag}(0,-d),\qquad d>0,\qquad0<\phi<\pi/2,\\
 V=\operatorname{span}\{(\cos\phi,\sin\phi)\},\qquad
 Z=\operatorname{span}\{(-\sin\phi,\cos\phi)\}.
 \end{gathered}
\]
The first coordinate is an exact radical and supplies a lift. For $B+\delta I$ the complement is $\delta-d\cos^2\phi$, while the unit-vector Schur scalar is
\begin{equation}\label{a:eq:model}
 S_\delta=\frac{\delta(\delta-d)}{\delta-d\cos^2\phi}.
\end{equation}
This follows either by the $2\times2$ Schur formula or by dividing the determinant $\delta(\delta-d)$ by the complement entry. When
$d\cos^2\phi<\delta<d$, the complement is positive but the Schur form is negative. Choosing $\cos^2\phi$ sufficiently small also gives a prescribed relative complement bound on a nonempty interval. Exact reference iterations converge there. The example is not the theta operator and is not a counterexample to its desired half bound. It shows why the abstract complement, radical, and residual algebra alone cannot assign the missing sign.

There is a corresponding conditional cost statement for the actual weighted form. Suppose that a normalized centered weighted eigenstate $u$ satisfies
$q_0(u,z)=-d\ip uz_{\Hh}$ for all $z\in X_c$, with $d>0$. It is weighted-orthogonal to $\Rcal$ by~\eqref{a:eq:radicalidentity}. If
$\ell_\delta|_Z\ge b\delta\norm{\cdot}_{\Hh}^2$ with $0<b<1$, put
$v=P_Vu$ and $w=u-Hv\in Z$. Since $u\perp_{\Hh} Hv$,
\[
 \ell_\delta[w]=(\delta-d)\norm u_{\Hh}^2
                   +\delta\norm{Hv}_{\Hh}^2,
 \qquad
 \norm w_{\Hh}^2=\norm u_{\Hh}^2+\norm{Hv}_{\Hh}^2.
\]
Complement coercivity forces
\begin{equation}\label{a:eq:conditionalcost}
 \norm{Hv}_{\Hh}^2\ge
 \left(\frac{d}{(1-b)\delta}-1\right)\norm u_{\Hh}^2.
\end{equation}
Thus a hypothetical negative state is compatible with the complement estimate: its cost can move into the interpolation norm as $\delta$ shrinks. This conditional inequality is not an exclusion of such a state.


\section{Curvature and positive local row obstructions}\label{sec:obstructions}
The generator of Section~\ref{b:sec:source} is actual on its invariant smooth algebra. On real profiles write $L=-A$ and
\begin{equation}\label{b:eq:carredef}
 \Gamma(f,h)=\tfrac12\{L(fh)-fLh-hLf\},\quad
 \Gamma(f)=\Gamma(f,f),\quad \Gamma_2(f)=\tfrac12L\Gamma(f)-\Gamma(f,Lf).
\end{equation}
For $f\in\dom A$ the integrated half defect is
\begin{equation}\label{b:eq:defect}
 \Def[f]=\norm{Af}_\mu^2-\tfrac12\E[f].
\end{equation}
Let $\Eloc=C^\infty_{\mathrm{even}}(\R;\R)$ have the compact-open smooth seminorms
\begin{equation}\label{b:eq:localtopology}
 p^{\mathrm{loc}}_{R,m}(f)=\max_{0\le q\le m}\sup_{|x|\le R}|f^{(q)}(x)|.
\end{equation}
A local row is a continuous real-linear map $\Eloc\to\C$, equivalently an even compactly supported distribution of finite order. A common row family is required in a factorization; its support or order need not be uniform over the family.


\begin{theorem}[Curvature obstruction]\label{b:thm:maincurvature}
For the actual operator $A$, the profile
\[
 g=\kappa''/\kappa-1/4
\]
belongs to $\dom(A)$, satisfies $Ag=g/2$ and $\mu(g)=0$, and obeys
\begin{equation}\label{b:eq:mainquarter}
 \frac{\Gamma_2(g)(x)}{\Gamma(g)(x)}\longrightarrow\frac14
 \quad(|x|\to\infty).
\end{equation}
The denominator is positive in both sufficiently distant tails. For every $K>1/4$ there are a real even $f\in C_c^\infty(\R)$ with $\mu(f)=0$ and a nonempty open interval on which $\Gamma_2(f)<K\Gamma(f)$. Nevertheless,
\[
 \int_\R[\Gamma_2(g)-\tfrac12\Gamma(g)]\,d\mu=0.
\]
\end{theorem}

\begin{theorem}[Local density and positive local rows]\label{b:thm:mainlocal}
The real span $\mathcal G$ of
\begin{equation}\label{b:eq:halfprofiles}
 g_j=\frac{\kappa^{(2j)}}\kappa-4^{-j},\qquad j\ge1,
\end{equation}
is dense in $\Eloc$. Every $g_j$ is an actual centered half-eigenstate, and has exactly centered even compact smooth approximants converging in the graph norm of $A$.

Let $(Z,\eta)$ be a positive measure space, and suppose that outside one null set $\ell_z:\Eloc\to\C$ is a continuous real-linear row, with $z\mapsto\ell_z(f)$ measurable for the profiles in question. There is no identity
\begin{equation}\label{b:eq:mainnoidentity}
 \Def[f]=\int_Z|\ell_z(f)|^2\,d\eta(z)
\end{equation}
for every real even $f\in C_c^\infty(\R)$ with $\mu(f)=0$, using one common row family and measure. If only the lower bound
$\Def[f]\ge\int_Z|\ell_z(f)|^2\,d\eta\ge0$ is asserted on that class, all rows must vanish on $\Eloc$ almost everywhere. Finite or countable collections of rows are included.
\end{theorem}

Sections~\ref{b:sec:curvature}--\ref{b:sec:cutoffs} prove Theorem~\ref{b:thm:maincurvature}. Sections~\ref{b:sec:density}--\ref{b:sec:nogo} prove Theorem~\ref{b:thm:mainlocal}, and Section~\ref{b:sec:paths} treats positive semidefinite finite-path matrices without discarding compatibility relations. The common foundation, explicit-formula bridge and actual half-eigenstates were proved in Sections~\ref{sec:foundations}--\ref{sec:spectral}. The classical explicit formula is used as a theorem, not as a hypothesis about zero locations.


\subsection{Carr\'e du champ and all mixed jump compositions}\label{b:sec:carre}

Define the symmetric displacement measure $J$ by
\[
 \int F(h)\,J(dh)=\int_{s>0}[F(s)+F(-s)]\,d\nu(s),
 \qquad r(x,h)=a(x)\kappa(x+h).
\]
The measure $d\mu(x)r(x,h)J(dh)=\kappa(x)\kappa(x+h)\,dx\,J(dh)$ is invariant under edge reversal $(x,h)\mapsto(x+h,-h)$. Thus
\begin{equation}\label{b:eq:energy-symmetric}
 \E[f]=\tfrac12\int_\R\int
       \kappa(x)\kappa(x+h)|f(x+h)-f(x)|^2\,J(dh)\,dx.
\end{equation}
The factor $1/2$ accounts for the two orientations.

\begin{proposition}[Pointwise and integrated identities]\label{b:prop:carre}
For real even $f,h\in\CE$, all terms in~\eqref{b:eq:carredef} belong to $\CE$, and
\begin{align}
 \Gamma(f,h)(x)&=\tfrac12\int r(x,z)
   [f(x+z)-f(x)][h(x+z)-h(x)]\,J(dz),\label{b:eq:carre-jump}\\
 \Gamma(f,h)&=\tfrac14fh+\tfrac12
   \{fT(\kappa h)+hT(\kappa f)-T(\kappa fh)\},\label{b:eq:carre-source}\\
 \int\Gamma(f)\,d\mu&=\E[f],\qquad
 \int\Gamma_2(f)\,d\mu=\norm{Af}_\mu^2.\label{b:eq:integrated}
\end{align}
\end{proposition}
\begin{proof}
The jump product rule gives~\eqref{b:eq:carre-jump}; its integral is~\eqref{b:eq:energy-symmetric}. Substitution of~\eqref{b:eq:generator-source} gives~\eqref{b:eq:carre-source} and membership in $\CE$. Proposition~\ref{b:prop:generator} admits every product and every subsequent application of $L$. Since $A1=0$ and $\Gamma(f)\in\dom(A)$,
$\int L\Gamma(f)\,d\mu=0$. Also $Lf\in\D$ and edge polarization gives
\[
 \int\Gamma(f,Lf)\,d\mu=\E(f,Lf)=-\norm{Af}_\mu^2.
\]
The bilinear edge integral is absolutely integrable by Cauchy--Schwarz in the edge measure; the other terms lie in $L^2(\mu)\subset L^1(\mu)$. This proves~\eqref{b:eq:integrated}.
\end{proof}

The formulas use the full operator, not its Gamma component alone. To record explicitly the two-step terms, put $G_\Gamma=D_\Gamma-\cg$. For $h\in\RR$,
\begin{align}\label{b:eq:mixed}
 T(\kappa Th)=aQ(WQh)
 &=aG_\Gamma(WG_\Gamma h)-aG_\Gamma(WPh)\notag\\
 &\quad-aP(WG_\Gamma h)+aP(WPh).
\end{align}
Every intermediate $WG_\Gamma h$ and $WPh$ belongs to $\RR$ by Lemma~\ref{b:lem:T}. Thus Gamma--Gamma, both mixed orders, and prime--prime are separately defined smooth $\SSS$ functions. Constants in $G_\Gamma$ and their products are retained.

There is also locally uniform absolute prime--prime summability. Let
\[
 U_h(x)=\sum_{n\ge2}a_n\bigl(|h(x+s_n)|+|h(x-s_n)|\bigr).
\]
Equation~\eqref{b:eq:primepoint} bounds it by $C_\epsilon e^{(1/2+\epsilon)|x|}$. Therefore $WU_h$ decreases faster than every exponential, and the absolute double sum is bounded by
\[
 \sum_{m\ge2}a_m\{W(x+s_m)U_h(x+s_m)+W(x-s_m)U_h(x-s_m)\},
\]
which converges locally uniformly by the same bound. Differentiating the inner prime series is legitimate; multiplication by $W$ places it in $\RR$ before the outer Gamma or prime operation.

For clarity, the corresponding jump integrals are absolutely controlled with their small-shift cancellations retained. For $\sigma\in\{1,-1\}$ define
\[
 I_f(z,t)=a(z)\kappa(z+\sigma t)[f(z+\sigma t)-f(z)],\qquad
 J_f(z,t)=\tfrac12a(z)\kappa(z+\sigma t)[f(z+\sigma t)-f(z)]^2.
\]
For $f\in\CE$, the $\nu$-integral of the absolute value of every fixed $z$ derivative has an exponential envelope in $z$. On $0<t\le1$ the two differentiated expressions are respectively $O(t)$ and $O(t^2)$ with such envelopes. Above one, expansion reduces the estimate to shifted rapid functions $\kappa f^j$ and exponential factors in $z$; the Gamma tail and~\eqref{b:eq:primepoint} apply to each channel.

In $L\Gamma(f)$ retain the difference $J_f(x+h,t)-J_f(x,t)$. Its absolute $t$ integral is $O(|h|)$ for $|h|\le1$ on compact $x$ intervals. In $\Gamma(f,Lf)$ retain
$[f(x+h)-f(x)][I_f(x+h,t)-I_f(x,t)]$, whose absolute $t$ integral is $O(h^2)$. For $|h|>1$, the outer factor $\kappa(x+h)$ controls the exponential envelope and the complete prime and Gamma tails. Fubini therefore applies to these cancelled integrands for every pair of channels, including large opposite shifts. It does not authorize separating the infinite Gamma diagonal rate.

For finite jump cutoffs in both instances of the generator, put
\[
 p=f(x+h)-f(x),\quad q=f(x+k)-f(x),\quad
 z=f(x+h+k)-f(x+h)-f(x+k)+f(x),
\]
\[
 C_{h,k}=a(x)\kappa(x+h)a(x+h)\kappa(x+h+k),\qquad
 D_{h,k}=a(x)^2\kappa(x+h)\kappa(x+k).
\]
Direct expansion gives
\begin{align}\label{b:eq:doublejump}
 4\Gamma_2^{\mathrm{cut}}(f)(x)
 =\iint\bigl\{&C_{h,k}(q+z)^2-D_{h,k}q^2\notag\\
 &-2p[C_{h,k}(q+z)-D_{h,k}q]\bigr\}
     J_{\mathrm{cut}}(dh)J_{\mathrm{cut}}(dk).
\end{align}
The preceding estimates give the full limit with the braces kept together. The first two terms cannot be independently integrated over the infinite activity in $h$. No general nonnegative-square interpretation of the bracket is asserted.

\subsection{The exact half state and the quarter curvature limit}\label{b:sec:curvature}

Proposition~\ref{f:prop:halfeigen} proves, using the full form domain, that every $g_j$ in~\eqref{f:eq:halfprofiles} satisfies
\begin{equation}\label{b:eq:actualhalf}
 g_j\in\dom(A),\qquad Ag_j=\tfrac12g_j,\qquad \mu(g_j)=0.
\end{equation}
The upper and lower theta estimates~\eqref{a:eq:derivativeupper}--\eqref{a:eq:thetalower} imply
\begin{equation}\label{b:eq:ratio}
 |\kappa^{(m)}(x)/\kappa(x)|\le C_me^{2m|x|}.
\end{equation}
Every derivative of a quotient in~\eqref{f:eq:halfprofiles} is a finite sum of products of these ratios; hence $g_j\in\CE$. In particular all ensuing identities concern actual operator-domain vectors, and~\eqref{b:eq:generator-source} gives the pointwise equality
\begin{equation}\label{b:eq:halfsource}
 T(\kappa g_j)=0.
\end{equation}
The smooth residual of the Hilbert-space identity vanishes pointwise because the density of $\mu$ is strictly positive.

For the first state $g=g_1$, let $t=\pi e^{2x}$ for $x\ge0$. The theta derivative recurrence is
\begin{equation}\label{b:eq:recurrence}
 \kappa^{(m)}(x)=e^{x/2}\sum_{n\ge1}P_m(n^2t)e^{-n^2t},\qquad
 P_{m+1}(t)=(1/2-2t)P_m(t)+2tP_m'(t),
\end{equation}
with $P_0(t)=4t^2-6t$. In particular,
\[
 P_2(t)=16t^4-112t^3+165t^2-\tfrac{75}{2}t.
\]
For $t\ge\pi$,
\begin{equation}\label{b:eq:thetatail}
 \sum_{n\ge2}n^{2m+4}e^{-(n^2-1)t}
 \le e^{-3t}\sum_{n\ge2}n^{2m+4}e^{-\pi(n^2-4)}.
\end{equation}
The last series is finite. Since $P_0(t)$ is bounded below by a positive multiple of $t^2$ on this range,
\[
 \frac{\kappa''(x)}{\kappa(x)}
  =\frac{P_2(t)}{P_0(t)}+O(t^2e^{-3t})=4t^2+O(t).
\]
Evenness gives the two-sided asymptotic
\begin{equation}\label{b:eq:gasymptotic}
 g(x)=4\pi^2e^{4|x|}+O(e^{2|x|}).
\end{equation}
No numerical theta truncation enters this estimate.

Define the full source corrections
\begin{equation}\label{b:eq:HK}
 h=\kappa g^2\in\RR,\qquad H=Th\in\SSS,\qquad
 K=T(\kappa H)=aQ(WQh)\in\SSS.
\end{equation}
Formula~\eqref{b:eq:mixed} expands $K$ with all four channel combinations.

\begin{proposition}[Exact signed identities]\label{b:prop:quarteridentities}
Pointwise on $\R$,
\begin{align}
 \Gamma(g)&=\tfrac14g^2-\tfrac12H,\label{b:eq:Gg}\\
 L\Gamma(g)&=-\tfrac18g^2+\tfrac12K,\label{b:eq:LGg}\\
 \Gamma_2(g)&=\tfrac1{16}g^2-\tfrac14H+\tfrac14K,\label{b:eq:G2g}\\
 \Gamma_2(g)-\tfrac12\Gamma(g)&=-\tfrac1{16}g^2+\tfrac14K.\label{b:eq:halfcurvature}
\end{align}
Consequently, as $r=|x|\to\infty$,
\begin{align}
 \Gamma(g)(x)&=4\pi^4e^{8r}(1+o(1)),\label{b:eq:Gtail}\\
 \Gamma_2(g)(x)&=\pi^4e^{8r}(1+o(1)),\label{b:eq:G2tail}\\
 \Gamma_2(g)(x)-\tfrac12\Gamma(g)(x)&=-\pi^4e^{8r}(1+o(1)).\label{b:eq:deftail}
\end{align}
\end{proposition}
\begin{proof}
Equations~\eqref{b:eq:carre-source} and~\eqref{b:eq:halfsource} give~\eqref{b:eq:Gg}. The source row gives $L(g^2)=-g^2/2-H$ and $LH=-H/2-K$, so the two $H$ terms cancel in $L\Gamma(g)$ and yield~\eqref{b:eq:LGg}. Since $Lg=-g/2$, bilinearity gives
$\Gamma_2(g)=L\Gamma(g)/2+\Gamma(g)/2$, proving~\eqref{b:eq:G2g}--\eqref{b:eq:halfcurvature}. Every application of $L$ is permitted by Proposition~\ref{b:prop:generator}. For every $\epsilon>0$, $H,K=O(e^{\epsilon r})$. Combining this with~\eqref{b:eq:gasymptotic} gives the three asymptotics.
\end{proof}

The quotient in~\eqref{b:eq:mainquarter} follows. More generally, for each fixed $K_0>1/4$,
\begin{equation}\label{b:eq:anycurvature}
 \Gamma_2(g)-K_0\Gamma(g)
  =(1/16-K_0/4)g^2+(K_0/2-1/4)H+K/4<0
\end{equation}
in both sufficiently distant tails. The number $1/4$ is an upper restriction on a possible pointwise lower bound; the calculation does not establish $\Gamma_2\ge\Gamma/4$.

\subsection{Compact witnesses and the integrated cancellation}\label{b:sec:cutoffs}

Fix real even $\chi\in C_c^\infty$ with $0\le\chi\le1$, equal to one on $[-1,1]$ and supported in $[-2,2]$, and put $\chi_R(x)=\chi(x/R)$ for $R\ge1$. Initially let $g_R=\chi_Rg$. These profiles and their squares have derivativewise exponential bounds uniform in $R$. The rapid theta decay therefore gives
\begin{equation}\label{b:eq:rapidcutoffs}
 \kappa g_R\longrightarrow\kappa g,\qquad
 \kappa g_R^2\longrightarrow\kappa g^2
 \quad\hbox{in }\RR.
\end{equation}
Indeed the profiles agree with their limits on every fixed compact set for large $R$, and outside that set stronger rapid seminorms control all cutoff derivative terms.

Set $F_R=T(\kappa g_R)$ and $H_R=T(\kappa g_R^2)$. Lemma~\ref{b:lem:T} gives
\begin{equation}\label{b:eq:sourcecutoffs}
 F_R\longrightarrow0,\qquad H_R\longrightarrow H\quad\hbox{in }\SSS.
\end{equation}
Using~\eqref{b:eq:carre-source}, write
\[
 B_R=\Gamma(g_R)=g_R^2/4+g_RF_R-H_R/2,
 \qquad B=\Gamma(g).
\]
Then $B_R\to B$ locally smoothly and $\kappa B_R\to\kappa B$ in $\RR$. For the product $g_RF_R$, use the uniform exponential bounds on $g_R$, convergence in one specified small-exponential seminorm at each derivative order, and the theta decay. Equation~\eqref{b:eq:generator-source} now yields $LB_R\to LB$ locally uniformly. Also
\[
 \Gamma(g_R,F_R)=g_RF_R/4+\tfrac12\{g_RT(\kappa F_R)
                  +F_RT(\kappa g_R)-T(\kappa g_RF_R)\}\longrightarrow0
\]
locally uniformly, since the two $T$ arguments containing $F_R$ converge to zero in $\RR$. Finally $Lg_R=-g_R/2-F_R$, so
\[
 \Gamma_2(g_R)-\tfrac12\Gamma(g_R)
       =\tfrac12LB_R+\Gamma(g_R,F_R).
\]
Thus
\begin{equation}\label{b:eq:curvaturecutoffs}
 \Gamma(g_R)\to\Gamma(g),\qquad \Gamma_2(g_R)\to\Gamma_2(g)
 \quad\hbox{locally uniformly}.
\end{equation}
This argument avoids interchanging an uncontrolled double jump sum.

Choose a fixed real even $\psi\in C_c^\infty$ with $\mu(\psi)=1$ and set
\begin{equation}\label{b:eq:centeredcutoff}
 f_R=\chi_Rg-\mu(\chi_Rg)\psi.
\end{equation}
Theta domination gives $\chi_Rg\to g$ in $\He$ and $\mu(\chi_Rg)\to0$. These centered compact profiles have the same uniform derivative envelopes and rapid-weighted limits. The proof above therefore also gives~\eqref{b:eq:curvaturecutoffs} with $f_R$ in place of $g_R$. Fix $K_0>1/4$ and a finite $x_0$ satisfying the strict tail inequality~\eqref{b:eq:anycurvature}. For sufficiently large finite $R$, that strict inequality persists at $x_0$ for $f_R$. Smoothness makes it hold on a nonempty open interval, hence not merely on a null set. This proves the compact-test assertion of Theorem~\ref{b:thm:maincurvature}.

There is actual graph convergence for these particular cutoffs: $g_R\to g$ and $Ag_R\to Ag$ in $\He$, by~\eqref{b:eq:sourcecutoffs} and theta domination. Likewise $B_R\to B$ and $AB_R\to AB$ in $\He$, using the uniform exponential envelopes and $\kappa B_R\to\kappa B$ in $\RR$. No assertion that compact smooth functions are an operator core for arbitrary vectors is required.

\begin{proposition}[Exact integrated cancellation]\label{b:prop:cancellation}
The curvature defect of $g$ integrates to zero, and
\begin{align}
 \int\Gamma(g)\,d\mu&=\tfrac12\norm g_\mu^2,
 &\int\Gamma_2(g)\,d\mu&=\tfrac14\norm g_\mu^2,\label{b:eq:halfintegrals}\\
 \int H\,d\mu&=-\tfrac12\norm g_\mu^2,
 &\int K\,d\mu&=\tfrac14\norm g_\mu^2.\label{b:eq:HKintegrals}
\end{align}
\end{proposition}
\begin{proof}
All terms are integrable by the admitted-class estimates. Equation~\eqref{b:eq:integrated} and $Ag=g/2$ prove~\eqref{b:eq:halfintegrals}; integrating~\eqref{b:eq:Gg} and~\eqref{b:eq:halfcurvature} proves~\eqref{b:eq:HKintegrals}.
\end{proof}

The negative tails in~\eqref{b:eq:deftail} must therefore be balanced by a positive part elsewhere. In particular $K$ cannot be discarded in an integrated argument. Pointwise curvature failure is not a counterexample to the integrated half inequality.

\subsection{Local smooth density of the half eigenstates}\label{b:sec:density}

\begin{theorem}[Local density]\label{b:thm:density}
The closure of $\mathcal G=\operatorname{span}_{\R}\{g_j:j\ge1\}$ in $\Eloc$ is all of $\Eloc$.
\end{theorem}
\begin{proof}
By Hahn--Banach separation it suffices to show that every continuous linear functional $\ell$ annihilating all $g_j$ is zero. Extend it to all smooth functions by even symmetrization, and then complexify. The extension is an even compactly supported distribution. Set $w=\ell/\kappa$, meaning multiplication by the smooth positive reciprocal on a neighborhood of its compact support.

By Lemma~\ref{d:lem:theta-tails}, $\kappa$ has an even holomorphic continuation to the strip $\mathcal S=\{z:|\operatorname{Im}z|<\pi/4\}$ with all the derivative tail bounds.

The function
\[
 F(t)=\ip{w(x)}{\kappa(x+t)},\qquad t\in\mathcal S,
\]
where brackets denote distributional pairing without conjugation, is holomorphic. Compact support and finite order justify differentiation under this pairing. Evenness gives $F(-t)=F(t)$, and annihilation of the $g_j$ gives
\[
 F(0)=\ell(1),\qquad F^{(2j)}(0)=4^{-j}\ell(1)\quad(j\ge1).
\]
Its Taylor series is therefore the series of $\ell(1)\cosh(t/2)$. The identity theorem on the connected strip gives
\begin{equation}\label{b:eq:Fcosh}
 F(t)=\ell(1)\cosh(t/2),\qquad t\in\mathcal S.
\end{equation}
On the real axis $F$ decreases faster than every exponential. Indeed, for $\supp(w)\subset[-R,R]$ and order at most $m$, the distribution bound with a fixed auxiliary cutoff gives
\[
 |F(t)|\le C\max_{0\le q\le m}\sup_{|x|\le R+1}|\kappa^{(q)}(x+t)|
       \le C_Be^{-B|t|}\quad(B>0).
\]
The derivative theta envelope justifies the last step. Hence~\eqref{b:eq:Fcosh} forces $\ell(1)=0$ and $F=0$.

On the real line, evenness identifies $F$ with $w*\kappa$. This convolution is Schwartz, because $w$ is compactly supported and $\kappa$ is Schwartz. In the convention $\widehat h(\xi)=\int e^{-i\xi x}h(x)\,dx$, Fourier transformation gives
\[
 \widehat w(\xi)\widehat\kappa(\xi)=0.
\]
The transform of a compact distribution is entire, directly by pairing it with $e^{-izx}$. Since $\widehat\kappa(0)=\int\kappa>0$, continuity makes $\widehat\kappa$ nonzero on a real interval around zero. Therefore $\widehat w$ vanishes on that interval, and the identity theorem makes it identically zero. Fourier injectivity for distributions yields $w=0$, hence $\ell=0$. Hahn--Banach proves the assertion.
\end{proof}

This proof does not divide by the completed zeta function at its zeros. Only the nonvanishing of the theta transform near zero is needed in its last step. Weighted centering is not continuous in the local smooth topology, so there is no contradiction between $\mu(g_j)=0$ and local approximation of the constant profile.

\begin{corollary}[Exact finite data interpolation]\label{b:cor:finiteimage}
For every continuous real-linear map $B:\Eloc\to V$ into a finite-dimensional space,
\[
 B(\mathcal G)=B(\Eloc).
\]
\end{corollary}
\begin{proof}
Density makes $B(\mathcal G)$ dense in $B(\Eloc)$. The former is a linear subspace of a finite-dimensional space and is closed, proving equality.
\end{proof}

Thus finite endpoint and jet data can be matched exactly by a finite linear combination of the half states, whenever the data are compatible with evenness and any repeated endpoints. This statement is local, not weighted.

\subsection{Graph cutoffs and a strictly positive compact defect}\label{b:sec:graph}

Use the fixed $\chi_R$ and $\psi$ of Section~\ref{b:sec:cutoffs}. For each fixed $j\ge1$ set
\begin{equation}\label{b:eq:fjR}
 e_{j,R}=\mu(\chi_Rg_j),\qquad
 f_{j,R}=\chi_Rg_j-e_{j,R}\psi.
\end{equation}
These profiles are real, even, compact smooth, and exactly centered.

\begin{lemma}[Special graph approximation]\label{b:lem:graph}
For each fixed $j$, as $R\to\infty$,
\[
 f_{j,R}\to g_j\text{ in }\Eloc\text{ and }\He,
 \qquad Af_{j,R}\to Ag_j\text{ in }\He,
 \qquad\Def[f_{j,R}]\to0.
\]
\end{lemma}
\begin{proof}
Theta domination and~\eqref{b:eq:actualhalf} give $e_{j,R}\to0$. Local smooth and weighted Hilbert convergence follow. Put $r_j=\kappa g_j=\kappa^{(2j)}-4^{-j}\kappa\in\RR$. Then
\[
 \kappa(f_{j,R}-g_j)=(\chi_R-1)r_j-e_{j,R}\kappa\psi\longrightarrow0
 \quad\hbox{in }\RR.
\]
For every seminorm this follows from Leibniz's rule and stronger rapid decay on the tail supporting the cutoff derivatives. Lemma~\ref{b:lem:T} gives convergence of its $T$ image in $\SSS$, hence in $\He$. Equation~\eqref{b:eq:generator-source} proves graph convergence. Finally~\eqref{b:eq:defect} is graph-continuous and $\Def[g_j]=0$.
\end{proof}

This is an approximation theorem for the specified half states. It does not assert a general compact smooth operator core.

\begin{lemma}[Logarithmic Gamma coercivity]\label{b:lem:fourier}
For the unitary Fourier transform,
\begin{equation}\label{b:eq:multiplier}
 \tg[u]=\int_\R d_\Gamma(\xi)|\widehat u(\xi)|^2\,d\xi,
 \qquad d_\Gamma(\xi)\ge\tfrac12\log(1+4\xi^2).
\end{equation}
\end{lemma}
\begin{proof}
This is Lemma~\ref{a:lem:fourier}, including its Plancherel--Tonelli identity for possibly infinite form values.
\end{proof}

\begin{proposition}[The defect is not identically zero]\label{b:prop:positive}
There is a real even $f\in C_c^\infty(\R)$ with $\mu(f)=0$ and $\Def[f]>0$.
\end{proposition}
\begin{proof}
Choose nonzero real even $\phi\in C_c^\infty$. For a real frequency $R>0$, put
\[
 m_R=\mu(\phi\cos(Rx)),\qquad h_R=\phi\cos(Rx)-m_R\psi,
 \qquad u_R=\kappa h_R,
\]
and write $b=\kappa\phi$, $d=\kappa\psi$. These are centered compact smooth profiles with fixed compact support. Riemann--Lebesgue gives
\begin{equation}\label{b:eq:HFmass}
 m_R\to0,\quad
 \norm{h_R}_\mu^2\to\tfrac12\int\phi^2\,d\mu>0,\quad
 \norm{u_R}_2^2\to\tfrac12\norm b_2^2>0.
\end{equation}
In the unitary Fourier convention,
\[
 \widehat u_R(\xi)=\tfrac12\{\widehat b(\xi-R)+\widehat b(\xi+R)\}
                         -m_R\widehat d(\xi).
\]
Because $\widehat b$ and $\widehat d$ are Schwartz, the Fourier mass in $[-M,M]$ tends to zero for every fixed $M$. A positive amount of mass therefore escapes to arbitrarily high frequencies. Lemma~\ref{b:lem:fourier} gives
\[
 \liminf_{R\to\infty}\tg[u_R]
 \ge\tfrac12\log(1+4M^2)\,\tfrac12\norm b_2^2.
\]
Letting $M\to\infty$ proves $\tg[u_R]\to\infty$.

The complete weighted prime estimate in Lemma~\ref{a:lem:prime} gives
$|p[u]|\le C_P\norm u_{\Hh}^2$ for a finite $C_P$, and
$\norm u_2^2\le\norm W_\infty\norm u_{\Hh}^2$. Since $\norm{\kappa h}_{\Hh}=\norm h_\mu$, the full bridge yields
\[
 \E[h_R]\ge\tg[u_R]-C_0\norm{h_R}_\mu^2\longrightarrow\infty,
 \qquad C_0=|\cg|\norm W_\infty+C_P.
\]
No nonzero prime tail is discarded. Proposition~\ref{b:prop:generator} places each $h_R$ in $\dom(A)$. Cauchy--Schwarz then gives
\[
 \Def[h_R]\ge\E[h_R]
       \left(\frac{\E[h_R]}{\norm{h_R}_\mu^2}-\frac12\right)>0
\]
for all sufficiently large $R$, by~\eqref{b:eq:HFmass}. Choose any such finite frequency.
\end{proof}

The witness is strictly positive. It is not a negative direction for the integrated inequality and does not construct an actual subthreshold state.

\subsection{No nontrivial positive local row factorization}\label{b:sec:nogo}

\begin{proof}[Proof of the factorization assertions in Theorem~\ref{b:thm:mainlocal}]
Suppose~\eqref{b:eq:mainnoidentity} holds on the exactly centered compact class. Fix $j$ and take integer $R\to\infty$ in~\eqref{b:eq:fjR}. Local smooth convergence and row continuity give
$\ell_z(f_{j,R})\to\ell_z(g_j)$ wherever the row is continuous. Fatou's lemma and Lemma~\ref{b:lem:graph} imply
\begin{align}\label{b:eq:Fatou}
 \int_Z|\ell_z(g_j)|^2\,d\eta
 &\le\liminf_{R\to\infty}\int_Z|\ell_z(f_{j,R})|^2\,d\eta\notag\\
 &=\liminf_{R\to\infty}\Def[f_{j,R}]=0.
\end{align}
Thus $\ell_z(g_j)=0$ almost everywhere. There are countably many $j$, so removal of their exceptional null sets and the common continuity exceptional set leaves one full-measure set on which every $g_j$ is annihilated. Theorem~\ref{b:thm:density} and continuity force $\ell_z=0$ on all of $\Eloc$ at every point of this set. The right side of~\eqref{b:eq:mainnoidentity} is then zero for every test, contradicting Proposition~\ref{b:prop:positive}.

For the lower-bound assertion, replace the equality in~\eqref{b:eq:Fatou} by the upper bound $\Def[f_{j,R}]\to0$. The same reasoning forces all rows to vanish, but gives no contradiction to $\Def\ge0$ itself: its lower bound may simply be zero. A countable row collection is handled by one further countable intersection of full-measure sets, or by the product with counting measure. Neither Fatou's lemma nor the zero-integral argument requires sigma-finiteness.
\end{proof}

The graph-cutoff step is essential. An identity asserted only on compact tests has not been substituted at an unbounded eigenstate without justification. The row family and positive measure must be fixed for all inputs; choosing rows separately for each $f$ is outside the theorem. Supports and distributional orders may vary without any uniform bound. Consequently the statement includes integrations over arbitrarily large Gamma and prime-power displacements, provided each individual row has compact dependence.

\subsection{Finite paths and the actual compatible data space}\label{b:sec:paths}

For each parameter $z$, let $B_z$ collect finitely many endpoint values, derivatives, or edge differences of a profile in $\Eloc$. Its components are continuous local rows. The image
\[
 V_z=B_z(\Eloc)
\]
is regarded as a real finite-dimensional space, even when the coordinates are complex. For edge differences it already contains all cycle relations, together with evenness and repeated-endpoint relations. It need not be the full redundant coordinate space.

\begin{corollary}[Positive finite-path matrices]\label{b:cor:paths}
Suppose $M_z$ is a finite Hermitian matrix whose quadratic form is nonnegative on the actual image $V_z$. Assume a common measurable family of these quadratic expressions and a positive measure $\eta$. There is no identity
\begin{equation}\label{b:eq:matrixidentity}
 \Def[f]=\int_Z(B_zf)^*M_z(B_zf)\,d\eta(z)
\end{equation}
for every real even exactly centered $f\in C_c^\infty(\R)$. If equality is replaced by a nonnegative lower bound, the quadratic form must vanish on $V_z$ almost everywhere.
\end{corollary}
\begin{proof}
Apply Lemma~\ref{b:lem:graph} and Fatou to each fixed $g_j$. On one common full-measure set,
$(B_zg_j)^*M_z(B_zg_j)=0$ for all $j$. On a real vector space a nonnegative quadratic form has zero real polarized pairing between any zero-norm vector and every other vector: expand its value on $v+tw$ for all real $t$. Hence the form vanishes on the real span of $\{B_zg_j:j\ge1\}$. Corollary~\ref{b:cor:finiteimage} identifies this span with $V_z$. The integrand is consequently zero for every even smooth profile, contradicting Proposition~\ref{b:prop:positive} in the equality case. The lower-bound conclusion follows by the same Fatou argument.
\end{proof}

This proof uses real polarization $\operatorname{Re}(v^*M_zw)$; it does not assume that $V_z$ is a complex subspace and does not require a measurable square-root selection. The support points can be
\[
 x,\quad x+h_1,\quad x+h_1+h_2,\quad\ldots,\quad x+h_1+\cdots+h_m
\]
for any finite $m$, with no bound on the individual shifts. The measure may include all Gamma and prime-power paths, both orientations, every mixed channel, and countably many finite path lengths. The obstruction comes from the finite local dependence of each nonnegative integrand, not from a bound on the total range or number of jumps.

\subsection{Weighted limits and the remaining integrated problem}\label{b:sec:scope}

Local density gives no control of the weighted cost of an approximation. A precise distinction is useful.

\begin{proposition}[Weighted cost of local approximation]\label{b:prop:blowup}
Let $f\in\He\cap\Eloc$ lie outside $\ker(A-\tfrac12I)$. If $v_n\in\mathcal G$ converges to $f$ in $\Eloc$, then $\norm{v_n}_\mu\to\infty$.
\end{proposition}
\begin{proof}
Otherwise there is a bounded subsequence in $\He$, with a further weakly convergent subsequence. The half-eigenspace is a closed linear subspace of $\He$, hence weakly closed, so the weak limit lies in it. Pairing with any compact smooth even test and using local convergence identifies the limit with $f$. Such tests are dense in $\He$, giving a contradiction. Thus no bounded subsequence exists, which is precisely divergence of the norms to infinity.
\end{proof}

The constant profile and the compact witness of Proposition~\ref{b:prop:positive} are examples. In particular even a centered compact target can have this divergent approximation cost. This local argument does not prove weighted completeness or identify the entire half-eigenspace. The weighted derivative closure, a different space from the full centered Hilbert space, is analyzed in Section~\ref{sec:synthesis-unified}.

For an admitted real profile, write $t=T(\kappa f)$. The actual source row gives the exact integrated identity
\begin{equation}\label{b:eq:signedsource-defect}
 \Def[f]=\norm t_\mu^2+\tfrac12\ip t f_\mu.
\end{equation}
The square is nonnegative, but no universal sign has been proved for the sum. Its global nature is visible without guessing a topology: for every fixed $x$, the row $f\mapsto T(\kappa f)(x)$ annihilates every $g_j$, while
\[
 T(\kappa)(x)=-\tfrac12.
\]
By Theorem~\ref{b:thm:density} it cannot be continuous in $\Eloc$. Thus this genuinely global row is outside the local-square obstruction, despite being evaluated at a single spatial point.

The integrated edge formulation is also exact but does not supply a sign. Equip edge functions with the norm
\[
 \norm F_{\mathrm{edge}}^2=\tfrac12\int_\R\int
     \kappa(x)\kappa(x+h)|F(x,h)|^2\,J(dh)\,dx,
 \qquad df(x,h)=f(x+h)-f(x).
\]
The difference map $d$ with domain $\D$ is closed, by the same almost-everywhere and Fatou argument used for the energy; its norm squared is $\E[f]$. The closed-form representation gives $A=d^*d$. For $f\in\CE$, invariance implies $Af\in\D$, so $df\in\dom(dd^*)$ and
\begin{equation}\label{b:eq:edge-defect}
 \Def[f]=\ip{df}{(dd^*-\tfrac12I)df}_{\mathrm{edge}}.
\end{equation}
This is a domain-correct reformulation, not a positivity proof. Corollary~\ref{b:cor:paths} shows that imposing the finite cycle relations alone cannot produce a nonzero positive finite-path density for this defect.

The obstruction permits genuinely global integral or frequency-space rows, nonlocal inverses, nonlinear constructions, and signed local identities with cancellation after integration. It also leaves open estimates using additional information from an actual candidate eigen-equation. No such method is proved positive here. In particular the explicit negative pointwise tails and the positive compact integrated witness establish neither a negative integrated direction nor a global half gap. The integrated half inequality and the Riemann hypothesis remain unresolved by the results of this paper.


\section{Weighted synthesis and full multiplicity zero jets}\label{sec:synthesis-unified}
The local density just proved is not weighted completeness. This section determines the weighted derivative closure through the exact theta transform, without a zero-location or simplicity hypothesis. The proofs of the actual Abel bound and its source recovery are included in Appendix~\ref{app:abel}; compact division and root-space structure are recorded in Appendix~\ref{app:compactmodel}.
Put $a=i/2$ and define
\begin{equation}\label{d:eq:center-intro}
 m(u)=F_u(a),\qquad Cu=u-2m(u)\kappa,\qquad \Hh_c=\ker m.
\end{equation}
Thus $C$ is exactly the weighted centering projection used earlier. Set
\begin{equation}\label{d:eq:N-intro}
 \Nn=\cl_{\Hh}\{r_j:j\ge1\},\qquad \Kk=\Hh_c\ominus\Nn,
 \qquad r_j=\kappa^{(2j)}-4^{-j}\kappa.
\end{equation}
If $\nu(w)$ is the full multiplicity of a zero $w$ of $\Xi$, define
\begin{align}
 \Ii&=\{u\in\Hh_c:F_u=\Xi Q\text{ for an even entire }Q\},\label{d:eq:I-intro}\\
 \Zz&=\cl_{\Hh}\{W(x)\partial_z^k\cos(zx)|_{z=w}:\Xi(w)=0,\ 0\le k<\nu(w)\}.
 \label{d:eq:Z-intro}
\end{align}
The algebraic span in the last line is denoted by $\mathcal J_{\mathrm{fin}}$. Its members are full zero jets, not only zero-value representatives.
The exact defect decomposition, proved below, is
\begin{equation}\label{d:eq:main-defect}
 \Kk=\Zz\oplus(\Ii\cap\Nn^\perp).
\end{equation}
The open translation lens is
\begin{equation}\label{d:eq:lens-intro}
 \Omega=\{b:|\operatorname{Im}b|<\pi/4,\quad
 2e^{-2|\operatorname{Re}b|}\cos(2\operatorname{Im}b)>1\}.
\end{equation}


\subsection{The full ideal and its annihilator}\label{d:sec:ideal}

The Fourier image $\Aa=\{F_u:u\in\Hh\}$ is a Hilbert space with the transported norm $\norm{F_u}_{\Aa}=\norm{u}_{\Hh}$. Injectivity makes this definition unambiguous. Equation~\eqref{d:eq:eval-bound} makes every derivative evaluation continuous. Since
\[
 \ip{u}{\kappa}_{\Hh}=2m(u),\qquad \norm{\kappa}_{\Hh}^2=1,
\]
$C$ is the orthogonal projection asserted above. Integration by parts, justified by the theta tails, gives
\begin{equation}\label{d:eq:derivative-transforms}
 F_{\kappa^{(2j)}}(z)=(-z^2)^j\Xi(z),\qquad
 F_{r_j}(z)=\bigl((-z^2)^j-4^{-j}\bigr)\Xi(z).
\end{equation}
Every derivative lies in $\Hh$, again by the tails.

It is useful to distinguish the raw wave from its conjugated Riesz representative:
\begin{equation}\label{d:eq:waves}
 J_{w,k}(x)=W(x)\partial_z^k\cos(zx)|_{z=w},\qquad
 E_{w,k}=J_{\overline w,k}.
\end{equation}
They belong to $\Hh$ for every fixed $w,k$, and the linear-first convention gives
\begin{equation}\label{d:eq:Riesz}
 F_u^{(k)}(w)=\ip{u}{E_{w,k}}_{\Hh}.
\end{equation}
The raw and Riesz families have the same closed span when $w$ ranges over all zeros. For $k<\nu(w)$,
\begin{align}
 m(J_{w,k})&=\tfrac12\Xi^{(k)}(w)=0,\label{d:eq:wave-center}\\
 \ip{J_{w,k}}{r_j}_{\Hh}
 &=\left.\partial_z^k\left[\bigl((-z^2)^j-4^{-j}\bigr)\Xi(z)\right]\right|_{z=w}=0.
 \label{d:eq:wave-radical}
\end{align}
Thus $\Zz\subset\Kk$.

\begin{proposition}[Closed ideal and exact defect]\label{d:prop:ideal}
The space $\Ii$ is closed and
\begin{equation}\label{d:eq:I-jets}
 \Ii=\{u\in\Hh_c:F_u^{(k)}(w)=0\text{ whenever }\Xi(w)=0,\ 0\le k<\nu(w)\}
 =\Hh_c\cap\Zz^\perp.
\end{equation}
It contains $\Nn$. For $u\in\Ii$, the entire quotient $Q=F_u/\Xi$ is even and satisfies $Q(a)=0$. Conversely these conditions on a quotient characterize $\Ii$. In particular~\eqref{d:eq:main-defect} holds and
\begin{equation}\label{d:eq:equivalence}
 \Zz=\Kk\quad\Longleftrightarrow\quad\Ii=\Nn.
\end{equation}
\end{proposition}
\begin{proof}
At each zero of $\Xi$, the conditions in~\eqref{d:eq:I-jets} are exactly the removable-singularity conditions for $F_u/\Xi$. This gives an entire quotient; its evenness follows off the zero set and then everywhere by continuation. The converse is immediate. Since $\Xi(a)=1/2$, centering is equivalent to $Q(a)=0$. The bounded evaluations make the intersection of all the specified kernels closed. Equations~\eqref{d:eq:Riesz} and~\eqref{d:eq:derivative-transforms} prove the orthogonal-complement description and $\Nn\subset\Ii$.

In $\Hh_c$, the decomposition $\Hh_c=\Zz\oplus\Ii$ is orthogonal. As $\Nn\subset\Ii$, taking its orthogonal complement within $\Hh_c$ gives~\eqref{d:eq:main-defect} and~\eqref{d:eq:equivalence}.
\end{proof}

This is full-multiplicity divisibility, not merely vanishing of the values at distinct zeros. Nothing in Proposition~\ref{d:prop:ideal} supplies a growth bound on $Q$. The full approximation statement proved in Theorem~\ref{d:thm:full-synthesis} is
\begin{equation}\label{d:eq:ideal-synthesis}
 \overline{\{\Xi p:p\text{ even polynomial},\ p(a)=0\}}^{\Aa}
 =\{\Xi Q\in\Aa:Q\text{ even entire},\ Q(a)=0\}.
\end{equation}

\subsection{The exact complex translation lens}\label{d:sec:translations}

For $|\operatorname{Im}b|<\pi/4$ define
\begin{equation}\label{d:eq:H-R}
 H_b(x)=\frac{\kappa(x+b)+\kappa(x-b)}2,\qquad
 R_b(x)=H_b(x)-\cosh(b/2)\kappa(x).
\end{equation}
These are even functions of $x$. Membership will restrict the parameter to the lens~\eqref{d:eq:lens-intro}.

\begin{theorem}[Exact admission and holomorphy]\label{d:thm:lens}
Within the continuation strip $|\operatorname{Im}b|<\pi/4$, each of the conditions $H_b\in\Hh$ and $R_b\in\Hh$ holds if and only if $b\in\Omega$. On $\Omega$ the maps $b\mapsto H_b$ and $b\mapsto R_b$ are $\Hh$-valued holomorphic. They satisfy
\begin{equation}\label{d:eq:translation-transform}
 F_{H_b}(z)=\cos(bz)\Xi(z),\qquad
 F_{R_b}(z)=\bigl(\cos(bz)-\cosh(b/2)\bigr)\Xi(z).
\end{equation}
In particular $R_b\in\Ii$.
\end{theorem}
\begin{proof}
Write $b=u+iv$. On a compact subset of $\Omega$,
$2e^{-2|u|}\cos(2v)-1$ has a positive lower bound. The theta upper bounds and the lower bound for $W$ show that $|H_b|^2/W$ and the corresponding expressions for every fixed parameter derivative are bounded by
\begin{equation}\label{d:eq:lens-majorant}
 C\exp(C'|x|-\epsilon e^{2|x|}),\qquad \epsilon>0,
\end{equation}
with constants uniform on that compact set. Differentiation under the integral and dominated difference quotients prove $\Hh$-holomorphy. Contour shifting within the theta strip is justified by the same rapid tails. Applying it first to the exponential Fourier transform and then using evenness gives~\eqref{d:eq:translation-transform}. Evaluation at $a$ proves centering.

For the converse, first suppose $u\ne0$ and set $A=e^{-2|u|}\cos(2v)>0$. In~\eqref{d:eq:H-R}, one summand has strictly slower exponential decay as $x\to+\infty$ than the other. Equation~\eqref{d:eq:complex-asymptotic} therefore gives
\begin{equation}\label{d:eq:nonadmission-u}
 \frac{|H_b(x)|^2}{W(x)}
 \sim C_b e^{5x-\pi(2A-1)e^{2x}},\qquad C_b>0.
\end{equation}
If $2A\le1$, the right side is not integrable. This proves nonadmission when the real part is nonzero.

If $u=0$, the leading terms have equal magnitude and must be treated together. Realness of the theta series gives
\begin{equation}\label{d:eq:oscillatory-tail}
 H_{iv}(x)=4\pi^2e^{9x/2-\pi\cos(2v)e^{2x}}
 \left[\cos\bigl(9v/2-\pi\sin(2v)e^{2x}\bigr)+O(e^{-2x})\right].
\end{equation}
Outside $\Omega$ but within the strip, $|v|\ge\pi/6$ and $\sin(2v)\ne0$. Set $t=e^{2x}$. On a fixed positive-length subinterval of every sufficiently late period in $t$, the absolute value of the cosine is at least $3/4$ and the error is at most $1/4$. The weighted integral over these intervals is bounded below by a positive constant times
\begin{equation}\label{d:eq:oscillatory-divergence}
 \int_{\mathcal P} t^{3/2}\exp\bigl(\pi[1-2\cos(2v)]t\bigr)\,dt,
\end{equation}
where $\mathcal P$ is that union of periodic subintervals. This integral diverges even when $\cos(2v)=1/2$. Hence the boundary points $b=\pm i\pi/6$, as well as the remaining excluded imaginary parameters, are not admitted. No pointwise positive lower bound for the oscillatory factor has been used.

Finally $\kappa\in\Hh$, so subtracting its scalar multiple cannot turn a nonmember into a member. This proves the identical admission criterion for $R_b$.
\end{proof}

\begin{theorem}[Translation span and derivative span]\label{d:thm:span}
For every real interval $U$ containing zero in its interior and contained in $\Omega$,
\begin{equation}\label{d:eq:span-equality}
 \Nn=\cl_{\Hh}\{R_b:b\in U\}
     =\cl_{\Hh}\{R_b:b\in\Omega\}.
\end{equation}
\end{theorem}
\begin{proof}
The lens is connected: its defining condition can be written
$2|u|-\log\cos(2v)<\log2$, a convex sublevel condition in $|v|<\pi/4$.
Holomorphic differentiation at zero gives
\[
 \left.\frac{d^{2j}}{db^{2j}}R_b\right|_{b=0}=r_j,
\]
while all odd derivatives and the zeroth derivative vanish. Hilbert-valued finite differences along $U$ therefore place every $r_j$ in the first closed span in~\eqref{d:eq:span-equality}.

Conversely, let a bounded linear functional on $\Hh$ annihilate all the $r_j$. Its value on $R_b$ is holomorphic on $\Omega$ with all Taylor coefficients at zero equal to zero. The scalar identity theorem makes it vanish throughout $\Omega$. Taking annihilators proves the reverse inclusion.
\end{proof}

\begin{lemma}[The nearest boundary]\label{d:lem:radius}
The distance from zero to $\partial\Omega$ is
\begin{equation}\label{d:eq:a-star}
 a_*=\frac{\log2}{2}.
\end{equation}
In particular every closed disk $|b|\le R<a_*$ lies in $\Omega$.
\end{lemma}
\begin{proof}
By symmetry it suffices to consider the first-quadrant boundary. Set $t=2v$, so $2u=\log(2\cos t)$ and $0\le t\le\pi/3$. Four times the squared distance is
$f(t)=\log^2(2\cos t)+t^2$. For $0<t<\pi/3$ its derivative has the sign of
\[
 h(t)=t\cot t-\log(2\cos t).
\]
Here $h(0+)=1-\log2>0$, and
\[
 h'(t)=\frac{\sin t-t\cos t}{\sin^2t\cos t}>0.
\]
Thus $f$ is increasing, and its minimum is at $t=0$.
\end{proof}

\begin{remark}[No global weighted convolution pairing]\label{d:rem:pairing}
The restriction on real translations is a genuine norm issue. Set
$u_*(x)=\sqrt{W(x)}/(1+|x|)$ and $v_*=u_*/W$. Then $u_*\in\Hh$. For any real $b\ge a_*$, the positive-tail product $v_*(x)\kappa(x-b)$ is comparable to
\[
 (1+x)^{-1}\exp\!\left(\tfrac52x-\pi[e^{-2b}-1/2]e^{2x}\right),
\]
so its integral diverges. Replacing $u_*$ by $Cu_*$ changes $v_*$ by a multiple of $2c$; that correction has a finite pairing with the translated theta function and does not remove the divergence. This is not a counterexample to~\eqref{d:eq:ideal-synthesis}: $Cu_*$ is not asserted to annihilate $\Nn$.
\end{remark}

\subsection{Strict type synthesis and the one ray barrier}\label{d:sec:type}

An entire function has exponential type strictly below $a_*$ here if
$|Q(z)|\le M e^{\tau|z|}$ on $\C$ for some $M<\infty$ and $0\le\tau<a_*$. This is a global bound, not merely a bound on one ray.

\begin{theorem}[Strict type partial synthesis]\label{d:thm:type}
Let $Q(z)=\sum_{j\ge0}a_{2j}z^{2j}$ be even entire of exponential type strictly below $a_*$. Then
\begin{equation}\label{d:eq:derivative-series}
 u_Q=\sum_{j\ge0}(-1)^j a_{2j}\kappa^{(2j)}
\end{equation}
converges absolutely in $\Hh$ and satisfies $F_{u_Q}=\Xi Q$. Its centered part belongs to $\Nn$ and is the $\Hh$-limit of
\begin{equation}\label{d:eq:centered-series}
 \sum_{j=1}^J(-1)^j a_{2j}r_j.
\end{equation}
Consequently $Q(a)=0$ implies $u_Q\in\Nn$. In particular every $u\in\Ii$ whose quotient has this strict global type bound belongs to $\Nn$.
\end{theorem}
\begin{proof}
Choose $0<\tau<R<a_*$ with $|Q(z)|\le M e^{\tau|z|}$, increasing an available exponent if necessary. Lemma~\ref{d:lem:radius} and the Banach-valued Cauchy estimate applied to $H_b$ give
\begin{equation}\label{d:eq:Banach-Cauchy}
 \norm{\kappa^{(2j)}}_{\Hh}\le M_R(2j)!R^{-2j},\qquad
 M_R=\sup_{|b|=R}\norm{H_b}_{\Hh}<\infty.
\end{equation}
Scalar Cauchy estimates at radius $n/\tau$, together with the positive-real Stirling estimate~\cite[\S5.11]{DLMF}, give
\begin{equation}\label{d:eq:scalar-Cauchy}
 |a_n|n!\le M'\sqrt{n+1}\,\tau^n\qquad(n\ge0).
\end{equation}
The $n=0$ term is included by enlarging $M'$. Combining the estimates yields a summable bound
$M'M_R\sqrt{2j+1}(\tau/R)^{2j}$ for the norms of the terms in~\eqref{d:eq:derivative-series}.

Bounded transform evaluations and~\eqref{d:eq:derivative-transforms} identify the transform of its sum with $\Xi Q$. Applying $C$ to its partial sums gives~\eqref{d:eq:centered-series}; equivalently their transforms are
$\Xi[P_J(z)-P_J(a)]$, where $P_J=\sum_{j=0}^J a_{2j}z^{2j}$. Since $C$ is bounded, they converge to $Cu_Q$. If $Q(a)=0$, then $Cu_Q=u_Q$. Fourier injectivity gives the last assertion.
\end{proof}

The same constant appears in a necessary bound for every divisible transform, but with no strict margin and only on the imaginary axis.

\begin{proposition}[Imaginary axis growth bound]\label{d:prop:ray}
As $y\to+\infty$,
\begin{equation}\label{d:eq:gamma-saddle}
 \int_\R W(x)\cosh^2(yx)\,dx
 =(1+o(1))\pi^{-y}\Gamma(y+2),
\end{equation}
and
\begin{equation}\label{d:eq:gamma-ratio}
 \frac{\sqrt{\pi^{-y}\Gamma(y+2)}}{\Xi(iy)}
 =(\sqrt2+o(1))\frac{2^{y/2}}{y}.
\end{equation}
Thus there is a finite constant $C$ such that every $u\in\Ii$, with $F_u=\Xi Q$, satisfies
\begin{equation}\label{d:eq:ray-bound}
 |Q(iy)|\le C\norm{u}_{\Hh}\frac{2^{y/2}}{y}
\end{equation}
for all sufficiently large $y$. The same bound holds at $-iy$.
\end{proposition}
\begin{proof}
By evenness and $\cosh^2(yx)=(e^{2yx}+2+e^{-2yx})/4$, the only leading term is
$\frac12\int_0^\infty W(x)e^{2yx}\,dx$. Substitute~\eqref{d:eq:W-asymptotic} and then $t=e^{2x}$. The leading expression becomes
\[
 \pi^2\int_1^\infty t^{y+1}e^{-\pi t}\,dt
 \sim\pi^{-y}\Gamma(y+2).
\]
Here the factor $1/2$ from evenness and $\cosh^2$, the factor $4\pi^2$ in $W$, and $dx=dt/(2t)$ give exactly $\pi^2$. To justify the asymptotic, fix a large $X$. On $x\ge X$ the relative error in~\eqref{d:eq:W-asymptotic} is uniformly small; on $[0,X]$ the integral is $O_X(e^{2yX})$, negligible against $\Gamma(y+2)\pi^{-y}$. The two remaining terms of $\cosh^2$ are also negligible. This proves~\eqref{d:eq:gamma-saddle}.

For $s=y+1/2>1$, Lemma~\ref{d:lem:Fourier} and its symmetry give
\[
 \Xi(iy)=\xi(s)=\frac{s(s-1)}2\pi^{-s/2}\Gamma(s/2)\zeta(s)>0.
\]
The Dirichlet series shows $\zeta(s)\to1$. Positive-real Stirling applied to the two Gamma factors gives~\eqref{d:eq:gamma-ratio}. Finally Cauchy--Schwarz in the defining transform gives
$|F_u(iy)|\le\norm{u}_{\Hh}(\int W\cosh^2(yx))^{1/2}$; division by $\Xi(iy)$ proves~\eqref{d:eq:ray-bound}.
\end{proof}

Equation~\eqref{d:eq:ray-bound} permits growth at the critical exponential rate $a_*$. It neither controls all directions nor characterizes the norm of the quotient. In particular it cannot replace the strict global bound in Theorem~\ref{d:thm:type}. The next section establishes a full directional lens bound, but not a strict circular type bound for every quotient.

\subsection{Global quotient growth and the closed lens}\label{d:sec:global-growth}

The one-ray estimate of Proposition~\ref{d:prop:ray} is only a necessary consequence of membership. We now obtain a global necessary estimate with the exact support of the closed translation lens. Set
\begin{equation}\label{d:eq:closed-lens}
 L=\overline\Omega,\qquad
 h_L(w)=\sup_{b\in L}\operatorname{Re}(bw),\qquad
 H_L(z)=h_L(iz)=\sup_{b\in L}\operatorname{Re}(ibz).
\end{equation}
The bilinear convention in $h_L$, with no hidden conjugation, will also fix the rotation in the Borel representation.

\begin{lemma}[Preliminary order and local quotient bounds]\label{d:lem:quotient-order}
For $u\in\Ii$ and $Q=F_u/\Xi$,
\begin{equation}\label{d:eq:quotient-order}
 \log^+M_Q(R)=O_u(R\log(R+2)),
 \qquad M_Q(R)=\max_{|z|=R}|Q(z)|.
\end{equation}
For every fixed $K<\infty$ there is a constant $C_K$, independent of $u$, such that
\begin{equation}\label{d:eq:quotient-local}
 \sup_{|z|\le K}|Q(z)|\le C_K\norm{u}_{\Hh}.
\end{equation}
\end{lemma}
\begin{proof}
The zero source is trivial. Cauchy--Schwarz, the theta upper bound, and the substitution $t=e^{2|x|}$ bound the squared evaluation norm for $|z|\le R$, $R\ge1$, by a constant times $\pi^{-R}\Gamma(R+2)$. Positive-real Stirling therefore gives
$\log^+M_{F_u}(R)=O_u(R\log(R+2))$. The same holds for $\Xi=F_\kappa$.

For an entire function $f$, let
\[
 m_R^\pm(f)=\frac1{2\pi}\int_0^{2\pi}\log^\pm|f(Re^{it})|\,dt.
\]
Jensen's formula, using $\Xi(0)\ne0$, gives
$m_R^-(\Xi)\le m_R^+(\Xi)-\log|\Xi(0)|$. Zeros on the circle cause only integrable logarithmic singularities; the same inequality follows by a limit through radii avoiding them. Almost everywhere on a circle,
\[
 \log^+|Q|\le\log^+|F_u|+\log^-|\Xi|.
\]
Consequently $m_R^+(Q)=O_u(R\log(R+2))$. Since $Q$ is entire, $\log^+|Q|$ is subharmonic. The Poisson kernel for the disk of radius $2R$ is at most $3$ at points of modulus at most $R$, so its Poisson majorant yields
\[
 \log^+M_Q(R)\le3m_{2R}^+(Q).
\]
This proves~\eqref{d:eq:quotient-order}. No lower bound for $\Xi$ near its moving zeros was used.

For~\eqref{d:eq:quotient-local}, choose one radius $R>K$ whose circle contains no zero of $\Xi$, possible by discreteness. Bounded Fourier evaluation and the positive minimum of $|\Xi|$ on that fixed circle give
$|F_u/\Xi|\le C_K\norm{u}_{\Hh}$ there. The maximum principle for the entire quotient gives the bound inside, including at the removable zero points.
\end{proof}

The order bound~\eqref{d:eq:quotient-order} is not yet finite exponential type. A uniform estimate away from the horizontal zero strip and a sector maximum-principle argument provide the additional information.

\begin{lemma}[Uniform off-strip estimate]\label{d:lem:offstrip}
Write $z=x+iy$, $X=|x|$, $Y=|y|$, and $R=(X^2+Y^2)^{1/2}$. For $Y>0$ set
\begin{equation}\label{d:eq:G-function}
 G(X,Y)=\frac Y2\log\frac{2Y}{R}+\frac X2\arctan\frac XY.
\end{equation}
There is a constant $C$ independent of $u\in\Ii$ such that, uniformly for $Y\ge1$,
\begin{equation}\label{d:eq:offstrip}
 |Q(x+iy)|\le C\norm{u}_{\Hh}Y^{3/4}R^{-7/4}e^{G(X,Y)}
 \le C\norm{u}_{\Hh}R^{-1}e^{G(X,Y)}.
\end{equation}
\end{lemma}
\begin{proof}
Since $|\cos((x+iy)t)|\le\cosh(Yt)$, the numerator estimate follows from the same Gamma integral used in Proposition~\ref{d:prop:ray}:
\begin{equation}\label{d:eq:offstrip-numerator}
 |F_u(x+iy)|\le C\norm{u}_{\Hh}
 Y^{Y/2+3/4}\pi^{-Y/2}e^{-Y/2}\qquad(Y\ge1).
\end{equation}
This is a uniform upper bound, with constants enlarged on bounded $Y$.

For the denominator, the symmetries of $\xi$ permit evaluation of its modulus at $s=v+1/2$, where $v=Y+iX$. Thus
\[
 |\Xi(x+iy)|
 =\left|\frac12(v^2-1/4)\pi^{-v/2-1/4}
       \Gamma(v/2+1/4)\zeta(v+1/2)\right|.
\]
The Dirichlet series and the absolutely convergent Euler product in $\operatorname{Re}s>1$ give
\begin{equation}\label{d:eq:zeta-uniform}
 \zeta(3/2)^{-1}\le|\zeta(s)|\le\zeta(3/2)
 \qquad(\operatorname{Re}s\ge3/2).
\end{equation}
Indeed the upper bound follows from the series, while
$|1/\zeta(s)|\le\prod_p(1+p^{-\operatorname{Re}s})
 \le\prod_p(1-p^{-\operatorname{Re}s})^{-1}
 =\zeta(\operatorname{Re}s)$.
Absolute convergence justifies these products and their nonvanishing; no prime enumeration is involved. These standard formulas are recorded in~\cite[Eq.~25.2.11]{DLMF}.

The uniform complex Stirling formula~\cite[\S5.11]{DLMF} gives the shifted form
\[
 \Gamma(v/2+1/4)
 =\sqrt{2\pi}(v/2)^{v/2-1/4}e^{-v/2}
       \bigl(1+O(R^{-1})\bigr),
\]
uniformly as $R\to\infty$ with $|\arg v|\le\pi/2$. Taking moduli and absorbing the remaining compact region, where the Gamma function is continuous and nonzero, yields
\[
 |\Gamma(v/2+1/4)|\asymp
 R^{-1/4}(R/2)^{Y/2}
 \exp\left(-Y/2-\frac X2\arctan(X/Y)\right).
\]
Also $|v^2-1/4|\asymp R^2$ for $Y\ge1$. Therefore
\begin{equation}\label{d:eq:offstrip-denominator}
 |\Xi(x+iy)|\ge c R^{7/4}
 \left(\frac R{2\pi}\right)^{Y/2}
 \exp\left(-Y/2-\frac X2\arctan(X/Y)\right).
\end{equation}
Dividing~\eqref{d:eq:offstrip-numerator} by this estimate gives~\eqref{d:eq:offstrip}. The final inequality uses $Y\le R$. The reflection and conjugation steps apply to the denominator only; no real-coefficient assumption on $u$ or $Q$ is needed.
\end{proof}

\begin{lemma}[Support function of the lens]\label{d:lem:lens-support}
For $X=|x|$ and $Y=|y|$,
\begin{equation}\label{d:eq:lens-support}
 H_L(x+iy)=
 \begin{cases}
 G(X,Y),&X\le\sqrt3Y,\ Y>0,\\
 \pi X/6,&X\ge\sqrt3Y.
 \end{cases}
\end{equation}
The two formulas agree at their common boundary, and the second includes $Y=0$.
\end{lemma}
\begin{proof}
Coordinate-reflection symmetry of $L$ and its first-quadrant boundary give
\[
 H_L(x+iy)=\max_{0\le\beta\le\pi/6}
 \left[\frac Y2\log(2\cos(2\beta))+X\beta\right].
\]
For $Y>0$, the interior critical point satisfies $\tan(2\beta)=X/Y$. It is available precisely when $X\le\sqrt3Y$, and substitution gives $G(X,Y)$. Otherwise the maximum is attained at $\beta=\pi/6$, where the real coordinate of the lens is zero. The case $Y=0$ follows from the same maximum.
\end{proof}

\begin{theorem}[Global closed-lens bound]\label{d:thm:global-growth}
There is a constant $C$, independent of $u\in\Ii$, such that its entire quotient satisfies
\begin{equation}\label{d:eq:global-growth}
 |Q(z)|\le C\norm{u}_{\Hh}\frac{e^{H_L(z)}}{1+|z|}
 \qquad(z\in\C).
\end{equation}
In particular every such quotient has finite exponential type.
\end{theorem}
\begin{proof}
In the vertical cones $X\le\sqrt3Y$, $R\ge2$ implies $Y\ge1$, so Lemmas~\ref{d:lem:offstrip} and~\ref{d:lem:lens-support} prove the assertion there. Lemma~\ref{d:lem:quotient-order} handles the compact disk.

For the right horizontal sector $|\arg z|\le\pi/6$, put
$f(z)=zQ(z)e^{-\pi z/6}$. On its boundary rays, the same off-strip estimate gives
$|f(z)|\le C\norm{u}_{\Hh}$ outside a fixed disk; the local quotient bound covers the remaining segments. By~\eqref{d:eq:quotient-order},
$\log^+|f(z)|=O_u(R\log(R+2))$.
For $\epsilon>0$ consider
\[
 f_\epsilon(z)=f(z)e^{-\epsilon z^2}.
\]
In this sector $\operatorname{Re}z^2\ge|z|^2/2$. Thus on the outer circular arc of a truncated sector the damping dominates the preliminary growth and the maximum tends to zero as the radius tends to infinity. On the two radial boundaries the damping has modulus at most one, preserving the bound $C\norm{u}_{\Hh}$. The maximum principle and then $\epsilon\downarrow0$ give that bound for $f$ throughout the sector. This proves~\eqref{d:eq:global-growth} there for $|z|\ge1$, because $H_L(z)=\pi|\operatorname{Re}z|/6$. Evenness covers the left horizontal sector, and~\eqref{d:eq:quotient-local} covers the remaining compact disk.
\end{proof}

The horizontal value $\pi|x|/6$ is essential: the raw off-strip expression $G$ is not the lens support near the real axis. The explicit damping argument supplies the needed improvement without a minimum-modulus assumption or an unproved quotient-order assertion.

\subsection{Directional Borel continuation and a sharp common carrier}\label{d:sec:Borel}

Set $P(w)=Q(-iw)$. Then~\eqref{d:eq:global-growth} becomes
\begin{equation}\label{d:eq:P-growth}
 |P(w)|\le C\norm{u}_{\Hh}\frac{e^{h_L(w)}}{1+|w|}.
\end{equation}
For $P(w)=\sum_{n\ge0}p_nw^n$, define its exponential Borel transform initially by
\begin{equation}\label{d:eq:B-series}
 B_P(b)=\sum_{n\ge0}n!p_n b^{-n-1}.
\end{equation}
The series converges outside a sufficiently large disk: scalar Cauchy estimates as in~\eqref{d:eq:scalar-Cauchy} apply to every finite exponential-type bound. A compact convex set is called a \emph{carrier} here when this Borel transform continues as a single-valued holomorphic function to its complement and the contour inversion below holds around that set.

\begin{theorem}[Borel continuation and uniform boundary estimate]\label{d:thm:Borel}
The function $B_P$ has a single-valued holomorphic continuation to $\C\setminus L$. For any positively oriented rectifiable Jordan contour $\Gamma$ in $\C\setminus L$ enclosing $L$ with winding number one,
\begin{equation}\label{d:eq:B-inversion}
 P(w)=\frac1{2\pi i}\int_\Gamma B_P(b)e^{bw}\,db.
\end{equation}
With $\delta=\operatorname{dist}(b,L)>0$,
\begin{equation}\label{d:eq:B-bound}
 |B_P(b)|\le C\norm{u}_{\Hh}
       \int_0^\infty\frac{e^{-\delta t}}{1+t}\,dt
 \le C'\norm{u}_{\Hh}\log(1+\delta^{-1}).
\end{equation}
The set $L$ is the smallest common closed convex carrier for all the rotated quotients $P$. Individual carriers may be smaller.
\end{theorem}
\begin{proof}
For a real direction $\theta$, the integral
\begin{equation}\label{d:eq:directional-Laplace}
 B_\theta(b)=e^{i\theta}\int_0^\infty
     P(te^{i\theta})e^{-bte^{i\theta}}\,dt
\end{equation}
converges and is holomorphic in the half-plane
\begin{equation}\label{d:eq:separation-halfplane}
 \operatorname{Re}(be^{i\theta})>h_L(e^{i\theta}).
\end{equation}
This follows directly from~\eqref{d:eq:P-growth}, with locally uniform majorants for derivatives. Far enough into that half-plane, termwise integration of the power series is absolutely justified and gives~\eqref{d:eq:B-series}.

Convex separation shows that these half-planes cover $\C\setminus L$. To verify compatibility, fix $b$ in an overlap. The separating directions $v$ satisfy
$\operatorname{Re}((b-k)v)>0$ for every $k\in L$. They form an open convex cone with no opposite pair. Two allowable unit directions can therefore be joined by an angular arc of length less than $\pi$ within that cone. On the compact arc the separation gap is uniformly positive. Applying Cauchy's theorem to the integrand $P(w)e^{-bw}$ between the two rays, its outer-arc integral is $O(e^{-cT})$ as the radius $T$ tends to infinity, by~\eqref{d:eq:P-growth}; the inner arc tends to zero with its radius. The two Laplace integrals agree. Thus the local functions patch to one holomorphic continuation.

On a large circle, multiplying~\eqref{d:eq:B-series} by the exponential series and integrating selects $p_nw^n$, proving~\eqref{d:eq:B-inversion}. Cauchy deformation through $\C\setminus L$ then proves it for the stated contours.

At a nearest point of the convex set $L$, choose a unit separating direction whose gap in~\eqref{d:eq:separation-halfplane} equals $\delta$. Formula~\eqref{d:eq:directional-Laplace} gives the first bound in~\eqref{d:eq:B-bound}. For $0<\delta\le1$, split the integral at $1$ and $\delta^{-1}$ to obtain $C(1+\log\delta^{-1})$. For $\delta\ge1$, it is at most $\delta^{-1}$. These are bounded by a constant times $\log(1+\delta^{-1})$, proving the second estimate.

For sharpness, take any nonzero $b_0\in\Omega$. The quotient of $R_{b_0}$ rotates to
$P_{b_0}(w)=\cosh(b_0w)-\cosh(b_0/2)$. Its Borel transform is
\[
 \frac{b}{b^2-b_0^2}-\frac{\cosh(b_0/2)}b,
\]
which has nonzero poles at both $b_0$ and $-b_0$. Any common closed convex carrier must contain all these pairs and therefore their closed convex hull $L$. Conversely the continuation and inversion just proved show that $L$ is a carrier for every quotient. A polynomial quotient can have carrier $\{0\}$, so no individual equality with $L$ is asserted.
\end{proof}

The rotation is important. In the original variable $Q(z)=P(iz)$, the exponential in~\eqref{d:eq:B-inversion} is $e^{ibz}$. Thus the translation parameters lie in $L$, whereas an exponential carrier expressed directly for $Q$ is rotated to $iL$.

\subsection{Compact interior Abel regularization}\label{d:sec:Abel-maps}

For $0<r<1$ the desired centered dilation is
$\Xi(z)[Q(rz)-Q(ra)]$. The Borel representation now proves its actual weighted admission for every $u\in\Ii$, not just for small circular type.

\begin{theorem}[Interior Abel maps]\label{d:thm:Abel-maps}
For every $0<r<1$ there is a unique linear operator $A_r:\Ii\to\Nn$ such that
\begin{equation}\label{d:eq:Abel-transform}
 F_{A_ru}(z)=\Xi(z)\bigl[Q(rz)-Q(ra)\bigr],
 \qquad F_u=\Xi Q.
\end{equation}
It is bounded, compact, and injective, and
\begin{equation}\label{d:eq:Abel-norm}
 \norm{A_r}_{\Ii\to\Hh}
 \le C(1-r)^{-5/4}\log\frac e{1-r}.
\end{equation}
Moreover
\begin{equation}\label{d:eq:Abel-semigroup}
 A_rA_s=A_{rs},\qquad A_rr_j=r^{2j}r_j,\qquad
 \overline{\Ran A_r}^{\Hh}=\Nn.
\end{equation}
The bound~\eqref{d:eq:Abel-norm} is a coarse interior estimate. Theorem~\ref{d:thm:uniform-endpoint} will replace its endpoint loss by a uniform bound using actual weighted admission.
\end{theorem}
\begin{proof}
Let $P_r(w)=P(rw)$, whose Borel transform is
$B_r(b)=r^{-1}B_P(b/r)$ outside $rL$. Put $\rho=(1+r)/2$ and let
$\Gamma_r=\rho\partial L$ with positive orientation. This is a piecewise smooth Jordan contour inside $\Omega$ and outside $rL$. Inversion and the evenness of $P_r$ give
\[
 Q(rz)=\frac1{2\pi i}\int_{\Gamma_r}B_r(b)\cos(bz)\,db.
\]
Define
\begin{equation}\label{d:eq:Abel-Bochner}
 A_ru=\frac1{2\pi i}\int_{\Gamma_r}B_r(b)R_b\,db.
\end{equation}
The contour is compactly contained in $\Omega$, so the integrand is continuous in the $\Hh$ norm and the Bochner integral exists. Every $R_b$ belongs to the closed linear subspace $\Nn$, so the integral does too. Bounded transform evaluations pass under the integral. Formula~\eqref{d:eq:translation-transform} then gives~\eqref{d:eq:Abel-transform}, including its exact centering. Fourier injectivity proves uniqueness and independence of the chosen intermediate contour.

We give quantitative estimates before proving compactness. For $b=u+iv\in\Omega$, put
\[
 \epsilon(b)=2e^{-2|u|}\cos(2v)-1.
\]
Uniform theta estimates on the compact closed lens give
\begin{equation}\label{d:eq:translation-blowup}
 \norm{R_b}_{\Hh}\le C\bigl(1+\epsilon(b)^{-5/4}\bigr).
\end{equation}
Indeed the squared positive-tail integrand is bounded by
$C\exp(5x-\pi\epsilon(b)e^{2x})$. With $t=e^{2x}$ its integral is at most a constant times
$\int_1^\infty t^{3/2}e^{-\pi\epsilon(b)t}\,dt
 \le C'\epsilon(b)^{-5/2}$.
The compact $x$ region and the centering term are uniformly bounded. Taking square roots proves~\eqref{d:eq:translation-blowup}.

The convex function $f(u+iv)=2|u|-\log\cos(2v)$ has $f(0)=0$ and equals $\log2$ on $\partial L$. Therefore, for $b\in\rho\partial L$,
\begin{equation}\label{d:eq:radial-epsilon}
 \epsilon(b)\ge 2^{1-\rho}-1\ge(1-\rho)\log2.
\end{equation}
At $b=\rho b_0$, choose a unit supporting direction at $b_0\in\partial L$. Since $L$ contains the closed disk of radius $a_*$, its support in that direction is at least $a_*$. The separation from $rL$ is consequently at least
\begin{equation}\label{d:eq:contour-separation}
 \delta=\operatorname{dist}(b,rL)\ge a_*(\rho-r).
\end{equation}

Applying the directional Laplace estimate directly to $P_r$ gives
\begin{equation}\label{d:eq:scaled-Borel}
 |B_r(b)|\le C\norm{u}_{\Hh}
       \int_0^\infty\frac{e^{-\delta t}}{1+rt}\,dt.
\end{equation}
For $r\ge1/2$,~\eqref{d:eq:contour-separation} bounds this by
$C\norm{u}_{\Hh}\log(e/(1-r))$. For $r\le1/2$, $\delta\ge a_*/4$, and dropping the denominator gives a uniform bound. This separate estimate avoids a spurious small-$r$ loss from $r^{-1}B_P(b/r)$. The contour lengths are uniformly bounded. Substituting~\eqref{d:eq:translation-blowup}--\eqref{d:eq:scaled-Borel} into~\eqref{d:eq:Abel-Bochner} proves~\eqref{d:eq:Abel-norm}.

For fixed $r$, partition the contour into short arcs and replace $R_b$ on each arc by one representative $R_{b_j}$. The resulting operator has finite rank: its coefficients are integrals of the bounded linear functionals $u\mapsto B_r(b)$. The uniform bound~\eqref{d:eq:scaled-Borel} and uniform $\Hh$ continuity of $R_b$ on the contour make the error tend to zero in operator norm. Thus $A_r$ is compact, not merely pointwise approximated by finite-rank maps.

For the composition law, the quotient of $A_su$ is
$Q_s(z)=Q(sz)-Q(sa)$. Hence
\[
 Q_s(rz)-Q_s(ra)=Q(sr z)-Q(sr a),
\]
and Fourier injectivity gives $A_rA_s=A_{rs}$. If $A_ru=0$, then $Q(rz)$ is constant. Since $r>0$, $Q$ itself is constant, and $Q(a)=0$ forces $u=0$. Finally~\eqref{d:eq:derivative-transforms} gives $A_rr_j=r^{2j}r_j$. The range is contained in $\Nn$ and contains every finite radical direction, proving its stated closure.
\end{proof}

For every $u\in\Ii$,~\eqref{d:eq:Abel-transform} tends compact-open to $F_u$ as $r\uparrow1$. Neither fixed-parameter compactness nor the semigroup law alone upgrades this to boundedness in $\Hh$. The actual-weight argument in Sections~\ref{d:sec:source-recovery}--\ref{d:sec:criterion} supplies the stronger estimate.

\subsection{A sharp boundary family and strong Abel convergence}\label{d:sec:boundary-family}

The strict-type theorem does not exhaust the synthesized quotients. A family whose carriers reach the boundary of $L$ has a sharp weighted admission threshold.

For $b_0\in\partial L$ and real $\alpha>0$, set
\begin{align}
 Q_{\alpha,b_0}(z)
 &=\int_0^1(1-t)^{\alpha-1}
       \bigl[\cos(tb_0z)-\cosh(tb_0/2)\bigr]\,dt,\label{d:eq:boundary-quotient}\\
 u_{\alpha,b_0}(x)
 &=\int_0^1(1-t)^{\alpha-1}R_{tb_0}(x)\,dt.\label{d:eq:boundary-source}
\end{align}
The second formula is initially pointwise, not an assumed $\Hh$ integral. For fixed $x$, the parameter segment lies inside the theta analytic strip, including its endpoint, so it defines a continuous integrand apart from the integrable scalar weight. The theta bounds uniformly on this compact segment also show that $u_{\alpha,b_0}$ is in $L^1(\R,e^{B|x|}dx)$ for every finite $B$. Indeed the slowest real exponential coefficient is $1/2$, still giving double-exponential ordinary decay. Fubini therefore proves
\begin{equation}\label{d:eq:boundary-transform}
 F_{u_{\alpha,b_0}}=\Xi Q_{\alpha,b_0},\qquad m(u_{\alpha,b_0})=0
\end{equation}
even before weighted admission is decided.

\begin{theorem}[Sharp boundary threshold]\label{d:thm:boundary-family}
For every $b_0\in\partial L$ and $\alpha>0$,
\begin{equation}\label{d:eq:boundary-threshold}
 u_{\alpha,b_0}\in\Hh\quad\Longleftrightarrow\quad\alpha>\frac54.
\end{equation}
In the admitted range, $u_{\alpha,b_0}\in\Nn$ and
\begin{equation}\label{d:eq:boundary-Abel}
 A_ru_{\alpha,b_0}\longrightarrow u_{\alpha,b_0}
 \quad\hbox{in }\Hh\quad(r\uparrow1).
\end{equation}
In the rotated variable $P_{\alpha,b_0}(w)=Q_{\alpha,b_0}(-iw)$, the segment $[-b_0,b_0]$ is a carrier and its endpoints are genuine boundary contacts.
\end{theorem}
\begin{proof}
By~\eqref{d:eq:translation-blowup} and~\eqref{d:eq:radial-epsilon},
\begin{equation}\label{d:eq:boundary-majorant}
 \norm{R_{tb_0}}_{\Hh}\le C(1-t)^{-5/4}\qquad(0\le t<1).
\end{equation}
When $\alpha>5/4$, multiplication by $(1-t)^{\alpha-1}$ gives an integrable norm majorant. Thus~\eqref{d:eq:boundary-source} is an $\Hh$ Bochner integral of vectors in $\Nn$, so it belongs to $\Nn$. The centered dilation of~\eqref{d:eq:boundary-quotient} is the same integral with $R_{rtb_0}$ in place of $R_{tb_0}$. Since $1-rt\ge1-t$, the same majorant is valid uniformly for $r\le1$. For fixed $t<1$, holomorphy gives $R_{rtb_0}\to R_{tb_0}$ in $\Hh$. Bochner dominated convergence and~\eqref{d:eq:Abel-transform} prove~\eqref{d:eq:boundary-Abel}.

For necessity, use $R_b=R_{-b}$ to reduce first to $b_0=u+iv$ with $u>0$. Put $q_0=e^{-2b_0}$. The boundary equation implies
\begin{equation}\label{d:eq:boundary-q}
 \operatorname{Re}q_0=\frac12,\qquad
 \operatorname{Re}(b_0q_0)
 =e^{-2u}\bigl[u\cos(2v)+v\sin(2v)\bigr]>0.
\end{equation}
In the branch $\kappa(x-tb_0)$ of~\eqref{d:eq:boundary-source}, set $T=e^{2x}$ and $\delta=1-t$. Uniformly near $\delta=0$,
\[
 e^{-2tb_0}=q_0+2b_0q_0\delta+O(\delta^2),\qquad
 e^{-9tb_0/2}=e^{-9b_0/2}\bigl(1+O(\delta)\bigr).
\]
The real part of $e^{-2tb_0}$ strictly decreases for $0\le t\le1$, since its derivative is
$-2e^{-2tu}[u\cos(2tv)+v\sin(2tv)]<0$.
Thus near the endpoint its increment above $\operatorname{Re}q_0$ is at least $c\delta$, and away from the endpoint it has a fixed positive gap.

Apply~\eqref{d:eq:complex-asymptotic}, split the integral at a small fixed $\delta$, and rescale $\delta=s/T$ in the endpoint piece. Dominated convergence against $s^{\alpha-1}e^{-cs}$ gives
\begin{equation}\label{d:eq:boundary-tail}
 u_{\alpha,b_0}(x)
 =C_{\alpha,b_0}e^{(9/2-2\alpha)x}
       e^{-\pi q_0e^{2x}}\bigl(1+o(1)\bigr),
\end{equation}
where
\begin{equation}\label{d:eq:boundary-constant}
 C_{\alpha,b_0}
 =2\pi^2 e^{-9b_0/2}
      \frac{\Gamma(\alpha)}{(2\pi b_0q_0)^\alpha}\ne0.
\end{equation}
The power is defined in the right half-plane from~\eqref{d:eq:boundary-q}. The Laplace identity
$\int_0^\infty s^{\alpha-1}e^{-As}\,ds
 =\Gamma(\alpha)A^{-\alpha}$ for $\operatorname{Re}A>0$
follows from the real Gamma integral by holomorphy.

No omitted term cancels~\eqref{d:eq:boundary-tail}. For the opposite branch, the real coefficient $e^{2tu}\cos(2tv)$ stays strictly above $1/2$ on the entire compact interval: for $t<1$ this follows from $tb_0\in\Omega$, and at $t=1$ it equals $e^{4u}/2>1/2$. The centered unshifted kernel is faster still. Consequently~\eqref{d:eq:W-asymptotic} yields
\begin{equation}\label{d:eq:boundary-weighted-tail}
 \frac{|u_{\alpha,b_0}(x)|^2}{W(x)}
 \sim c_{\alpha,b_0}e^{(5-4\alpha)x},
 \qquad c_{\alpha,b_0}>0.
\end{equation}
This is not integrable for $\alpha\le5/4$, including equality.

It remains to treat $b_0=\pm i\pi/6$. The two endpoint branches now have conjugate coefficients of equal magnitude. The same rescaling is valid because $\operatorname{Re}(b_0e^{-2b_0})>0$. Their sum has the form
\begin{equation}\label{d:eq:boundary-tip}
 2|C|e^{(9/2-2\alpha)x}e^{-\pi T/2}
 \left[\cos\left(\frac{\sqrt3\pi}{2}T+\phi\right)+o(1)\right],
 \qquad C\ne0.
\end{equation}
This is an additive error inside the bracket, not a relative estimate at the zeros of the cosine. On a fixed positive proportion of each sufficiently late $T$-period the cosine has magnitude at least $1/2$ and the error is at most $1/4$. After division by $W$ and $dx=dT/(2T)$, the weighted integral over these subintervals dominates a constant times
\[
 \int_{\mathcal P} T^{3/2-2\alpha}\,dT.
\]
It diverges for $\alpha\le5/4$, with logarithmic divergence at equality. This proves the threshold at the interference tips as well.

Finally, $P_{\alpha,b_0}$ is an integral of $\cosh(tb_0w)$ and a constant with an integrable scalar weight. Its Borel transform can therefore be represented, outside $[-b_0,b_0]$, by integrating the corresponding rational transforms. This proves that segment is a carrier. In every fixed direction with $\operatorname{Re}(b_0w)>0$, endpoint rescaling gives a nonzero leading term
\[
 \frac{\Gamma(\alpha)}2\,\frac{e^{b_0w}}{(b_0w)^\alpha}
 \quad(|w|\to\infty),
\]
where the opposite exponential and the centered constant are smaller. Symmetry gives the analogous term at the other endpoint. The directional exponential growth therefore attains the support of the segment, so its endpoints cannot be removed from the convex carrier. In particular, taking $b_0=i\pi/6$ gives type $\pi/6>a_*$ in the original variable, beyond the strict-type subclass of Theorem~\ref{d:thm:type}.
\end{proof}

\begin{proposition}[Carrier growth is not an admission criterion]\label{d:prop:sinc-warning}
There exists an even entire $Q_*$ with $Q_*(a)=0$ and
\begin{equation}\label{d:eq:sinc-growth}
 |Q_*(z)|\le C\frac{e^{H_L(z)}}{1+|z|}
\end{equation}
whose rotated carrier lies in $L$, but $\Xi Q_*$ is not the transform of any vector in $\Hh$.
\end{proposition}
\begin{proof}
Take $b_0=a_*$ and $\alpha=1$ in~\eqref{d:eq:boundary-quotient}. Then
\begin{equation}\label{d:eq:sinc-quotient}
 Q_*(z)=\frac{\sin(a_*z)}{a_*z}
        -\frac{\sinh(a_*/2)}{a_*/2},
\end{equation}
with its removable value at zero. Its carrier is the segment $[-a_*,a_*]$ in the rotated variable. The sinc term satisfies~\eqref{d:eq:sinc-growth} by
$|\sin(a_*z)|\le e^{a_*|\operatorname{Im}z|}$ and
$a_*|\operatorname{Im}z|\le H_L(z)$, with a compact-disk adjustment. The constant term is also absorbed: Lemma~\ref{d:lem:radius} implies $H_L(z)\ge a_*|z|$.

The unique ordinary $L^1$ source in~\eqref{d:eq:boundary-source} is not in $\Hh$, since~\eqref{d:eq:boundary-weighted-tail} gives squared weighted tail comparable to $e^x$. If an $\Hh$ source existed, it would also be in $L^1$ and have the same Fourier transform, contradicting Fourier uniqueness.

For each $r<1$, the centered function $Q_*(rz)-Q_*(ra)$ has strict exponential type $ra_*<a_*$, so Theorem~\ref{d:thm:type} supplies an $\Nn$ source. These sources cannot have a bounded sequence as $r\uparrow1$: a weakly convergent subsequence and bounded evaluations would provide the forbidden $\Hh$ inverse of $\Xi Q_*$.
\end{proof}

This example disproves only sufficiency of the carrier and growth condition for weighted admission. It is not an element of $\Ii$ and not a counterexample to $\Ii=\Nn$ or to endpoint boundedness for sources already in $\Ii$.


\subsection{Uniform Abel bounds and full weighted synthesis}\label{d:sec:criterion}

\begin{theorem}[The actual Hilbert endpoint]\label{d:thm:uniform-endpoint}
The centered Abel operators satisfy
\begin{equation}\label{d:eq:uniform-endpoint}
 \sup_{3/4\le r<1}\norm{A_r}_{\Ii\to\Hh}<\infty.
\end{equation}
\end{theorem}
\begin{proof}
For $u\in\Ii$, use its fixed scalar functional $\mu$ to define the uncentered source
\[
 v_r(x)=\mu(\kappa(x+rb)).
\]
The parameter set $r\Gamma$, $3/4\le r<1$, remains compactly inside the theta strip. Hence $v_r$ is even and has all ordinary exponential moments, by the same argument as Lemma~\ref{d:lem:scalar-source}. Its transform is $\Xi(z)Q(rz)$ and its leading Laplace source is exactly $f_r$ from~\eqref{d:eq:fr-definition}. The exact theta identity~\eqref{d:eq:U-H-identity} applies with $f_r$.

Theorem~\ref{d:thm:weighted-kernel}, Theorem~\ref{d:thm:recovery}, and Lemma~\ref{d:lem:forward-reconstruction} give
\begin{equation}\label{d:eq:vr-bound}
 \norm{v_r}_{\Hh}
 \le C\norm{f_r}_{E_\pi}
 \le C'\norm{f}_E
 \le C''\norm{u}_{\Hh}.
\end{equation}
In particular this proves actual weighted admission of $v_r$; it was not inferred from a formal transform or from an exterior $\Hh$-valued contour integral.

Centering is exact:
\[
 2m(v_r)=Q(ra),\qquad Cv_r=v_r-Q(ra)\kappa.
\]
The scalar coefficient is bounded uniformly by $C\norm{u}_{\Hh}$ from~\eqref{d:eq:quotient-local}, since $|ra|\le1/2$, and $\norm{\kappa}_{\Hh}=1$. Thus~\eqref{d:eq:vr-bound} pays for the centering correction and bounds $Cv_r$ uniformly. Its transform is precisely~\eqref{d:eq:Abel-transform}, so Fourier injectivity identifies it with the already admitted $A_ru\in\Nn$. This proves~\eqref{d:eq:uniform-endpoint}.

The argument used actual membership $u\in\Hh$ in~\eqref{d:eq:actual-weight-comparison} and Theorem~\ref{d:thm:recovery}. It assumed neither positivity of $u$ nor a positive measure representation of its Borel functional. Only the auxiliary fractional-power kernel is positive.
\end{proof}

\begin{theorem}[Full multiplicity weighted synthesis]\label{d:thm:full-synthesis}
The full divisible ideal equals the derivative span closure:
\begin{equation}\label{d:eq:full-synthesis}
 \Ii=\Nn.
\end{equation}
For every $u\in\Ii$,
\begin{equation}\label{d:eq:strong-Abel}
 A_ru\longrightarrow u\quad\hbox{strongly in }\Hh
 \qquad(r\uparrow1).
\end{equation}
Consequently
\begin{equation}\label{d:eq:zerojet-completeness}
 \Kk=\Zz
 =\cl_{\Hh}\{J_{w,k}:\Xi(w)=0,\ 0\le k<\nu(w)\}.
\end{equation}
Equivalently, the generalized nullspace of $T|_{\Ii}$ is dense in $\Ii$, and its quotient $\Ii/\Nn$ is zero.
\end{theorem}
\begin{proof}
Fix $u\in\Ii$ and any sequence $r_n\uparrow1$ with $r_n\ge3/4$. Theorem~\ref{d:thm:uniform-endpoint} makes $A_{r_n}u$ bounded in $\Hh$. A subsequence has a weak limit $v$. Since $\Nn$ is a closed linear subspace, it is weakly closed, so $v\in\Nn$. Bounded transform evaluations and~\eqref{d:eq:Abel-transform} give
\[
 F_v(z)=\lim_n\Xi(z)[Q(r_nz)-Q(r_na)]
       =\Xi(z)[Q(z)-Q(a)]=F_u(z).
\]
Fourier injectivity yields $v=u$, so $\Ii\subset\Nn$. The reverse inclusion was proved in Proposition~\ref{d:prop:ideal}; this establishes~\eqref{d:eq:full-synthesis}.

On a finite radical combination $p$, the identity $A_rr_j=r^{2j}r_j$ gives $A_rp\to p$ in $\Hh$. This finite span is dense in $\Ii=\Nn$. With the uniform bound already proved,
\[
 \norm{A_ru-u}_{\Hh}
 \le C\norm{u-p}_{\Hh}
       +\norm{A_rp-p}_{\Hh}+\norm{p-u}_{\Hh}.
\]
First let $r\uparrow1$, then let $p$ approach $u$. This proves strong convergence, not merely compact-open convergence of transforms.

Taking orthogonal complements in $\Hh_c$ in Proposition~\ref{d:prop:ideal} gives~\eqref{d:eq:zerojet-completeness}. The generalized-nullspace statement follows from Theorem~\ref{d:thm:backward}, and the quotient assertion is immediate.
\end{proof}

The uniform bound precedes both limiting arguments. Neither fixed-$r$ compactness, the semigroup identity, nor a carrier estimate was substituted for it. Proposition~\ref{d:prop:sinc-warning} is not contradicted: its source is outside $\Hh$, so the actual-weight recovery hypothesis fails. The carrier condition is still not a sufficient admission criterion.


\subsection{Stronger-topology duals and the same closed carrier}\label{d:sec:Gamma-dual-carrier}

Let
\begin{equation}\label{d:eq:jet-fin}
 \mathcal J_{\rm fin}=\operatorname{span}\{J_{w,k}:\Xi(w)=0,\ 0\le k<\nu(w)\},
 \qquad \Nn_0=\C\kappa\oplus\Nn.
\end{equation}
The sum is $\Hh$-orthogonal and closed, and $\Kk=\Nn_0^\perp$. Fix either $E={X}$ or $E=\mathcal G_{\mathsf D}$, with $\mathsf D$ fixed as above.

The theta tails and smooth positivity on compact intervals give
\begin{equation}\label{d:eq:W-derivative-bound}
 W^2\le C W,\qquad |W'|^2\le C W\quad\hbox{on }\R.
\end{equation}
For the second inequality, differentiating $W=\kappa/(2c)$ adds at most a fixed exponential factor to the theta derivative bound; dividing its square by the lower bound for $W$ leaves
$C\exp(C'|x|-\pi e^{2|x|})$ outside a compact interval. It is bounded there, and compact positivity handles the rest. A bounded logarithmic derivative of $W$ is neither used nor asserted.

For $z\in\C$ define
\[
 e_z(x)=W(x)\cos(zx),\qquad
 M_y^2=\int_\R W(x)\cosh^2(yx)\,dx\quad(y\ge0).
\]
Since $|\cos(zx)|,|\sin(zx)|\le\cosh(|\operatorname{Im}z|x)$, equations~\eqref{d:eq:W-derivative-bound} and~\eqref{d:eq:Gamma-E-upper} give
\begin{equation}\label{d:eq:dual-evaluation}
 \norm{e_z}_{\Hh}\le M_{|\operatorname{Im}z|},\qquad
 \norm{e_z}_{H^1}\le C(1+|z|)M_{|\operatorname{Im}z|},\qquad
 \norm{e_z}_E\le C_E(1+|z|)M_{|\operatorname{Im}z|}.
\end{equation}
The map $z\mapsto e_z$ is entire into $E$: on compact parameter sets, all fixed parameter derivatives add only fixed powers of $x$, dominated by the same theta tails in $\Hh$ and ordinary $H^1$. Difference quotients and~\eqref{d:eq:Gamma-E-upper} give $E$ holomorphy. In particular every jet in~\eqref{d:eq:jet-fin} belongs to both stronger spaces.

Let $\ell$ be a continuous conjugate-linear functional on $E$ annihilating $\mathcal J_{\rm fin}$, and put $L_\ell=\norm{\ell}_{E^*}$. Define the even tempered distribution
\begin{equation}\label{d:eq:dual-distribution}
 v(\varphi)=\ell(W\overline{\varphi_{\rm even}}),\qquad
 \varphi_{\rm even}(x)=\frac{\varphi(x)+\varphi(-x)}2,
 \quad\varphi\in\mathcal S(\R).
\end{equation}
It is linear in the test function. Multiplication by $W$ maps Schwartz functions continuously into $\Hh\cap H^1$, and~\eqref{d:eq:Gamma-E-upper} gives
$\norm{W\varphi}_E\le C_E(\sup|\varphi|+\sup|\varphi'|)$.
Thus~\eqref{d:eq:dual-distribution} genuinely defines a tempered distribution; no weighted function representative is yet assumed.

Set
\begin{equation}\label{d:eq:dual-entire}
 F(z)=\ell(e_{\overline z}).
\end{equation}
Conjugate linearity and the entire $E$-valued family make $F$ even entire. Its derivatives obey
$F^{(k)}(w)=\ell(J_{\overline w,k})$.
Conjugation invariance of the zero set, with multiplicities, therefore implies
\begin{equation}\label{d:eq:dual-quotient}
 F=\Xi Q,\qquad Q\text{ even entire},\qquad
 |F(z)|\le C_E L_\ell(1+|z|)M_{|\operatorname{Im}z|}.
\end{equation}
There is no centering assumption: $Q(a)$ may be nonzero, and the $\kappa$ direction must be retained in the eventual annihilator.

\begin{theorem}[Closed-lens bound for stronger-topology duals]\label{d:thm:dual-lens}
The quotient in~\eqref{d:eq:dual-quotient} satisfies
\begin{equation}\label{d:eq:dual-lens}
 |Q(z)|\le C_E L_\ell e^{H_L(z)}\qquad(z\in\C).
\end{equation}
This has the same exponential carrier as~\eqref{d:eq:global-growth}, but not its reciprocal-linear prefactor. No weighted admission of $v$ is used.
\end{theorem}
\begin{proof}
The zero functional is trivial. The theta evaluation bound and~\eqref{d:eq:dual-quotient} give
$\log^+M_F(R)=O_\ell(R\log(R+2))$, with the same bound for $\Xi$.
Jensen, using $\Xi(0)\ne0$, bounds the circle average of $\log^-|\Xi|$ by its positive average minus $\log|\Xi(0)|$. Since
$\log^+|Q|\le\log^+|F|+\log^-|\Xi|$, Poisson majorization from radius $2R$ to $R$ yields
\begin{equation}\label{d:eq:dual-order}
 \log^+M_Q(R)=O_\ell(R\log(R+2)).
\end{equation}
Circle zeros have integrable logarithmic singularities; limiting radii avoiding them give the same result. On any fixed compact disk, choose a larger zero-free circle for $\Xi$ and divide~\eqref{d:eq:dual-quotient} by its positive minimum there. The maximum principle gives
\begin{equation}\label{d:eq:dual-local}
 \sup_{|z|\le R_0}|Q(z)|\le C_{E,R_0}L_\ell.
\end{equation}

For $z=x+iy$, write $X=|x|$, $Y=|y|\ge1$, and $R=(X^2+Y^2)^{1/2}$. The Gamma integral used in Proposition~\ref{d:prop:ray} gives
\[
 M_Y\le C Y^{Y/2+3/4}\pi^{-Y/2}e^{-Y/2}.
\]
The same symmetry, reciprocal Euler product in $\operatorname{Re}s\ge3/2$, and uniform complex Stirling calculation used in~\eqref{d:eq:offstrip-denominator} give
\[
 |\Xi(z)|\ge cR^{7/4}\left(\frac R{2\pi}\right)^{Y/2}
       \exp\left(-\frac Y2-\frac X2\arctan\frac XY\right).
\]
These estimates do not require an $\Hh$ representative for $F$; the numerator bound needed here is exactly~\eqref{d:eq:dual-quotient}. With $G(X,Y)$ defined in~\eqref{d:eq:G-function}, division gives
\[
 |Q(z)|\le C_E L_\ell(1+R)Y^{3/4}R^{-7/4}e^{G(X,Y)}
 \le C'_E L_\ell e^{G(X,Y)}.
\]
The last inequality uses $Y\le R$ and $(1+R)/R\le2$. Thus the single derivative loss removes the reciprocal-linear gain, without enlarging the exponential support.

The support calculation in Lemma~\ref{d:lem:lens-support} gives $H_L=G$ in the vertical cones and $H_L(z)=\pi|\operatorname{Re}z|/6$ in the horizontal cones. For the right sector $|\arg z|\le\pi/6$, put $f(z)=Q(z)e^{-\pi z/6}$. The preceding estimate and~\eqref{d:eq:dual-local} bound $f$ on both radial boundaries by $C_E L_\ell$. Apply the maximum principle on truncated sectors to $f(z)e^{-\epsilon z^2}$. There $\operatorname{Re}z^2\ge|z|^2/2$, so~\eqref{d:eq:dual-order} makes the outer circular boundary tend to zero. On the rays the damping has modulus at most one. Let the outer radius tend to infinity and then $\epsilon\downarrow0$. Evenness covers the left sector, and~\eqref{d:eq:dual-local} covers the compact remainder. This proves~\eqref{d:eq:dual-lens} with a constant linear in $L_\ell$.
\end{proof}

\begin{lemma}[Ordinary Sobolev source from the exterior contour]\label{d:lem:dual-source}
There is a smooth even ordinary $H^1$ function $v_0$, with all fixed ordinary exponential moments, such that
\begin{equation}\label{d:eq:dual-source}
 F_{v_0}=F=\Xi Q,\qquad
 \norm{v_0}_{H^1(\R)}\le C_E L_\ell.
\end{equation}
This conclusion initially makes no assertion that $v_0\in\Hh$.
\end{lemma}
\begin{proof}
Set $P(w)=Q(-iw)$, so~\eqref{d:eq:dual-lens} gives
$|P(w)|\le C_E L_\ell e^{h_L(w)}$.
The exponential Borel series $B_\ell(b)=\sum_{n\ge0}n!p_nb^{-n-1}$ converges outside a disk by scalar Cauchy estimates. Its directional continuation is
\[
 e^{i\theta}\int_0^\infty
       P(te^{i\theta})e^{-bte^{i\theta}}\,dt,
 \qquad\operatorname{Re}(be^{i\theta})>h_L(e^{i\theta}).
\]
Expansion sufficiently far out recovers the series. The half-planes cover $\C\setminus L$ by convex separation. At a point in an overlap, separating directions form an open convex cone without opposite rays; a compact angular arc joining two such directions has a uniform positive gap. Contour rotation is therefore justified by an outer-arc bound $O(e^{-\delta T})$, and the inner arc tends to zero. The local continuations patch to a single holomorphic function. Choosing a normal at a nearest point of $L$ gives
\begin{equation}\label{d:eq:dual-Borel}
 |B_\ell(b)|\le\frac{C_E L_\ell}{\operatorname{dist}(b,L)}
 \qquad(b\notin L).
\end{equation}
The inverse-distance bound replaces the earlier logarithmic one; a fixed exterior contour only requires its finiteness.

Choose a fixed positively oriented rectifiable Jordan contour $\Gamma\subset\C\setminus L$ with $L\subset\operatorname{int}\Gamma$, winding number one, and contour and interior contained in the theta strip. Series inversion on a large circle, followed by deformation in $\C\setminus L$, yields
\[
 P(w)=\frac1{2\pi i}\int_\Gamma B_\ell(b)e^{bw}\,db.
\]
Define
\begin{equation}\label{d:eq:dual-contour-source}
 v_0(x)=\frac1{2\pi i}\int_\Gamma B_\ell(b)H_b(x)\,db,
 \qquad H_b(x)=\frac{\kappa(x+b)+\kappa(x-b)}2.
\end{equation}
The theta tails and their fixed derivatives are uniform on this compact parameter contour. Consequently this is an ordinary $H^1$-valued Bochner integral, with norm bounded as in~\eqref{d:eq:dual-source}; it is smooth, even, and has all ordinary exponential moments. It is not an $\Hh$-valued contour integral. Ordinary Fourier interchange and contour shifting in the strip give $F_{H_b}(z)=\Xi(z)\cos(bz)$. Since $P$ is even, integrating $\cos(bz)$ gives $P(iz)=Q(z)$, proving the transform identity in~\eqref{d:eq:dual-source}.
\end{proof}

\subsection{Dual recovery and full zero-jet closure in both topologies}\label{d:sec:Gamma-jet-core}

The ordinary source must now be identified with the dual distribution, and only then admitted to the actual weighted space.

\begin{lemma}[Distribution identity and weighted dual recovery]\label{d:lem:dual-recovery}
The distribution $v$ in~\eqref{d:eq:dual-distribution} equals the ordinary function $v_0$ in~\eqref{d:eq:dual-contour-source}. Moreover
\begin{equation}\label{d:eq:dual-H-recovery}
 v_0\in\Hh,\qquad \norm{v_0}_{\Hh}\le C_E L_\ell,\qquad
 \ell(u)=\ip{v_0}{u}_{\Hh}\quad(u\in E).
\end{equation}
\end{lemma}
\begin{proof}
For this test-level identification write the nonunitary Fourier transform explicitly as
$\widehat\varphi_0(t)=\int_\R\varphi(x)e^{-itx}\,dx$.
For even Schwartz $\varphi$, its conjugate inversion formula multiplied by $W$ is the $E$-valued identity
\[
 W\overline\varphi
 =\frac1{2\pi}\int_\R\overline{\widehat\varphi_0(t)}\,e_t\,dt.
\]
The Bochner integral converges because~\eqref{d:eq:dual-evaluation} on the real axis is $O(1+|t|)$ and $\widehat\varphi_0$ is Schwartz. Its $\Hh$ component identifies it with the displayed pointwise product by ordinary inversion. Applying the conjugate-linear $\ell$ gives
\begin{equation}\label{d:eq:dual-test-Fourier}
 v(\varphi)=\frac1{2\pi}\int_\R\widehat\varphi_0(t)F(t)\,dt.
\end{equation}
The ordinary $L^1$ function $v_0$ has the same identity by Fubini and~\eqref{d:eq:dual-source}. For non-even tests both sides equal their even part. Thus $v=v_0$ as tempered distributions, with no weighted admission premise.

Take the real even compact smooth cutoffs $\chi_R$ above, for $R\ge1$, and define the finite nonnegative number
\[
 b_R=\int_\R\chi_R(x)\frac{|v_0(x)|^2}{W(x)}\,dx.
\]
The compact smooth test $\varphi=\chi_R\overline{v_0}/W$ is legitimate because $v_0$ is smooth and $W$ is smooth and positive on its support. Equation~\eqref{d:eq:dual-distribution}, now identified with $v_0$, gives exactly
\begin{equation}\label{d:eq:dual-cutoff-test}
 b_R=\ell(\chi_Rv_0),\qquad
 \norm{\chi_Rv_0}_{\Hh}^2\le b_R.
\end{equation}
The inequality uses $\chi_R^2\le\chi_R$, not equality of the two cutoff masses. The ordinary estimate already proved in Lemma~\ref{d:lem:dual-source} gives, uniformly in $R\ge1$,
\[
 \norm{\chi_Rv_0}_{H^1}\le C\norm{v_0}_{H^1}\le C_E L_\ell.
\]
Every $\chi_Rv_0$ is compact smooth, so~\eqref{d:eq:Gamma-E-upper} applies before any assertion about weighted admission of $v_0$. Thus
\begin{equation}\label{d:eq:dual-quadratic}
 b_R\le L_\ell\norm{\chi_Rv_0}_E
 \le C_E L_\ell(b_R+L_\ell^2)^{1/2}.
\end{equation}
For the form norm this uses the full Gamma energy bound in~\eqref{d:eq:Gamma-first-order}; for the graph norm it uses the actual operator admission~\eqref{d:eq:Gamma-H1}. The bounded perturbation costs only $\norm{\mathsf D}\sqrt{b_R}$ and is absorbed by the same inequality.

If $L_\ell=0$ the conclusion is immediate. Otherwise divide~\eqref{d:eq:dual-quadratic} by $L_\ell^2$, square, and solve the resulting quadratic inequality to obtain $b_R\le C'_E L_\ell^2$. Since $\chi_R\to1$ pointwise, Fatou's lemma gives $v_0\in\Hh$ and the norm bound in~\eqref{d:eq:dual-H-recovery}. The ordinary $H^1$ estimate was independent of this recovery, so the argument is noncircular.

For any even compact smooth $u$, the test $\varphi=\overline u/W$ in~\eqref{d:eq:dual-distribution} yields
$\ell(u)=\int v_0\overline u/W=\ip{v_0}{u}_{\Hh}$.
The compact smooth form core, or respectively graph core, now extends this equality to every $u\in E$, because both sides are $E$-continuous. This last extension uses the graph core in the graph case, not mere form density.
\end{proof}

\begin{theorem}[Full zero jets in the exact Gamma and ambient graph norms]\label{d:thm:Gamma-jet-core}
For every bounded self-adjoint $\mathsf D$ on $\Hh$,
\begin{align}
 \overline{\mathcal J_{\rm fin}}^{{X}}
   &={X}\cap\Kk,\label{d:eq:Gamma-form-jet-core}\\
 \overline{\mathcal J_{\rm fin}}^{\mathcal G_{\mathsf D}}
   &=\operatorname{Dom}T_\Gamma\cap\Kk.
 \label{d:eq:Gamma-graph-jet-core}
\end{align}
More precisely, for either $E={X}$ or $E=\mathcal G_{\mathsf D}$, its continuous conjugate-linear annihilator is
\begin{equation}\label{d:eq:Gamma-annihilator}
 \{\ell\in E^*: \ell(\mathcal J_{\rm fin})=0\}
 =\{u\mapsto\ip{n_0}{u}_{\Hh}:n_0\in\Nn_0\}.
\end{equation}
The graph closure is in the ambient domain norm of $T_\Gamma-\mathsf D$. It does not assert that $\Kk$ reduces this operator or that its restriction to $\Kk$ is self-adjoint.
\end{theorem}
\begin{proof}
For an annihilating $\ell$, Lemma~\ref{d:lem:dual-recovery} gives an actual $\Hh$ representative $v_0$ with $F_{v_0}=\Xi Q$. Only at this stage apply the Hilbert synthesis theorem. Since $\Xi(a)=1/2$,
\[
 Cv_0=v_0-Q(a)\kappa,\qquad
 F_{Cv_0}(z)=\Xi(z)[Q(z)-Q(a)].
\]
Thus $Cv_0\in\Ii=\Nn$ by Theorem~\ref{d:thm:full-synthesis}, and $v_0\in\Nn_0$. This proves one inclusion in~\eqref{d:eq:Gamma-annihilator}. Conversely, each $n_0\in\Nn_0$ represents a continuous functional on $E$ because $E\hookrightarrow\Hh$. It annihilates every full jet by the $\kappa$ and radical pairings in~\eqref{d:eq:wave-center}--\eqref{d:eq:wave-radical}, followed by $\Hh$ continuity. This proves the reverse inclusion.

The common kernel of the functionals on the right of~\eqref{d:eq:Gamma-annihilator} is $E\cap\Nn_0^\perp=E\cap\Kk$. The closed-linear-span form of Hahn--Banach therefore gives the two closure identities. Explicitly, an element of $E\cap\Kk$ outside the $E$-closure of $\mathcal J_{\rm fin}$ would admit a continuous separating functional; its conjugate-linear form is classified by~\eqref{d:eq:Gamma-annihilator} and must vanish on that element, a contradiction. The reverse inclusion follows from $\mathcal J_{\rm fin}\subset\Kk$ and continuous inclusion $E\hookrightarrow\Hh$.

The representer in~\eqref{d:eq:Gamma-annihilator} is unique because compact smooth vectors make $E$ dense in $\Hh$. Its $\Hh$ norm and its $E$-dual norm are equivalent on this particular annihilator by~\eqref{d:eq:dual-H-recovery} and the continuous inclusion. No equivalence of the stronger and weaker vector norms on $\Kk$ is asserted.
\end{proof}

The theorem is a topology-lifting argument, not an inference from $\Hh$ density alone or from individual jet smoothness. The distributional quotient, its weakened prefactor, the ordinary exterior-contour source, compact weighted recovery, and the appropriate compact smooth core are all needed. It supplies existence of finite-jet approximations in these two precise norms, but no uniform rate, coefficient bound, unconditional basis, finite-section conditioning, or algebraic invariance of the finite jet span.


The preceding graph closure is in the ambient graph norm for an arbitrary bounded self-adjoint perturbation. It does not assert that $\Kk$ reduces such a perturbation. The actual physical operator has additional radical identities, which provide precisely that missing reduction in the next section.


\section{Actual physical reduction and the exact finite sign}\label{sec:actual-jets}
The stronger-topology theorem applies to the actual physical correction $D_b$ from~\eqref{eq:Db-physical}. The conclusion here is more specific than ambient graph density: the radical identities make the actual complementary sector reducing. No corresponding claim is made for an arbitrary bounded self-adjoint perturbation.

\begin{theorem}[Actual reducing sector and full zero jet cores]\label{thm:actual-core}
Let $B=T_\Gamma-D_b=UAU^{-1}-I/2$ on $\Hh$, and let $\Nn$ and $\Kk$ be as in~\eqref{d:eq:N-intro}. Then
\begin{equation}\label{eq:actual-radical}
 \Nn\subset\dom B,\qquad B|_{\Nn}=0,\qquad B\kappa=-\kappa/2.
\end{equation}
The orthogonal decomposition
\begin{equation}\label{eq:actual-decomposition}
 \Hh=\C\kappa\oplus\Nn\oplus\Kk
\end{equation}
reduces $B$. In particular $B_{\Kk}=B|_{\dom B\cap\Kk}$ is self-adjoint in $\Kk$, and
\begin{equation}\label{eq:actual-cores}
 \overline{\mathcal J_{\mathrm{fin}}}^{X}=X\cap\Kk,
 \qquad
 \overline{\mathcal J_{\mathrm{fin}}}^{\,\operatorname{graph}(B)}
       =\dom B\cap\Kk.
\end{equation}
Thus the restriction of $B_{\Kk}$ to the algebraic full-zero-jet span is essentially self-adjoint and has closure $B_{\Kk}$. Its finite span need not be invariant under $B$.
\end{theorem}
\begin{proof}
Every finite radical $r$ belongs to $\dom B$ and satisfies $Br=0$, by Lemma~\ref{a:lem:radical} and the representative criterion. If $r_n\to n\in\Nn$ in $\Hh$, then $(r_n,Br_n)\to(n,0)$ in $\Hh\oplus\Hh$. Closedness of $B$ proves $n\in\dom B$ and $Bn=0$. The source identity $q_0(\kappa,z)=-\overline{m(z)}$ and $\ip\kappa z_{\Hh}=2\overline{m(z)}$ give $B\kappa=-\kappa/2$ by the same criterion.

For $u\in\dom B$, self-adjointness gives
\[
 \ip{Bu}{n}_{\Hh}=0\quad(n\in\Nn),\qquad
 \ip{Bu}{\kappa}_{\Hh}=-\tfrac12\ip u\kappa_{\Hh}.
\]
The projections onto $\Nn$ and $\C\kappa$ map $\dom B$ into itself, since both target spaces are in the domain, and commute with $B$ by these identities. Their complementary projection onto $\Kk$ therefore also preserves the domain and commutes with $B$. This proves reduction and self-adjointness of the part.

We also record form-domain preservation explicitly. For $n\in\Nn$,
\[
 \tg[n]=\ip{D_bn}{n}_{\Hh}\le\norm{D_b}\norm n_{\Hh}^2,
 \qquad \norm n_X\le\sqrt{1+\norm{D_b}}\norm n_{\Hh}.
\]
Hence $P_{\Nn}:\Hh\to X$ is bounded. The one-dimensional projection onto $\kappa\in X$ is bounded into $X$ as well, so $P_{\Kk}=I-P_{\Nn}-P_{\kappa}$ preserves $X$ continuously. This does not assert that $T_\Gamma$ alone reduces these spaces.

Theorem~\ref{d:thm:Gamma-jet-core}, with $\mathsf D=D_b$, gives both closures in~\eqref{eq:actual-cores}; its hypotheses hold by the common Gamma and prime-domain construction. Every jet is in the actual operator domain by Lemma~\ref{d:lem:Gamma-H1}: the theta derivative estimates give $W\partial_z^k\cos(zx)|_{z=w}\in\Hh\cap H^1(\R)$. The graph closure is therefore an operator core for the self-adjoint part $B_{\Kk}$. This proves essential self-adjointness of its symmetric core restriction, without assuming invariance of the finite span.
\end{proof}

Under $U^{-1}$, the finite profile family is
\begin{equation}\label{eq:actual-profile-jets}
 \frac{J_{w,k}(x)}{\kappa(x)}
 =\frac{\partial_z^k\cos(zx)|_{z=w}}{2\cosh(x/2)},
 \qquad \Xi(w)=0,\quad 0\le k<\nu(w).
\end{equation}
It is a form core and an $A$-graph core for the reducing profile sector $U^{-1}\Kk$. Unitarity and the bounded graph shift by $I/2$ preserve these closure statements. No backward-division operator has been commuted through $B$ or through the prime and Gamma rows.

\subsection{The exact Hermitian matrices}
For a finite list of actual zero jets indexed by $\alpha=(w,k)$, put $j_\alpha=J_{w,k}$ and define
\begin{align}\label{eq:actual-matrix}
 M_{\alpha\beta}
 ={}&\int_0^\infty\rho(s)\int_\R
 [j_\alpha(x+s)-j_\alpha(x)]
 \overline{[j_\beta(x+s)-j_\beta(x)]}\,dx\,ds\notag\\
 &-\cg\int_\R j_\alpha(x)\overline{j_\beta(x)}\,dx\notag\\
 &-\sum_{n\ge2}a_n\left\{\int_\R j_\alpha(x+s_n)\overline{j_\beta(x)}\,dx
       +\int_\R j_\alpha(x)\overline{j_\beta(x+s_n)}\,dx\right\}.
\end{align}
All entries exist. Gamma pairings follow by Cauchy--Schwarz from admitted finite Gamma energies; prime pairings are absolutely convergent by Lemma~\ref{a:lem:prime}. The matrix is Hermitian, and
\begin{equation}\label{eq:actual-quadratic}
 q_0\!\left[\sum_\alpha c_\alpha j_\alpha\right]
 =\sum_{\alpha,\beta}c_\alpha\overline{c_\beta}M_{\alpha\beta}.
\end{equation}
No Gamma shift cutoff, prime head, reality hypothesis on $w$, or simplicity assumption appears in this definition. Redundant columns can make the matrix singular without affecting the sign criterion.

\begin{theorem}[Exact remaining sign equivalence]\label{thm:actual-sign}
The following four assertions are equivalent:
\begin{enumerate}[label=\textup{(\roman*)}]
\item $q_0[u]\ge0$ for every $u\in X_c$.
\item $q_0[u]\ge0$ for every $u\in X\cap\Kk$.
\item Every finite exact full-zero-jet matrix~\eqref{eq:actual-matrix} is positive semidefinite.
\item $\E[f]\ge\tfrac12\Var_\mu(f)$ for every $f\in\D$.
\end{enumerate}
The theorem proves the equivalence, not any one of these assertions.
\end{theorem}
\begin{proof}
The full-domain physical bridge after subtracting the original mean gives (i)$\Leftrightarrow$(iv). For $u\in X_c$, Theorem~\ref{thm:actual-core} gives $u=n+k$ with $n=P_{\Nn}u\in\Nn$ and $k=P_{\Kk}u\in X\cap\Kk$. Since $Bn=0$, the operator-to-form identity yields $q_0(n,z)=\ip{Bn}{z}_{\Hh}=0$ for every $z\in X$. Thus $q_0[u]=q_0[k]$, proving (i)$\Leftrightarrow$(ii).

The form $q_0$ is continuous on $X\times X$: the Gamma term is controlled by Cauchy--Schwarz in its form norm, and the complete correction is bounded on $\Hh$. The form-core equality in~\eqref{eq:actual-cores}, together with~\eqref{eq:actual-quadratic}, therefore proves (ii)$\Leftrightarrow$(iii). This passage is in the exact form norm, not merely the Hilbert norm.
\end{proof}

In particular, if a centered negative-form vector exists, some finite full-zero-jet combination already has strictly negative $q_0$. The same is true if an actual centered eigenstate of $A_E$ has eigenvalue in $(0,1/2)$, by the bridge. This follows from form convergence and strict negativity; it supplies no bound on the number, height, multiplicity or conditioning of the necessary jets. The constant ground state is not in the centered sector. Establishing all the finite signs remains open.


\section{Uniform coercivity at fixed exponential complexity}\label{sec:complexity-unified}
The finite-jet sign equivalence concerns every finite complexity. A complementary analytic result controls the high-frequency region when that complexity is fixed. We state it for general rapidly decreasing positive weights and then specialize to the same theta form. In this section and its cluster appendix only, $\HH=L^2(\R,dx/W)$ is the full space: neither weights nor functions need be even unless this is stated. The physical theta specialization then restricts to the even sector. The proof uses compact ODE control through collisions and finite clusters at arbitrary separations, not a lower bound on frequency spacing.


\begin{definition}[Admissible weight]\label{e:def:weight}
A function $W:\R\to(0,\infty)$ is admissible if it is $C^1$ and, for every
$a>0$, there is $C_a<\infty$ such that
\begin{equation}\label{e:eq:weight}
 W(x)+|W'(x)|\le C_a e^{-a|x|}\qquad(x\in\R).
\end{equation}
Set $\HH=L^2(\R,dx/W)$, with
$\norm{u}_W^2=\int_\R |u(x)|^2/W(x)\,dx$.
\end{definition}

The assumptions imply that $W$ is bounded, has a positive minimum on every
compact interval, and tends to zero at both ends. In particular,
\begin{equation}\label{e:eq:embedding}
 \norm{u}_2^2\le \norm W_\infty\norm{u}_W^2\qquad(u\in\HH).
\end{equation}
The converse below is specific to a nonlinear class, not to all of $\HH$.

\begin{definition}[Exponential complexity]\label{e:def:complexity}
For $N\in\N$ and $\beta\ge0$, let $\EE$ consist of zero and all functions
\begin{equation}\label{e:eq:exp-poly}
 P(x)=\sum_{j=1}^r p_j(x)e^{i\lambda_jx},\qquad
 \sum_{j=1}^r(\deg p_j+1)\le N,\qquad |\Ima\lambda_j|\le\beta,
\end{equation}
where the $\lambda_j\in\C$ are distinct and the $p_j$ are nonzero
polynomials. Repeated frequencies are merged first. Let $\EE(T)$ be the
subclass for which $|\Rea\lambda_j|\ge T$ for every displayed frequency.
\end{definition}

The cost of a degree-$k$ polynomial at one frequency is $k+1$, regardless
of how many of its coefficients vanish. Distinct exponential-polynomial
terms are linearly independent, as recalled in Section~\ref{e:sec:cosine},
so the merged nonzero representation is unique. These classes need not be
closed under addition at fixed $N$.

For $s>0$ define
\begin{equation}\label{e:eq:gamma-form}
 \rho(s)=\frac{e^{-s/2}}{1-e^{-2s}},\qquad
 \tg[u]=\int_0^\infty\rho(s)\norm{u(\cdot+s)-u}_2^2\,ds,
 \qquad \XX=\{u\in\HH:\tg[u]<\infty\}.
\end{equation}
The nonnegative integral is interpreted without cancellation. If $D$ is
any bounded self-adjoint operator on $\HH$, set
\begin{equation}\label{e:eq:q}
 \qD(u,v)=\tg(u,v)-\ip{Du}{v}_W,\qquad
 \qD[u]=\qD(u,u).
\end{equation}
Here $\tg(u,v)$ is the polarized difference form. Every $WP$ in
Definition~\ref{e:def:complexity} belongs to $\XX$, as proved below.

\begin{theorem}[Uniform finite-complexity coercivity]\label{e:thm:main}
Fix an admissible weight $W$, $N\in\N$, and $\beta\ge0$.
\begin{enumerate}[label=\textup{(\roman*)},leftmargin=2.1em]
\item There is $C=C(N,\beta,W)<\infty$ such that
\begin{equation}\label{e:eq:comparison}
 \norm{WP}_W^2\le C\norm{WP}_2^2\qquad(P\in\EE).
\end{equation}
\item For every fixed $0\le R<\infty$,
\begin{equation}\label{e:eq:escape}
 \varepsilon_{N,\beta,W}(T;R):=
 \sup_{0\ne P\in\EE(T)}
 \frac{\norm{\mathbf1_{[-R,R]}\widehat{WP}}_2}{\norm{WP}_2}
 \longrightarrow0\qquad(T\to\infty).
\end{equation}
\item For every bounded self-adjoint $D$ on $\HH$,
\begin{equation}\label{e:eq:coercivity}
 \inf_{0\ne P\in\EE(T)}\frac{\qD[WP]}{\norm{WP}_2^2}
 \longrightarrow+\infty,\qquad
 \inf_{0\ne P\in\EE(T)}\frac{\qD[WP]}{\norm{WP}_W^2}
 \longrightarrow+\infty.
\end{equation}
\end{enumerate}
No lower bound on frequency separation or bound on coefficient
conditioning is required. The constants and height thresholds may depend
on $N$, $\beta$, $W$, and, in \textup{(iii)}, $\norm D$.
\end{theorem}

In particular, for every prescribed $L>0$ there is $T_L<\infty$ such that
$\qD[WP]\ge L\norm{WP}_W^2$ for every $P\in\EE(T_L)$.
This is uniform in the represented function norm, not in the Euclidean
norm of raw exponential coefficients. The theorem gives no effective
value of $T_L$ and no uniformity as $N\to\infty$.

Section~\ref{e:sec:gamma} fixes the exact Gamma form and its domain.
Sections~\ref{e:sec:ode}--\ref{e:sec:proof} prove the theorem.
Section~\ref{e:sec:theta} verifies the theta and full-prime specialization
without importing an energy-realization theorem. The remaining sections
record finite cosine-jet consequences and the limitations of the
complexity hypothesis. No historical-priority claim is made.


\subsection{The exact Gamma form}\label{e:sec:gamma}

\begin{lemma}[Multiplier and domain]\label{e:lem:gamma}
For $u\in\HH$,
\begin{equation}\label{e:eq:multiplier}
 \tg[u]=\int_\R d_\Gamma(t)|\widehat u(t)|^2\,dt,
 \qquad
 d_\Gamma(t)=2\int_0^\infty\rho(s)(1-\cos(ts))\,ds.
\end{equation}
The multiplier is nonnegative and satisfies
\begin{equation}\label{e:eq:logbound}
 d_\Gamma(t)\ge\frac12\log(1+4t^2).
\end{equation}
Moreover, for a constant $C_\Gamma<\infty$,
\begin{equation}\label{e:eq:h1bound}
 \tg[u]\le C_\Gamma\bigl(\norm u_2^2+\norm{u'}_2^2\bigr)
 \qquad(u\in H^1(\R)).
\end{equation}
Thus $WP\in\HH\cap H^1(\R)\subset\XX$ for every exponential polynomial
$P$.
\end{lemma}
\begin{proof}
By~\eqref{e:eq:embedding}, $u\in L^2$. Plancherel for each shift and Tonelli
for the nonnegative integrand give~\eqref{e:eq:multiplier}, including when
either side is infinite. Near zero, $\rho(s)=O(s^{-1})$; at infinity,
$\rho(s)=O(e^{-s/2})$. Thus the integral defining $d_\Gamma(t)$ converges
for each $t$. The estimates
\[
 \norm{u(\cdot+s)-u}_2\le s\norm{u'}_2,
 \qquad \norm{u(\cdot+s)-u}_2\le2\norm u_2
\]
prove~\eqref{e:eq:h1bound} by splitting at $s=1$.

Since $1-e^{-2s}\le2s$,
\[
 d_\Gamma(t)\ge I(t):=\int_0^\infty
 \frac{e^{-s/2}(1-\cos(ts))}{s}\,ds.
\]
Differentiation under the integral on compact $t$ intervals gives
$I'(t)=\int_0^\infty e^{-s/2}\sin(ts)\,ds=t/(t^2+1/4)$.
Integrating from $0$ yields $I(t)=\tfrac12\log(1+4t^2)$.
Finally $P$ and $P'$ grow at most like a polynomial times an exponential.
Equation~\eqref{e:eq:weight} makes
$\int W|P|^2$, $\int W^2|P|^2$, and
$\int |W'P+WP'|^2$ finite. This proves admission.
\end{proof}

For reference, the exact digamma normalization is
\begin{equation}\label{e:eq:digamma}
 d_\Gamma(t)=\Rea\psi\left(\frac14+\frac{it}{2}\right)-\psi(1/4),
 \qquad \psi=\Gamma'/\Gamma.
\end{equation}
Indeed, subtract the integral representations
\[
 \psi(z)+\gamma=\int_0^\infty
 \frac{e^{-r}-e^{-zr}}{1-e^{-r}}\,dr\qquad(\Rea z>0)
\]
at $z=1/4+it/2$ and $z=1/4$, take real parts, and set $r=2s$
\cite[Eq.~5.9.16]{DLMF}. This identity names the form; only the directly
proved bound~\eqref{e:eq:logbound} is needed for coercivity.

\begin{proposition}[Closed form]\label{e:prop:closed}
The space $\XX$, with norm
$\norm u_X^2=\norm u_W^2+\tg[u]$, is complete. The form $\tg$ is densely
defined and nonnegative on $\HH$. The form $\qD$ is closed and bounded
below by $-\norm D\norm u_W^2$.
\end{proposition}
\begin{proof}
If $u_n$ is Cauchy in $\norm{\cdot}_X$, it converges to $u$ in $\HH$,
hence in ordinary $L^2$. In Fourier variables,
$d_\Gamma^{1/2}\widehat u_n$ is Cauchy in $L^2$. A subsequence of the two
convergent $L^2$ sequences converges almost everywhere, which identifies
the second limit as $d_\Gamma^{1/2}\widehat u$. Therefore $u\in\XX$ and
$u_n\to u$ in $\norm{\cdot}_X$. Compactly supported smooth functions lie
in $\XX$ by Lemma~\ref{e:lem:gamma} and are dense in $\HH$: truncate first,
then use the positive continuous weight bounds on a compact interval and
ordinary smooth approximation. Finally
$|\ip{Du}{u}_W|\le\norm D\norm u_W^2$ and
\[
 \norm u_X^2\le \qD[u]+(\norm D+1)\norm u_W^2
 \le \tg[u]+(2\norm D+1)\norm u_W^2.
\]
The shifted form norm is equivalent to the complete $X$ norm.
\end{proof}

No ordinary Fourier projection is asserted to preserve $\HH$. Fourier
cutoffs below are used only to estimate integrals in ordinary $L^2$.


The complete proof of Theorem~\ref{e:thm:main} is in Appendix~\ref{app:clusters}. Its constants are qualitative and may depend on the total confluent complexity.


\subsection{Specialization to the actual theta form}\label{e:sec:theta}
Lemmas~\ref{a:lem:theta} and~\ref{d:lem:theta-tails}, together with Leibniz's rule for $W=\kappa/(2c)$, show that $W$ and every fixed derivative decrease faster than every exponential. Thus the theta weight is admissible. On the full space $\HH$, put $K_su(x)=W(x)u(x+s)$ and
\[
 D_\theta=\cg M_W+\sum_{n\ge2}a_n(K_{s_n}+K_{s_n}^*),\qquad
 (K_s^*u)(x)=W(x)u(x-s).
\]
Weighted Cauchy--Schwarz gives $\norm{K_s}\le\sup_x\sqrt{W(x)W(x+s)}$. The complete norm-summable majorant proved in Lemma~\ref{a:lem:prime} did not require evenness, so this is a bounded self-adjoint operator on the full $\HH$. Reflection interchanges $K_s$ and $K_s^*$, hence the sum preserves the even sector and its restriction is $D_b$. Define $\qth=\tg-\ip{D_\theta\,\cdot}{\cdot}_W$. Its even restriction is precisely $q_0$ on $X$; no physical term is discarded.
\begin{corollary}[Complete theta specialization]\label{e:cor:theta}
For each fixed $N$ and $\beta$, both Rayleigh quotients in~\eqref{e:eq:coercivity} for $\qth$ tend uniformly to $+\infty$. After restriction to even functions, this is the same conclusion for the actual form $q_0$.
\end{corollary}
\begin{proof}
Apply Theorem~\ref{e:thm:main} with $D=D_\theta$. Restricting the class cannot decrease the infimum.
\end{proof}


\subsection{Confluent cosine families and finite matrices}\label{e:sec:cosine}

Let $W$ be even and admissible. For $w\in\C$ and $k\ge0$, put
\begin{equation}\label{e:eq:jets}
 J_{w,k}(x)=W(x)\left.\partial_z^k\cos(zx)\right|_{z=w}.
\end{equation}
These are even functions; here the word jet means only a parameter
derivative and does not impose a zero equation. Suppose that
$w_1,\ldots,w_r$ are distinct, $\Rea w_j\ge T>0$,
$|\Ima w_j|\le\beta$, and $m_j\ge1$. Set $d=\sum_jm_j$ and
\begin{equation}\label{e:eq:cosine-combination}
 u=\sum_{j=1}^r\sum_{k=0}^{m_j-1}b_{j,k}J_{w_j,k}.
\end{equation}
The identity
\begin{equation}\label{e:eq:jet-exp}
 \partial_z^k\cos(zx)=\frac12\bigl((ix)^ke^{izx}+(-ix)^ke^{-izx}\bigr)
\end{equation}
shows that $u/W\in\mathcal E_{2d,\beta}(T)$.

\begin{corollary}[Uniform finite cosine matrices]\label{e:cor:cosine}
Fix $N$, $\beta$, and a bounded self-adjoint $D$ on $\HH$. For every
$L>0$ there is $T_N(L)<\infty$ such that every combination
\eqref{e:eq:cosine-combination} with $d\le N$ and
$\min_j\Rea w_j\ge T_N(L)$ satisfies
\begin{equation}\label{e:eq:cosine-coercivity}
 \qD[u]\ge L\norm u_W^2.
\end{equation}
For the nonredundant prefix list displayed above, the form matrix is
positive definite at those heights. More precisely it dominates $L$ times
the weighted Gram matrix as a Hermitian form on the coefficient space.
\end{corollary}
\begin{proof}
The inequality is Theorem~\ref{e:thm:main} at complexity $2N$. To justify
strictness of the coefficient form, recall exponential-polynomial
independence. If $\sum_jp_j(x)e^{i\lambda_jx}=0$ with distinct
$\lambda_j$, apply
$\prod_{h\ne j}(\partial_x-i\lambda_h)^{\deg p_h+1}$.
Every term except the $j$th vanishes. After removing $e^{i\lambda_jx}$,
the remaining operator on $p_j$ is a product of powers of
$\partial_x+i(\lambda_j-\lambda_h)$.
Each factor is invertible on the finite polynomial space: in the monomial
basis its matrix is triangular with a nonzero constant diagonal.
Therefore $p_j=0$ for every $j$.

The frequencies $w_1,\ldots,w_r,-w_1,\ldots,-w_r$ are all distinct
because the first group has positive real parts and the $w_j$ are distinct.
Equation~\eqref{e:eq:jet-exp} and independence imply that the prefix jets
are linearly independent. Hence a nonzero coefficient vector represents a
nonzero $u$, and~\eqref{e:eq:cosine-coercivity} gives $\qD[u]>0$.
\end{proof}

For repeated or redundant columns the same inequality holds, but the
coefficient matrix can be singular. For a sparse selection at $w_j$, the
correct complexity cost is one plus the largest derivative order selected,
not the number of selected columns. Near collisions the unnormalized Gram
matrix can be ill-conditioned; no uniform lower eigenvalue in raw
Euclidean coordinates is asserted.

This corollary applies to arbitrary frequencies meeting its hypotheses.
It neither uses nor proves that a given frequency is a zero of a special
function. In particular, no conclusion about all actual zeta-zero jets,
a centered half inequality, or the Riemann hypothesis follows from it.

\subsection{A fixed width cosine packet}\label{e:sec:packet}

Although the general theorem has no stated rate, a common bounded offset
window permits a useful explicit form of estimate with unevaluated
constants. Fix $R_0\ge0$ and $\beta\ge0$. Let $P_+$ have complexity at most $N$ and frequencies in
$|\Rea\lambda|\le R_0$, $|\Ima\lambda|\le\beta$. Put $A=WP_+$ and
\begin{equation}\label{e:eq:packet}
 u_T(x)=\frac12\bigl(e^{iTx}A(x)+e^{-iTx}A(-x)\bigr),\qquad T>0,
\end{equation}
where $W$ is even. By Lemma~\ref{e:lem:ode}, choose a constant $C\ge1$
depending only on $N,R_0,\beta,W$ such that
$\norm A_W,\norm{A'}_2\le C\norm A_2$.

\begin{proposition}[Carrier estimate]\label{e:prop:packet}
For $T\ge2C$,
\begin{equation}\label{e:eq:packet-bound}
 \qD[u_T]\ge\left\{
 \frac12\log(1+T^2)\left(1-\frac{16C^2}{T^2}\right)
 -4\norm D C^2\right\}\norm{u_T}_2^2.
\end{equation}
\end{proposition}
\begin{proof}
Integration by parts gives
\[
 \left|\int_\R e^{2iTx}A(x)\overline{A(-x)}\,dx\right|
 \le\frac{\norm{A'}_2\norm A_2}{T}
 \le\frac C T\norm A_2^2.
\]
Expanding~\eqref{e:eq:packet} therefore yields
$\norm{u_T}_2^2\ge\norm A_2^2/4$ for $T\ge2C$.
The weighted triangle inequality gives
$\norm{u_T}_W^2\le C^2\norm A_2^2\le4C^2\norm{u_T}_2^2$.
Using~\eqref{e:eq:low-band} with the two carriers $\pm T$ and band
$[-T/2,T/2]$, including their factors $1/2$, gives
\[
 \norm{\mathbf1_{[-T/2,T/2]}\widehat{u_T}}_2
 \le\frac{2\norm{A'}_2}{T}
 \le\frac{4C}{T}\norm{u_T}_2.
\]
Now apply~\eqref{e:eq:logbound} outside that band and subtract
$\norm D\norm{u_T}_W^2$.
\end{proof}

The right side tends to a positive and then arbitrarily large coefficient
as $T\to\infty$ with the window and complexity fixed. The constants have
not been evaluated, so this is not an effective numerical threshold.
If the displayed lower bound is negative at a smaller carrier, that fact
is not a negative value of the form.

\subsection{Why complexity cannot be removed}\label{e:sec:limits}

The quantifiers in Theorem~\ref{e:thm:main} have the order
\[
 \text{for each fixed }N,\beta,L\text{ there is a height }T(N,\beta,L).
\]
They do not provide a single height for all complexities. Taking the span
of all high-frequency jets allows unbounded complexity and is not covered
by the theorem. The following constructions demonstrate that two key
uniform estimates genuinely fail on the larger arbitrary-frequency class.

\begin{proposition}[Failure of complexity independent norm comparison]\label{e:prop:norm-counter}
For any admissible $W$, let $P_m(x)=x^{2m}$. Then
\begin{equation}\label{e:eq:norm-counter}
 \frac{\norm{WP_m}_W^2}{\norm{WP_m}_2^2}\longrightarrow\infty.
\end{equation}
\end{proposition}
\begin{proof}
Define the probability measure
\[
 d\nu_m(x)=\frac{W(x)^2x^{4m}\,dx}{\int_\R W(y)^2y^{4m}\,dy}.
\]
For each fixed $R\ge1$,
\[
 \nu_m([-R,R])\le
 \left(\frac R{R+1}\right)^{4m}
 \frac{\int_{-R}^R W(x)^2\,dx}{\int_{R+1}^{R+2}W(x)^2\,dx}
 \longrightarrow0.
\]
The denominator in the fixed ratio is strictly positive. Since
$W(x)\to0$ as $|x|\to\infty$, given $B>0$ choose $R$ so that
$1/W\ge B$ outside $[-R,R]$. Then
\[
 \frac{\norm{WP_m}_W^2}{\norm{WP_m}_2^2}
 =\int_\R\frac1W\,d\nu_m
 \ge B\bigl(1-\nu_m([-R,R])\bigr).
\]
Take the lower limit and then let $B\to\infty$.
\end{proof}

\begin{proposition}[Two high carrier locations with growing degree]\label{e:prop:escape-counter}
Suppose $W$ is even and admissible. There exist $T_n\to\infty$ and real
even exponential polynomials $P_n$ whose only frequency locations are
$\pm T_n$, but whose complexities are unbounded, such that
\begin{equation}\label{e:eq:taylor-limit}
 WP_n\longrightarrow W\quad\text{in }\HH\text{ and in }H^1(\R),
 \quad\text{hence in }\XX.
\end{equation}
Consequently uniform low-frequency escape fails when no complexity bound
is imposed. For every fixed bounded self-adjoint $D$,
\begin{equation}\label{e:eq:q-limit}
 \frac{\qD[WP_n]}{\norm{WP_n}_2^2}
 \longrightarrow\frac{\qD[W]}{\norm W_2^2}<\infty.
\end{equation}
\end{proposition}
\begin{proof}
For fixed $T>0$, put
\[
 S_{m,T}(x)=\sum_{k=0}^m\frac{(-iTx)^k}{k!},\qquad
 P_{m,T}(x)=\frac12\bigl(e^{iTx}S_{m,T}(x)
 +e^{-iTx}S_{m,T}(-x)\bigr).
\]
These functions are real and even on $\R$, have frequencies $\pm T$, and
have complexity $2(m+1)$. For fixed $T$, the Taylor series and its derivative
converge locally uniformly, so $P_{m,T}\to1$ and $P_{m,T}'\to0$ pointwise.
For all $m$,
\[
 |S_{m,T}(x)|\le e^{T|x|},\qquad
 |S_{m,T}'(x)|\le T e^{T|x|},
\]
and hence $|P_{m,T}|\le e^{T|x|}$ and
$|P_{m,T}'|\le2T e^{T|x|}$. Equation~\eqref{e:eq:weight} supplies integrable
dominating functions for the weighted norm, ordinary norm, and derivative
of $W(P_{m,T}-1)$. Dominated convergence proves
$WP_{m,T}\to W$ in $\HH\cap H^1$ at fixed $T$.

Choose any $T_n\to\infty$ and then choose $m_n$ so large that the sum of
the two norm distances is below $1/n$. Set $P_n=P_{m_n,T_n}$.
Equation~\eqref{e:eq:h1bound} proves~\eqref{e:eq:taylor-limit} in $\XX$.
Because $W\in L^1$, $\widehat W$ is continuous and
$\widehat W(0)=(2\pi)^{-1/2}\int W>0$.
Thus some fixed low Fourier band has positive limiting mass, contradicting
any complexity independent escape statement. Finally $\qD$ is continuous
on $\XX$: Cauchy--Schwarz holds for $\tg$, and $D$ is bounded.
This proves~\eqref{e:eq:q-limit}. The complexities must be unbounded,
otherwise Theorem~\ref{e:thm:main}\textup{(ii)} would contradict the
nonzero limiting low-band mass.
\end{proof}

These are auxiliary arbitrary-frequency examples. The first has frequency
zero; the second prescribes arbitrarily high polynomial order at two
chosen frequencies. Neither is asserted to use actual zeta zeros or their
actual multiplicities. Neither proves that $\qth$ is negative at any
vector, and~\eqref{e:eq:q-limit} does not determine the sign of $\qth[W]$.
They rule out deleting the complexity hypothesis from this general
mechanism, not a possible theorem using additional zero-specific structure.


\section{The remaining sign}\label{sec:remaining}
All reductions in this paper retain the complete Gamma integral, every prime power, the original mean, and the relevant minimum domains. The local-sign branch gives $49307/100000$; the additional-response implication gives $2861149/5760000$. Neither is an endpoint theorem. The analytic seed is independent of these inputs but is not the best coefficient among the branches.

The exact half-eigenstates and essential bottom $1/2$ leave open a lower discrete spectrum. The delta-scaled complement controls an infinite block but does not assign the finite Schur sign as $\delta\downarrow0$. The full zero-jet core closes the approximation issue for the actual operator, but the entries of the resulting finite Hermitian forms still require a signed estimate. Fixed-complexity high-frequency coercivity does not provide a threshold uniform in complexity and does not control arbitrary combinations mixing low and high frequencies. The local-row and curvature obstructions rule out certain sufficient mechanisms, not the desired global inequality.

Thus the remaining mathematical assertion is exactly the nonnegativity of every finite full-zero-jet form in Theorem~\ref{thm:actual-sign}. This assertion, the global variance half bound, and the Riemann hypothesis are not proved here. No zero-simplicity, historical-priority, journal-acceptance, human-peer-review or proof-assistant-verification claim is made.

\appendix
\section{Complete scalar and interface estimates}\label{app:costs}
The estimates below are part of the proofs of Section~\ref{sec:bounds}. They retain the full theta series, every prime tail and the endpoint singularities of the Gamma kernels.


\subsection{Elementary full-series bounds}\label{c:app:scalar}
For $r>0$ set $S_m(r)=\sum_{k=0}^m r^k/k!$.
The positive exponential series and its geometric tail give
\[
 S_m(r)<e^r\le S_m(r)+\frac{r^{m+1}}{(m+1)!}
                  \frac1{1-r/(m+2)}\qquad(0<r<m+2).
\]
The positive logarithm identity is
\begin{equation}\label{c:eq:atanh}
 \log\frac{1+z}{1-z}=2\sum_{k\ge0}\frac{z^{2k+1}}{2k+1},
 \qquad0<z<1.
\end{equation}
Omitted denominators may be bounded below and powers summed geometrically.
For example
\begin{gather*}
 693/1000<842/1215<\log2<25/36<7/10,\\
 1098/1000<7379/6720<\log3<1+1/12+1/60=11/10,\\
 1609/1000<\log5<13724/8505<81/50.
\end{gather*}
The lower bound for $\log5$ follows from the first eight terms at $z=2/3$,
which exceed respectively
$1.33333,.19753,.05267,.01672,.00578,.00210,.00079,.00030$.
These are exact finite decimals summing to $1.60922$.
The upper bound uses the first three terms and their geometric remainder.
Squaring proves
$\sqrt2>140/99$, $\sqrt3>55/32$, $\sqrt5>324/145$;
these imply \eqref{c:eq:coefficientcaps}.
Also $\log n/\sqrt n\le2/e<1$, so $a_n\le1$.
The inequality $\Lambda(n)\le n/2$ follows from $\log n<n/2$ for $n\ge2$.

The arctangent addition formula and alternating series give
\[
 \pi>4\sum_{d\in\{2,3\}}\sum_{k=0}^3
       \frac{(-1)^k}{(2k+1)d^{2k+1}}
 =\frac{1538665}{489888}>\frac{157}{50},
\]
where the last difference is $10417/12247200$.
Also
$\frac{22}{7}-\pi=\int_0^1z^4(1-z)^4/(1+z^2)\,dz>0$.
Thus $\pi^2<10$.
Useful Taylor consequences are
\[
 e^{2/5}<3/2,\quad e^{1/5}<5/4,\quad
 e^{8/5}<5,\quad e^3>401/20>20,\quad e^2>73/10.
\]
The bound at $8/5$ follows from the displayed tail formula with $m=6$.

\subsubsection{The exterior theta envelope}
For $g(t)=(4t^2-6t)e^{-t}$ and $t\ge3$,
\begin{equation}\label{c:eq:thetatotal}
 \sum_{j\ge2}g(j^2t)<128t^2e^{-4t}.
\end{equation}
Indeed bound each term by $4t^2j^4e^{-tj^2}$.
Consecutive ratios for $j\ge2$ are at most
$(3/2)^4e^{-5t}<1/2$, so twice the first majorant pays the complete tail.
Both $g(t)$ and $t^2e^{-4t}$ decrease for $t>3$.

At ${B_0}=4/5$, $S_7(8/5)>99/20$ gives
$t_{B_0}>T_0=15543/1000>31/2$.
The logarithm bounds imply $e^{B_0}<\sqrt5<9/4$, so $1+e^{-{B_0}}>13/9$.
The degree-nine exponential sum at $3$ is
$7280604/362880>401/20$; the degree-four sum at $543/1000$ is
$41297535960801/24000000000000>43/25$.
Hence $e^{15}>20^5(1+5/400)=3240000$ and $e^{T_0}>5572800$.
Since $4T_0^2-6T_0<8731/10$, the leading part of $A({B_0})$ is below
$78579/724464000$.
The entire tail is below $32768/20^{20}<10^{-12}$, while
\[
 \frac1{9216}-\frac{78579}{724464000}
 =\frac{279936}{9216\cdot724464000}>10^{-8}.
\]
This proves $A<1/9216$.
The logarithmic derivatives of its two summands are at most
\[
 \frac1{1+e^x}+\frac{2(8t-6)}{4t-6}-2t\le6-2t<-25,
 \qquad 4-8t<-25
\]
on $x\ge {B_0}$.
Therefore $W({B_0}+h)\le A({B_0}+h)\le Ae^{-25h}$ and
$\int_{B_0}^\infty W\le A/25$.

\subsubsection{A global theta cap and the exterior probability}
Termwise differentiation is justified by locally uniform summable majorants, and gives
\[
 \kappa'(x)=e^{x/2}\sum_{j\ge1}t_jP(t_j)e^{-t_j},
 \qquad P(t)=-8t^2+30t-15.
\]
At zero the first theta term is below $46/49$:
$4\pi^2-6\pi<4(22/7)^2-6(22/7)$, while
$e^\pi>e^{31/10}>22$.
The remaining terms are bounded by $40j^4e^{-3j^2}$,
whose first term is below $1/250$ and successive ratio below $1/10000$.
Thus $\kappa(0)<46/49+40/9999$.
The polynomial $P$ decreases on $[3,\infty)$, with
$P(31/10)=28/25$ and $P(16/5)=-23/25$.
All $j\ge2$ derivative terms are negative.
A positive derivative is therefore possible only on
$0<x<\frac12\log(32/31)<1/62$.
There,
\[
 \kappa'_+(x)<\frac{60}{59}\frac{16}{5}\frac{28}{25}\frac1{22}
 =\frac{2688}{16225}<\frac16.
\]
Hence
\[
 0<\kappa(x)<\frac{46}{49}+\frac{40}{9999}+\frac1{372}
 =\frac{57440653}{60753924}<\frac{19}{20}.
\]
Evenness gives the global assertion and $\sup W<1/2$.
The normalization proved in Section~\ref{a:sec:energy} gives
\[
 \int_0^\infty W
 =\int_0^\infty\frac{\kappa}{2c}\le\frac12\int_0^\infty c\kappa=\frac18.
\]

For completeness, the endpoint cap $\kappa({B_0})<1/4000$ follows from
$t_{B_0}>31/2$, $e^{{B_0}/2}<3/2$, $e^{31/2}>5300000$:
\[
 \kappa({B_0})<\frac{1302}{5300000}
 +\frac{23064}{5300000^4}\frac{10000}{9999}
 <\frac{1302}{5300000}+\frac1{10^9}<\frac1{4000}.
\]
The exponential lower bound is supplied by $S_{24}(31/2)$;
its excess over $5300000$ is
\[
 \frac{37076883949467729399970752217216753}
 {10409396852733332453861621760000}>0.
\]
The tail ratio in this estimate is below $1/10000$.
Furthermore
\[
 -\kappa'-25\kappa
 =e^{x/2}\sum_{j\ge1}t_j(8t_j^2-130t_j+165)e^{-t_j}>0
 \quad(x\ge {B_0}),
\]
since at $t=31/2+u$ the polynomial is $8u^2+118u+72$.
It follows that
\[
 \int_{B_0}^\infty e^{x/2}\kappa(x)\,dx
 <\frac{(3/2)(1/4000)}{49/2}=\frac3{196000},
 \qquad \mu(O)<\frac3{49000}.
\]

\subsubsection{All needed interior point estimates}
Since $t(1/10)>19/5>16/5$, the differentiated polynomial proves that
$\kappa$ and $W$ decrease to the right of $1/10$.
Every tail below is the complete bound \eqref{c:eq:thetatotal}, with the factor $e^{x/2}$ when estimating $\kappa$.

At $x=1/10$,
$t>(157/50)(61/50)>19/5$,
$1+e^{-x}>19/10$, and $e^{19/5}>20(11/5)=44$.
The leading bound for $W$ is $437/1045$; its tail is below
$46208/96000000<1/1000$.
The margin $21/50-437/1045=19/10450$ pays this tail.

At $x=1/4$, $t>(157/50)(79/48)>103/20$,
$e^{1/4}<9/7$ gives $1+e^{-x}>16/9$, and
$e^{103/20}>20(73/10)(29/25)=4234/25$.
The leading bound is $67671/270976$, the full tail is below
$3456/(20^6\cdot7)<1/100000$, and the margin to $1/4$ is
$73/270976$.

At $x=9/20$,
$(157/50)S_3(9/10)=761921/100000>15/2$,
$1+e^{-9/20}>31/20$, and $e^{15/2}>20^2\cdot4=1600$.
The leading bound is $9/124$, the tail is below $7200/20^{10}<1/1000$,
and the margin to $3/40$ is $3/1240$.

For $\kappa(2/5)$, use $t>69/10$, $e^{1/5}<5/4$,
$e^{69/10}>900$.
The leading bound is $207/1000$, and the tail is below
$\frac{19044}{5\cdot900^4}\frac{10000}{9999}<10^{-6}$,
proving $\kappa(2/5)<1/4$.
The first inequality follows from $(157/50)S_3(4/5)>69/10$.

At $x=1/2$ take $t>17/2$, since $(157/50)S_5(1)>17/2$.
Use $e^{1/4}<9/7<4/3$ and $e^{17/2}>20^2\cdot7\cdot(79/48)$.
The leading bound is $952/13825$ and the tail is below
$36992/(3\cdot20^{11})<1/1000$.
The margin $7/100-952/13825=63/55300$ pays the tail.

For $\kappa(3/5)$,
$S_6(6/5)-520/157=88083/12265625>0$ and
$S_5(7/5)-4=1979/46875>0$ give $t>52/5$ and
$e^{52/5}>32000$; also $e^{3/10}<10/7$.
The leading bound is $(10/7)(9256/25)/32000$.
The full tail is below $19840/20^{13}$, by the majorants
$620j^4e^{-10j^2}$ and their ratio below $1/2$.
Their sum is below $1/50$.

At $x=7/10$, take $t>25/2$, using $e^{7/5}>4$.
Also $e^{7/20}<47/33<10/7$ and
$e^{25/2}>20^4(79/48)$.
The leading bound is $33/11060$ and the tail is below
$200000/(7\cdot20^{16})<10^{-6}$.
The margin to $3/1000$ is $9/553000$.

At $x=3/4$,
$(157/50)S_5(3/2)=896627/64000>14$,
$e^{3/8}<3/2$, and $e^{14}>20^4(73/10)=1168000$.
The leading bound is $21/23360$ and the tail is below
$37632/20^{18}<10^{-5}$.
The margin to $1/1000$ is $59/584000$.
All exponential comparisons above follow from the stated positive sums or their geometric remainder formula; no evaluated theta value enters.

\subsubsection{Gamma constants and all overlap attenuations}
Since $\rho(s)=e^{s/2}/(2\sinh s)$, the inequality
$\rho(s)\le1/(2s)+1/4$ is equivalent to
$2se^{s/2}\le(2+s)\sinh s$.
For odd degree $n$ coefficient comparison reduces to $n\le2^{n-1}$;
for even $n\ge2$ it reduces to $2^{2-n}\le1$.
Thus the power-series difference is nonnegative.
The lower bound $\rho({B_0})>5/6$ uses $e^{2/5}<3/2$, $e^{8/5}<5$;
the upper bound $\rho({B_0})<9/10$ uses $e^{2/5}>10/7$, $e^{8/5}>9/2$.

The harmonic error formula gives
\[
 H_n-\log n-\gamma_E
 =\sum_{k=n}^\infty\int_k^{k+1}\left(\frac1x-\frac1{k+1}\right)\,dx
 >\frac1{2(n+1)}.
\]
At $n=5$ and $\log5>8/5$ this yields $\gamma_E<3/5$.
Together with $\log\pi<6/5$, $\pi/2<11/7$ and $\log2<7/10$,
\[
 \cg<\frac65+\frac35+\frac{11}{7}+\frac{21}{10}
 =\frac{383}{70}<\frac{11}{2}.
\]

For the six overlap caps let $\sigma=25/2$.
The reflection exponents exceed $19/8,47/8,35/4$, from the stated logarithm lower bounds.
The positive exponential sums give
\begin{align*}
 e^{19/8}&>7(185/128)>10,\\
 e^{47/8}&>20\cdot7(289/128)>300,\\
 e^{35/4}&>400(73/10)(21/10)=6132>6000.
\end{align*}
For the translations, use
$\sqrt{3/4}<7/8$, $\sqrt{3/5}<4/5$, $\sqrt{4/5}<9/10$, and
\[
 28\cdot3^{12}<4^{12},\qquad
 2000\cdot3^{12}<5^{13},\qquad
 27\cdot4^{12}<2\cdot5^{12}.
\]
These give $(3/4)^\sigma<1/32$, $(3/5)^\sigma<1/500$ and
$(4/5)^\sigma<1/15$, paying every listed pair.


\subsection{Coercivity of the rank 64 lift}\label{c:sec:lift}
Let
\[
 e_0=(2{B_0})^{-1/2},\qquad e_n={B_0}^{-1/2}\cos(k_nx),\quad
 k_n=n\pi/{B_0}\quad(n\ge1),
\]
be the physical even Neumann basis. Let $P_{64}$ project onto
$e_0,\ldots,e_{63}$ and $H_{64}=L_I+9P_{64}$. Write $H=H_{64}$ in this appendix. We prove the lift estimate
without using the stored response data.

\subsubsection{The entire Gamma form dominates a diagonal form}
For $a>0$ the restricted whole-line Green kernel
$e^{-a\abs{x-y}}/(2a)$ is the inverse $R_{R,a}$ of
$-\partial_x^2+a^2$ with Robin form boundary contribution
$a(\abs{v(-{B_0})}^2+\abs{v({B_0})}^2)$. Removing this nonnegative term gives
the Neumann inverse $R_{N,a}$, so $R_{R,a}\preceq R_{N,a}$.
For every zero-extended local vector,
\[
 \int_0^\infty e^{-as}\norm{v(\cdot+s)-v}_2^2\dd s
 =\frac2a\norm v_2^2-2a\ip v{R_{R,a}v}
 \ge\frac2a\norm v_2^2-2a\ip v{R_{N,a}v}.
\]
Using $\rho(s)=\sum_{j\ge0}e^{-\lambda_js}$, where
$\lambda_j=2j+1/2$, and monotone summation of nonnegative forms, yields
\begin{equation}\label{c:eq:DN}
 t_\Gamma\succeq D_N,\qquad
 D_Ne_n=d_ne_n,\qquad
 d_n=2\sum_{j\ge0}\frac{k_n^2}{\lambda_j(\lambda_j^2+k_n^2)}.
\end{equation}
This is a form inequality on the whole local minimum domain. It does
not impose Neumann boundary conditions on an original vector. Lower
semicontinuity extends the core comparison, with inclusion in the
form domain of $D_N$.

The numbers $d_n$ are nonnegative and increasing. We give a short
rational proof that $d_{64}>9$. Set
\[
 f_t(a)=\frac{2t^2}{a(a^2+t^2)},\qquad G(t)=\sum_{j\ge0}f_t(2j+1/2).
\]
Each summand increases in $t$, and $k_{64}=80\pi>240$.
Differentiation gives
\[
 f_t''(a)=\frac{4t^2(6a^4+3a^2t^2+t^4)}{a^3(a^2+t^2)^3}>0.
\]
The mesh-two trapezoidal inequality for this convex decreasing function
therefore gives
\begin{align*}
 G(240)\ge{}&f_{240}(1/2)+f_{240}(5/2)+f_{240}(9/2)
 +\tfrac12f_{240}(13/2)\\
 &+\tfrac12\log\bigl(1+(480/13)^2\bigr),
\end{align*}
since $\int_A^\infty f_t(a)\dd a=\log(1+t^2/A^2)$.
The inequality $(1+x)^{-1}>1-x$ proves that the four rational terms
exceed, respectively,
\[
 \frac{3999}{1000},\qquad\frac{799}{1000},\qquad
 \frac{444}{1000},\qquad\frac{153}{1000}.
\]
For example they exceed $4-1/57600$, $4/5-5/57600$,
$4/9-1/6400$ and $2/13-13/115200$, respectively.
Also \eqref{c:eq:atanh} gives
\[
 \log2>\frac{842}{1215}>\frac{693}{1000},\qquad
 \log(15/13)>\frac17.
\]
It follows that
\[
 \tfrac12\log\bigl(1+(480/13)^2\bigr)
 >\log(480/13)=5\log2+\log(15/13)
 >\frac{693}{200}+\frac17.
\]
Consequently
\begin{equation}\label{c:eq:d64}
 d_{64}>G(240)>\frac{3151}{350}=9+\frac1{350}>9.
\end{equation}

\subsubsection{A whole operator bound for the local prime term}
Write $s=\log2$, $t=\log3$, $d={B_0}-s$ and
$a=s/\sqrt2$, $b=s/2$, $a_3=t/\sqrt3$.
The elementary log bounds give $.69<s<.70$, $1.09<t<1.10$ and
$\log5>2{B_0}$. The lower bound for $s$ follows from $842/1215>.693$;
for $t$, the first three terms in \eqref{c:eq:atanh} give
$\log3>1+1/12+1/80>1.09$.
Under the unitary fold on $(0,{B_0})$ partition that interval
(up to endpoints) into
\[
 C=(0,d),\quad L_-=(s-d,s),\quad L_+=(s,s+d),\quad
 J_0=(d,s-d).
\]
Prime 2 connects $y\in C$ to $s-y$ and $s+y$, with weight $a$,
and reflects $J_0$. Prime 4 connects the two leaves $L_-,L_+$
with weight $b$. Prime 3 reflects $S=(t-{B_0},{B_0})$ with weight $a_3$.
The inequalities
$2s>t$, $t+s>2{B_0}$, $t+2{B_0}<4s$ show that the leaves lie in $S$,
that their prime-3 images lie in $J_0$, and that $C$ does not meet $S$.

Define the positive number
\[
 R_p=\frac{b+a_3+\sqrt{(b+a_3)^2+8a^2}}2,
 \qquad\eta_0=2a/R_p,
 \qquad w_0=\eta_0 1_C+1_{J_0\cup L_-\cup L_+}.
\]
The weighted row ratio $(\PP_Iw_0)/w_0$ equals $R_p$ on the center
and leaves and is at most $a+a_3$ on $J_0$. Here
$\PP_I$ is the complete local prime operator in the definition of $L_I$.
One has $R_p>5/4$: its defining quadratic at $5/4$ is negative using
$b>1/3$, $a_3>4/7$, $2a^2>4/9$.
Moreover $a<1/2$, $a_3<121/190$ and hence $a+a_3<5/4$.
Every edge map is measure-preserving and every edge weight is
nonnegative. Applying
$2\abs{uv}\le(w_y/w_x)u^2+(w_x/w_y)v^2$ on each symmetric edge,
and summing, proves
$\abs{\ip v{\PP_Iv}}\le R_p\norm v_2^2$.
This retains all low/high interactions of the full operator.
Finally
\[
 b+a_3<\frac7{20}+\frac{121}{190}=\frac{75}{76},\quad
 2a^2<\frac{49}{100},\quad
 \left(\frac{27}{20}\right)^2-\frac{75}{76}\frac{27}{20}
 -\frac{49}{100}=\frac1{3800}>0.
\]
Thus
\begin{equation}\label{c:eq:primeop}
 \norm{\PP_I}<\frac{27}{20}.
\end{equation}

\subsubsection{The lift constant and a residual majorant}
The elementary estimate
\begin{equation}\label{c:eq:cgamma}
 c_\Gamma<\frac65+\frac35+\frac{11}{7}+\frac{21}{10}
 =\frac{383}{70}
\end{equation}
uses $e^{6/5}>401/125>22/7$ and $\gamma_E<3/5$.
To check the latter without evaluating Euler's constant, the harmonic
error identity gives
\[
 H_n-\log n-\gamma_E
 =\sum_{k=n}^\infty\int_k^{k+1}
       \left(\frac1x-\frac1{k+1}\right)\dd x
 >\frac1{2(n+1)}.
\]
Take $n=5$ and $\log5>8/5$ to get
$\gamma_E<H_5-8/5-1/12=3/5$.

\begin{theorem}[Rank-64 lift coercivity]\label{c:thm:lift}
On the complete mean-inclusive local minimum domain,
\begin{equation}\label{c:eq:lift}
 H_{64}=L_I+9P_{64}\succeq\frac{61}{28}I,
 \qquad H^{-1}\preceq\frac{28}{61}I.
\end{equation}
This is a coercivity statement for $H$, not for the unlifted $L_I$.
\end{theorem}
\begin{proof}
Dropping only the nonnegative pole in a lower estimate, equations
\eqref{c:eq:DN} and \eqref{c:eq:primeop} give
\[
 L_I\succeq D_N-(c_\Gamma+27/20)I.
\]
Below mode 64 the lifted diagonal is at least $9-c_\Gamma-27/20$;
from mode 64 onward the same bound follows from \eqref{c:eq:d64}.
Since $9-383/70-27/20=61/28$, \eqref{c:eq:lift} follows.
The pole has not been removed from the operator being inverted.
\end{proof}
For the residual contract define the larger diagonal majorant
\begin{equation}\label{c:eq:Dold}
 D_{\rm old}e_n=\beta_ne_n,\qquad
 \beta_n=\begin{cases}(d_n+1/2)^{-1},&n<64,\\
 (d_n-17/2)^{-1},&n\ge64.
 \end{cases}
\end{equation}
Translation Cauchy--Schwarz gives
$\norm{\PP_I}\le2(a_2+a_3+a_4)<3$, and
$c_\Gamma<11/2$, so $H\succeq D_N-17I/2+9P_{64}$.
Equation~\eqref{c:eq:d64} proves
$0<H^{-1}\preceq D_{\rm old}\preceq2I$.
Contract (a) uses this larger majorant. The sharper scalar bound for $H^{-1}$ is used separately in the inverse-cost estimate.

\subsection{An exact signed exterior prime identity}\label{c:sec:signed}
In this appendix write $A_{B_0}=A$.
This section is analytic and requires no stored numerical input or
exterior moment constraint. Set $g=u/\alpha$ and
$E_n=\{x:x,x+s_n\in O\}$. Define the nonnegative full-prime slack
\begin{equation}\label{c:eq:slack}
 D_p(g)=\sum_{n\ge2}a_n\int_{E_n}
 \alpha(x)\alpha(x+s_n)[g(x+s_n)-g(x)]^2\dd x.
\end{equation}
The sum includes all prime powers. The two row orientations arise from
the two diagonal terms in each edge square, as the proof below shows.

\subsubsection{The complete exterior row and its one reflection}
On the positive exterior fold $f=\sqrt2u$, the full prime expression is
\begin{equation}\label{c:eq:exteriorfold}
 (\PP_Of)(x)=\sum_{n\ge2}a_n\big[
 f(x+s_n)+1_{x>{B_0}+s_n}f(x-s_n)
 +1_{{B_0}<x<s_n-{B_0}}f(s_n-x)\big],\quad x>{B_0}.
\end{equation}
The first two terms are same-side orientations. The third is one
reflected opposite-side input, with one coefficient $a_n$. The cases
for $x-s_n$ are disjoint: positive exterior, killed interval and
negative exterior. Since $\log4<2{B_0}<\log5$, no reflection occurs for
$n=2,3,4$. The unitary fold already absorbs the factor two from the
two half-lines.

The relation $p(u,u)=\int_{B_0}^\infty f\PP_Of$ is understood as an
absolutely convergent form or weighted-duality pairing, first for
truncated sums. It is not an assertion that $\PP_Of$ belongs to
physical $L^2$ for every finite-$\wo$ vector. Equivalently it is the
$L^2$ pairing with the bounded conjugate
$M_\alpha\PP_OM_\alpha$, justified by the following row bound.
Define the nonnegative potential on all of $O$ by
\begin{equation}\label{c:eq:primepotential}
 V_p(x)=\alpha(x)\sum_{n\ge2}a_n
 \big[\alpha(x+s_n)1_O(x+s_n)
      +\alpha(x-s_n)1_O(x-s_n)\big].
\end{equation}
For $x>{B_0}$, either existing same-side product is at most
$A_{B_0}n^{-25/2}$: the sum of the two distances from ${B_0}$ is at least
$s_n$. On the reflected branch it is at most
$A_{B_0}e^{-(25/2)(s_n-2{B_0})}$, because those distances sum exactly to
$s_n-2{B_0}$. There is only one such branch in a row.

The two same-side budgets together satisfy
\[
 2\sum_{n\ge2}a_nn^{-25/2}\le2\sum_{n\ge2}n^{-12}
 \le2\left(2^{-12}+\int_2^\infty t^{-12}\dd t\right)
 =\frac{13}{22528}<\frac1{1000}.
\]
At $n=5$, equations \eqref{c:eq:coefficientcaps} and
$e^{-(25/2)(\log5-2{B_0})}<9/10$ give cost less than $261/400$.
There is no $n=6$ coefficient. Every remaining reflection is covered by
\eqref{c:eq:bothtails}, since
$a_ne^{-(25/2)(\log n-2{B_0})}\le(5/n)^{12}$ for $n\ge7$.
Consequently, for all $x\in O$,
\begin{equation}\label{c:eq:primecap}
 V_p(x)<\left(\frac1{1000}+\frac{261}{400}+\frac3{100}\right)A_{B_0}
 =\frac{1367}{2000}A_{B_0}<\frac7{10}A_{B_0}.
\end{equation}
Evenness extends the pointwise bound to the negative half-line without
doubling it. The summable majorants give uniform row convergence and
boundedness of the conjugated operator, including every prime tail.

\subsubsection{The signed identity and the analytic cap}
\begin{theorem}[Full-prime slack and mass bound]\label{c:thm:signed}
For every real even exterior $u$ of finite $\wo$, one has
\begin{equation}\label{c:eq:slackidentity}
 p(u,u)+D_p(u/\alpha)
 =\int_O V_p(x)|u(x)/\alpha(x)|^2\dd x,
\end{equation}
with $0\le D_p(u/\alpha)\le(7/5)A_{B_0}\wo(u)$.
Put
\[
 V(x)=c_\Gamma\alpha(x)^2+V_p(x),\qquad N=\frac{31}{46080}.
\]
Then $0<V(x)<N$ and
\begin{equation}\label{c:eq:negativeanalytic}
 c_\Gamma\norm u_2^2+p(u,u)
 \le N \wo(u)-D_p(u/\alpha).
\end{equation}
Whenever $t_\Gamma(u)$ is finite, the exact spatial identity is
\begin{equation}\label{c:eq:signedq0}
 q_0(u)=t_\Gamma(u)+D_p(u/\alpha)
              -\int_OV(x)|u(x)/\alpha(x)|^2\dd x.
\end{equation}
The cap is an analytic conclusion, not a numerical premise.
\end{theorem}
\begin{proof}
For a finite sum, expansion of one edge square in \eqref{c:eq:slack}
gives cross term $-2a_nu(x)u(x+s_n)$. The two diagonal terms,
after a change of variable in one, give the two orientations in
\eqref{c:eq:primepotential}. This proves \eqref{c:eq:slackidentity} for
truncated sums with exactly the factor in the definition of $p$.
The inequality $(a-b)^2\le2(a^2+b^2)$ gives
\[
 0\le D_p(g)\le2\int_OV_pg^2\le\frac75 A_{B_0}\norm g_2^2.
\]
The same row bound controls the absolute prime bilinear form by
weighted Cauchy--Schwarz. Monotone convergence of the slack series and
absolute convergence of the prime pairing extend the identity to all
$g\in L^2(O)$. Thus negative prime correlation increases the retained
slack; its sign has not been replaced by an absolute correlation.

Equation~\eqref{c:eq:cgamma} gives $c_\Gamma<11/2$. Together with
\eqref{c:eq:envelope} and \eqref{c:eq:primecap},
\[
 0<V(x)<\left(\frac{11}{2}+\frac7{10}\right)A_{B_0}
 =\frac{31}{5}A_{B_0}<\frac{31}{46080}.
\]
Adding the mass term proves \eqref{c:eq:negativeanalytic}, and
subtracting the exact identity from $t_\Gamma$ proves
\eqref{c:eq:signedq0}. The slack and spatial potential are $\wo$-bounded,
so the same identities extend in the original form topology, in
particular on the fixed minimum domain $\D_O^0$.
\end{proof}

The slack has trivial kernel on the even weighted exterior space.
Indeed its $n=2$ same-side edges have strictly positive weight, so
zero slack forces $g(x+\log2)=g(x)$ almost everywhere on $({B_0},\infty)$.
The $L^2$ norm is then an infinite sum of equal nonnegative period
integrals and can be finite only if each vanishes. Evenness supplies
the other half-line. This trivial-kernel statement does not assert a
uniform inequality $D_p\ge\delta\wo$ with $\delta>0$.


\subsection{Sharp cuts and the centered exterior minimum domain}\label{c:app:cuts}
We record the domain argument, rather than identifying a minimum
closure with an arbitrary maximal finite-energy class.
Write $W_\R(h)=\int_\R |h|^2/\alpha^2$. Under $h=\kappa f$ it equals
$\norm f_{L^2(\mu)}^2$. Proposition~\ref{a:prop:domain}
shows that the original topology and
$(t_\Gamma(h)+W_\R(h))^{1/2}$ are equivalent on numerators.

For a smooth numerator $h$, expansion gives
\[
 t_\Gamma(h1_I)+t_\Gamma(h1_O)
 =t_\Gamma(h)+2\ip{h1_I}{\J_\Gamma(h1_O)}.
\]
For $0<s\le2$, \eqref{c:eq:rhoelementary} gives
$\rho(s)\le1/(2s)+1/4\le1/s$. For $s\ge2$, the function
$se^{-s/2}$ decreases, so
$s\rho(s)\le(2/e)/(1-e^{-4})<1$; the last inequality follows already
from $e>8/3$. Thus the unweighted Carleman bound for the Gamma crossing
is finite: $\rho(s)\le1/s$ and the two boundary Carleman maps give,
for example, $\norm{\J_\Gamma}_{L^2(O)\to L^2(I)}\le2\pi$.
Thus the right-hand correction is bounded by a constant times
$\norm h_2^2$, and $W_\R$ splits exactly. Since $\alpha^2$ is bounded,
both cut maps extend continuously in the numerator form topology.

For a fixed smooth compact numerator, approximate either step cut from
its own side by smooth killed-core cutoffs. The error $e_\delta$ is
supported in strips of width $\delta$ at $\pm {B_0}$ and satisfies
\[
 \norm{e_\delta}_\infty+\operatorname{TV}(e_\delta)\le C,
 \qquad\norm{e_\delta}_2^2+W_\R(e_\delta)\le C\delta.
\]
The bounded-variation translation estimate gives
\[
 \norm{e_\delta(\cdot+s)-e_\delta}_2^2\le C\min(s,\delta).
\]
Integrating against $\rho(s)=O(1/s)$ near zero and its integrable
large-shift tail proves
$t_\Gamma(e_\delta)=O(\delta(1+|\log\delta|))\to0$.
Consequently the cuts are in the corresponding killed minimum closures.
Apply continuity to a general minimum-domain numerator to obtain the
same conclusion. On the bounded local interval $W_\R$ is equivalent
to physical $L^2$, because $\alpha$ is continuous and strictly positive.

The constant function belongs to $\D$. Indeed even smooth cutoffs
$\chi_R$ of the constant satisfy $\kappa\chi_R\to\kappa$ in physical
$H^1$ and in $W_\R$ by theta decay. The topology equivalence and
$t_\Gamma(h)\le C\norm h_{H^1}^2$ make them Cauchy in the original
energy topology. Subtracting a constant is an admissible gauge change,
preserving energy and variance.

For an exterior limit with zero moment, choose even compact smooth
$\psi$ in $O$ with $m_O(\psi)=1$. If $u_n$ are killed-core approximants,
replace them by $u_n-m_O(u_n)\psi$. Moment continuity makes the
correction vanish in form norm. Hence the centered cut lies in
precisely $\D_O^0$, not just in the moment kernel of a larger domain.


\section{Complete response quantities and contracts}\label{app:response}
These quantities specify the additional response branch. Their definitions and analytic error propagation do not constitute a new verification of the stated numerical enclosures.


\subsection{The exact centered-row quantities}\label{c:app:rows}
This appendix fixes the analytic interpretation of contract (b).
Use $K=64$, $J_0=512$, $d_0=6$ and $x_*=2$.
Let $\varphi_0={B_0}^{-1/2}$ and
$\varphi_m=(2/{B_0})^{1/2}\cos(k_mx)$ for $m>0$, where $k_m=m\pi/{B_0}$.
They are the unitary folds of the physical cosines.
All rows below are formed from the same matrix $Z$ and the same fixed rational vector $a$ in Hypothesis~\ref{c:hyp:response}.
The displayed formulas define quantities to be enclosed; their small numerical bounds remain assumptions.
For the physical columns $w_i$ of $\V$, let $Q_i=\J^*w_i$ denote
their weighted-dual kernels. This notation means
$\ip{\J u}{w_i}=\int_OuQ_i$ and
$\alpha Q_i\in L^2(O)$; it does not assert $Q_i\in L^2(O,\dd x)$.
If $m_O(u)=0$, then for every fixed vector $a$,
\begin{equation}\label{c:eq:centering}
 \V^*\J u=\int_Ou(x)[Q(x)-ac(x)]\dd x,\qquad
 \norm{\V^*\J u}\le\sqrt{\wo(u)}\norm{\alpha(Q-ac)}_{L^2(O;\R^{64})}.
\end{equation}
This is exact moment cancellation, not an approximate mean repair.
The choice of $a$ need not be accurate as a numerical minimizer.
The direction $c$ is subtracted in the weighted dual, not by physical
$L^2$ projection; $c$ is not a physical exterior $L^2$ vector.

Split the trial response into cosine rows below $K=64$ and rows
64 through 1023. Duality and Theorem~\ref{c:thm:source} give
\begin{equation}\label{c:eq:highrow}
 \norm{\alpha\J^*\V_{\ge64}}_{L^2(O;\R^{64})}
 \le j\norm{Z_{\ge64}}_{\rm HS}.
\end{equation}
This pays every discarded row of the finite trial response. The error
of the trial response as an approximation to $H_{64}^{-1}E$ belongs instead
to \eqref{c:eq:contracta}; it is not silently removed.

Let $b_m=\varphi_m({B_0})$ and $\lambda_j=2j+1/2$.
For $x>{B_0}$, direct cosine integration of the two folded exponential
kernels gives the low-response Gamma adjoint
\begin{equation}\label{c:eq:Qgamma}
 Q_{\Gamma,<K}(x)=\sum_{j\ge0}g_je^{-\lambda_j(x-{B_0})},\qquad
 g_j=\frac{1-e^{-2\lambda_jB_0}}{\sqrt2}
       \sum_{m<K}\frac{b_m\lambda_j}{\lambda_j^2+k_m^2}Z_m.
\end{equation}
The factor $1/\sqrt2$ converts folded coefficients to physical
full-line columns. Put $A_J=2J_0+1/2$, and define signed moments
\[
 M_r=\sum_{m<K}b_m(-k_m^2)^rZ_m\quad(0\le r<d_0),\qquad
 H_{d_0}=\sum_{m<K}|b_m|\norm{Z_m}k_m^{2d_0}.
\]
The exact finite geometric identity for
$\lambda/(\lambda^2+k^2)$ has remainder of modulus at most
$k^{2d_0}/\lambda^{2d_0+1}$, without requiring $k<\lambda$.
Therefore
\[
 \norm{\sum_{m<K}\frac{b_m\lambda}{\lambda^2+k_m^2}Z_m}
 \le\sum_{r<d_0}\frac{\norm{M_r}}{\lambda^{2r+1}}
       +\frac{H_{d_0}}{\lambda^{2d_0+1}}.
\]
With $\sigma_p(A)=A^{-p}+A^{1-p}/[2(p-1)]$, decreasing integral
comparison gives $\sum_{j\ge J_0}\lambda_j^{-p}\le\sigma_p(A_J)$.
Minkowski's inequality and
$\norm{e^{-\lambda h}}_{L^2(0,\infty)}=(2\lambda)^{-1/2}$ prove
\begin{equation}\label{c:eq:egamma}
 e_\Gamma=\sqrt{A/2}\left[
 \sum_{r<d_0}\norm{M_r}\sigma_{2r+3/2}(A_J)
 +H_{d_0}\sigma_{2d_0+3/2}(A_J)\right]
\end{equation}
as a bound for the complete omitted weighted Gamma norm. Both
exterior halves are included. The signed moments are summed before
their vector norms are taken.

The remaining finite kernel is
\begin{equation}\label{c:eq:K0}
 \begin{split}
 K_0(x)={}&\sum_{j<J_0}g_je^{-\lambda_j(x-{B_0})}-ac(x)\\
 &+\sum_{\substack{2\le n\le16\\|x-\log n|<{B_0}}}
  \frac{a_n}{\sqrt2}\sum_{m<K}|b_m|Z_m\cos(k_m(x-\log n)).
 \end{split}
\end{equation}
All ten prime powers through 16 are included:
$2,3,4,5,7,8,9,11,13,16$. Their Gamma, prime and mean contributions
are combined before vector norms. Boundary values at support switches
have measure zero.

The rough bound
\begin{equation}\label{c:eq:roughalpha}
 \alpha(x)^2\le80e^{4|x|-\pi e^{2|x|}}\quad(x\in O)
\end{equation}
follows already from the envelope: $4\pi^2<40$ and the tail-to-leading
coarse ratio $32e^{-3t}<1$ permit the constant 80.
Define
\[
 C_X=\sum_{j<J_0}\norm{g_j}
 +\left(\sum_{n\le16}a_n\right)
      \sum_{m<K}\frac{|b_m|\norm{Z_m}}{\sqrt2}+\norm a.
\]
For $x\ge x_*=2$, $\norm{K_0(x)}\le C_Xe^{x/2}$, including any finite
prime supports continuing past $x_*$. Since
$5-2\pi e^{2x}\le5-2\pi e^{2x_*}<0$, the full spatial tail is bounded by
\begin{equation}\label{c:eq:eX}
 e_X^2=\frac{160C_X^2e^{5x_*-\pi e^{2x_*}}}{2\pi e^{2x_*}-5}.
\end{equation}
The factor 160 includes both exterior half-lines.

On $[{B_0},2]$ the nine support switches, in increasing order, are
\[
 \begin{gathered}\log5-{B_0},\ \log7-{B_0},\ \log8-{B_0},\ \log9-{B_0},\ \log2+{B_0},\\
 \log11-{B_0},\ \log13-{B_0},\ \log3+{B_0},\ \log16-{B_0}.\end{gathered}
\]
The ten intervals have the fixed subdivision counts
\[
 (5,169,67,59,48,53,84,67,37,14),
\]
totaling 603.
On an exact coarse interval $[L_c,R_c]$ with count $N$, the exact finer
endpoints are $l_i=L_c+(R_c-L_c)i/N$, $r_i=L_c+(R_c-L_c)(i+1)/N$.
This algebraic identity proves adjacency and coverage. Overlap of
arbitrary interval boxes alone would not do so. Each open cell lies
between consecutive support switches and has a constant active-prime
set.

For a cell $[l,r]$, center $x_0$ and half-width $h$, define
$v_s=K_0^{(s)}(x_0)/s!$, $0\le s<16$. Taylor's theorem gives the
following explicit uniform vector remainder bounds:
\begin{align*}
 R_\Gamma&=\sum_{j<J_0}\norm{g_j}e^{-\lambda_j(l-{B_0})}
                   \frac{(\lambda_jh)^{16}}{16!},\\
 R_p&=\left(\sum_{n\ {\rm active}}a_n\right)
       \sum_{m<K}\frac{|b_m|\norm{Z_m}}{\sqrt2}
                   \frac{(k_mh)^{16}}{16!},\\
 R_c&=\norm a\cosh(r/2)\frac{(h/2)^{16}}{16!}.
\end{align*}
Thus $M_{[l,r]}=\sum_{s=0}^{15}\norm{v_s}h^s+R_\Gamma+R_p+R_c$
bounds $\norm{K_0(x)}$ on the whole cell, and
\begin{equation}\label{c:eq:Ifin}
 I_{\rm fin}=\sum_{[l,r]}2(r-l)A(l)M_{[l,r]}^2
\end{equation}
bounds its full even weighted integral up to $x_*$.
Contract (b) requires outward enclosures of these
quantities; the analytical formula is a supremum-and-remainder bound,
not a quadrature accuracy inferred from mesh size. The subdivision counts do not replace these exact endpoint definitions.

\subsubsection{The complete prime tail in the response budget}
For $y\in I$, $|y\pm\log n|\ge\log n-{B_0}$ and each of the two
translated restrictions has weighted norm at most
$\sqrt{80}e^{-2{B_0}}n^2e^{-\zeta n^2}$, with
$\zeta=(\pi/2)e^{-2{B_0}}>3/10$. For $N\ge16$, integration by parts and
$\int_N^\infty e^{-\zeta x^2}\dd x\le
 e^{-\zeta N^2}/(2\zeta N)$ give
\[
 \norm{\J_{p,>N}}_{\wo\to L^2(I)}
 \le2\sqrt{80}e^{-2{B_0}}
 \left(\frac{N}{2\zeta}+\frac1{4\zeta^2N}\right)e^{-\zeta N^2}.
\]
At $N=16$, $2\sqrt{80}e^{-2{B_0}}<18$, the parenthesis is less than 27,
and $\zeta N^2>75$; hence the bound is smaller than
$b_{16}=500/20^{25}$. These entire integer tails dominate all omitted
prime powers and retain both orientations.


The finite part has weighted norm at most $\sqrt{I_{\rm fin}}+e_X$,
the complete omitted Gamma part has norm at most $e_\Gamma$,
and the complete omitted prime part has norm at most
$b_{16}\norm{Z_{<64}}_{\mathrm{HS}}$.
The high trial rows have norm at most $j\norm{Z_{64:1023}}_{\mathrm{HS}}$ by
\eqref{c:eq:highrow}. Exact exterior centering therefore gives
\[
 \norm{V^*Ju}\le
 \left(\sqrt{I_{\rm fin}}+e_X+e_\Gamma+
 b_{16}\norm{Z_{<64}}_{\mathrm{HS}}
 +j\norm{Z_{64:1023}}_{\mathrm{HS}}\right)\sqrt{\wo(u)}.
\]
This proves the decomposition used in \eqref{c:eq:responsebound};
the numerical caps themselves remain contract (b).


\section{Actual weight source recovery and the Abel endpoint}\label{app:abel}
This appendix proves the weighted estimates invoked by the synthesis theorem. In particular, the source is recovered in the actual weight before the endpoint conclusion is drawn. Carrier growth alone is not used as a weighted admission criterion.


\subsection{A scalar Borel source and exact theta reduction}\label{d:sec:scalar-source}

The coarse estimate~\eqref{d:eq:Abel-norm} uses quotient growth but not the full strength of actual weighted admission of $u$. We now use that admission to recover an auxiliary Laplace source in its correct weighted norm. This is the additional input that makes the Abel bound uniform.

The exponential coordinate $q=e^{2b}$ maps $L$ bijectively onto
\begin{equation}\label{d:eq:q-domain}
 \mathcal D_q=\{q:\operatorname{Re}q\ge1/2,\ |q-1|\le1\}.
\end{equation}
Indeed the two lens inequalities are
$\operatorname{Re}(e^{2b})\ge1/2$ and
$\operatorname{Re}(e^{-2b})\ge1/2$. The latter becomes
$\operatorname{Re}q/|q|^2\ge1/2$, equivalent to $|q-1|\le1$.
Conversely $\mathcal D_q$ lies in the right half-plane, and
$b=\frac12\operatorname{Log}q$ satisfies the two inequalities in the principal theta strip. Fractional powers of $q$ below always use this principal branch.

Fix once and for all a positively oriented, centrally symmetric rectifiable Jordan contour $\Gamma\subset\mathbb C\setminus L$ with $L\subset\operatorname{int}\Gamma$, winding number one about every point of $L$, with its interior in $|\operatorname{Im}b|<\pi/4$, and such that
\begin{equation}\label{d:eq:fixed-contour-margin}
 \operatorname{Re}(e^{2b})\ge\frac14\qquad(b\in\Gamma).
\end{equation}
A sufficiently small exterior parallel contour has these properties by compactness and the strict margin between $L$ and the strip boundary. Symmetry also gives the same margin for $e^{-2b}$. The contour remains fixed independently of $u$ and $r$.

For $u\in\Ii$, with $P(w)=Q(-iw)$ and Borel transform $B_P$, define the scalar analytic functional
\begin{equation}\label{d:eq:mu-functional}
 \mu(h)=\frac1{2\pi i}\int_\Gamma B_P(b)h(b)\,db.
\end{equation}
All functions $h$ used below are holomorphic on a neighborhood of the contour and its interior. Theorem~\ref{d:thm:Borel} gives
\begin{equation}\label{d:eq:mu-bound}
 |\mu(h)|\le C_\Gamma\norm{u}_{\Hh}\sup_{b\in\Gamma}|h(b)|.
\end{equation}
The Borel series of an even $P$ is odd, and continuation gives
$B_P(-b)=-B_P(b)$. A change of variable on the symmetric contour therefore yields
\begin{equation}\label{d:eq:mu-symmetry}
 \mu(h(-b))=\mu(h(b)),\qquad \mu(e^{bw})=P(w).
\end{equation}

\begin{lemma}[Scalar source identification]\label{d:lem:scalar-source}
The function
\begin{equation}\label{d:eq:scalar-v}
 v(x)=\mu(\kappa(x+b))
\end{equation}
is even, belongs to $L^1(\R,e^{B|x|}dx)$ for every finite $B$, and equals $u$ almost everywhere.
\end{lemma}
\begin{proof}
The contour lies in a compact subset of the theta strip. The two-ended theta bounds therefore give uniform ordinary exponential moments for $\kappa(\cdot+b)$, $b\in\Gamma$. Equation~\eqref{d:eq:mu-bound} and ordinary dominated integration prove the asserted $L^1$ properties. Symmetry~\eqref{d:eq:mu-symmetry} and evenness of $\kappa$ give evenness of $v$.

For each $b$ on the contour, a shift within the strip gives
\[
 \int_\R\kappa(x+b)\cos(zx)\,dx=\Xi(z)\cos(bz).
\]
The vertical sides vanish by the theta tails. Uniform exponential moments justify passing the scalar contour through the Fourier integral, so
\[
 F_v(z)=\Xi(z)\mu(\cos(bz))=\Xi(z)P(iz)=\Xi(z)Q(z)=F_u(z).
\]
The given $u\in\Hh$ also lies in ordinary $L^1$ by~\eqref{d:eq:eval-bound}. Fourier uniqueness there gives $v=u$ almost everywhere.
\end{proof}

The integral~\eqref{d:eq:scalar-v} is a scalar contour integral, not an $\Hh$-valued Bochner integral over an exterior or boundary contour. No such weighted contour admission is being assumed.

Define for $t\ge0$
\begin{equation}\label{d:eq:fg-source}
 f(t)=\mu(q^{9/4}e^{-tq}),\qquad
 g(t)=\mu(q^{5/4}e^{-tq}),\qquad q=e^{2b}.
\end{equation}
These are continuous, and the fixed contour gives
\begin{equation}\label{d:eq:fg-crude}
 |f(t)|+|g(t)|\le C\norm{u}_{\Hh}e^{-t/4}.
\end{equation}
Ordinary integration on the contour, justified by~\eqref{d:eq:fixed-contour-margin}, gives
\begin{equation}\label{d:eq:g-tail}
 g(t)=\int_t^\infty f(s)\,ds.
\end{equation}

For $x\ge0$, put
\begin{equation}\label{d:eq:t-U-H}
 t=\pi e^{2x},\qquad
 U(t)=\frac{u(x)}{4\pi^2e^{9x/2}},\qquad
 H_f(t)=\sum_{n\ge1}n^4f(n^2t).
\end{equation}
Here $u$ is represented by the continuous scalar source of Lemma~\ref{d:lem:scalar-source}. The precise theta identity on the strip is
\begin{align}
 \kappa(x+b)
 &=4\pi^2e^{9x/2}\sum_{n\ge1}n^4q^{9/4}e^{-n^2tq}\notag\\
 &\quad-6\pi e^{5x/2}\sum_{n\ge1}n^2q^{5/4}e^{-n^2tq}.
 \label{d:eq:exact-theta-q}
\end{align}
It follows by expanding~\eqref{a:eq:kappa}, with analytic continuation as in Lemma~\ref{d:lem:theta-tails}. The uniform exponential bound on the fixed contour justifies all sums and integrals.

\begin{proposition}[Exact reduction in the actual weight]\label{d:prop:exact-reduction}
For $t\ge\pi$,
\begin{equation}\label{d:eq:U-H-identity}
 U(t)=H_f(t)-\frac3{2t}\int_t^\infty H_f(s)\,ds.
\end{equation}
Writing
\begin{equation}\label{d:eq:E-tail}
 E_T=L^2([T,\infty),t^{3/2}e^t\,dt),\qquad T\ge\pi,
\end{equation}
one has two-sided norm comparability
\begin{equation}\label{d:eq:actual-weight-comparison}
 \norm{u}_{\Hh}\asymp\norm{U}_{E_\pi}.
\end{equation}
The comparison holds as an identity of extended finite or infinite integrals for every even scalar source with the displayed normalization.
\end{proposition}
\begin{proof}
Insert~\eqref{d:eq:exact-theta-q} into~\eqref{d:eq:mu-functional} and divide by the normalization in~\eqref{d:eq:t-U-H}. The lower-degree coefficient is exactly
$6\pi e^{5x/2}/(4\pi^2e^{9x/2})=3/(2t)$. Thus
\[
 U(t)=\sum_{n\ge1}n^4f(n^2t)
       -\frac3{2t}\sum_{n\ge1}n^2g(n^2t).
\]
By~\eqref{d:eq:fg-crude} and~\eqref{d:eq:g-tail}, the latter sum equals
$\int_t^\infty H_f(s)\,ds$, proving~\eqref{d:eq:U-H-identity}. No theta remainder has been discarded.

Finally use~\eqref{d:eq:W-two-sided}, $u(x)=4\pi^2e^{9x/2}U(t)$, and $dx=dt/(2t)$. The squared source normalization contributes $t^{9/2}$, the weight contributes $t^{-2}e^t$, and the change of variables contributes $t^{-1}$. Their product is $t^{3/2}e^t$. Evenness covers both physical half-lines and proves~\eqref{d:eq:actual-weight-comparison}.
\end{proof}

\subsection{Noncircular recovery of the Laplace source}\label{d:sec:source-recovery}

The contour bound alone does not imply the weighted membership needed for the endpoint. The actual membership $u\in\Hh$, through~\eqref{d:eq:actual-weight-comparison}, supplies it.

\begin{theorem}[Recovery from actual weighted admission]\label{d:thm:recovery}
For every $u\in\Ii$, the function $f$ in~\eqref{d:eq:fg-source} belongs to
\begin{equation}\label{d:eq:E-global}
 E=L^2((0,\infty),(1+t)^{3/2}e^t\,dt),
\end{equation}
and
\begin{equation}\label{d:eq:recovery-bound}
 \norm{f}_E\le C\norm{u}_{\Hh}.
\end{equation}
\end{theorem}
\begin{proof}
First recover $H_f$ without assuming it belongs to $E_\pi$. The crude bound~\eqref{d:eq:fg-crude} gives
$H_f(t)=O(\norm{u}_{\Hh}e^{-t/4})$ for $t\ge\pi$. Set
$J_0(t)=\int_t^\infty H_f(s)\,ds$. Equation~\eqref{d:eq:U-H-identity} becomes
\[
 J_0'(t)+\frac3{2t}J_0(t)=-U(t).
\]
The known exponential bound already gives $t^{3/2}J_0(t)\to0$. Multiplying by $t^{3/2}$ and integrating to infinity yields the pointwise identities
\begin{align}
 J_0(t)&=t^{-3/2}\int_t^\infty s^{3/2}U(s)\,ds,\notag\\
 H_f(t)&=U(t)+\frac32t^{-5/2}\int_t^\infty s^{3/2}U(s)\,ds.
 \label{d:eq:Volterra-recovery}
\end{align}
This derivation has not used weighted membership of $H_f$ or $f$.

The second operator in~\eqref{d:eq:Volterra-recovery} is bounded on $E_\pi$. Conjugate by the isometry
$V(s)=s^{3/4}e^{s/2}U(s)$ to ordinary $L^2$. Apart from its constant $3/2$, the resulting kernel is
\[
 \boldsymbol1_{s\ge t}\,t^{-1}(s/t)^{3/4}e^{-(s-t)/2}.
\]
For $t\ge\pi$ it is bounded by the integrable function
$\pi^{-1}(1+(s-t)/\pi)^{3/4}e^{-(s-t)/2}$ of $s-t\ge0$.
The elementary $L^1$ convolution bound, after extension by zero, proves boundedness. Therefore
\begin{equation}\label{d:eq:Hf-weighted}
 \norm{H_f}_{E_\pi}\le C\norm{U}_{E_\pi}
 \le C'\norm{u}_{\Hh}.
\end{equation}

To remove the higher theta terms, define
\[
 (Sh)(t)=\sum_{n\ge2}n^4h(n^2t).
\]
On $E_T$, a change of variables gives
\begin{equation}\label{d:eq:dilation-E}
 \norm{n^4h(n^2\,\cdot)}_{E_T}
 \le n^{3/2}e^{-(n^2-1)T/2}\norm{h}_{E_T}.
\end{equation}
Indeed the squared polynomial prefactor is $n^8n^{-3}n^{-2}=n^3$, and on $s\ge n^2T$ the exponential ratio
$e^{s/n^2-s}$ is at most $e^{-(n^2-1)T}$. Hence the norm of $S$ on $E_T$ is less than one for all sufficiently large fixed $T$.

One must not apply this inverse to $f$ before proving $f\in E_T$. Use simultaneously the Banach space
\[
 \mathcal B_T=\{h\in C([T,\infty)):
       \norm{h}_{\mathcal B_T}:=\sup_{t\ge T}e^{t/4}|h(t)|<\infty\}.
\]
There
\begin{equation}\label{d:eq:dilation-B}
 \norm{S}_{\mathcal B_T\to\mathcal B_T}
 \le\sum_{n\ge2}n^4e^{-(n^2-1)T/4}.
\end{equation}
The sums in~\eqref{d:eq:dilation-E} and~\eqref{d:eq:dilation-B} tend to zero as $T$ tends to infinity, by domination by any fixed convergent Gaussian series. Choose one $T$ for which both operator norms are less than one. This is an analytic choice; no numerical tail calculation is needed.

The known contour bound places $f$ and $H_f=(I+S)f$ in $\mathcal B_T$, while~\eqref{d:eq:Hf-weighted} places $H_f$ in $E_T$. The same finite Neumann sums
\[
 f_M=\sum_{k=0}^M(-S)^kH_f
\]
converge to $f$ in $\mathcal B_T$ and to a weighted inverse in $E_T$. The first convergence is uniform on compact subintervals, while the second implies local $L^2$ convergence. Their limits therefore agree almost everywhere. Consequently $f\in E_T$, with
\[
 \norm{f}_{E_T}\le
 (1-\norm{S}_{E_T\to E_T})^{-1}\norm{H_f}_{E_T}
 \le C\norm{u}_{\Hh}.
\]
Finally~\eqref{d:eq:fg-crude} controls the fixed interval $[0,T]$ with a constant times $\norm{u}_{\Hh}$, and the weights $t^{3/2}$ and $(1+t)^{3/2}$ are comparable on $[T,\infty)$. This proves~\eqref{d:eq:recovery-bound}.
\end{proof}

\begin{lemma}[Forward weighted reconstruction]\label{d:lem:forward-reconstruction}
Suppose an even scalar theta source has the exact identity~\eqref{d:eq:U-H-identity} with $H_f=(I+S)f$, and $f\in E_\pi$. Then its physical numerator is admitted to $\Hh$, with
\begin{equation}\label{d:eq:forward-reconstruction}
 \norm{u}_{\Hh}\le C\norm{f}_{E_\pi}.
\end{equation}
\end{lemma}
\begin{proof}
The dilation estimate~\eqref{d:eq:dilation-E} at $T=\pi$ proves boundedness of $I+S$, without requiring contraction. The tail operator
$Jh(t)=t^{-1}\int_t^\infty h(s)\,ds$ is bounded on $E_\pi$: after the same isometry, its kernel is
\[
 \boldsymbol1_{s\ge t}\,t^{-1}(t/s)^{3/4}e^{-(s-t)/2}
 \le\pi^{-1}e^{-(s-t)/2}.
\]
Its norm is at most $2/\pi$. Thus~\eqref{d:eq:U-H-identity} bounds $U$ in $E_\pi$, and the extended-integral comparison~\eqref{d:eq:actual-weight-comparison} proves admission and~\eqref{d:eq:forward-reconstruction}. The scalar source is already defined; the argument does not assume the weighted conclusion.
\end{proof}

\subsection{Construction of the positive fractional power kernel}\label{d:sec:positive-kernel}

For the endpoint analysis restrict
\begin{equation}\label{d:eq:kernel-parameters}
 \frac34\le r<1,\qquad
 \eta=\eta_r=\frac94(1-r),\qquad
 c_0=\frac12,\qquad d_r=c_0^{1/r}.
\end{equation}
Thus $0<\eta\le9/16$ and $2^{-4/3}\le d_r\le c_0$. The symbol $a=i/2$ remains the centering point and is not reused as a fractional-order parameter.

Let $I_\eta$ denote the positive measure
$s^{\eta-1}ds/\Gamma(\eta)$ on $(0,\infty)$, with $I_0=\delta_0$ if the limiting parameter is used. Its Laplace transform is $q^{-\eta}$ in the right half-plane~\cite[Eq.~5.9.1]{DLMF}.

\begin{lemma}[Positive probability and convolution kernels]\label{d:lem:positive-kernel}
For $0<r<1$ and $t\ge0$ there is a positive probability measure $\rho_{r,t}$ on $[0,\infty)$ with
\begin{equation}\label{d:eq:rho-transform}
 \int_0^\infty e^{-qs}\rho_{r,t}(ds)=e^{-tq^r},
 \qquad \operatorname{Re}q>0.
\end{equation}
The family is weakly continuous in $t$. The positive locally finite kernel
\begin{equation}\label{d:eq:k-definition}
 k_r(t,ds)=(I_\eta*\rho_{r,t})(ds)
\end{equation}
is measurable in $t$ and satisfies
\begin{equation}\label{d:eq:k-Laplace}
 \int_0^\infty e^{-qs}k_r(t,ds)=q^{-\eta}e^{-tq^r}.
\end{equation}
\end{lemma}
\begin{proof}
The Gamma integral and integration by parts give, first for positive $q$,
\begin{equation}\label{d:eq:Levy-power}
 q^r=\frac r{\Gamma(1-r)}
      \int_0^\infty(1-e^{-qs})s^{-1-r}\,ds.
\end{equation}
Both endpoints converge for $0<r<1$. For $\epsilon>0$ let
\[
 \nu_\epsilon(ds)=\frac r{\Gamma(1-r)}
           \boldsymbol1_{s\ge\epsilon}s^{-1-r}\,ds.
\]
This measure is finite. Its compound-Poisson probability measure
\[
 e^{-t\nu_\epsilon([0,\infty))}
 \sum_{j=0}^\infty\frac{t^j}{j!}\nu_\epsilon^{*j},
 \qquad \nu_\epsilon^{*0}=\delta_0,
\]
has Laplace transform
$\exp[-t\int(1-e^{-qs})\nu_\epsilon(ds)]$.

These probabilities are tight as $\epsilon\downarrow0$. If $L_\epsilon$ denotes their Laplace transform, their mass above $R>0$ is at most
\[
 (1-e^{-1})^{-1}(1-L_\epsilon(1/R))
 \le C tR^{-r},
\]
because the truncated exponent is at most $R^{-r}$ by~\eqref{d:eq:Levy-power}. Every weak subsequential limit therefore has transform $e^{-tq^r}$ for positive $q$. Uniqueness of Laplace transforms identifies all such limits. For clarity, uniqueness for finite measures follows by the change of variable $y=e^{-s}$ and polynomial approximation on $[0,1]$; for positive locally finite measures with a finite exponential tilt, first apply that tilt. This also gives the uniqueness used below.

The same tightness argument, locally uniformly in $t$, and continuity of $e^{-tq^r}$ prove weak continuity in $t$. The right-half-plane identity~\eqref{d:eq:rho-transform} follows by holomorphy and the identity theorem. Convolution with $I_\eta$ gives~\eqref{d:eq:k-Laplace}. Its finite exponentially weighted mass proves local finiteness. Weak continuity of the probability kernel and integration against the locally finite gamma measure give measurability of the convolution kernel, first on bounded continuous compactly supported tests and then on Borel sets. Positivity permits all these convolution integrations by Tonelli.
\end{proof}

\begin{proposition}[Fractional contraction of the actual Laplace source]\label{d:prop:fractional-source}
For the $f$ of~\eqref{d:eq:fg-source}, define
\begin{equation}\label{d:eq:fr-definition}
 f_r(t)=\mu(q^{9r/4}e^{-tq^r}).
\end{equation}
Then
\begin{equation}\label{d:eq:fr-kernel}
 f_r(t)=\int_0^\infty f(s)k_r(t,ds).
\end{equation}
The integral is absolutely convergent. This identity applies to arbitrary complex sources, without positivity of $f$ or $\mu$.
\end{proposition}
\begin{proof}
The contour lies in $\operatorname{Re}q\ge1/4$, so
$|q^{9/4}e^{-sq}|\le C e^{-s/4}$ there. The integral of this bound against $k_r(t,ds)$ is finite by~\eqref{d:eq:k-Laplace} at $q=1/4$. Thus absolute Fubini permits interchange with the scalar contour. The result is
\[
 \mu\left(q^{9/4}\int_0^\infty e^{-sq}k_r(t,ds)\right)
 =\mu(q^{9/4-\eta}e^{-tq^r})=f_r(t).
\]
The branches agree because $\operatorname{Log}q=2b$ in the strip, so $q^r=e^{2rb}$ and $9/4-\eta=9r/4$ exactly.
\end{proof}

\subsection{A positive resolvent and uniform weighted Schur control}\label{d:sec:weighted-Schur}

We require a bound uniform as $r\uparrow1$, including small values of the input variable. It can be obtained without a pointwise density formula for the probability measures in Lemma~\ref{d:lem:positive-kernel}.

\begin{lemma}[Keyhole resolvent formula]\label{d:lem:keyhole}
For the parameters~\eqref{d:eq:kernel-parameters} and real $q>d_r$,
\begin{equation}\label{d:eq:resolvent-Stieltjes}
 \frac1{q^r-c_0}
 =\frac{d_r^{1-r}/r}{q-d_r}
   +\int_0^\infty\frac{\varrho_r(v)}{q+v}\,dv,
\end{equation}
where
\begin{equation}\label{d:eq:cut-density}
 \varrho_r(v)=\frac{\sin(\pi r)}{\pi}
 \frac{v^r}{v^{2r}-2c_0v^r\cos(\pi r)+c_0^2}\ge0.
\end{equation}
The positive function
\begin{equation}\label{d:eq:pr-positive}
 p_r(s)=\frac{d_r^{1-r}}r e^{d_rs}
             +\int_0^\infty e^{-vs}\varrho_r(v)\,dv
\end{equation}
has Laplace transform $(q^r-c_0)^{-1}$ for $q>d_r$, and satisfies uniformly for $s>0$ and $3/4\le r<1$,
\begin{equation}\label{d:eq:pr-bound}
 p_r(s)\le C e^{d_rs}(1+s^{-1/4}).
\end{equation}
\end{lemma}
\begin{proof}
Use the principal branch of $z^r$ on $\C\setminus(-\infty,0]$. The function $(z^r-c_0)^{-1}$ has one simple pole there at $z=d_r$, of residue $d_r^{1-r}/r$. Apply a keyhole Cauchy contour to
$(z^r-c_0)^{-1}/(z-q)$. The large-circle integral tends to zero as $O(R^{-r})$, while the small circle tends to zero because $(z^r-c_0)^{-1}$ remains bounded at zero. On the upper bank of the negative axis its negative imaginary part is
\[
 \frac{v^r\sin(\pi r)}
 {v^{2r}-2c_0v^r\cos(\pi r)+c_0^2}.
\]
The lower bank is its conjugate. Their oriented difference and the simple pole give~\eqref{d:eq:resolvent-Stieltjes}, including its positive sign. The cut integral converges at both endpoints. Integrating the nonnegative expression~\eqref{d:eq:pr-positive} against $e^{-qs}$ and using Tonelli proves its claimed Laplace transform.

Since $\cos(\pi r)\le0$ in the specified range, the denominator in~\eqref{d:eq:cut-density} is at least $v^{2r}+c_0^2$. For $0<v\le1$, $\varrho_r(v)\le C v^r$; for $v\ge1$,
$\varrho_r(v)\le C\sin(\pi r)v^{-r}$. Thus the cut term in~\eqref{d:eq:pr-positive} is bounded by
\[
 C+C\sin(\pi r)\Gamma(1-r)s^{r-1}
 \le C'(1+s^{-1/4}).
\]
The last step uses
$\sin(\pi r)\Gamma(1-r)=\pi/\Gamma(r)$, uniformly bounded on $[3/4,1]$, by the Gamma reflection formula~\cite[Eq.~5.5.3]{DLMF}. For $s\le1$, $s^{r-1}\le s^{-1/4}$, and for $s\ge1$ it is at most one. The residue coefficient is uniformly bounded and $e^{d_rs}\ge1$, proving~\eqref{d:eq:pr-bound}.
\end{proof}

\begin{lemma}[Integrated potentials and their densities]\label{d:lem:potentials}
Uniformly in the same parameter range,
\begin{equation}\label{d:eq:cubic-potential}
 p_r^{*3}(s)\le C e^{d_rs}(s^{5/4}+s^2).
\end{equation}
For every integer $j\ge0$, the positive locally finite measure
\[
 M_j(ds)=\int_0^\infty t^je^{c_0t}k_r(t,ds)\,dt
\]
has the density
\begin{equation}\label{d:eq:potential-identity}
 M_j(ds)=j!\,[I_\eta*p_r^{*(j+1)}](s)\,ds.
\end{equation}
In particular these measures have no atom at zero.
\end{lemma}
\begin{proof}
For~\eqref{d:eq:cubic-potential}, expand the triple convolution of the bound
$C e^{d_rs}(1+s^{-1/4})$. If $k$ of the three factors are singular, the beta integral gives a constant times $e^{d_rs}s^{2-k/4}$, for $0\le k\le3$. The beta parameters are positive; explicitly
\[
 \int_0^s y^{\alpha-1}(s-y)^{\beta-1}\,dy
 =s^{\alpha+\beta-1}
       \frac{\Gamma(\alpha)\Gamma(\beta)}{\Gamma(\alpha+\beta)}
 \quad(\alpha,\beta>0),
\]
as in~\cite[Eq.~5.12.1]{DLMF}. The intermediate powers are bounded by a constant times $s^{5/4}+s^2$.

For real $q>d_r$, Tonelli and~\eqref{d:eq:k-Laplace} give
\[
 \int_0^\infty e^{-qs}M_j(ds)
 =q^{-\eta}\int_0^\infty t^je^{-t(q^r-c_0)}\,dt
 =\frac{j!\,q^{-\eta}}{(q^r-c_0)^{j+1}}.
\]
This is finite, so $M_j$ is locally finite. The right side is also the Laplace transform of the measure on the right of~\eqref{d:eq:potential-identity}, by Lemma~\ref{d:lem:keyhole}. Laplace uniqueness after a fixed finite exponential tilt proves the measure identity. The convolution on its right is locally integrable and is a density, proving the no-atom assertion.
\end{proof}

We will also use, for $\beta\ge0$,
\begin{equation}\label{d:eq:gamma-convolution-bound}
 [I_\eta*(e^{d_r\,\cdot}\,\cdot^\beta)](s)
 \le e^{d_rs}s^\beta\int_0^\infty e^{-d_ry}I_\eta(dy)
 =d_r^{-\eta}e^{d_rs}s^\beta.
\end{equation}
The factor $d_r^{-\eta}$ is uniformly bounded under~\eqref{d:eq:kernel-parameters}. This formulation preserves the cancellation between the singular gamma integral and $1/\Gamma(\eta)$ as $\eta\downarrow0$.

\begin{theorem}[Uniform weighted kernel estimate]\label{d:thm:weighted-kernel}
There is a constant independent of $3/4\le r<1$ such that
\begin{align}
 &\int_\pi^\infty(1+t)^{3/2}e^t
       \left|\int_0^\infty f(s)k_r(t,ds)\right|^2dt\notag\\
 &\hspace{35mm}\le C\int_0^\infty(1+s)^{3/2}e^s|f(s)|^2\,ds.
 \label{d:eq:uniform-kernel}
\end{align}
For general $E$ elements the integral is interpreted by the bounded extension; for the continuous source in~\eqref{d:eq:fg-source} it agrees with the absolutely convergent integral~\eqref{d:eq:fr-kernel}.
\end{theorem}
\begin{proof}
Set
\[
 K_r(t,ds)=e^{c_0t-c_0s}k_r(t,ds).
\]
Its mass and first moment follow from~\eqref{d:eq:k-Laplace} and one derivative at $q=c_0$:
\begin{align}
 \int K_r(t,ds)&=c_0^{-\eta}e^{-t(c_0^r-c_0)},\notag\\
 \int sK_r(t,ds)&=
 c_0^{-\eta}e^{-t(c_0^r-c_0)}
        [\eta/c_0+trc_0^{r-1}].
 \label{d:eq:kernel-moments}
\end{align}
The derivative is legitimate by exponential integrability. Since $c_0^r\ge c_0$ and the remaining parameters stay in bounded ranges, Cauchy--Schwarz for the positive finite measure gives
\begin{equation}\label{d:eq:Schur-row}
 \int(1+s)^{1/2}K_r(t,ds)\le C(1+t)^{1/2}.
\end{equation}

The required column estimate is a measure domination on the whole input half-line:
\begin{equation}\label{d:eq:Schur-column}
 \int_\pi^\infty(1+t)^2K_r(t,ds)\,dt
 \le C(1+s)^2\,ds.
\end{equation}
To prove it, use $(1+t)^2\le Ct^2$ for $t\ge\pi$, and then the $j=2$ case of~\eqref{d:eq:potential-identity}. By~\eqref{d:eq:cubic-potential} and~\eqref{d:eq:gamma-convolution-bound}, the resulting density is bounded by
\[
 C e^{-c_0s}[I_\eta*p_r^{*3}](s)
 \le C'e^{-(c_0-d_r)s}(s^{5/4}+s^2)
 \le C''(1+s)^2.
\]
This holds also for $0<s<1$, and the potential measure has no atom at zero. Thus no small-input region has been discarded. The restriction $t\ge\pi$ matches the complete physical half-line $x\ge0$, not an omitted physical tail.

Write $G(s)=e^{c_0s}f(s)$ and $G_r(t)=e^{c_0t}f_r(t)$. Weighted Cauchy--Schwarz gives
\[
 |G_r(t)|^2
 \le\left[\int(1+s)^{1/2}K_r(t,ds)\right]
       \left[\int |G(s)|^2(1+s)^{-1/2}K_r(t,ds)\right].
\]
Multiply by $(1+t)^{3/2}$, apply~\eqref{d:eq:Schur-row}, integrate over $t\ge\pi$, and use Tonelli and~\eqref{d:eq:Schur-column}. The result is
\[
 \int_\pi^\infty(1+t)^{3/2}|G_r(t)|^2dt
 \le C\int_0^\infty(1+s)^{3/2}|G(s)|^2ds,
\]
which is~\eqref{d:eq:uniform-kernel}. The argument starts with bounded compactly supported measurable $G$ and extends by completion. The column domination ensures that input null sets are ignored for almost every output parameter. For the actual continuous exponentially bounded $f$, absolute convergence from Proposition~\ref{d:prop:fractional-source} identifies the extension with its ordinary kernel integral. Cauchy--Schwarz applies without change to complex inputs.
\end{proof}


\section{Compact division and supporting spectral distinctions}\label{app:compactmodel}
The following compact-model results describe exact root spaces and explain why a compact spectral list alone would not establish synthesis. They are supporting results, not substitutes for Appendix~\ref{app:abel}.


\subsection{Compact division and the origin boundary condition}\label{d:sec:division}

For $w\in\C$, write $S_w(s)=\sin(ws)/w$ with $S_0(s)=s$. On the positive half-line set
\begin{equation}\label{d:eq:V-definition}
 (V_wu)(x)=\int_x^\infty S_w(t-x)u(t)\,dt,\qquad x\ge0,
\end{equation}
and extend evenly.

\begin{theorem}[Hilbert Schmidt division]\label{d:thm:division}
For every $w$, $V_w$ is a Hilbert--Schmidt operator on $\Hh$. The map $w\mapsto V_w$ is entire with values in $\HS(\Hh)$ and is bounded there on each horizontal strip $|\operatorname{Im}w|\le B$. If $F_u(w)=0$, then $g=V_wu$ satisfies, distributionally on $\R$,
\begin{equation}\label{d:eq:division-equations}
 g''+w^2g=u,\qquad
 F_g(z)=\frac{F_u(z)}{w^2-z^2},
\end{equation}
where all apparent singularities are removable. Thus division in~\eqref{d:eq:division-equations} is bounded and compact on the closed subspace $\{u:F_u(w)=0\}$. Without that condition the even extension has a derivative jump at zero and the displayed division identity need not hold.
\end{theorem}
\begin{proof}
The isometry $u\mapsto\sqrt2u/\sqrt W$ from $\Hh$ to $L^2(0,\infty)$ transforms $V_w$ into the kernel
\[
 \boldsymbol1_{t\ge x}S_w(t-x)\sqrt{W(t)/W(x)}.
\]
For $x,s\ge0$,~\eqref{d:eq:W-two-sided} gives
\begin{equation}\label{d:eq:weight-ratio}
 \frac{W(x+s)}{W(x)}\le C\exp\bigl(4s-\pi e^{2x}(e^{2s}-1)\bigr).
\end{equation}
Using $S_w(s)=\int_0^s\cos(wr)\,dr$, we have $|S_w(s)|\le s e^{Bs}$ whenever $|\operatorname{Im}w|\le B$. Thus the squared kernel is bounded by
\begin{equation}\label{d:eq:HS-majorant}
 Cs^2\exp\bigl((4+2B)s-\pi e^{2x}(e^{2s}-1)\bigr).
\end{equation}
For $x$ in a bounded interval, the integral in $s$ is finite by double-exponential decay. Choose $X$ so that $\pi e^{2X}\ge4+2B$. For $x\ge X$, the inequality $e^{2s}-1\ge2s$ makes the exponent at most $-\pi e^{2x}s$. The $s^2$ integral is then $2\pi^{-3}e^{-6x}$, integrable in $x$. This proves both the Hilbert--Schmidt property and the uniform strip bound. Parameter derivatives introduce only additional fixed powers of $s$; the same estimates justify holomorphy in the Hilbert--Schmidt norm.

The exponential integrability following~\eqref{d:eq:eval-bound} justifies differentiating the tail integral in the weak sense. On $x>0$,
$g''+w^2g=u$, while
\begin{equation}\label{d:eq:jump}
 g'(0+)=-\int_0^\infty\cos(wt)u(t)\,dt=-\tfrac12F_u(w).
\end{equation}
The even extension therefore satisfies, on all of $\R$,
\begin{equation}\label{d:eq:jump-distribution}
 g''+w^2g=u-F_u(w)\delta_0.
\end{equation}
Exactly the stipulated vanishing removes the delta term. Since both $g,u\in\Hh$ have all exponential moments, their distributional Fourier identity extends from the real axis to the entire functions in~\eqref{d:eq:division-equations}. This proves the quotient identity and removability. At $w=0$, evenness makes $F_u(0)=0$ a zero of order at least two, as required by the denominator $-z^2$.

There is no nonzero even solution of $g''+w^2g=0$ in $\Hh$: such a solution is a scalar cosine, or a constant when $w=0$, and cannot be square integrable against the reciprocal double-exponential weight. Thus the admitted solution is unique. Restricting a Hilbert--Schmidt operator to a closed subspace is compact, completing the proof.
\end{proof}

Every fixed finite sequence of divisions is bounded whenever the remaining zero orders permit it. No bound uniform in the number of divisions, no uniform spectral-projection bound, and no frame estimate has been established. Compact division is not a polynomial-density theorem.

\subsection{A centered compact backward operator}\label{d:sec:model}

Put
\begin{equation}\label{d:eq:d-delta-T}
 d=\frac14,\qquad a=\frac i2,\qquad
 \Delta(z)=a^2-z^2=-\frac14-z^2,\qquad
 T=CV_a|_{\Hh_c}.
\end{equation}
Because $F_u(a)=0$ for $u\in\Hh_c$, $V_a$ is genuine division on this whole subspace. Write
\begin{equation}\label{d:eq:sj}
 s_j=(D^2-d)^j\kappa\qquad(j\ge1).
\end{equation}

\begin{theorem}[Dense range and all generalized nullspaces]\label{d:thm:backward}
The operator $T$ is compact on $\Hh_c$, its range is dense, and for every $n\ge1$,
\begin{equation}\label{d:eq:kernel-exact}
 \ker T^n=\operatorname{span}\{s_1,\ldots,s_n\}
 =\operatorname{span}\{r_1,\ldots,r_n\}.
\end{equation}
These spaces have dimension $n$. The chain action is
\begin{equation}\label{d:eq:shift-chain}
 Ts_1=0,\qquad Ts_j=s_{j-1}\quad(j\ge2),
\end{equation}
and consequently
\begin{equation}\label{d:eq:N-generalized}
 \Nn=\overline{\bigcup_{n\ge1}\ker T^n}^{\Hh}.
\end{equation}
\end{theorem}
\begin{proof}
Compactness follows from Theorem~\ref{d:thm:division} and boundedness of $C$. Since $F_{s_j}=\Delta^j\Xi$, all $s_j$ are centered. Division and centering give~\eqref{d:eq:shift-chain}. The polynomial $\Delta^j$ has leading term $(-z^2)^j$ and vanishes at $a$, so a triangular change of basis gives the two finite spans in~\eqref{d:eq:kernel-exact}. Linear independence follows by dividing any transform relation by the nonzero entire function $\Xi$.

If $Tu=0$, then $V_au=\alpha\kappa$ for some scalar $\alpha$. Applying $D^2-d$ gives $u=\alpha s_1$. More generally, if $Tu$ lies in $\operatorname{span}\{s_1,\ldots,s_n\}$, then $V_au$ is that vector plus a multiple of $\kappa$. Applying $D^2-d$ puts $u$ in $\operatorname{span}\{s_1,\ldots,s_{n+1}\}$. Induction, together with~\eqref{d:eq:shift-chain}, proves the exact kernel formula rather than only one inclusion.

For every even $g\in C_c^\infty(\R)$, let $u=(D^2-d)g$. Integration by parts against $c$ gives $m(u)=0$. The uniqueness part of Theorem~\ref{d:thm:division} gives $V_au=g$ and hence $Tu=Cg$. Even compact smooth functions are dense in $\Hh$: first truncate in the weighted norm, then approximate on a compact set where $W$ is positive and continuous, and symmetrize. Their images under $C$ are therefore dense in $\Hh_c$. This proves dense range and the last assertion.
\end{proof}

For later use the transform action is
\begin{equation}\label{d:eq:T-transform}
 F_{Tu}(z)=\frac{F_u(z)}{\Delta(z)}-\beta(u)\Xi(z),\qquad
 \beta(u)=2F_{V_au}(a).
\end{equation}
In particular $T$ is not self-adjoint: its nontrivial kernel and dense range would be incompatible with self-adjointness.

\begin{proposition}[The invariant divisible ideal]\label{d:prop:T-ideal}
The spaces $\Nn$ and $\Ii$ are $T$-invariant. If $u\in\Ii$ and $F_u=\Xi Q$, then
\begin{equation}\label{d:eq:T-quotient}
 F_{Tu}=\Xi\bigl(R(z)-R(a)\bigr),\qquad
 R(z)=\frac{Q(z)}{\Delta(z)}.
\end{equation}
The function $R$ is even entire.
\end{proposition}
\begin{proof}
Invariance of $\Nn$ follows from~\eqref{d:eq:shift-chain} and continuity. For $\Ii$, the condition $Q(a)=0$ and evenness give zeros at both $a$ and $-a$, so division by $\Delta$ is removable. Since $2\Xi(a)=1$,~\eqref{d:eq:T-transform} becomes~\eqref{d:eq:T-quotient}, whose quotient again vanishes at $a$.
\end{proof}

\subsection{Exact nonzero root spaces and multiplicities}\label{d:sec:roots}

A root space of a bounded operator at $\lambda$ means the union of the kernels of all powers of that operator minus $\lambda$. For a nonzero eigenvalue of a compact operator this is finite dimensional and the union stabilizes. Put
\begin{equation}\label{d:eq:lambda-w}
 \lambda(w)=\frac1{\Delta(w)}=-\frac1{w^2+1/4}.
\end{equation}
It is defined at every zero of $\Xi$ by~\eqref{d:eq:Xi-xi}.

\begin{theorem}[All nonzero adjoint root spaces]\label{d:thm:adjoint}
For each distinct zero pair $\{w,-w\}$ of $\Xi$, the root space of $T^*$ at $\lambda(w)$ is exactly
\begin{equation}\label{d:eq:adjoint-root}
 \operatorname{span}\{J_{w,k}:0\le k<\nu(w)\}.
\end{equation}
It has dimension $\nu(w)$ and consists of one Jordan block. These are all the nonzero root spaces of $T^*$. For the Riesz waves $E_{w,k}$ the corresponding eigenvalue is $\overline{\lambda(w)}$. The sign pair yields the same space, not two copies. Consequently $\Zz$ is the closed span of all nonzero adjoint root spaces.
\end{theorem}
\begin{proof}
For $k<\nu(w)$, differentiate~\eqref{d:eq:T-transform}. All derivatives of the correction term that occur vanish at $w$. Using~\eqref{d:eq:Riesz} and the linear-first convention gives the raw-wave formula
\begin{equation}\label{d:eq:adjoint-action}
 T^*J_{w,k}=\sum_{\ell=0}^k\binom{k}{\ell}
 \lambda^{(\ell)}(w)J_{w,k-\ell}.
\end{equation}
The vectors in~\eqref{d:eq:adjoint-root} are linearly independent. Indeed their exponential-polynomial expressions have distinct leading polynomial degrees at the two distinct frequencies $iw$ and $-iw$; such expressions are independent for $w\ne0$. One may verify this by successively applying powers of $D-iw$ and $D+iw$ to isolate each frequency. Here $w\ne0$ by~\eqref{d:eq:Xi-xi}. Also
\[
 \lambda'(w)=\frac{2w}{\Delta(w)^2}\ne0.
\]
Thus the first adjacent coefficient in the triangular action~\eqref{d:eq:adjoint-action} never vanishes. It is one block of length $\nu(w)$. The identity $J_{-w,k}=(-1)^kJ_{w,k}$ handles the sign pair.

It remains to exhaust the root spaces. Let $h$ be a root vector of $T^*$ at some $\lambda\ne0$, and put $v=h/W$. This is an even locally square-integrable distribution. For even compact smooth $g$, Theorem~\ref{d:thm:backward} gives $T(D^2-d)g=Cg$. Since $h\in\Hh_c$,
\begin{equation}\label{d:eq:root-test}
 \ip{(D^2-d)g}{T^*h}_{\Hh}=\ip{g}{h}_{\Hh}.
\end{equation}
For a chain $T^*h=\lambda h+h_-$, conjugating the tested distributional identity gives
\[
 (D^2-d)(\lambda v+v_-)=v,
 \qquad (D^2-d-\lambda^{-1})v=-\lambda^{-1}(D^2-d)v_-.
\]
The constant-coefficient factors commute. Iterating along a chain proves
\begin{equation}\label{d:eq:root-ODE}
 (D^2-d-\lambda^{-1})^n v=0
\end{equation}
for some $n$. Even test functions suffice because each distribution involved is even.

Choose $w$ such that $\Delta(w)=\lambda^{-1}$. Then~\eqref{d:eq:root-ODE} is $(D^2+w^2)^nv=0$. Distributional solutions of a constant-coefficient ordinary differential equation are the usual exponential polynomials: factor the operator into first-order factors and solve each distributional equation after multiplication by its integrating exponential. For $w\ne0$, its even solutions are therefore finite linear combinations of $\partial_z^k\cos(zx)|_{z=w}$. For $w=0$, they are even polynomials. It remains to determine which of these possibilities survive the Hilbert-space orthogonality constraints.

Every nonzero adjoint root vector annihilates $\ker T^m$ for every $m$. To see this, $T^*$ is invertible on its finite root space, so write $h=(T^*)^mh_m$ within that space. Then $\ip{u}{h}=\ip{T^mu}{h_m}=0$ for $u\in\ker T^m$. In particular $h\perp\Nn$.

Suppose $w\ne0$ and write $h=\sum_{k=0}^M c_kJ_{w,k}$ with $c_M\ne0$. Let $\nu$ be the zero order of $\Xi$ at $w$, using $\nu=0$ if it is not a zero. If $M\ge\nu$, choose the even polynomial
\[
 p(z)=\Delta(z)(z^2-\overline{w}^{\,2})^{M-\nu}.
\]
It vanishes at $a$, so the source $u_p$ with transform $p\Xi$ is a finite radical combination by~\eqref{d:eq:derivative-transforms}. Since $\Delta(\overline w)\ne0$, this transform has exact order $M$ at $\overline w$. But
\[
 \ip{u_p}{h}_{\Hh}
 =\sum_{k=0}^M\overline{c_k}(p\Xi)^{(k)}(\overline w)
 =\overline{c_M}(p\Xi)^{(M)}(\overline w)\ne0,
\]
contradicting $h\perp\Nn$. Hence $M<\nu$. This proves that $w$ must be a zero and that no jet above the asserted order is present.

If $w=0$, write $v=\sum_{j=0}^M c_jx^{2j}$ with $c_M\ne0$. The centered polynomial $p(z)=\Delta(z)z^{2M}$ makes $p\Xi$ have exact zero order $2M$ at zero, since $\Xi(0)\ne0$. The same pairing, using $F_u^{(2j)}(0)=(-1)^j\int x^{2j}u(x)\,dx$, is nonzero. This excludes $w=0$. The case $\Delta(w)=0$ is impossible because $\lambda$ is finite and nonzero. Exhaustion, multiplicity, and the final span statement follow.
\end{proof}

\begin{corollary}[Primal divided chains]\label{d:cor:primal}
For a zero $w$ of multiplicity $\nu$, set $s=w^2+1/4$ and define $h_j\in\Hh$ by
\begin{equation}\label{d:eq:primal-source}
 F_{h_j}(z)=\frac{\Xi(z)}{(w^2-z^2)^j},\qquad
 \phi_j=h_j-s^{-j}\kappa,\qquad 1\le j\le\nu.
\end{equation}
The vectors $h_j$ exist, and the vectors $\phi_j$ are centered and constitute the full root space of $T$ at $-s^{-1}$, of dimension $\nu$. Their action is
\begin{equation}\label{d:eq:primal-action}
 T\phi_j=-\sum_{r=1}^j s^{-(j-r+1)}\phi_r.
\end{equation}
If $j<\nu$, $F_{\phi_j}$ vanishes in value at every zero of $\Xi$, but $\phi_j\notin\Ii$.
\end{corollary}
\begin{proof}
Repeated application of Theorem~\ref{d:thm:division}, starting with $\kappa$, gives $h_j$ for $j\le\nu$. At $a$, its transform is $s^{-j}/2$, giving centering. With $y=w^2-z^2$ and $\Delta=y-s$, the identity
\[
 \frac{y^{-j}-s^{-j}}{y-s}
 =-\sum_{r=1}^j s^{-(j-r+1)}y^{-r}
\]
and centering prove~\eqref{d:eq:primal-action}. The vectors are independent by the resulting rational-function identity after division by $\Xi$. Their adjacent block coefficients are nonzero. For a compact operator the nonzero algebraic multiplicities of $T$ and $T^*$ at conjugate eigenvalues agree; this follows by restricting the two adjoint Riesz projections to their finite-dimensional ranges. Theorem~\ref{d:thm:adjoint}, applied at $\overline w$, therefore shows there is no larger primal root space.

At the pair $\{w,-w\}$, the first term in~\eqref{d:eq:primal-source} has zero order $\nu-j$, whereas the subtracted multiple of $\Xi$ has order $\nu$. The resulting order is exactly $\nu-j$. For $j<\nu$ this is positive but less than $\nu$; all other zeros retain their values. This proves the last assertion without identifying the values-only kernel with $\Ii$.
\end{proof}

The root-space theorems use only multiplicities that actually occur. They neither assume nor prove simplicity. Their algebraic description alone does not establish completeness; that conclusion follows from Theorem~\ref{d:thm:full-synthesis}.

\subsection{The compact quasinilpotent model}\label{d:sec:gap}

\begin{theorem}[Quasinilpotent restriction and synthesis criterion]\label{d:thm:quasinilpotent}
The compact operator $T|_{\Ii}$ is quasinilpotent. Its generalized nullspace is precisely the algebraic span of all the $r_j$. Hence global synthesis~\eqref{d:eq:ideal-synthesis} is equivalent to
\begin{equation}\label{d:eq:remaining-density}
 \overline{\bigcup_{n\ge1}\ker(T|_{\Ii})^n}^{\Hh}=\Ii.
\end{equation}
Before using synthesis, the quotient $\Ii/\Nn$ carries a compact quasinilpotent induced operator and identifies as a Hilbert space with $\Ii\cap\Nn^\perp$. Theorem~\ref{d:thm:full-synthesis} proves this quotient is zero; compactness alone does not.
\end{theorem}
\begin{proof}
Suppose $Tu=\lambda u$ with $u\in\Ii$ and $\lambda\ne0$. For some scalar $\alpha$,
$V_au=\lambda u+\alpha\kappa$. Applying $D^2-d$ and transforming gives
\[
 [1-\lambda\Delta(z)]F_u(z)=\alpha\Delta(z)\Xi(z).
\]
Write $F_u=\Xi Q$ and cancel $\Xi$ away from its discrete zero set, then continue analytically:
\begin{equation}\label{d:eq:no-eigen}
 [1-\lambda\Delta(z)]Q(z)=\alpha\Delta(z).
\end{equation}
At either root of $1-\lambda\Delta(z)$, the right side is $\alpha/\lambda$. Thus $\alpha=0$, including when the root is the double point zero. Equation~\eqref{d:eq:no-eigen} then forces $Q=0$. There are no nonzero eigenvalues of the restriction. A compact operator has no nonzero spectral points except eigenvalues; hence its spectrum is $\{0\}$. The space $\Ii$ is infinite dimensional because it contains the independent radical chains, so zero does belong to this spectrum.

The exact kernels follow from Theorem~\ref{d:thm:backward}, since every $s_j$ belongs to $\Ii$. Their closed span is $\Nn$, proving the equivalence. Both $\Nn$ and $\Ii$ are invariant, so the quotient operator is defined and compact. Its $n$th power has norm at most $\norm{(T|_{\Ii})^n}$, proving spectral radius zero. The usual orthogonal representative identifies the quotient Hilbert space with $\Ii\ominus\Nn$ and represents its action by compression. These structural facts do not by themselves prove that the space is zero; the uniform endpoint estimate will do so.
\end{proof}

Compactness, dense range, and the exact dimensions in~\eqref{d:eq:kernel-exact} do not imply~\eqref{d:eq:remaining-density}. The following elementary example records the precise logical obstruction.

\begin{proposition}[A compact model with an invisible residual part]\label{d:prop:countermodel}
There is a compact quasinilpotent operator with dense range, a one-dimensional kernel, and an $n$-dimensional $n$th kernel for every $n$, whose generalized nullspace is not dense.
\end{proposition}
\begin{proof}
On $\ell^2(\mathbb N_0)$ define $Be_0=0$ and $Be_n=n^{-2}e_{n-1}$ for $n\ge1$. The weights tend to zero, so $B$ is compact, and every basis vector lies in its range. Moreover
\[
 \ker B^n=\operatorname{span}\{e_0,\ldots,e_{n-1}\},\qquad
 \norm{B^n}=(n!)^{-2}.
\]
Thus it has dense range and spectral radius zero.

On $L^2(0,1)$ let $(Vf)(x)=\int_0^x f(t)\,dt$. Its triangular kernel is square integrable, so it is Hilbert--Schmidt. It is injective. Its adjoint is $V^*f(t)=\int_t^1f(x)\,dx$, also injective, so $V$ has dense range. The iterated kernel is
$\boldsymbol1_{t<x}(x-t)^{n-1}/(n-1)!$. Its row and column integrals are bounded by $1/n!$, so $\norm{V^n}\le1/n!$ and $V$ is quasinilpotent.

The direct sum $B\oplus V$ has all the claimed properties, but
\[
 \overline{\bigcup_{n\ge1}\ker(B\oplus V)^n}
 =\ell^2(\mathbb N_0)\oplus\{0\}
 \ne\ell^2(\mathbb N_0)\oplus L^2(0,1).
\]
\end{proof}

This is not a counterexample for the theta space. It shows that the compact-model properties alone cannot exclude a residual quotient. The actual theta-space conclusion uses the additional uniform weighted endpoint estimate of Theorem~\ref{d:thm:uniform-endpoint}, rather than this invalid generic inference.

\subsection{Failure of the inherited de Branges norm axiom}\label{d:sec:debranges}

Every even entire transform can be written uniquely as $f(\lambda)=F_u(z)$ with $\lambda=z^2$. Transport the $\Hh$ norm to this entire-function space in $\lambda$. A necessary de Branges axiom says that, at a nonreal zero $\alpha$ of $f$, the entire function
\begin{equation}\label{d:eq:zero-flip}
 f(\lambda)\frac{\lambda-\overline\alpha}{\lambda-\alpha}
\end{equation}
has the same norm as $f$. This is axiom (H1) in de Branges's original formulation~\cite{deBranges1961}. Bounded division alone does not provide this isometric statement.

\begin{theorem}[Explicit zero flip obstruction]\label{d:thm:debranges}
The Fourier image of $\Hh$, written in the variable $\lambda=z^2$ with its inherited norm, fails the axiom~\eqref{d:eq:zero-flip}. The centered image of $\Hh_c$ fails it as well. The counterexample can be chosen so that both physical numerators are even and compactly supported smooth functions.
\end{theorem}
\begin{proof}
Write $\omega=1/W$ and take even $g\in C_c^\infty(\R;\C)$. For nonreal $\alpha$, define
\[
 u=g''+\alpha g,\qquad \widetilde u=g''+\overline\alpha g.
\]
If $G(\lambda)=F_g(z)$, their transforms are $(\alpha-\lambda)G(\lambda)$ and $(\overline\alpha-\lambda)G(\lambda)$, exactly related by~\eqref{d:eq:zero-flip}. Both are in the space. Direct expansion in the linear-first convention gives
\begin{align}
 \norm{u}_{\Hh}^2-\norm{\widetilde u}_{\Hh}^2
 &=4\operatorname{Im}\alpha\,\operatorname{Im}\int_\R g''\overline g\,\omega\,dx\notag\\
 &=-4\operatorname{Im}\alpha\int_\R\omega'(x)
           \operatorname{Im}(g'(x)\overline{g(x)})\,dx.
 \label{d:eq:flip-norm}
\end{align}
The integration by parts has no endpoint term.

Differentiating the theta asymptotic gives
$(\log\omega)'(x)=2\pi e^{2x}-4+o(1)>0$ for sufficiently large positive $x$. Choose two reflected compact intervals in these tails. On their union $x\omega'(x)>0$. Let $\eta$ be a nonzero real even compact smooth function supported there and put $g(x)=\eta(x)e^{itx^2}$ with real $t\ne0$. Then
\[
 \operatorname{Im}(g'\overline g)=2tx\eta^2,
\]
so~\eqref{d:eq:flip-norm} becomes
\begin{equation}\label{d:eq:flip-nonzero}
 \norm{u}_{\Hh}^2-\norm{\widetilde u}_{\Hh}^2
 =-8t\operatorname{Im}\alpha\int_\R x\omega'(x)\eta(x)^2\,dx\ne0.
\end{equation}
This disproves the axiom for the full image.

For the centered image, choose $\eta$ additionally in the common kernel of the two real linear functionals
\[
 \eta\longmapsto\int_\R c(x)\eta(x)\cos(tx^2)\,dx,
 \qquad
 \eta\longmapsto\int_\R c(x)\eta(x)\sin(tx^2)\,dx.
\]
Such a nonzero choice exists because the real even compact test space on the chosen intervals is infinite dimensional and these impose at most two linear constraints. Now $m(g)=0$, and integration by parts gives $m(g'')=dm(g)=0$. Hence $u,\widetilde u\in\Hh_c$, while the strict sign in~\eqref{d:eq:flip-nonzero} persists. The transform is nonzero by Fourier injectivity and the absence of a nonzero compactly supported solution of $g''+\alpha g=0$. This supplies a genuine nonzero function with a zero at $\alpha$ in the centered space as well.
\end{proof}

The conclusion concerns the actual inherited norm. It does not exclude every equivalent norm, auxiliary entire-function space, nonisometric embedding, or other completeness theorem. Any such approach must establish its own identification and hypotheses before drawing a synthesis conclusion.


\begin{proposition}[Semigroup structure alone is insufficient]\label{d:prop:semigroup-model}
On $\ell^2(\mathbb N)\oplus\C$, define
\[
 B_r(x,c)=\bigl((r^{2n}(x_n+c))_{n\ge1},0\bigr),\qquad 0<r<1.
\]
Each $B_r$ is bounded, compact, and injective, and $B_rB_s=B_{rs}$. Its range closure is $\ell^2(\mathbb N)\oplus\{0\}$, while $B_re_n=r^{2n}e_n$ there. Nevertheless
\[
 \norm{B_r(0,1)}^2=\frac{r^4}{1-r^4}\longrightarrow\infty.
\]
\end{proposition}
\begin{proof}
The map is a compact diagonal plus a rank-one operator. A zero image requires the constant sequence $x_n=-c$, which is square summable only when $c=0$. Direct substitution gives the composition law. Finite sequences in the first summand belong to the range, while the second output coordinate is always zero. The geometric series gives the displayed norm.
\end{proof}

This abstract example is not a theta counterexample: it has no analogue of the proved recovery estimate~\eqref{d:eq:recovery-bound} and uniform kernel bound~\eqref{d:eq:uniform-kernel}. It explains why those estimates, rather than semigroup structure alone, are the decisive steps.


\section{Finite clusters and the fixed complexity proof}\label{app:clusters}


\subsection{Uniform envelopes in a bounded frequency window}\label{e:sec:ode}

\begin{lemma}[Compact ODE estimate]\label{e:lem:ode}
Fix $N\in\N$ and $0\le M<\infty$. If $P$ has complexity at most $N$ and all
its frequencies obey $|\lambda_j|\le M$, then, with $A=WP$,
\begin{equation}\label{e:eq:envelope}
 \norm A_W+\norm{A'}_2\le C_{N,M,W}\norm A_2.
\end{equation}
The estimate is uniform when distinct frequencies approach one another
and when polynomial coefficients vary arbitrarily.
\end{lemma}
\begin{proof}
Fix the total order $m=\sum_j(\deg p_j+1)\le N$. The function $P$
satisfies the monic differential equation
\begin{equation}\label{e:eq:annihilator}
 \prod_j(\partial_x-i\lambda_j)^{\deg p_j+1}P=0.
\end{equation}
The coefficients of the monic polynomial on the left lie in a compact set:
they are continuous symmetric polynomials in $m$ roots from a closed disk.
Repeated roots are included in that set. Let $B_a$ be the corresponding
companion matrix and set $Y=(P,P',\ldots,P^{(m-1)})^{\mathsf T}$. Then
$Y'=B_aY$, and $\norm{B_a}\le L$ uniformly in this compact family.

We claim the uniform observability estimate
\begin{equation}\label{e:eq:observable}
 |Y(0)|^2\le C_m\int_{-1}^1|P(x)|^2\,dx.
\end{equation}
Otherwise there would be coefficient vectors $a_n$ in the compact family
and unit initial vectors $v_n$ for which the integrals tend to zero.
Extract $a_n\to a$ and $v_n\to v$, with $|v|=1$. The power series of the
matrix exponential converges uniformly for $|x|\le1$ and bounded matrices,
so the first components of $e^{xB_{a_n}}v_n$ converge uniformly to that of
$e^{xB_a}v$. The latter component vanishes identically on $[-1,1]$.
Its derivatives of orders $0,\ldots,m-1$ at zero are the components of $v$,
so $v=0$, a contradiction.

The matrix exponential estimate gives
$|Y(x)|\le e^{L|x|}|Y(0)|$ for all real $x$. The same type of bound holds
for $P'$, including when $m=1$ by its differential equation. Since
$w_*:=\min_{[-1,1]}W>0$,
\[
 \int_{-1}^1|P|^2\,dx\le w_*^{-2}\norm{WP}_2^2.
\]
Combining this with~\eqref{e:eq:observable} yields
\begin{equation}\label{e:eq:growth}
 |P(x)|+|P'(x)|\le C e^{L|x|}\norm{WP}_2.
\end{equation}
The integrals of $e^{2L|x|}W$, $e^{2L|x|}W^2$, and
$e^{2L|x|}|W'|^2$ are finite by~\eqref{e:eq:weight}.
Integrating~\eqref{e:eq:growth} and using $(WP)'=W'P+WP'$ proves the lemma
for order $m$. Take the maximum of the finitely many constants for
$1\le m\le N$.
\end{proof}

The compact family is formed only after frequencies are placed in a
bounded window. The companion matrices of an original configuration with
unbounded real frequencies need not be bounded. No inverse Vandermonde
matrix or lower bound on root spacing is used.

\subsection{Finite clusters at arbitrary real separations}\label{e:sec:clusters}

\begin{lemma}[Subsequential cluster decomposition]\label{e:lem:clusters}
Let $P_n\in\EE$ be nonzero. There is a subsequence, still indexed by $n$,
and a fixed finite collection of clusters $\mathcal C$ such that
\begin{equation}\label{e:eq:cluster}
 WP_n=\sum_{C\in\mathcal C}e^{it_{C,n}x}A_{C,n},\qquad
 A_{C,n}=WQ_{C,n},
\end{equation}
where each $t_{C,n}$ is the real part of an original frequency,
all frequencies of each $Q_{C,n}$ lie in a fixed bounded disk depending
on $C$, and $|t_{C,n}-t_{D,n}|\to\infty$ for $C\ne D$.
There are constants $L_C<\infty$, independent of $n$, such that
\begin{equation}\label{e:eq:cluster-envelope}
 \norm{A_{C,n}}_W+\norm{A_{C,n}'}_2\le L_C\norm{A_{C,n}}_2.
\end{equation}
Moreover,
\begin{equation}\label{e:eq:decouple}
 \norm{WP_n}_2^2=(1+o(1))
 \sum_{C\in\mathcal C}\norm{A_{C,n}}_2^2,
\end{equation}
where the relative error is independent of the cluster amplitudes.
\end{lemma}
\begin{proof}
Extract a subsequence on which the number of distinct frequencies and all
polynomial degrees are fixed; there are only finitely many such choices.
Label the frequencies along this subsequence. For every pair $j,k$, extract
further so that $\Rea\lambda_{j,n}-\Rea\lambda_{k,n}$ converges in the
extended real line. There are finitely many pairs. Declare $j\sim k$ when
this limit is finite. Reflexivity and symmetry are immediate; transitivity
follows by adding the two convergent real differences. The equivalence
classes are the clusters.

Choose a member $j_C$ in each cluster and put
\[
 t_{C,n}=\Rea\lambda_{j_C,n},\qquad
 Q_{C,n}(x)=\sum_{j\in C}p_{j,n}(x)
 e^{i(\lambda_{j,n}-t_{C,n})x}.
\]
Within a cluster the real offsets are bounded, while the imaginary parts
remain in $[-\beta,\beta]$. Between distinct clusters the carrier differences
tend in modulus to infinity. This proves the stated decomposition and,
by Lemma~\ref{e:lem:ode},~\eqref{e:eq:cluster-envelope}.

For distinct clusters, integration by parts gives
\begin{align}\label{e:eq:cross}
 \left|\ip{e^{it_{C,n}x}A_{C,n}}{e^{it_{D,n}x}A_{D,n}}_2\right|
 &\le\frac{\norm{A_{C,n}'}_2\norm{A_{D,n}}_2+
 \norm{A_{C,n}}_2\norm{A_{D,n}'}_2}{|t_{C,n}-t_{D,n}|}\notag\\
 &\le\frac{L_C+L_D}{|t_{C,n}-t_{D,n}|}
 \norm{A_{C,n}}_2\norm{A_{D,n}}_2.
\end{align}
The product $A_{C,n}\overline{A_{D,n}}$ belongs to $W^{1,1}(\R)$ by
Cauchy--Schwarz. Its endpoint limits are zero: the derivative is integrable,
so the limits exist, and its integrability forces those limits to vanish.
Thus the integration by parts is valid.

Write $s=|\mathcal C|\le N$ and $h_{C,n}=\norm{A_{C,n}}_2$. If $s>1$,
let $\delta_n$ be the largest coefficient on the right of~\eqref{e:eq:cross}.
Then $\delta_n\to0$ and
\[
 \left|\norm{WP_n}_2^2-\sum_C h_{C,n}^2\right|
 \le2\delta_n\sum_{C<D}h_{C,n}h_{D,n}
 \le(s-1)\delta_n\sum_C h_{C,n}^2.
\]
For $s=1$ the identity is exact. This proves~\eqref{e:eq:decouple}; in
particular, for all sufficiently large $n$,
\begin{equation}\label{e:eq:amplitude}
 \sum_C\norm{A_{C,n}}_2^2\le2\norm{WP_n}_2^2.
\end{equation}
\end{proof}

The constants $L_C$ depend on the bounded offset windows found along the
subsequence. They are not yet universal constants for the original class.
This is sufficient for the contradiction arguments that follow. Crucially,
\eqref{e:eq:amplitude} is established without any prior bound on the separate
envelope norms or on the raw coefficients.

\subsection{Proof of uniform comparison and coercivity}\label{e:sec:proof}

\begin{proof}[Proof of Theorem~\ref{e:thm:main}\textup{(i)}]
Suppose no constant in~\eqref{e:eq:comparison} existed. Choose nonzero
$P_n\in\EE$ with $\norm{WP_n}_W^2>n\norm{WP_n}_2^2$. Apply
Lemma~\ref{e:lem:clusters}. Multiplication by a real modulation is an
isometry in both weighted and ordinary norms. If $s=|\mathcal C|$ and
$L=\max_C L_C$, the triangle inequality and Cauchy--Schwarz give
\[
 \norm{WP_n}_W^2\le s\sum_C\norm{A_{C,n}}_W^2
 \le sL^2\sum_C\norm{A_{C,n}}_2^2
 \le2sL^2\norm{WP_n}_2^2
\]
for all sufficiently large $n$, a contradiction. Thus a single finite
constant exists for the entire class.
\end{proof}

\begin{proof}[Proof of Theorem~\ref{e:thm:main}\textup{(ii)}]
If~\eqref{e:eq:escape} failed, there would be $T_n\to\infty$ and
$P_n\in\EE(T_n)$ such that $\norm{WP_n}_2=1$ and the low-band Fourier
norm remained bounded below by a positive constant. Extract the clusters
of Lemma~\ref{e:lem:clusters}. Each $t_{C,n}$ is the real part of an original
frequency, so $|t_{C,n}|\ge T_n\to\infty$.

For $|t_{C,n}|>R$, Fourier translation and Plancherel imply
\begin{align}\label{e:eq:low-band}
 \norm{\mathbf1_{[-R,R]}
 \widehat{e^{it_{C,n}x}A_{C,n}}}_2^2
 &=\int_{|t|\le R}|\widehat A_{C,n}(t-t_{C,n})|^2\,dt\notag\\
 &\le\frac{\norm{A_{C,n}'}_2^2}{(|t_{C,n}|-R)^2}
 \le\frac{L_C^2\norm{A_{C,n}}_2^2}{(|t_{C,n}|-R)^2}.
\end{align}
Equation~\eqref{e:eq:amplitude} bounds every envelope norm by $\sqrt2$.
There are finitely many clusters, so the triangle inequality makes the
low-band norm of their sum tend to zero, a contradiction.
\end{proof}

\begin{proof}[Proof of Theorem~\ref{e:thm:main}\textup{(iii)}]
Let $C$ be the constant in~\eqref{e:eq:comparison} and put
$\varepsilon(T;R)=\varepsilon_{N,\beta,W}(T;R)$. For every fixed $R$,
Lemma~\ref{e:lem:gamma} and the bounded perturbation estimate give
\begin{equation}\label{e:eq:master-lower}
 \qD[WP]\ge\left\{
 \frac12\log(1+4R^2)\bigl(1-\varepsilon(T;R)^2\bigr)
 -\norm D C\right\}\norm{WP}_2^2
 \qquad(P\in\EE(T)).
\end{equation}
All Gamma frequencies remain in the original form; its nonnegative
low-frequency part was only discarded in this lower estimate. Given
$L>0$, first choose $R$ so that
$\tfrac14\log(1+4R^2)-\norm D C\ge L$. Then choose $T$ large enough
that $\varepsilon(T;R)^2\le1/2$. This proves the first divergence in
\eqref{e:eq:coercivity}. For a prescribed weighted coefficient $L_W>0$,
apply that result with $L=L_W C$ and use~\eqref{e:eq:comparison} to get
$\qD[WP]\ge L_W\norm{WP}_W^2$.
\end{proof}

The order of choices is important: first fix a Fourier band $R$, then
send the minimum real carrier height $T$ to infinity. The extraction proof
does not yield a logarithmic-in-$T$ rate for arbitrary configurations.


\subsection{An entire two frequency formula}\label{e:sec:kernel}

For the theta weight and $D_\theta$ of Section~\ref{e:sec:theta}, define
\begin{equation}\label{e:eq:overlap}
 e_z(x)=W(x)\cos(zx),\qquad
 H_s(t)=\int_\R W(x+s/2)W(x-s/2)\cos(tx)\,dx\quad(s\ge0).
\end{equation}
The following exact formula records the complete form on cosine vectors.
It is not a zero-sum representation and makes no zero-location assumption.

\begin{proposition}[Entire bilinear kernel]\label{e:prop:kernel}
The function
$\mathcal K(z,v)=\qth(e_z,e_{\overline v})$ is entire on $\C^2$ and is
symmetric in $z,v$, even in each variable, and real under simultaneous
complex conjugation. It satisfies
\begin{align}\label{e:eq:kernel}
 \mathcal K(z,v)
 &=\int_0^\infty\rho(s)\Bigl[
 H_0(z-v)+H_0(z+v)\notag\\[-2pt]
 &\hspace{3.5em}-\cos\bigl((z+v)s/2\bigr)H_s(z-v)
 -\cos\bigl((z-v)s/2\bigr)H_s(z+v)\Bigr]\,ds\notag\\
 &\quad-\frac{c_\Gamma}{2}\bigl[H_0(z-v)+H_0(z+v)\bigr]\notag\\
 &\quad-\sum_{n=2}^\infty a_n\Bigl[
 \cos\bigl((z+v)s_n/2\bigr)H_{s_n}(z-v)
 +\cos\bigl((z-v)s_n/2\bigr)H_{s_n}(z+v)\Bigr].
\end{align}
The Gamma bracket is kept combined. The integral and series converge
locally uniformly, and every fixed parameter derivative can be taken under
them. In particular,
\begin{equation}\label{e:eq:jet-matrix}
 \qth(J_{w,k},J_{v,\ell})
 =\left.\partial_z^k\partial_\eta^\ell\mathcal K(z,\eta)
 \right|_{z=w,\,\eta=\overline v}.
\end{equation}
\end{proposition}
\begin{proof}
For a compact parameter set, $e_z$ and each of its fixed parameter
derivatives have uniformly bounded $\HH$ and $H^1$ norms by the theta
tails. The same estimates, with one additional derivative, justify
complex differentiation by dominated convergence in both norms. Thus
$z\mapsto e_z$ is entire into $\XX$ by~\eqref{e:eq:h1bound}.

Set $C_s(z,v)=\int_\R e_z(x+s)e_v(x)\,dx$. With $y=x+s/2$, the weight
product is even in $y$. The product-to-sum identity, whose odd sine terms
integrate to zero, yields
\begin{equation}\label{e:eq:Cs}
 C_s(z,v)=\frac12\Bigl[
 \cos\bigl((z+v)s/2\bigr)H_s(z-v)
 +\cos\bigl((z-v)s/2\bigr)H_s(z+v)\Bigr].
\end{equation}
In particular $C_s$ is symmetric in $z,v$. Translation invariance gives
\[
 \int_\R(e_z(x+s)-e_z(x))(e_v(x+s)-e_v(x))\,dx
 =2C_0(z,v)-2C_s(z,v).
\]
The two prime orientations also both equal $C_{s_n}(z,v)$.
Substitution into the exact forms proves~\eqref{e:eq:kernel}.

On compact parameter sets, the Gamma bracket is $O(s^2)$ near zero,
bounded by $s^2\norm{e_z'}_2\norm{e_v'}_2$; for large $s$ it is bounded
by $4\norm{e_z}_2\norm{e_v}_2$. Both bounds are integrable against
$\rho(s)\,ds$. They remain valid for fixed parameter derivatives, with
the corresponding differentiated vectors. For the prime part,
Lemma~\ref{a:lem:prime} and its norm-summable weighted translation proof give absolute
locally uniform convergence, also after every fixed parameter derivative.
Each $H_s$ is entire by the same theta domination. This proves all
convergence and analyticity assertions. Evenness, symmetry and the real
conjugation identity follow from the formulas; differentiation gives
\eqref{e:eq:jet-matrix}. Notice that the sesquilinear kernel is
$\qth(e_z,e_v)=\mathcal K(z,\overline v)$, not $\mathcal K(z,v)$.
\end{proof}


\section{Computational evidence and public reproducibility}\label{app:provenance-unified}
The analytic arguments are proved in this single source file. The following evidence identifies computational premises without treating a hash as a proof, and without citing an earlier manuscript in place of an analytic argument.


\subsection{Computational input identity and verification boundary}\label{c:app:provenance}
The proof has two different computational input classes.
Hypothesis~\ref{c:hyp:local} is an accepted whole-domain local theorem at ${B_0}=4/5$.
Its supplied packet gives $q_I[u]\ge10^{-22}\norm u_2^2$.
The records below identify that input and its computational evidence.
They are data-provenance identifiers, not citations to an earlier manuscript as a substitute for an analytic proof.
A cryptographic identity does not establish an inequality.

\begingroup\small
\begin{longtable}{@{}p{0.96\textwidth}@{}}
\toprule Local input artifact and SHA-256\\\midrule\endhead
Local theorem source\\*
\nolinkurl{ac04a393e8b6112d81bfb598c2dd00450b2a22eca36c298e3bbe82f7ef78b4c6}\\[3pt]
Local certificate archive version 1.1\\*
\nolinkurl{b805a431c4c77d955d5a6d4ed7397a01faf12cd727c1051af607ed14f77ab74b}\\[3pt]
Radius-four-fifths source diagnostic\\*
\nolinkurl{4c67f0c3f37bb8ce3c4f7c6da166050f2a1760c62bd47227c7709fb3a0b2ea05}\\[3pt]
Radius-four-fifths producer success record\\*
\nolinkurl{19156f6635bb25e853b9a1aecd650672043113f85d279c562e1a090269b0ca91}\\[3pt]
Exact-rational replay source\\*
\nolinkurl{3c5e006df82fbc985ce9b7ed73e26d62a9679086cc399dcb022224325c9fde13}\\[3pt]
Exact-rational replay result\\*
\nolinkurl{f2d4eef6e3efcd632ee252b2cca56be120e5ce5b3c59cdcc499277bc4ee528e2}\\[3pt]
Separate MPFR input archive\\*
\nolinkurl{4e57e261db7c8cfebdc5dc82b23d8aecfab3b639aa7fa5c95901b90f3aa18f8b}\\[3pt]
Separate MPFR evidence\\*
\nolinkurl{2a6fe3e8dfa04308ed555271bb19441b7d3ce783b6cc427419840a2d39a231d8}\\
\bottomrule
\end{longtable}\endgroup
The full local minimum domain, ordinary physical $L^2$ norm, positive mean pole, exact local prime powers $2,3,4$, and complete Gamma remainder are part of the input statement.
Finite matrix positivity alone is not the input theorem.
The supplied exact-rational replay and separately supplied source-enclosure evidence represent different validation layers.
No source generation, quadrature, local certificate, rational replay, or numerical sign gate was executed in preparing the present analytic manuscript.

The second class consists of Hypothesis~\ref{c:hyp:response}(a),(b).
One fixed data source for those contracts is the archive
\texttt{centered\_exterior\_cost\_sources.zip}, SHA-256
\begin{center}\small
\nolinkurl{4984450bc9f3174d354fdfe9731919dfb4c134f2db1ef0b55c40bc701b04792e}.
\end{center}
The relevant member identifiers are
\begingroup\small
\begin{longtable}{@{}p{0.96\textwidth}@{}}
\toprule Response input member and SHA-256\\\midrule\endhead
\texttt{centering\_vector.json}\\*
\nolinkurl{c04a7d7ac467971ae101290adbaec719826af343d7c0f589bdb7789c8ca8b4ed}\\[3pt]
\texttt{run/evidence.json}\\*
\nolinkurl{6c404a912f5986ae08080f2811a5d797d9aed22fc5683e24c9f2241525e4d3cd}\\[3pt]
\texttt{run/cells.json}\\*
\nolinkurl{d93657d28a399250f87f0d12f5421ced5c72178ff874c67ec5a3087a387ea6c7}\\[3pt]
\texttt{local\_response/run1024/response\_coefficients.json}\\*
\nolinkurl{b40fff3249a60dbb1224c9099360c96333d19fe22b7bd10bda126e9ece44f7a9}\\[3pt]
\texttt{local\_response/run1024/exact\_matrices.json}\\*
\nolinkurl{792718e3d299d74ca5757e5f17942608ab005468759673bd2f36388cf102dfba}\\[3pt]
\texttt{local\_response/run1024/evidence.json}\\*
\nolinkurl{62232dc8890083bfd967ad7fa3e854fe4ca00925b7b493a45dc7a45ecf234cbb}\\
\bottomrule
\end{longtable}\endgroup
The archive also identifies 60 finite local Gram blocks covering rows 1024 through 16383 and the infinite-tail enclosure.
Its local parameters are 1024 trial modes, cutoff 16384, 192-bit precision, 64 Gamma labels and tail order six.
The exterior parameters are fixed in Appendix~\ref{c:app:rows}.
These identify the declared inputs, not a new certification.
File identity, finite rational replay, and full source-enclosure validation remain distinct.
The clean source distribution contains no certificate data or executable research code; its build compiles the document only.


\subsection{Public supporting archive and verification boundary}
The immutable public supporting archive is
\begin{center}\small
\href{https://github.com/serabing-hash/theta-weighted-certificate/tree/f539e52fe08379c0639ff94a812ec73a40d29a0e}{\texttt{serabing-hash/theta-weighted-certificate}}\\
commit \nolinkurl{f539e52fe08379c0639ff94a812ec73a40d29a0e}.
\end{center}
The accepted radius-$4/5$ local certificate package is available in that commit at
\begin{center}\small
\href{https://github.com/serabing-hash/theta-weighted-certificate/tree/f539e52fe08379c0639ff94a812ec73a40d29a0e/research/2026-10-28-four-fifths-local-certificate}{\texttt{research/2026-10-28-four-fifths-local-certificate}}.
\end{center}
Its complete local-domain theorem, source-enclosure layer and exact-rational replay layer must be distinguished. No numerical research program, prime or zero enumeration, local certificate replay, response solve, quadrature or source enclosure was executed in preparing this unified manuscript.
The exact source archive with SHA-256 beginning \texttt{b805a431} in the table above is directly linked as
\href{https://github.com/serabing-hash/theta-weighted-certificate/blob/f539e52fe08379c0639ff94a812ec73a40d29a0e/research/2026-10-28-four-fifths-local-certificate/prime_gamma_local_certificate_sources.zip}{\texttt{prime\_gamma\_local\_certificate\_sources.zip}} in the same immutable commit.

The additional response archive identified above is not asserted to be publicly downloadable from that commit. The response theorem is an implication from the explicitly displayed complete-operator contracts, not a claim that their enclosure premises were newly certified or made public. This distinction does not affect the analytic implication, but it limits independent computational reproduction of that branch until the corresponding source archive is published.

This manuscript builds from one \texttt{.tex} file with standard packages. Every proof appendix and the bibliography are inline; no figures, bibliography database, external input file or certificate data are required to compile it. Document compilation, PDF rendering and file hashes verify artifact production only, not the mathematical correctness of computational input statements.

\section*{License}
Copyright \textcopyright\ 2026 Jongmin Choi. This work is licensed under the Creative Commons Attribution-NonCommercial 4.0 International License (CC BY-NC 4.0): \url{https://creativecommons.org/licenses/by-nc/4.0/}. External references retain their respective rights.

\begin{thebibliography}{9}
\bibitem{Guinand}
A.~P. Guinand, A summation formula in the theory of prime numbers,
\emph{Proceedings of the London Mathematical Society} (2) \textbf{50} (1948), 107--119.
\url{https://doi.org/10.1112/plms/s2-50.2.107}.
\bibitem{vonMangoldt}
H.~von Mangoldt, Zur Verteilung der Nullstellen der Riemannschen Funktion,
\emph{Mathematische Annalen} \textbf{60} (1905), 1--19.
\url{https://doi.org/10.1007/BF01447494}.
\bibitem{Kato}
T.~Kato, \emph{Perturbation Theory for Linear Operators}, second edition,
Classics in Mathematics, Springer, Berlin, 1995, Chapter~VI.
\url{https://doi.org/10.1007/978-3-642-66282-9}.
\bibitem{deBranges1961}
L.~de Branges, Some Hilbert spaces of entire functions.~II,
\emph{Transactions of the American Mathematical Society} \textbf{99} (1961), 118--152.
\url{https://doi.org/10.1090/S0002-9947-1961-0133456-2}.
\bibitem{DLMF}
National Institute of Standards and Technology, \emph{Digital Library of Mathematical Functions}.
Jacobi transformation, Eq.~20.7.32; Gamma integral, reflection, beta integral and uniform complex Stirling formulas, Sections~5.5, 5.9, 5.11 and 5.12; completed zeta function and Euler product, Sections~25.2 and 25.4.
\url{https://dlmf.nist.gov/20.7}; \url{https://dlmf.nist.gov/5.9};
\url{https://dlmf.nist.gov/5.11}; \url{https://dlmf.nist.gov/5.12};
\url{https://dlmf.nist.gov/25.2}; \url{https://dlmf.nist.gov/25.4}.
\end{thebibliography}
\end{document}
